\beginnumbering
\eledsection*{Adhyāya 1}
\pstart \ledsidenote{\passref{6.1.1}}\edtext{\skt{अथातः}}{\Afootnote{K; \skt{अथात}\ A H}}\skt{ औपद्रविकमध्यायं व्याख्यास्यामः ॥}\pend

\pstart\versepar \ledsidenote{\passref{6.1.3}}\edtext{\skt{सविंशे तु}}{\Afootnote{K H; \emph{om.}\ A}}\skt{ }\edtext{\skt{शते}}{\Afootnote{H; \emph{om.}\ A; \skt{शत}\ K}}\skt{ }\edtext{\skt{पूर्वमध्यायानां}}{\lemma{\skt{पूर्वम्}}\Afootnote{K H; \emph{om.}\ A}}\skt{ }\edtext{\skt{तु}}{\Afootnote{K; \skt{शते विंशे यदु}\ A; \emph{om.}\ H}}\skt{ }\edtext{\skt{कीर्त्तिते}}{\Afootnote{K; \skt{क्तमस}\ A; \skt{तु}\ H}}\skt{ । }\pend
\pstart\versepar \skt{}\edtext{\skt{वक्ष्ये}}{\Afootnote{K; \skt{कृन्मया । वक्ष्यामि}\ A; \skt{कीर्त्तितः । वक्ष्ये}\ H}}\skt{ }\edtext{\skt{बहुविधां}}{\Afootnote{K; \skt{बहुधा}\ A; \skt{वहुविधा न}\ H}}\skt{ }\edtext{\skt{सम्यगुत्तरेऽर्थानतः}}{\lemma{\skt{ऽर्थानतः}}\Afootnote{H; \skt{ऽर्थानि}\ A; \skt{र्थानतः}\ K}}\skt{ }\edtext{\skt{परं}}{\Afootnote{K H; \skt{मानिति ॥}\ A}}\skt{ ।}\pend

\pstart\versepar \ledsidenote{\passref{6.1.4}}\edtext{\skt{इदानीन्तत्प्रवक्ष्यामि}}{\lemma{\skt{तत्प्रवक्ष्यामि}}\Afootnote{A K; \skt{सम्प्रवक्ष्यामि}\ H}}\skt{ }\edtext{\skt{तन्त्रमुत्तरसंज्ञितं}}{\lemma{\skt{उत्तरसंज्ञितं}}\Afootnote{K H; \skt{उत्तरमुत्तमम् ।}\ A}}\skt{ । }\pend
\pstart\versepar \skt{}\edtext{\skt{निखिलेनोपदृष्यन्ते}}{\Afootnote{K; \skt{निखिलेनोपदिश्यन्ते}\ A; \skt{निखिलेनोपदेक्ष्यन्ते}\ H}}\skt{ यत्र रोगाः पृथग्विधाः ॥}\pend

\pstart\versepar \ledsidenote{\passref{6.1.5}}\edtext{\skt{शालाक्यशास्त्रा विहिता}}{\Afootnote{K; \skt{शालाक्यतन्त्राभिहिता}\ A; \skt{शालाक्यशास्त्रार्भिहिता}\ H}}\skt{ विदेहाधियकीर्त्तिता । }\pend
\pstart\versepar \skt{ये च }\edtext{\skt{विस्तरशो}}{\Afootnote{K\textsuperscript{pc} H; \skt{विस्तरतो}\ A; \skt{विरस्तरशो}\ K\textsuperscript{ac}}}\skt{ दृष्टाः }\edtext{\skt{कुमाराबाधहेतवः}}{\Afootnote{A K; \skt{कुमारावाधहेतवः ॥}\ H}}\skt{ ।}\edtext{}{\lemma{\skt{विदेहाधियकीर्त्तिता । ये च}}\Afootnote{K; \skt{विदेहाधिपकीर्तिताः ये च}\ A H}}\pend

\pstart\versepar \ledsidenote{\passref{6.1.6}}\skt{षट्सु कायचिकित्सासु ये चोक्ताः }\edtext{\skt{परमर्षिणा}}{\Afootnote{K; \skt{परमर्षिभिः ।}\ A; \skt{परर्षिणा ।}\ H}}\skt{ । }\pend
\pstart\versepar \skt{उपसर्गादयो }\edtext{\skt{दोषा}}{\Afootnote{K H; \skt{रोगा}\ A}}\skt{ }\edtext{\skt{येनाप्यागन्तवः}}{\lemma{\skt{येनाप्य्}}\Afootnote{K; \skt{ये चाप्य्}\ A; \skt{यो चाप्य्}\ H}}\skt{ स्मृताः ॥}\pend

\pstart\versepar \ledsidenote{\passref{6.1.7}}\edtext{\edtext{\skt{त्रिषष्टिरससंसर्ग्गाः}}{\lemma{\skt{त्रिषष्टिर्}}\Afootnote{K; \skt{त्रिषष्टी}\ A; \skt{त्रिषष्टि}\ H}}}{\lemma{\skt{अससंसर्ग्गाः}}\Afootnote{K; \skt{रससंसर्गाः}\ A H}}\skt{ स्वस्थवृत्तन्तथैव च । }\pend
\pstart\versepar \skt{युक्तार्था }\edtext{\skt{युक्तयश्चापि}}{\lemma{\skt{चापि}}\Afootnote{K H; \skt{चैव}\ A}}\skt{ }\edtext{\skt{दोषभेदाश्च}}{\Afootnote{K H; \skt{दोषभेदास्त}\ A}}\skt{ }\edtext{\skt{तत्वतः}}{\Afootnote{K H; \skt{थैव च ॥}\ A}}\skt{ ॥}\pend

\pstart\versepar \ledsidenote{\passref{6.1.8}}\skt{यत्रोक्ता }\edtext{\edtext{\skt{विविधाश्चार्था}}{\lemma{\skt{विविधाश्}}\Afootnote{K H; \skt{विविधा}\ A}}}{\lemma{\skt{चार्था}}\Afootnote{K H; \skt{अर्था}\ A}}\skt{ }\edtext{\skt{रोगसाधनहेतवः}}{\Afootnote{A H; \skt{रोगसाधणहेतवः ।}\ K}}\skt{ । }\pend
\pstart\versepar \skt{}\edtext{\skt{अस्य}}{\Afootnote{K H; \skt{महतस्तस्य}\ A}}\skt{ तन्त्रस्य }\edtext{\skt{महतः}}{\Afootnote{K H; \skt{दुर्गाधस्याम्बुधेरिव ॥}\ A}}\skt{ }\edtext{\skt{सागराकारवर्चसः}}{\Afootnote{K H; \emph{om.}\ A}}\skt{ ।}\pend

\pstart\versepar \ledsidenote{\passref{6.1.9}}\edtext{\skt{आदावेवोत्तमांगस्थां}}{\lemma{\skt{एवोत्तमांगस्थां}}\Afootnote{K; \skt{एवोत्तमाङ्गस्था}\ A H}}\skt{ }\edtext{\skt{रोगा}}{\Afootnote{K; \skt{न्रोगा}\ A H}}\skt{ }\edtext{\skt{समभिधास्यते}}{\Afootnote{K; \skt{नभिदधाम्यहम् ।}\ A; \skt{न्समभिधास्यते ।}\ H}}\skt{ । }\pend
\pstart\versepar \skt{संख्यया }\edtext{\skt{लक्षणैश्चापि}}{\lemma{\skt{लक्षणैश्}}\Afootnote{A K\textsuperscript{pc} H; \skt{लणैश्}\ K\textsuperscript{ac}}}\skt{ साध्यासाध्यक्रमेण च ।}\pend

\pstart\versepar \ledsidenote{\passref{6.1.10}}\edtext{\skt{विद्याद्यंगुलबाहुल्यं}}{\Afootnote{K; \skt{विद्याद्द्वयङ्गुलबाहुल्यं}\ A; \skt{विद्याद्व्यङ्गुलवाहुल्यं}\ H}}\skt{ }\edtext{\edtext{\skt{स्वांगुष्ठोदरमायतं}}{\lemma{\skt{स्वांगुष्ठोदरम्}}\Afootnote{K; \skt{स्वाङ्गुष्ठोदर}\ A H}}}{\lemma{\skt{आयतं}}\Afootnote{K; \skt{संमितम् ।}\ A; \skt{मायतं ।}\ H}}\skt{ । }\pend
\pstart\versepar \skt{द्व्यङ्गुलं सर्वतः सार्धं }\edtext{\skt{भिषङ्नयनबुद्बुदं}}{\lemma{\skt{नयनबुद्बुदं}}\Afootnote{A K; \skt{नयनवुद्वुदं ॥}\ H}}\skt{ ।}\pend

\pstart\versepar \ledsidenote{\passref{6.1.11}}\edtext{\skt{वृत्तं}}{\Afootnote{K H; \skt{सुवृत्तं}\ A}}\skt{ }\edtext{\skt{गोस्तनकाकारं}}{\Afootnote{K H; \skt{गोस्तनाकारं}\ A}}\skt{ सर्वभूतगुणोद्भवं । }\pend
\pstart\versepar \skt{}\edtext{\skt{वलं}}{\Afootnote{K H; \skt{पलं}\ A}}\skt{ }\edtext{\skt{मृद्भ्योग्नितो}}{\Afootnote{K; \skt{भुवोऽग्नितो}\ A; \skt{मृड्भ्योऽग्नितो}\ H}}\skt{ रक्तं वातात्कृष्णं सितं जलात् ।}\pend

\pstart\versepar \ledsidenote{\passref{6.1.12}}\edtext{\skt{आकाशादश्रुमार्गस्तु}}{\Afootnote{K H; \skt{आकाशादश्रुमार्गाश्च}\ A}}\skt{ जायन्ते नेत्र }\edtext{\skt{बुद्बुदे}}{\Afootnote{A K; \skt{वुद्वुदे ।}\ H}}\skt{ । }\pend
\pstart\versepar \skt{}\edtext{\skt{दृष्टिश्चात्र}}{\lemma{\skt{दृष्टिश्}}\Afootnote{K; \skt{दृष्टिं}\ A; \skt{दृष्टि}\ H}}\skt{ }\edtext{\skt{यथा}}{\Afootnote{K H; \skt{तथा}\ A}}\skt{ }\edtext{\skt{वक्ष्ये}}{\Afootnote{A K; \skt{वक्ष्यो}\ H}}\skt{ }\edtext{\skt{तथा}}{\Afootnote{K H; \skt{यथा}\ A}}\skt{ }\edtext{\skt{विद्याद्विचक्षणः}}{\Afootnote{K H; \skt{ब्रूयाद्विशारदः ॥}\ A}}\skt{ ।}\pend

\pstart\versepar \ledsidenote{\passref{6.1.13}}\edtext{\skt{नेत्रायामविभागस्तु}}{\Afootnote{K; \skt{नेत्रायामत्रिभागं तु}\ A; \skt{नेत्रायामविभागस्तु}\ H}}\skt{ }\edtext{\edtext{\skt{कृष्णम्मण्डलमुच्यते}}{\lemma{\skt{कृष्णम्मण्डलम्}}\Afootnote{K H; \skt{कृष्णमण्डल}\ A}}}{\lemma{\skt{उच्यते}}\Afootnote{K H; \skt{मुच्यते ।}\ A}}\skt{ । }\pend
\pstart\versepar \skt{}\edtext{\skt{कृष्णात्सप्तममिच्छन्ति}}{\lemma{\skt{कृष्णात्सप्तमम्}}\Afootnote{A K; \skt{कृत्सप्तमम्}\ H}}\skt{ दृष्टिं }\edtext{\skt{दृष्टिविदो}}{\Afootnote{K H; \skt{दृष्टिविशा}\ A}}\skt{ }\edtext{\skt{जनाः}}{\Afootnote{K H; \skt{रदाः ॥}\ A}}\skt{ ।}\pend

\pstart\versepar \ledsidenote{\passref{6.1.14}}\edtext{\skt{मण्डलानि च}}{\Afootnote{A K H\textsuperscript{ac}; \skt{लानि च}\ H\textsuperscript{pc}}}\skt{ }\edtext{\skt{सन्धीश्च}}{\Afootnote{K; \skt{सन्धींश्च}\ A H}}\skt{ पटलानि च लोचने । }\pend
\pstart\versepar \skt{यथाक्रमं }\edtext{\skt{विजानीयात्पञ्चषट्च}}{\Afootnote{A K; \skt{विजानीयात्पञ्चषट्चा}\ H}}\skt{ षडेव च ।}\pend

\pstart\versepar \ledsidenote{\passref{6.1.15}}\edtext{\skt{पक्ष्मवर्त्मसितश्यावदृष्टीनाम्मण्डलानि}}{\lemma{\skt{पक्ष्मवर्त्मसितश्यावदृष्टीनाम्}}\Afootnote{K; \skt{पक्ष्मवर्त्मश्वेतकृष्णदृष्टीनां}\ A; \skt{पक्ष्मवर्त्मसितश्यावदृष्टीनां}\ H}}\skt{ तु । }\pend
\pstart\versepar \skt{}\edtext{\skt{अनुपूर्वन्तु}}{\Afootnote{K H; \skt{अनुपूर्वं तु ते}\ A}}\skt{ }\edtext{\skt{चत्वारो}}{\Afootnote{K H; \skt{मध्याश्चत्वारो}\ A}}\skt{ }\edtext{\skt{मध्यादुक्ता}}{\Afootnote{H; \skt{ऽन्त्या}\ A; \skt{मध्यादुक्ता}\ K}}\skt{ }\edtext{\skt{यथाक्रमं}}{\Afootnote{K H; \skt{यथोत्तरम् ॥}\ A}}\skt{ ।}\pend

\pstart\versepar \ledsidenote{\passref{6.1.16}}\edtext{\skt{पक्ष्मवर्मगतः}}{\Afootnote{K; \skt{पक्ष्मवर्त्मगतः}\ A; \skt{पक्ष्मवर्त्म्मगतः}\ H}}\skt{ }\edtext{\edtext{\skt{सन्धिर्वर्त्मशुक्लगतोपरः}}{\lemma{\skt{सन्धिर्}}\Afootnote{A H; \skt{सन्धि}\ K}}}{\lemma{\skt{वर्त्मशुक्लगतोपरः}}\Afootnote{\skt{वर्त्मशुक्लगतोऽपरः}\ A K; \skt{व्वर्त्म्मशुक्लगतः परः ।}\ H}}\skt{ । }\pend
\pstart\versepar \skt{}\edtext{\skt{शुक्लकृष्णगतस्त्वन्यः}}{\Afootnote{\skt{शुक्लकृष्णगतस्त्वन्यः}\ A K; \skt{शुक्लदृष्टिगतस्त्वन्यः}\ H}}\skt{ }\edtext{\skt{कृष्णदृष्टिगतस्तथा}}{\Afootnote{H; \skt{कृष्णदृष्टिगतोऽपरः ।}\ A; \skt{कृष्णदृष्टिगतस्तथा ।}\ K}}\skt{ । }\pend
\pstart\versepar \skt{ततः }\edtext{\skt{कनीनिकगतः}}{\Afootnote{K H; \skt{कनीनकगतः}\ A}}\skt{ }\edtext{\edtext{\skt{षष्ठश्चापांगगः}}{\lemma{\skt{षष्ठश्}}\Afootnote{A K; \skt{षष्ठ}\ H}}}{\lemma{\skt{चापांगगः}}\Afootnote{A K; \skt{श्चापाङ्गगः}\ H}}\skt{ स्मृतः ।}\pend

\pstart\versepar \ledsidenote{\passref{6.1.17}}\edtext{\skt{द्वे वर्त्मपटले}}{\Afootnote{A K; \skt{द्वे वर्त्म्मपटले}\ H}}\skt{ }\edtext{\edtext{\skt{विद्याच्चत्वार्यन्यानि}}{\lemma{\skt{विद्याच्चत्वार्य्}}\Afootnote{A K; \skt{विद्याच्चत्वा}\ H}}}{\lemma{\skt{अन्यानि}}\Afootnote{A K; \skt{र्यन्यानि}\ H}}\skt{ }\edtext{\skt{चाक्ष्णि तु}}{\Afootnote{K; \skt{चाक्षिणि ।}\ A; \skt{चाक्षि तु ॥}\ H}}\skt{ । }\pend
\pstart\versepar \skt{जायते तिमिरं येषु }\edtext{\skt{व्याधिपरमदारुणः}}{\Afootnote{K; \skt{व्याधिः परमदारुणः}\ A H}}\skt{ ।}\pend

\pstart\versepar \ledsidenote{\passref{6.1.18}}\edtext{\skt{तेजो}}{\Afootnote{A K; \skt{तेजो}\ H}}\skt{ जलाश्रितं }\edtext{\skt{बाह्यं}}{\Afootnote{A K; \emph{om.}\ H}}\skt{ }\edtext{\edtext{\skt{तेष्वन्यपिशिताश्रितं}}{\lemma{\skt{तेष्व्}}\Afootnote{A K; \skt{…}\ H}}}{\lemma{\skt{अन्यपिशिताश्रितं}}\Afootnote{K; \skt{अन्यत्पिशिताश्रितम् ।}\ A; \skt{ष्वन्यत्पिशिताश्रितं ॥}\ H}}\skt{ । }\pend
\pstart\versepar \skt{}\edtext{\edtext{\edtext{\edtext{\skt{मेदस्तृतीयम्पटलमाश्रितं}}{\lemma{\skt{मेदस्}}\Afootnote{A K; \skt{मेद}\ H}}}{\lemma{\skt{तृतीयम्}}\Afootnote{A K; \skt{स्तृतीय}\ H}}}{\lemma{\skt{पटलम्}}\Afootnote{A K; \skt{पटल}\ H}}}{\lemma{\skt{आश्रितं}}\Afootnote{A K; \skt{माश्रितं}\ H}}\skt{ त्वस्थि चापरं ।}\pend

\pstart\versepar \ledsidenote{\passref{6.1.19}}\edtext{\skt{पञ्चमाशं}}{\Afootnote{K; \skt{पञ्चमांश}\ A; \skt{पञ्चमांस}\ H}}\skt{ }\edtext{\edtext{\skt{समन्दृष्टेस्तेषां}}{\lemma{\skt{समन्दृष्टेस्}}\Afootnote{A K; \skt{समन्दृष्टे}\ H}}}{\lemma{\skt{तेषां}}\Afootnote{A K; \skt{स्तेषां}\ H}}\skt{ }\edtext{\skt{बाहुल्यमिष्यते}}{\lemma{\skt{बाहुल्यम्}}\Afootnote{A K; \skt{वाहल्यम्}\ H}}\skt{ । }\pend
\pstart\versepar \skt{सिराणां कण्डराणाञ्च मेदसः }\edtext{\skt{कृष्णकस्य}}{\Afootnote{K H; \skt{कालकस्य}\ A}}\skt{ च}\pend

\pstart\versepar \ledsidenote{\passref{6.1.20}}\edtext{\skt{गुणा}}{\Afootnote{K; \skt{गुणाः}\ A H}}\skt{ }\edtext{\skt{कृष्णात्परं}}{\Afootnote{K H; \skt{कालात्परः}\ A}}\skt{ }\edtext{\skt{श्लेष्मा}}{\Afootnote{A K; \skt{श्ल्लेष्मा}\ H}}\skt{ }\edtext{\skt{बन्धनेक्ष्णोः}}{\Afootnote{K; \skt{बन्धनेऽक्ष्णोः}\ A; \skt{वन्धनेऽक्ष्णोः}\ H}}\skt{ सिरायुतः । }\pend
\pstart\versepar \skt{}\edtext{\edtext{\skt{सिरानुसारिभिर्दोषैर्विगुणैरूर्ध्वमाश्रितैः}}{\lemma{\skt{सिरानुसारिभिर्}}\Afootnote{A K; \skt{शिरानुसारिभिर्}\ H}}}{\lemma{\skt{विगुणैरूर्ध्वमाश्रितैः}}\Afootnote{\skt{विगुणैरूर्ध्वमागतैः ॥}\ A; \skt{विगुणैरूर्ध्वमाश्रितैः ।}\ K; \skt{व्विगुणैरूर्द्ध्वमास्थितैः ।}\ H}}\skt{ ।}\pend

\pstart\versepar \ledsidenote{\passref{6.1.21}}\skt{जायन्ते नेत्रभागेषु }\edtext{\edtext{\skt{व्याधयश्चास्तु}}{\lemma{\skt{व्याधयश्}}\Afootnote{K H; \skt{रोगाः प}\ A}}}{\lemma{\skt{चास्तु}}\Afootnote{K; \skt{रम}\ A; \skt{त्वाशु}\ H}}\skt{ दारुणाः । }\pend
\pstart\versepar \skt{}\edtext{\skt{तत्राविलं}}{\Afootnote{A K; \skt{तत्रापि…}\ H}}\skt{ }\edtext{\skt{ससंरम्भमश्रुकण्डूपदेहवत्}}{\Afootnote{K; \skt{ससंरम्भमश्रुकण्डूपदेहवत् ॥}\ A; \skt{मस्रुकण्डूपदेहवत् ।}\ H}}\skt{ ।}\pend

\pstart\versepar \ledsidenote{\passref{6.1.22}}\edtext{\edtext{\skt{गुरूषाचोपरागाद्यैर्जुष्टञ्चाव्यक्तलक्षणैः}}{\lemma{\skt{गुरूषाचोपरागाद्यैर्}}\Afootnote{K; \skt{गुरूषातोदरागाद्यैर्}\ A; \skt{गुरूषाचोषरागाद्यै}\ H}}}{\lemma{\skt{जुष्टञ्चाव्यक्तलक्षणैः}}\Afootnote{A K; \skt{र्जुष्टं चाव्यक्तलक्षणैः ॥}\ H}}\skt{ । }\pend
\pstart\versepar \skt{}\edtext{\skt{सशूलं}}{\Afootnote{A K; \skt{सशूल}\ H}}\skt{ }\edtext{\skt{वर्त्मकोशेषु}}{\Afootnote{K; \skt{वर्त्मकोषेषु}\ A; \skt{वर्त्यकोशेषु}\ H}}\skt{ }\edtext{\skt{शुकपूर्ण्णाभमेव}}{\lemma{\skt{शुकपूर्ण्णाभम्}}\Afootnote{K; \skt{शूकपूर्णाभम्}\ A; \skt{शूककर्ण्णाभम्}\ H}}\skt{ }\edtext{\skt{च}}{\Afootnote{A K; \skt{वा ।}\ H}}\skt{ ।}\pend

\pstart\versepar \ledsidenote{\passref{6.1.23}}\skt{विहन्यमानं }\edtext{\skt{रूपेषु}}{\Afootnote{K H; \skt{रूपे वा}\ A}}\skt{ }\edtext{\skt{क्रियास्वक्षि}}{\Afootnote{H; \skt{क्रियास्वक्षि}\ A K}}\skt{ }\edtext{\skt{यथापरं}}{\Afootnote{K; \skt{यथा पुरा ।}\ A; \skt{यथा पुरं ॥}\ H}}\skt{ । }\pend
\pstart\versepar \skt{}\edtext{\skt{दृष्ट्वैवम्धीमाञ्जानीयाद्दोषैणाधिष्ठितन्तु}}{\Afootnote{\skt{दृष्ट्वैव धीमान्बुध्येत दोषेणाधिष्ठितं तु}\ A; \skt{दृष्ट्वैवम्धीमाञ्जानीयाद्दोषैणाधिष्ठितन्तु}\ K; \skt{दृष्ट्वैव धीमान्जानीयाद्दोषेणाधिष्टितं तु}\ H}}\skt{ तत् ।}\pend

\pstart \ledsidenote{\passref{6.1.24}}\skt{तत्र }\edtext{\skt{संचयमासाद्य}}{\Afootnote{\skt{संभवमासाद्य यथादोषं}\ A; \skt{संचयमासाद्य}\ K H}}\skt{ }\edtext{\skt{विदध्यात्तत्र}}{\Afootnote{K H; \skt{भिषग्जितम् । विदध्यान्नेत्र}\ A}}\skt{ }\edtext{\skt{भेषजं}}{\Afootnote{K H; \skt{जा रोगा बलवन्तः स्युरन्यथा ॥}\ A}}\skt{ ।}\pend

\pstart\versepar \ledsidenote{\passref{6.1.25}}\skt{संक्षेपतः क्रियायोगो निदानपरिवर्जनं । }\pend
\pstart\versepar \skt{वातादीनां }\edtext{\skt{प्रतीकारःप्रोक्तो}}{\Afootnote{K H; \skt{प्रतीघातः प्रोक्तो}\ A}}\skt{ }\edtext{\skt{विस्तरतः}}{\Afootnote{A K; \skt{विस्तरणः}\ H}}\skt{ }\edtext{\skt{पुरा}}{\Afootnote{K H; \skt{पुनः ॥}\ A}}\skt{ ।}\pend

\pstart \ledsidenote{\passref{6.1.26}}\edtext{\skt{उष्णाभितप्तस्य}}{\Afootnote{A K; \skt{उष्णाभिघातादतिमै}\ H}}\skt{ }\edtext{\skt{जलप्रवेशाद्दूरेक्षणात्स्वप्नविपर्ययाच्च}}{\Afootnote{\skt{जलप्रवेशाद्दूरेक्षणात्स्वप्नविपर्ययाच्च । प्रसक्तसंरोदनकोपशोकक्लेशाभिघातादतिमैथुनाच्च ॥}\ A; \skt{जलप्रवेशाद्दूरेक्षणात्स्वप्नविपर्ययाच्च ।}\ K; \skt{थुनाच्च वाष्पग्रहात्सूक्ष्मनिरीक्षणाच्च स्वेदाद्रजो धूमनिषेवणाद्वा छर्द्दोऽभिघाताद्वमनातियोगात् ॥}\ H}}\skt{ ।}\pend

\pstart\versepar \ledsidenote{\passref{6.1.27}}\edtext{\skt{स्वेसद्रजोधूमनिषेवणाद्वा}}{\Afootnote{\skt{शुक्तारनालाम्लकुलत्थमाषनिषेवणाद्वे}\ A; \skt{स्वेसद्रजोधूमनिषेवणाद्वा}\ K}}\skt{ छर्दोभिघाताद्वमनातियोगात् । }\pend
\pstart\versepar \skt{तथा }\edtext{\skt{द्रवान्नातिनिषेवणाच्च}}{\Afootnote{\skt{मनिषेवणाच्च}\ A; \skt{द्रवान्नातिनिषेवणाच्च}\ K\textsuperscript{pc}; \skt{द्रवान्नातिमिषेवणाच्च}\ H\textsuperscript{pc}}}\skt{ }\edtext{\skt{विण्मूत्रवातक्रमनिग्रहाच्च}}{\Afootnote{\skt{छर्देर्विघाताद्वमनातियोगा}\ A; \skt{विण्मूत्रवातक्रमनिग्रहाच्च}\ K\textsuperscript{pc} H\textsuperscript{pc}}}\skt{ । }\pend
\pstart\versepar \skt{प्रसक्तसंरोदनशोककोपाच्छिराभिघातादतिमद्ययोगात् । }\pend
\pstart\versepar \skt{नाञ्च विपर्ययेण }\edtext{\skt{क्लेशाभिघातादतिमैथुनाच्च}}{\Afootnote{\skt{क्लेशाभिघातादतिमैथुनाच्च}\ K H}}\skt{ । }\pend
\pstart\versepar \skt{}\edtext{\skt{बाष्पग्रहात्सूक्ष्मनिरीक्षणाच्च}}{\Afootnote{K; \skt{वाष्पग्रहात्सूक्ष्मनिरीक्षणाच्च}\ H}}\skt{ नेत्रे }\edtext{\skt{विकाराञ्जनयन्ति}}{\Afootnote{K; \skt{विकारान्जनयन्ति}\ A H}}\skt{ दोषाः ।}\edtext{}{\lemma{\skt{… प्र}}\Afootnote{\emph{om.}\ H\textsuperscript{ac}}}\edtext{}{\lemma{\skt{… त}}\Afootnote{\emph{om.}\ H\textsuperscript{pc}}}\edtext{}{\lemma{\skt{छर्दोभिघाताद्वमनातियोगात् । तथा}}\Afootnote{\skt{गविनिग्रहाच्च । स्वेदादथो धू}\ A; \skt{छर्दोभिघाताद्वमनातियोगा}\ K\textsuperscript{ac}; \skt{छर्दोभिघाताद्वमनातियोगात् । तथा}\ K\textsuperscript{pc}}}\edtext{}{\lemma{\skt{त्… प्र}}\Afootnote{\emph{om.}\ K\textsuperscript{ac}}}\edtext{}{\lemma{\skt{प्रसक्तसंरोदनशोककोपाच्छिराभिघातादतिमद्ययोगात्… नाञ्}}\Afootnote{\emph{om.}\ A}}\edtext{}{\lemma{\skt{प्रसक्तसंरोदनशोककोपाच्छिराभिघातादतिमद्ययोगात् । नाञ्च}}\Afootnote{\skt{त् । प्रसक्तसंरोदनशोककोपाच्छिराभिघातादतिमद्ययोगात् । नाञ्च}\ K\textsuperscript{ac}; \skt{प्रसक्तसंरोदनशोककोपाच्छिराभिघातादतिमद्ययोगात् । नाञ्च}\ K\textsuperscript{pc}; \skt{प्रसक्तसंरोदनशोककोपाच्छिरोऽभिघातादतिमद्यपानात्तथाह क्रूराश्च}\ H}}\edtext{}{\lemma{\skt{च… बाष्प}}\Afootnote{\emph{om.}\ A}}\pend

\pstart\versepar \ledsidenote{\passref{6.1.28}}\edtext{\skt{तैस्त्रिभिर्दशदश च}}{\Afootnote{\skt{वाताद्दश तथा पित्तात्कफाच्चैव त्रयोदश ।}\ A; \skt{तैस्त्रिभिर्दशद च}\ K\textsuperscript{ac}; \skt{तैस्त्रिभिर्दशदश च}\ K\textsuperscript{pc}; \skt{तैस्त्रिभिर्द्दश दश च}\ H}}\skt{ }\edtext{\skt{कफेनाप्यधिकास्त्रयः}}{\Afootnote{\skt{कफेनाप्यधिकास्त्रयः ।}\ K; \skt{कफेनात्यधिकास्त्रयः ।}\ H}}\skt{ । }\pend
\pstart\versepar \skt{}\edtext{\skt{रक्तजा}}{\Afootnote{K; \skt{रक्ता}\ A; \skt{रक्तजाः}\ H}}\skt{ }\edtext{\skt{षोडश}}{\Afootnote{K H; \skt{त्षोडश}\ A}}\skt{ }\edtext{\skt{ज्ञेयाः}}{\Afootnote{K H; \skt{विज्ञेयाः}\ A}}\skt{ }\edtext{\skt{सर्वजाःपञ्चविंशति}}{\Afootnote{K; \skt{सर्वजाः पञ्चविंशतिः}\ A H}}\skt{ ।}\edtext{}{\lemma{\skt{च… र}}\Afootnote{\emph{om.}\ A}}\pend

\pstart\versepar \ledsidenote{\passref{6.1.29}}\edtext{\skt{बाह्यौ}}{\Afootnote{K; \skt{तथा बाह्यौ}\ A; \skt{वाह्यै}\ H}}\skt{ }\edtext{\skt{पुनर्द्वौ च}}{\Afootnote{\skt{पुनर्द्वौ च}\ A K; \skt{र्पुनर्द्वौ च}\ H}}\skt{ }\edtext{\skt{तथा}}{\Afootnote{K H; \emph{om.}\ A}}\skt{ रोगाः }\edtext{\skt{षट्सप्तति}}{\Afootnote{K; \skt{षट्सप्ततिः}\ A; \skt{षड्सप्ततिः}\ H}}\skt{ स्मृताः । }\pend
\pstart\versepar \skt{हताधिमन्थो }\edtext{\skt{निमिशो}}{\Afootnote{K; \skt{निमिषो}\ A H}}\skt{ }\edtext{\skt{दृष्टिर्गम्भीरिका च}}{\Afootnote{\skt{दृष्टिर्गम्भीरिका च}\ A K; \skt{दृष्टिगंभीरिका}\ H}}\skt{ या ।}\pend

\pstart\versepar \ledsidenote{\passref{6.1.30}}\skt{यच्च }\edtext{\skt{वाताहतं}}{\Afootnote{K H; \skt{वातहतं}\ A}}\skt{ }\edtext{\skt{वर्त्म ते न}}{\Afootnote{K H; \skt{वर्त्म न ते}\ A}}\skt{ सिध्यन्ति }\edtext{\skt{वातजाः}}{\Afootnote{K H; \skt{वात}\ A}}\skt{ । }\pend
\pstart\versepar \skt{}\edtext{\skt{काचोरुणौ}}{\Afootnote{K; \emph{om.}\ A; \skt{काचोरुणो}\ H}}\skt{ }\edtext{\skt{मारुतजो}}{\Afootnote{K H; \skt{जाः ।}\ A}}\skt{ याप्यो }\edtext{\skt{भवति}}{\Afootnote{K; \skt{ऽथ तन्}\ A; \skt{भवन्ति}\ H}}\skt{ }\edtext{\skt{जानता}}{\Afootnote{K H; \skt{मयः काचः साध्याः स्युः सान्यमारुताः ॥}\ A}}\skt{ ।}\pend

\pstart\versepar \ledsidenote{\passref{6.1.31}}\edtext{\skt{शुष्काक्षिपाकाभिष्यन्दो यश्च}}{\Afootnote{K; \skt{शुष्काक्षिपाकाधीमन्थस्यन्द}\ A; \skt{शुष्काक्षिपाकोभिष्यन्दो यश्च}\ H}}\skt{ }\edtext{\skt{मारुतपर्ययः}}{\Afootnote{K H; \skt{मारुतपर्य}\ A}}\skt{ । }\pend
\pstart\versepar \skt{}\edtext{\skt{अधिमन्थान्यवातश्च}}{\Afootnote{\skt{याः । अ}\ A; \skt{अधिमन्थान्यवातश्च}\ K\textsuperscript{pc}; \skt{अधिमन्थोन्यवातश्च}\ H\textsuperscript{pc}}}\skt{ }\edtext{\skt{साध्या}}{\Afootnote{K\textsuperscript{pc} H\textsuperscript{pc}; \skt{साध्यो}\ A}}\skt{ }\edtext{\skt{मारुतजेषु ते}}{\Afootnote{K\textsuperscript{pc}; \emph{om.}\ A; \skt{मारुतजेषु वै ॥}\ H\textsuperscript{pc}}}\skt{ । }\pend
\pstart\versepar \skt{}\edtext{\skt{ह्रस्वजाड्यजलस्रावावसाध्यौ}}{\Afootnote{\skt{ह्रस्वजाड्यो यो जलस्रा}\ A; \skt{ह्रस्वजाड्यजलस्रावावसाध्यौ}\ K\textsuperscript{ac} K\textsuperscript{pc}; \skt{ह्रस्वजाड्यजलश्रावावसाध्यौ}\ H\textsuperscript{ac} H\textsuperscript{pc}}}\skt{ }\edtext{\skt{पित्तसम्भवौ}}{\Afootnote{K H; \skt{वश्च पैत्तिकः ॥}\ A}}\skt{ ।}\edtext{}{\lemma{\skt{अधिमन्थान्यवातश्… ह्र}}\Afootnote{\emph{om.}\ K\textsuperscript{ac}}}\edtext{}{\lemma{\skt{अधिमन्थान्यवातश्… ह्र}}\Afootnote{\emph{om.}\ H\textsuperscript{ac}}}\pend

\pstart\versepar \ledsidenote{\passref{6.1.32}}\skt{परिम्लायी च }\edtext{\skt{नीली च}}{\Afootnote{K H; \skt{नीलश्च}\ A}}\skt{ }\edtext{\skt{काचो}}{\Afootnote{K H; \skt{याप्यः}\ A}}\skt{ }\edtext{\edtext{\skt{याप्यश्च}}{\lemma{\skt{याप्यश्}}\Afootnote{K H; \skt{काचो}\ A}}}{\lemma{\skt{च}}\Afootnote{K H; \skt{ऽथ}\ A}}\skt{ तन्मयः । }\pend
\pstart\versepar \skt{}\edtext{\skt{शुक्तिस्यन्दोधिमन्थश्च}}{\Afootnote{K; \skt{अभिष्यन्दोऽधिमन्थो}\ A; \skt{शुक्तिष्यन्दोधिमन्थश्च}\ H}}\skt{ }\edtext{\skt{तथाम्लाध्युषितञ्च}}{\Afootnote{K; \skt{ऽम्लाध्युषितं शुक्तिका च}\ A; \skt{तथाम्लध्यूषितञ्च}\ H}}\skt{ }\edtext{\skt{यत्}}{\Afootnote{K H; \skt{या ॥}\ A}}\skt{ ।}\pend

\pstart\versepar \ledsidenote{\passref{6.1.33}}\skt{दृष्टिः पित्तविदग्धा च धूमदर्शी च सिध्यति । }\pend
\pstart\versepar \skt{}\edtext{\skt{सुस्रावः}}{\Afootnote{K; \skt{असाध्यः}\ A; \skt{आस्रावः}\ H}}\skt{ }\edtext{\skt{कफजे}}{\Afootnote{K; \skt{कफजः}\ A; \skt{कफजो}\ H}}\skt{ }\edtext{\skt{साध्यो}}{\Afootnote{K H; \skt{स्रावो}\ A}}\skt{ याप्यः }\edtext{\skt{काचस्तु}}{\Afootnote{\skt{काचश्च}\ A; \skt{काचस्तु}\ K H}}\skt{ तन्मयः ।}\pend

\pstart\versepar \ledsidenote{\passref{6.1.34}}\edtext{\skt{अभिष्यन्दोधिवातश्च}}{\Afootnote{K; \skt{अभिष्यन्दोऽधिमन्थश्च}\ A; \skt{अभिष्यन्दोधिमन्थश्च}\ H}}\skt{ }\edtext{\skt{बलासग्रथितन्तथा}}{\Afootnote{K; \skt{बलासग्रथितं च य}\ A; \skt{वलासग्रथितन्तथा ॥}\ H}}\skt{ ॥ }\pend
\pstart\versepar \skt{}\edtext{\skt{दृष्टिः}}{\Afootnote{K H; \skt{त् । दृष्टिः}\ A}}\skt{ श्लेष्मविदग्धा च पोथक्यो }\edtext{\skt{लगणश्च}}{\Afootnote{A K; \skt{लगणञ्च}\ H}}\skt{ }\edtext{\skt{यत्}}{\Afootnote{K H; \skt{यः ॥}\ A}}\skt{ ।}\pend

\pstart\versepar \ledsidenote{\passref{6.1.35}}\edtext{\skt{क्रिमिग्रन्थिपरिक्लिन्नं}}{\Afootnote{K H; \skt{क्रिमिग्रन्थिपरिक्लिन्न}\ A}}\skt{ }\edtext{\skt{वर्त्मशुक्लार्मपिष्टकः}}{\Afootnote{K; \skt{वर्त्मशुक्लार्मपिष्टकाः ।}\ A; \skt{वर्त्म्मशुक्लार्म्मपिष्टकः ॥}\ H}}\skt{ । }\pend
\pstart\versepar \skt{}\edtext{\skt{श्लेष्मोपनाहाः}}{\Afootnote{K H; \skt{श्लेष्मोपनाहः}\ A}}\skt{ }\edtext{\skt{साध्यास्ते}}{\Afootnote{H; \skt{साध्यास्तु}\ A; \skt{साध्यास्ते}\ K}}\skt{ }\edtext{\skt{कथिता}}{\Afootnote{K; \skt{कथिताः}\ A H}}\skt{ श्लेष्मजेषु तु ।}\pend

\pstart\versepar \ledsidenote{\passref{6.1.36}}\edtext{\skt{रक्तस्रावोजकाजातं}}{\Afootnote{K; \skt{रक्तस्रावोऽजकाजातं}\ A; \skt{रक्तास्रावोजकाजातं}\ H}}\skt{ }\edtext{\skt{शोणितार्शोव्रणान्वितं}}{\Afootnote{A K; \skt{शोणेतासोवणान्वितं ॥}\ H}}\skt{ । }\pend
\pstart\versepar \skt{शुक्रं न साध्यं }\edtext{\skt{काचन्तु}}{\Afootnote{K; \skt{काचश्च}\ A; \skt{काचस्तु}\ H}}\skt{ }\edtext{\skt{याप्यं}}{\Afootnote{K H; \skt{याप्य}\ A}}\skt{ }\edtext{\skt{कृष्णश्च}}{\Afootnote{K; \skt{स्तज्जः प्र}\ A; \skt{कृष्णञ्च}\ H}}\skt{ }\edtext{\skt{तन्मयः}}{\Afootnote{K H; \skt{कीर्तितः ॥}\ A}}\skt{ ।}\pend

\pstart\versepar \ledsidenote{\passref{6.1.37}}\skt{मन्थस्यन्दौ }\edtext{\skt{क्लिष्टवर्त्म}}{\Afootnote{A K; \skt{क्लिष्टवर्त्म्म}\ H}}\skt{ }\edtext{\skt{हर्षोत्पत्तौ}}{\Afootnote{K H; \skt{हर्षोत्पातौ}\ A}}\skt{ तथैव च । }\pend
\pstart\versepar \skt{}\edtext{\skt{सिराजाताञ्जनाख्या च}}{\Afootnote{A K; \skt{सिरापाकाञ्जनाख्या च}\ H}}\skt{ }\edtext{\skt{सिराजालश्च}}{\Afootnote{K; \skt{सिराजालं च}\ A H}}\skt{ यत्स्मृतं ।}\pend

\pstart\versepar \ledsidenote{\passref{6.1.38}}\edtext{\edtext{\skt{पर्व्वण्यथाव्रणं}}{\lemma{\skt{पर्व्वण्य्}}\Afootnote{K; \skt{पर्व}\ A H}}}{\lemma{\skt{अथाव्रणं}}\Afootnote{K; \skt{ण्यथाव्रणं}\ A H}}\skt{ शुक्रं }\edtext{\skt{लोहितार्म}}{\Afootnote{K H; \skt{शोणितार्मा}\ A}}\skt{ }\edtext{\skt{तथार्जुनः}}{\Afootnote{K H; \skt{र्जुनश्च यः ।}\ A}}\skt{ । }\pend
\pstart\versepar \skt{एते साध्या विकारेषु रक्तजेषु भवन्ति हि ॥}\pend

\pstart\versepar \ledsidenote{\passref{6.1.39}}\skt{पूयास्रावो }\edtext{\skt{नाकुलान्ध्यमक्षीपाकात्ययोजली}}{\Afootnote{K; \skt{नाकुलान्ध्यमक्षिपाकात्य}\ A; \skt{नाकुलान्ध्यं मक्षिपाकात्ययोलजी ॥}\ H}}\skt{ । }\pend
\pstart\versepar \skt{}\edtext{\skt{रोगाअसाध्याश्चत्वारः}}{\Afootnote{K; \skt{रोगा असाध्याश्चत्वारः}\ H}}\skt{ }\edtext{\skt{सर्वजेषु}}{\Afootnote{K H; \skt{योऽ}\ A}}\skt{ }\edtext{\skt{भवन्ति}}{\Afootnote{K H; \skt{ल}\ A}}\skt{ }\edtext{\skt{हि}}{\Afootnote{K H\textsuperscript{ac}; \skt{जी । असाध्याः सर्वजा याप्यः काचः कोपश्च पक्ष्मणः ॥}\ A; \skt{हि । याप्यो वर्त्म्माववन्धस्तु पक्ष्मकोपः सविप्लुतः ॥}\ H\textsuperscript{pc}}}\skt{ ।}\edtext{}{\lemma{\skt{योजली… स}}\Afootnote{\emph{om.}\ A}}\pend

\pstart\versepar \ledsidenote{\passref{6.1.40}}\edtext{\edtext{\skt{याप्या}}{\Afootnote{K; \emph{om.}\ A}}\skt{ }\edtext{\skt{वर्त्मावबन्धस्तु}}{\Afootnote{K; \skt{वर्त्मावबन्}\ A}}\skt{ }\edtext{\skt{पक्ष्मकोपः}}{\Afootnote{K; \skt{धो}\ A}}\skt{ भविप्लुतः}}{\lemma{\skt{… भवि}}\Afootnote{\emph{om.}\ H}}\skt{ । }\pend
\pstart\versepar \skt{}\edtext{\skt{पिण्डाक्षस्तु यो}}{\Afootnote{K; \skt{क्लिन्न वर्त्मासु}\ H}}\skt{ }\edtext{\skt{व्याधिः}}{\Afootnote{A K; \emph{om.}\ H}}\skt{ सिरासु }\edtext{\skt{पिटका च या}}{\Afootnote{K H; \skt{पिडका च या ।}\ A}}\skt{ । }\pend
\pstart\versepar \skt{}\edtext{\skt{प्रस्तार्यर्माधीमांसार्म}}{\Afootnote{K; \skt{प्रस्तार्यर्माधिमांसार्म}\ A H}}\skt{ स्नाय्वर्मोत्सङ्गिनी च या ।}\edtext{}{\lemma{\skt{भविप्लुतः… यो}}\Afootnote{\emph{om.}\ A}}\pend

\pstart\versepar \ledsidenote{\passref{6.1.41}}\edtext{\edtext{\skt{पूयालसश्चार्बुदश्च}}{\lemma{\skt{पूयालसश्}}\Afootnote{A K; \skt{पूयालसञ्}\ H}}}{\lemma{\skt{चार्बुदश्च}}\Afootnote{K; \skt{चार्बुदं च}\ A; \skt{चार्व्वुदञ्च}\ H}}\skt{ }\edtext{\skt{श्यावकर्दमवर्त्मनी}}{\Afootnote{A K; \skt{श्यावकन्दमवर्त्म्मनी ।}\ H}}\skt{ । }\pend
\pstart\versepar \skt{}\edtext{\skt{तथार्शोवर्त्मशुष्कार्शः}}{\Afootnote{A K; \skt{तथार्शोवर्त्म्मशुष्कार्शः}\ H}}\skt{ }\edtext{\skt{शर्करावर्त्म}}{\Afootnote{A K; \skt{सर्करावर्त्म्}\ H}}\skt{ }\edtext{\skt{एव}}{\Afootnote{K; \skt{यच्च वै ॥}\ A; \skt{म एव}\ H}}\skt{ }\edtext{\skt{च}}{\Afootnote{K H; \emph{om.}\ A}}\skt{ ।}\pend

\pstart\versepar \ledsidenote{\passref{6.1.42}}\edtext{\skt{अशोफशोफै}}{\Afootnote{K; \skt{सशोफश्चाप्यशोफश्}\ A; \skt{अशोफशोफौ}\ H}}\skt{ }\edtext{\skt{महतौ}}{\Afootnote{K H; \skt{च}\ A}}\skt{ }\edtext{\skt{पाकौ}}{\Afootnote{K H; \skt{पाको}\ A}}\skt{ }\edtext{\skt{बहलवर्त्म च}}{\Afootnote{A K; \skt{वहलव}\ H}}\skt{ । }\pend
\pstart\versepar \skt{}\edtext{\skt{कुम्भिकिनी}}{\Afootnote{K; \skt{अक्लिन्नवर्त्म कुम्भी}\ A; \skt{र्त्म्म च । कुम्भीकिनी}\ H}}\skt{ }\edtext{\skt{वर्त्मगता}}{\Afootnote{K; \skt{का}\ A; \skt{वर्त्म्मगता}\ H}}\skt{ }\edtext{\skt{बिसवर्त्म च}}{\Afootnote{A K; \skt{विसवर्त्म्म च}\ H}}\skt{ सिध्यति ।}\pend

\pstart\versepar \ledsidenote{\passref{6.1.43}}\edtext{\skt{सनिमित्तो निमित्तश्च}}{\Afootnote{K; \skt{सनिमित्तोऽनिमित्तश्च}\ A; \skt{सनिमित्तोनिमित्तश्च}\ H}}\skt{ }\edtext{\skt{द्वावसाध्यौ च}}{\Afootnote{K; \skt{द्वावसाध्यौ तु}\ A; \skt{द्वावसाध्यौ}\ H\textsuperscript{ac}; \skt{द्वावसाध्यौतु}\ H\textsuperscript{pc}}}\skt{ }\edtext{\skt{बाह्यतः}}{\Afootnote{K; \skt{बाह्य}\ A; \skt{वाह्यतः ।}\ H}}\skt{ । }\pend
\pstart\versepar \skt{}\edtext{\skt{रोगाः}}{\Afootnote{K H; \skt{जौ ।}\ A}}\skt{ }\edtext{\skt{षट्सप्ततिस्त्वेते}}{\Afootnote{K; \skt{षट्सप्ततिर्विकाराणामेषा}\ A; \skt{षट्सप्ततिस्वेते}\ H}}\skt{ }\edtext{\skt{सङ्ग्रहेण}}{\Afootnote{K H; \skt{संग्रह}\ A}}\skt{ }\edtext{\skt{प्रकीर्त्तिताः}}{\Afootnote{K H; \skt{कीर्तिता ॥}\ A}}\skt{ ।}\pend

\pstart \ledsidenote{\passref{6.1.45}}\edtext{\edtext{\skt{भूयस्तेषां}}{\lemma{\skt{भूयस्}}\Afootnote{\skt{सर्वाश्रयाः सप्तदश दृष्टिजा द्वादशैव तु । बाह्यजौ द्वौ समाख्यातौ रोगौ परमदारुणौ । भूय}\ A; \skt{भूय}\ K H}}}{\lemma{\skt{तेषां}}\Afootnote{\skt{एता}\ A; \skt{स्तेषां}\ K; \skt{स्त्वेषाम्}\ H}}\skt{ }\edtext{\skt{प्रवक्ष्यामि}}{\Afootnote{K H; \skt{न्प्रवक्ष्यामि}\ A}}\skt{ }\edtext{\edtext{\skt{संख्यारूपचिकित्सितैरिति}}{\lemma{\skt{संख्यारूपचिकित्सितैर्}}\Afootnote{A H; \skt{संख्यारूपचिकित्सितै}\ K}}}{\lemma{\skt{इति}}\Afootnote{H; \emph{om.}\ A; \skt{रिति ॥}\ K}}\skt{ ॥ }\edtext{\skt{शालाकृ}}{\Afootnote{K; \emph{om.}\ A H}}\skt{ १ ॥}\pend

\endnumbering
\beginnumbering
\eledsection*{Adhyāya 2}
\pstart \ledsidenote{\passref{6.2.1}}\skt{अथातः }\edtext{\edtext{\skt{सन्धिगतरोगविज्ञानीयम्व्याख्यास्यामः}}{\lemma{\skt{सन्धिगतरोगविज्ञानीयम्}}\Afootnote{H; \skt{सन्धिगतरोगविज्ञानीयमध्यायं}\ A; \skt{सन्धिगतरोगविज्ञानीयम्व्या व्याख्यास्यामः ॥}\ K}}}{\lemma{\skt{व्याख्यास्यामः}}\Afootnote{A H; \emph{om.}\ K}}\skt{ ॥}\pend

\pstart\versepar \ledsidenote{\passref{6.2.2-i}}\edtext{\skt{नवसन्ध्यश्रयश्च}}{\Afootnote{\emph{om.}\ A; \skt{नवसन्ध्यश्रजश्च}\ K; \skt{नवसन्धाग्र स्}\ H\textsuperscript{ac}; \skt{नव सन्ध्याश्रयास्}\ H\textsuperscript{pc}}}\skt{ }\edtext{\skt{वर्त्मजाताचेक विंशति}}{\Afootnote{K; \emph{om.}\ A; \skt{वर्त्मजातेकविंशति ।}\ H}}\skt{ । }\pend
\pstart\versepar \skt{}\edtext{\skt{शुक्लभागैदसेकश्च}}{\Afootnote{K; \emph{om.}\ A; \skt{शुक्लभादलैकश्च}\ H}}\skt{ }\edtext{\skt{चत्वारः}}{\Afootnote{K H; \emph{om.}\ A}}\skt{ }\edtext{\skt{कृष्णभागजाः}}{\Afootnote{K H; \emph{om.}\ A}}\skt{ ।}\pend

\pstart\versepar \ledsidenote{\passref{6.2.2-ii}}\edtext{\skt{सर्वाश्रयाः}}{\Afootnote{K H; \emph{om.}\ A}}\skt{ सपुदमदृष्टिद्वादशैव त्तु । }\pend
\pstart\versepar \skt{द्वौ तु }\edtext{\skt{सर्वाश्रया}}{\Afootnote{K H; \emph{om.}\ A}}\skt{ }\edtext{\skt{वन्यावनिमित्तनिमित्तजौ}}{\Afootnote{K H; \emph{om.}\ A}}\skt{ ॥}\edtext{}{\lemma{\skt{सपुदमदृष्टिद्वादशैव त्तु । द्वौ तु}}\Afootnote{\emph{om.}\ A; \skt{सपुदम दृष्टिद्वादशैवत्तु । द्वौ तु}\ K; \skt{सपुदल दृष्टिद्वादशैवत्तु । द्वौ तु}\ H}}\pend

\pstart\versepar \ledsidenote{\passref{6.2.3}}\skt{पूयालसः }\edtext{\skt{सोपोनाहः}}{\Afootnote{K; \skt{सोपनाहः}\ A H}}\skt{ स्रावाः पर्वणिकालजी । }\pend
\pstart\versepar \skt{क्रिमिग्रन्थिश्च विज्ञेया रोगाः सन्धिगता नव ॥}\pend

\pstart\versepar \ledsidenote{\passref{6.2.4}}\skt{पक्वः शोफः }\edtext{\skt{सन्धिजेयः}}{\Afootnote{K; \skt{सन्धिजः सं}\ A; \skt{सन्धिजो यः}\ H}}\skt{ }\edtext{\skt{सतोदः}}{\Afootnote{K H; \skt{स्रवेद्यः}\ A}}\skt{ }\pend
\pstart\versepar \skt{}\edtext{\skt{स्रवेत्पूयं स हि}}{\Afootnote{K; \skt{सान्द्रं पूयं पूति}\ A; \skt{स्रवेत्पूयं}\ H}}\skt{ }\edtext{\skt{पूयालसाख्यः}}{\Afootnote{K H; \skt{पूयालसः सः ।}\ A}}\skt{ । }\pend
\pstart\versepar \skt{}\edtext{\edtext{\skt{ग्रन्थिर्महान्दृष्टिसन्धावपाकः}}{\lemma{\skt{ग्रन्थिर्}}\Afootnote{K; \skt{ग्रन्थि}\ A; \skt{ग्रन्थिम्}\ H}}}{\lemma{\skt{महान्दृष्टिसन्धावपाकः}}\Afootnote{K H; \skt{र्नाल्पो दृष्टिसन्धावपाकः}\ A}}\skt{ }\pend
\pstart\versepar \skt{}\edtext{\skt{कण्डुस्रावोनीरूजाः}}{\Afootnote{K H; \skt{कण्डूप्रायो नीरुज}\ A}}\skt{ }\edtext{\skt{सोपनाः}}{\Afootnote{K H; \skt{स्तूपनाहः ॥}\ A}}\skt{ ।}\pend

\pstart\versepar \ledsidenote{\passref{6.2.5}}\skt{गत्वा सन्धीनश्रुमार्गेण दोषाः }\pend
\pstart\versepar \skt{}\edtext{\skt{श्रवन्ती}}{\Afootnote{K; \skt{कुर्युः स्रा}\ A; \skt{श्रवन्तीम्}\ H}}\skt{ }\edtext{\skt{वास्वामिम्विकारां}}{\Afootnote{K; \skt{वान्रुग्विहीना}\ A; \skt{वास्वामिम्विकारा}\ H}}\skt{ कनीनात् । }\pend
\pstart\versepar \skt{तद्विश्रावं }\edtext{\skt{नेत्रनाडीति}}{\Afootnote{K H; \skt{नेत्रनाडीम}\ A}}\skt{ }\edtext{\skt{चैके}}{\Afootnote{K H; \skt{थैके}\ A}}\skt{ }\pend
\pstart\versepar \skt{}\edtext{\skt{तस्यालिङ्गं}}{\Afootnote{A K; \skt{तस्यालिङ्गा}\ H}}\skt{ }\edtext{\skt{कीर्तियेष्ये}}{\Afootnote{K; \skt{कीर्तयिष्ये}\ A; \skt{न्कीर्त्तयिष्ये}\ H}}\skt{ }\edtext{\skt{चतुर्धा}}{\Afootnote{A K; \skt{चतुधा ॥}\ H}}\skt{ ॥}\edtext{}{\lemma{\skt{कनीनात् । तद्विश्रावं}}\Afootnote{K; \skt{न्कनीनात् । तान्वै स्रावान्}\ A; \skt{न्कनीनात् । तद्विश्रावं}\ H}}\pend

\pstart\versepar \ledsidenote{\passref{6.2.6}}\edtext{\skt{पाकात्सन्धौ}}{\Afootnote{K H; \skt{पाकः सन्धौ}\ A}}\skt{ }\edtext{\skt{संश्रवेद्यम्स}}{\Afootnote{\skt{संस्रवेद्यश्च}\ A; \skt{संश्रवेद्यम्स}\ K; \skt{संश्रवेद्यः स}\ H}}\skt{ पूयं }\pend
\pstart\versepar \skt{}\edtext{\skt{पूयाश्रावः}}{\Afootnote{K H; \skt{पूयास्रावो}\ A}}\skt{ }\edtext{\skt{सविकारो}}{\Afootnote{K H; \skt{नैकरूपः}\ A}}\skt{ }\edtext{\skt{मतस्तु}}{\Afootnote{K H; \skt{प्रदिष्टः ।}\ A}}\skt{ । }\pend
\pstart\versepar \skt{}\edtext{\skt{श्वेतं}}{\Afootnote{A K; \skt{शेषं}\ H}}\skt{ }\edtext{\skt{स्कत्रं}}{\Afootnote{K H; \skt{सान्द्रं}\ A}}\skt{ }\edtext{\skt{पिच्छिलं यः}}{\Afootnote{K H; \skt{पिच्छिलं सं}\ A}}\skt{ }\edtext{\skt{श्रवेत्तु}}{\Afootnote{K H; \skt{स्रवेद्यः}\ A}}\skt{ }\pend
\pstart\versepar \skt{}\edtext{\skt{श्लेष्माश्रावः}}{\Afootnote{K H; \skt{श्लेष्मास्रावो}\ A}}\skt{ }\edtext{\skt{सविकारः}}{\Afootnote{K H; \skt{नीरुजः स}\ A}}\skt{ प्रदिष्टः ।}\pend

\pstart\versepar \ledsidenote{\passref{6.2.7}}\edtext{\skt{रक्ताश्रावः}}{\Afootnote{K H; \skt{रक्तास्रावः}\ A}}\skt{ }\edtext{\skt{शोणितजोविकारः}}{\Afootnote{K H; \skt{शोणितो}\ A}}\skt{ }\pend
\pstart\versepar \skt{}\edtext{\skt{श्रवदूष्ण}}{\Afootnote{K H; \emph{om.}\ A}}\skt{ }\edtext{\skt{तत्तु}}{\Afootnote{K H; \skt{त्थः स}\ A}}\skt{ }\edtext{\edtext{\skt{रक्तम्प्रभूतं}}{\lemma{\skt{रक्तम्}}\Afootnote{K H; \skt{रक्त}\ A}}}{\lemma{\skt{प्रभूतं}}\Afootnote{K; \skt{मुष्णं नाल्पं संस्रवेन्}\ A; \skt{प्रभू}\ H}}\skt{ । }\pend
\pstart\versepar \skt{}\edtext{\skt{हारिद्राभं}}{\Afootnote{K; \skt{नातिसान्द्रम् । पीताभासं}\ A; \skt{तं । हारिद्राभं}\ H}}\skt{ }\edtext{\edtext{\skt{पीतमुष्णं}}{\lemma{\skt{पीतम्}}\Afootnote{K H; \skt{नील}\ A}}}{\lemma{\skt{उष्णं}}\Afootnote{K H; \skt{मुष्णं}\ A}}\skt{ }\edtext{\skt{जलम्वा}}{\Afootnote{K H; \skt{जलाभं}\ A}}\skt{ }\pend
\pstart\versepar \skt{}\edtext{\skt{पित्ताश्रावः}}{\Afootnote{K H; \skt{पित्तास्रावः}\ A}}\skt{ }\edtext{\skt{संश्रवेत्सन्धिमध्यात्}}{\Afootnote{K H; \skt{संस्रवेत्सन्धिमध्यात् ॥}\ A}}\skt{ ॥}\pend

\pstart\versepar \ledsidenote{\passref{6.2.8}}\edtext{\skt{ताम्रां}}{\Afootnote{K; \skt{ताम्रा}\ A H}}\skt{ }\edtext{\edtext{\skt{तन्वीन्दाहशूलोपपन्नां}}{\lemma{\skt{तन्वीन्}}\Afootnote{K; \skt{तन्वी}\ A; \skt{नन्वीं}\ H}}}{\lemma{\skt{दाहशूलोपपन्नां}}\Afootnote{K H; \skt{दाहशूलोपपन्ना}\ A}}\skt{ }\pend
\pstart\versepar \skt{}\edtext{\skt{ब्रूयाद्वैद्यः}}{\Afootnote{\skt{रक्ताज्ज्ञेया}\ A; \skt{व्रूयाद्वैड्यः}\ K H}}\skt{ }\edtext{\skt{पिटकां}}{\Afootnote{K H; \emph{om.}\ A}}\skt{ }\edtext{\skt{पर्वणीकां}}{\Afootnote{K H; \skt{पर्वणी वृत्तशोफा ।}\ A}}\skt{ । }\pend
\pstart\versepar \skt{जाता सन्धौ }\edtext{\skt{कृष्णशुक्लेलजी}}{\Afootnote{K H; \skt{कृष्णशुक्लेऽलजी}\ A}}\skt{ स्यात् }\pend
\pstart\versepar \skt{तस्मिन्नेव }\edtext{\skt{प्रथिताः}}{\Afootnote{K H; \skt{ख्यापिता}\ A}}\skt{ पूर्वलिङ्गैः ॥}\pend

\pstart\versepar \ledsidenote{\passref{6.2.9}}\edtext{\skt{क्रिमिग्रन्थैवर्त्मतः}}{\Afootnote{K; \skt{क्रिमिग्रन्थिर्वर्त्मनः}\ A; \skt{क्रिमिग्रन्थैवर्त्म्मतः}\ H}}\skt{ }\edtext{\skt{पक्ष्मतश्च}}{\Afootnote{K H; \skt{पक्ष्मणश्च}\ A}}\skt{ }\pend
\pstart\versepar \skt{}\edtext{\skt{कूर्युः}}{\Afootnote{K H; \emph{om.}\ A}}\skt{ कण्डूं }\edtext{\skt{क्रिमयः}}{\Afootnote{K H; \skt{कुर्युः क्रिमयः}\ A}}\skt{ सन्धिजाताः । }\pend
\pstart\versepar \skt{नानारूपा }\edtext{\skt{वर्त्मशुक्लान्तसन्धौ}}{\Afootnote{K H; \skt{वर्त्मशुक्लस्य सन्धौ}\ A}}\skt{ }\pend
\pstart\versepar \skt{}\edtext{\skt{चरन्त्यन्तर्नयनं}}{\lemma{\skt{चरन्त्य्}}\Afootnote{K H; \skt{चरन्तो}\ A}}\skt{ }\edtext{\skt{दूषयन्त्य}}{\Afootnote{K H; \skt{दूषयन्ति ॥}\ A}}\skt{ }\edtext{\skt{इति}}{\Afootnote{K H; \emph{om.}\ A}}\skt{ ॥ ० ॥}\pend

\pstart \ledsidenote{\passref{6.2.trailer}}\edtext{\skt{॥ शालाक्ये}}{\Afootnote{K H; \skt{इति सुश्रुतसंहितायामुत्तरतन्त्रे सन्धिगतरोगविज्ञानीयो नाम}\ A}}\skt{ }\edtext{\edtext{\skt{द्वितीयोऽध्यायः}}{\lemma{\skt{द्वितीयो}}\Afootnote{A H; \emph{om.}\ K}}}{\lemma{\skt{ऽध्यायः}}\Afootnote{A; \emph{om.}\ K; \skt{ऽध्यायः ॥ {\symfont ❈} ॥}\ H}}\skt{ ॥}\pend

\endnumbering
\beginnumbering
\eledsection*{Adhyāya 3}
\pstart \ledsidenote{\passref{6.3.1}}\skt{अथातो }\edtext{\edtext{\skt{वर्त्मगतरोगविञ्जानीयम्व्याख्यास्यामः}}{\lemma{\skt{वर्त्मगतरोगविञ्जानीयम्}}\Afootnote{H; \skt{वर्त्मगतरोगविज्ञानीयमध्यायं}\ A; \skt{वर्त्मगतरोगविञ्जानीयम्व्या व्याख्यास्यामः ॥}\ K}}}{\lemma{\skt{व्याख्यास्यामः}}\Afootnote{A H; \emph{om.}\ K}}\skt{ ॥}\pend

\pstart\versepar \ledsidenote{\passref{6.3.5}}\edtext{\skt{उत्सङ्न्योथ}}{\Afootnote{K; \skt{उत्सङ्गिन्यथ}\ A; \skt{उत्सङ्न्योऽथ}\ H}}\skt{ कुम्भीका }\edtext{\skt{पोथक्यो}}{\Afootnote{A K; \skt{पोत्छक्यो}\ H}}\skt{ }\edtext{\skt{वर्त्मशर्कराः}}{\Afootnote{K H; \skt{वर्त्मशर्करा ।}\ A}}\skt{ । }\pend
\pstart\versepar \skt{}\edtext{\edtext{\skt{तथार्शावर्मशुष्कार्शस्तथैवाञ्जनामिका}}{\lemma{\skt{तथार्शावर्मशुष्कार्शस्}}\Afootnote{K H; \skt{तथार्शोवर्त्म शुष्कार्श}\ A}}}{\lemma{\skt{तथैवाञ्जनामिका}}\Afootnote{K H; \skt{स्तथैवाञ्जननामिका ॥}\ A}}\skt{ ।}\pend

\pstart\versepar \ledsidenote{\passref{6.3.6}}\edtext{\skt{वहलं}}{\Afootnote{K H; \skt{बहलं}\ A}}\skt{ }\edtext{\skt{वर्त्मयच्चापि}}{\Afootnote{K; \skt{वर्त्म यच्चापि}\ A; \skt{वर्त्मीयच्चापि}\ H}}\skt{ }\edtext{\edtext{\skt{व्याधिर्वर्त्माववन्धकः}}{\lemma{\skt{व्याधिर्}}\Afootnote{K H; \skt{व्याधि}\ A}}}{\lemma{\skt{वर्त्माववन्धकः}}\Afootnote{K H; \skt{र्वर्त्मावबन्धकः ।}\ A}}\skt{ । }\pend
\pstart\versepar \skt{}\edtext{\skt{क्लिष्टकर्दमवर्त्माख्यः}}{\Afootnote{K H; \skt{क्लिष्टकर्दमवर्त्माख्यौ}\ A}}\skt{ }\edtext{\skt{श्याववर्त्मस्तथैव}}{\lemma{\skt{श्याववर्त्मस्}}\Afootnote{K H; \skt{श्याववर्त्म}\ A}}\skt{ च ।}\pend

\pstart\versepar \ledsidenote{\passref{6.3.7}}\edtext{\edtext{\skt{प्रक्लिन्नमपरिक्लिन्न}}{\lemma{\skt{प्रक्लिन्नम्}}\Afootnote{K H; \skt{प्रक्लिन्न}\ A}}}{\lemma{\skt{अपरिक्लिन्न}}\Afootnote{K H; \skt{मपरिक्लिन्नं}\ A}}\skt{ }\edtext{\skt{वर्त्मवातहरञ्च}}{\Afootnote{K H; \skt{वर्त्म वातहतं तु}\ A}}\skt{ यत् । । }\pend
\pstart\versepar \skt{अर्वूद }\edtext{\skt{निमिषं}}{\Afootnote{K H; \skt{निमिषश्}\ A}}\skt{ चापि शोणितार्शश्च }\edtext{\skt{यस्मृताः}}{\Afootnote{K; \skt{यत्स्मृतम् ॥}\ A; \skt{यः स्मृतः ॥}\ H}}\skt{ ॥}\edtext{}{\lemma{\skt{यत् । । अर्वूद}}\Afootnote{K; \skt{यत् । अर्बुदं}\ A; \skt{यत् । । अर्वूदन्}\ H}}\pend

\pstart\versepar \ledsidenote{\passref{6.3.8}}\skt{लगणो }\edtext{\skt{विसनामा च}}{\Afootnote{K H; \skt{बिशनामा च}\ A}}\skt{ }\edtext{\skt{पक्ष्मकोपस्तथैव च}}{\lemma{\skt{तथैव च}}\Afootnote{A K; \skt{सपुतः ।}\ H}}\skt{ । }\pend
\pstart\versepar \skt{}\edtext{\edtext{\skt{एकविंशतिरित्येते}}{\lemma{\skt{एकविंशतिर्}}\Afootnote{K H; \skt{एकविंशति}\ A}}}{\lemma{\skt{इत्येते}}\Afootnote{K H; \skt{रित्येते}\ A}}\skt{ विकारा }\edtext{\skt{वर्त्मसम्भवाः}}{\Afootnote{K H; \skt{वर्त्मसंश्रयाः ॥}\ A}}\skt{ ।}\pend

\pstart\versepar \ledsidenote{\passref{6.3.9}}\skt{नामभिस्ते }\edtext{\skt{समूद्दिष्टाच्सांजनैः}}{\Afootnote{K; \skt{समुद्दिष्टा लक्षणै}\ A; \skt{समूद्दिष्टाः सांजनैः}\ H}}\skt{ }\edtext{\skt{सम्प्रवच्यते}}{\Afootnote{K H; \skt{स्तान्प्रचक्ष्महे ।}\ A}}\skt{ ॥ }\pend
\pstart\versepar \skt{}\edtext{\skt{पिडकाभ्यन्तरमूखी}}{\Afootnote{K H; \skt{अभ्यन्तरमुखी}\ A}}\skt{ }\edtext{\skt{वाह्याधोवर्मसंश्रया}}{\Afootnote{K H; \skt{बाह्योत्सङ्गेऽधो वर्त्मनश्च या ॥}\ A}}\skt{ ।}\pend

\pstart\versepar \ledsidenote{\passref{6.3.10}}\edtext{\skt{ज्ञेया}}{\Afootnote{K H; \skt{विज्ञे}\ A}}\skt{ }\edtext{\skt{सोत्साङ्गिनी}}{\Afootnote{K; \skt{योत्सङ्गिनी}\ A; \skt{सोत्सङ्गिनी}\ H}}\skt{ नाम }\edtext{\skt{पिडाका}}{\Afootnote{K; \skt{तद्रूपपिडका}\ A; \skt{पिडका तु}\ H}}\skt{ }\edtext{\skt{कृतिः}}{\Afootnote{K; \skt{चिता ।}\ A; \skt{तदा कृतिः ।}\ H}}\skt{ । }\pend
\pstart\versepar \skt{}\edtext{\skt{कुम्भीकवीजप्रतिमा}}{\Afootnote{K H; \skt{कुम्भीकबीजप्रतिमाः}\ A}}\skt{ पिडका }\edtext{\skt{यस्य}}{\Afootnote{K H; \skt{यास्तु}\ A}}\skt{ }\edtext{\skt{वर्त्मनि}}{\Afootnote{K H; \skt{वर्त्मजाः ॥}\ A}}\skt{ ।}\pend

\pstart\versepar \ledsidenote{\passref{6.3.11}}\edtext{\skt{आध्मायन्ते}}{\Afootnote{K H; \skt{आध्मापयन्}\ A}}\skt{ }\edtext{\skt{भिद्यमानाः}}{\Afootnote{K H; \skt{ति भिन्ना याः}\ A}}\skt{ }\edtext{\skt{कुम्भीकिन्यस्तु ताः}}{\Afootnote{K H; \skt{कुम्भीकपिडकास्}\ A}}\skt{ }\edtext{\skt{स्मृताः}}{\Afootnote{K H; \skt{तु ताः ।}\ A}}\skt{ ॥ }\pend
\pstart\versepar \skt{स्राविण्यः }\edtext{\skt{कण्डूरा}}{\Afootnote{K H; \skt{कण्डुरा}\ A}}\skt{ }\edtext{\skt{गुर्व्यो}}{\Afootnote{A K; \skt{गुर्व्यो}\ H}}\skt{ }\edtext{\skt{रक्तसर्षपसम्मिताः}}{\Afootnote{K; \skt{रक्तसर्षपसन्निभाः । पिडकाश्च}\ A; \skt{चक्तसर्षपसम्मिताः ।}\ H}}\skt{ ॥ }\edtext{\skt{रुजावन्त्यश्च}}{\Afootnote{K H; \skt{रुजावत्यः}\ A}}\skt{ }\edtext{\skt{पिडकाः}}{\Afootnote{K; \emph{om.}\ A; \skt{पिडकाः}\ H}}\skt{ }\edtext{\skt{पोथक्य}}{\Afootnote{A K; \skt{योत्छक्य}\ H}}\skt{ इति संज्ञिताः ॥}\pend

\pstart\versepar \ledsidenote{\passref{6.3.12}}\edtext{\skt{पिडका या}}{\Afootnote{K H; \skt{पिडकाभिः सुसूक्ष्माभिर्घनाभिरभिसंवृता । पिडका या}\ A}}\skt{ खरा }\edtext{\skt{स्थूलाः}}{\Afootnote{K H; \skt{स्थूला}\ A}}\skt{ }\edtext{\edtext{\skt{शूक्ष्माभिरभिसम्वृताः}}{\lemma{\skt{शूक्ष्माभिरभिस्}}\Afootnote{K; \skt{सा ज्ञे}\ A; \skt{शूक्ष्माभिरभिस्}\ H}}}{\lemma{\skt{अम्वृताः}}\Afootnote{K H; \skt{या}\ A}}\skt{ । }\pend
\pstart\versepar \skt{}\edtext{\skt{वर्त्मस्थाशर्करानाम}}{\Afootnote{K; \skt{वर्त्मशर्करा ॥}\ A; \skt{वर्त्मास्थाशर्करा नाम}\ H}}\skt{ }\edtext{\skt{सरोगावर्त्मदूषणाः}}{\Afootnote{K H; \emph{om.}\ A}}\skt{ ।}\pend

\pstart \ledsidenote{\passref{6.3.13}}\edtext{\skt{एर्वाचूवीजप्रतिमापिडका}}{\Afootnote{K; \skt{एर्वारुबीजप्रतिमाः पिडका}\ A; \skt{एर्व्वाचूवीजप्रतिमाः पिडका}\ H}}\skt{ }\edtext{\skt{मृदुवेदनाः}}{\Afootnote{K H; \skt{मन्दवेदनाः ।}\ A}}\skt{ । }\edtext{\skt{शुक्ष्माखरावर्त्मस्थाः}}{\Afootnote{K; \skt{सूक्ष्माः खराश्च वर्त्मस्था}\ A; \skt{शुक्ष्माः खराश्च वर्त्मस्था}\ H}}\skt{ }\edtext{\skt{अर्शावर्त्मस्तदुच्यते}}{\Afootnote{K; \skt{स्तदर्शोवर्त्म कीर्त्यते ॥}\ A; \skt{अर्शावर्त्म्मस्तदुच्यते ॥}\ H}}\skt{ ।}\pend

\pstart\versepar \ledsidenote{\passref{6.3.14}}\edtext{\skt{दीर्घाण्कुरः}}{\Afootnote{K H; \skt{दीर्घोऽङ्कुरः}\ A}}\skt{ }\edtext{\skt{खरःस्तथा}}{\lemma{\skt{तथा}}\Afootnote{K H; \skt{स्तब्धो}\ A}}\skt{ }\edtext{\skt{दारूणैस्यक्तरोद्भवः}}{\Afootnote{K; \skt{दारुणो वर्त्मसंभवः ।}\ A; \skt{दारूणास्यक्तरोद्भवः ।}\ H}}\skt{ । }\pend
\pstart\versepar \skt{}\edtext{\skt{व्याधिरेषौति}}{\Afootnote{K; \skt{व्याधिरेष स}\ A; \skt{व्याधिरेषाति}\ H}}\skt{ }\edtext{\skt{विख्यातः}}{\Afootnote{K; \skt{माख्यातः}\ A; \skt{विख्याताः}\ H}}\skt{ }\edtext{\skt{शुष्कार्शा}}{\Afootnote{K H; \skt{शुष्कार्श}\ A}}\skt{ }\edtext{\skt{नाशनामतः}}{\Afootnote{K H; \skt{इति संज्ञितः ॥}\ A}}\skt{ ॥}\pend

\pstart\versepar \ledsidenote{\passref{6.3.15}}\skt{दाहतोदवती }\edtext{\skt{स्रवि}}{\Afootnote{K; \skt{ताम्रा}\ A; \skt{स्रावि}\ H}}\skt{ }\edtext{\skt{पिडकावर्त्मसम्भवाः}}{\Afootnote{A K; \skt{पिडकावर्त्म्मसम्भवाः ।}\ H}}\skt{ । }\pend
\pstart\versepar \skt{मृद्वी }\edtext{\skt{मन्दरूजा}}{\Afootnote{K H; \skt{मन्दरुजा}\ A}}\skt{ }\edtext{\skt{शुक्ष्मा}}{\Afootnote{K H; \skt{सूक्ष्मा}\ A}}\skt{ ज्ञेया साञ्जननामिका ।}\pend

\pstart\versepar \ledsidenote{\passref{6.3.16}}\skt{वर्त्मोपचीयते यस्य पिडकाभिः समन्ततः । }\pend
\pstart\versepar \skt{सवर्ण्णाभिः }\edtext{\skt{स्थिराभिश्च}}{\Afootnote{K H; \skt{समाभिश्च}\ A}}\skt{ }\edtext{\skt{विद्याद्वहलवर्त्मवत्}}{\Afootnote{K; \skt{विद्याद्बहलवर्त्म तत् ॥}\ A; \skt{विद्याद्वहलवर्त्म्मवत् ॥}\ H}}\skt{ ॥}\pend

\pstart\versepar \ledsidenote{\passref{6.3.17}}\edtext{\skt{कण्डूरोगाल्पतोदेन}}{\Afootnote{K H; \skt{कण्डूमताल्पतोदेन}\ A}}\skt{ }\edtext{\skt{वर्त्माशोफेन यो}}{\Afootnote{K H; \skt{वर्त्मशोफेन यो}\ A}}\skt{ नरः । }\pend
\pstart\versepar \skt{}\edtext{\skt{नभः}}{\Afootnote{K H; \skt{न स}\ A}}\skt{ }\edtext{\skt{सञ्छादयेदक्षि}}{\Afootnote{K; \skt{मं छादयेदक्षि}\ A; \skt{सञ्जा}\ H}}\skt{ }\edtext{\skt{भवेद्वन्धः स}}{\Afootnote{K H; \skt{भवेद्बन्धः स}\ A}}\skt{ }\edtext{\skt{वर्त्मनाः}}{\Afootnote{K H; \skt{वर्त्मनः ॥}\ A}}\skt{ ॥}\pend

\pstart\versepar \ledsidenote{\passref{6.3.18}}\edtext{\skt{मृद्वल्पवेदनाताम्रं}}{\Afootnote{K; \skt{मृद्वल्पवेदनं ताम्रं}\ A H}}\skt{ }\edtext{\skt{यद्वर्ण्णसममेव}}{\lemma{\skt{यद्वर्ण्णसमम्}}\Afootnote{K H; \skt{यद्वर्त्म समम्}\ A}}\skt{ च । }\pend
\pstart\versepar \skt{अकस्माच्च }\edtext{\skt{भवेद्रक्तं}}{\Afootnote{A K; \skt{भवेदेकं}\ H}}\skt{ }\edtext{\skt{क्लिष्टवर्त्मति }\edtext{\skt{तम्विदुः}}{\lemma{\skt{विदुः}}\Afootnote{K H; \skt{दिशेत् ॥}\ A}}}{\lemma{\skt{क्लिष्टवर्त्मति तम्}}\Afootnote{K H; \skt{क्लिष्टवर्त्म तदा}\ A}}\skt{ ।}\pend

\pstart\versepar \ledsidenote{\passref{6.3.19}}\skt{क्लिष्टम्पुनः }\edtext{\skt{पित्तयुतं}}{\Afootnote{K H; \skt{पित्तयु}\ A}}\skt{ }\edtext{\skt{शोणितं}}{\Afootnote{K; \skt{तं}\ A H}}\skt{ }\edtext{\skt{विदहेद्यदा}}{\Afootnote{K; \skt{विदहेच्छोणितं यदा ।}\ A; \skt{विदहे यदा ।}\ H}}\skt{ । }\pend
\pstart\versepar \skt{}\edtext{\skt{ततः}}{\Afootnote{K H; \skt{तदा}\ A}}\skt{ }\edtext{\skt{क्लिष्टत्वमापन्नमुच्यते}}{\lemma{\skt{क्लिष्टत्वम्}}\Afootnote{K H; \skt{क्लिन्नत्वम्}\ A}}\skt{ }\edtext{\skt{वर्त्मकर्दमम्}}{\Afootnote{K H; \skt{वर्त्मकर्दमः ॥}\ A}}\skt{ ॥}\pend

\pstart\versepar \ledsidenote{\passref{6.3.20}}\edtext{\skt{वर्त्मवद्वाह्यतो}}{\Afootnote{K H; \skt{यद्वर्त्म बाह्यतो}\ A}}\skt{ }\edtext{\skt{तश्च}}{\Afootnote{K; \skt{ऽन्यश्च}\ A; \skt{यश्च}\ H}}\skt{ श्यावं }\edtext{\edtext{\skt{शूनाञ्च}}{\lemma{\skt{शूनाञ्}}\Afootnote{K H; \skt{शूनं}\ A}}}{\lemma{\skt{च}}\Afootnote{K H; \skt{सवेदनम् ।}\ A}}\skt{ }\edtext{\skt{जायते}}{\Afootnote{K H; \skt{दाहकण्डू}\ A}}\skt{ । }\pend
\pstart\versepar \skt{}\edtext{\skt{तदाहुः}}{\Afootnote{K H; \skt{परिक्लेदि}\ A}}\skt{ }\edtext{\skt{श्याववर्मेति}}{\Afootnote{K H; \skt{श्याववर्त्मेति}\ A}}\skt{ }\edtext{\skt{वर्त्मरोगविशोरदाः}}{\Afootnote{K; \skt{तन्मतम् ॥}\ A; \skt{तर्त्मरोविशोरदाः ॥}\ H}}\skt{ ॥}\pend

\pstart\versepar \ledsidenote{\passref{6.3.21}}\edtext{\skt{अरुजम्वाह्यतः}}{\lemma{\skt{वाह्यतः}}\Afootnote{K H; \skt{बाह्यतः}\ A}}\skt{ }\edtext{\skt{शूनं}}{\Afootnote{K H; \skt{शून}\ A}}\skt{ }\edtext{\skt{वर्त्मयस्य}}{\Afootnote{K; \skt{मन्तः क्लिन्नं स्रवत्यपि । कण्डूनिस्}\ A; \skt{वर्त्मर्यस्य न हि । प्र}\ H}}\skt{ }\edtext{\skt{नरम्ये हि}}{\Afootnote{K; \skt{तोदभूयि}\ A}}\skt{ । }\pend
\pstart\versepar \skt{}\edtext{\edtext{\edtext{\edtext{\edtext{\skt{प्रक्लिन्नवर्त्मतम्विद्यात्क्लिन्नमत्यर्थम्न्ततः}}{\lemma{\skt{प्रक्लिन्नवर्त्मतम्}}\Afootnote{K; \skt{ष्ठं क्लिन्नवर्त्म त}\ A}}}{\lemma{\skt{विद्यात्क्लिन्नम्}}\Afootnote{K H; \skt{दुच्यते ॥}\ A}}}{\lemma{\skt{अत्य्}}\Afootnote{K H; \emph{om.}\ A}}}{\lemma{\skt{अर्थम्न्}}\Afootnote{K H; \emph{om.}\ A}}}{\lemma{\skt{ततः}}\Afootnote{K H; \emph{om.}\ A}}\skt{ ॥}\edtext{}{\lemma{\skt{र्त्मयस्य… प्रक्लिन्}}\Afootnote{\emph{om.}\ H}}\pend

\pstart\versepar \ledsidenote{\passref{6.3.22}}\skt{यस्य धौतानि धौतानि }\edtext{\skt{सन्दिद्यन्ते}}{\Afootnote{K H; \skt{संबध्यन्ते}\ A}}\skt{ पुनः पुनः । }\pend
\pstart\versepar \skt{}\edtext{\skt{वर्त्मान्यपरिपक्वानि}}{\Afootnote{A K; \skt{वर्त्मान्यपरिपक्वानि}\ H}}\skt{ }\edtext{\skt{निद्यपरि}}{\Afootnote{K H; \skt{विद्याद}\ A}}\skt{ }\edtext{\skt{किन्नवर्त्म}}{\Afootnote{K H; \skt{क्लिन्नवर्त्म}\ A}}\skt{ तत् ।}\pend

\pstart\versepar \ledsidenote{\passref{6.3.23}}\edtext{\edtext{\skt{विदग्धसन्धिर्निश्चेष्टं}}{\lemma{\skt{विदग्धसन्धिर्}}\Afootnote{K H; \skt{विमुक्तसन्धि}\ A}}}{\lemma{\skt{निश्चेष्टं}}\Afootnote{K H; \skt{निश्चेष्टं}\ A}}\skt{ वर्त्म यस्य }\edtext{\skt{निर्मील्यते}}{\Afootnote{K; \skt{न मील्यते ।}\ A; \skt{निमील्यते ।}\ H}}\skt{ । }\pend
\pstart\versepar \skt{}\edtext{\skt{एतद्वातहतन्नाम}}{\lemma{\skt{एतद्वातहतन्}}\Afootnote{K H; \skt{एतद्वातहतं}\ A}}\skt{ }\edtext{\edtext{\skt{जानीयादक्षिचिन्तकः}}{\Afootnote{K; \skt{जानीयादक्षिचिन्तकैः ॥}\ H}}\skt{ ।}}{\lemma{\skt{नाम… अक्षिचिन्तकः}}\Afootnote{\emph{om.}\ A}}\pend

\pstart\versepar \ledsidenote{\passref{6.3.24}}\skt{वर्त्मान्तरस्थम्विषमं }\edtext{\edtext{\skt{ग्रन्थिभूतमवेदनं}}{\lemma{\skt{ग्रन्थिभूतम्}}\Afootnote{K; \skt{ग्रन्थिभूत}\ A H}}}{\lemma{\skt{अवेदनं}}\Afootnote{K; \skt{मवेदनम् ।}\ A; \skt{मवे । विलम्}\ H}}\skt{ । }\pend
\pstart\versepar \skt{}\edtext{\edtext{\skt{तार्वुदमिति}}{\lemma{\skt{तार्वुदम्}}\Afootnote{K; \skt{विज्ञेयमर्बुदं}\ A; \skt{वि च ॥}\ H}}}{\lemma{\skt{इति}}\Afootnote{K; \skt{पुंसां}\ A; \emph{om.}\ H}}\skt{ }\edtext{\edtext{\skt{सरक्षमविलम्वि}}{\lemma{\skt{सरक्षम्}}\Afootnote{K; \skt{सरक्त}\ A; \emph{om.}\ H}}}{\lemma{\skt{अविलम्वि}}\Afootnote{K; \skt{मवलम्बि}\ A; \emph{om.}\ H}}\skt{ }\edtext{\skt{च}}{\Afootnote{K; \skt{तम् ॥}\ A; \emph{om.}\ H}}\skt{ ॥}\pend

\pstart\versepar \ledsidenote{\passref{6.3.25}}\edtext{\skt{निमेषिणीः}}{\Afootnote{K; \skt{निमेषणीः}\ A; \skt{निमेषिण्यः}\ H}}\skt{ सिरा वायुः }\edtext{\skt{संविष्टः}}{\Afootnote{K H; \skt{प्रविष्टो}\ A}}\skt{ }\edtext{\skt{सन्धिसंश्रयः}}{\Afootnote{K; \skt{वर्त्मसंश्रयाः ।}\ A; \skt{सन्धिसंश्रितः ।}\ H}}\skt{ । }\pend
\pstart\versepar \skt{}\edtext{\edtext{\skt{चालयत्यथ}}{\lemma{\skt{चालयत्य्}}\Afootnote{K H; \skt{चालय}\ A}}}{\lemma{\skt{अथ}}\Afootnote{K H; \skt{त्यति}\ A}}\skt{ वर्त्मानि }\edtext{\skt{निमेष}}{\Afootnote{K H; \skt{निमेषः}\ A}}\skt{ }\edtext{\skt{इति}}{\Afootnote{K H; \skt{स ग}\ A}}\skt{ }\edtext{\skt{तम्विदुः}}{\Afootnote{K H; \skt{दो मतः ॥}\ A}}\skt{ ॥}\pend

\pstart\versepar \ledsidenote{\passref{6.3.26}}\edtext{\skt{यः श्रितौ}}{\Afootnote{K; \skt{छिन्नाश्छिना विवर्धन्ते}\ A; \skt{शोणितार्शश्चिन्नं चाभिप्रवद्धते ॥}\ H}}\skt{ }\edtext{\skt{वर्त्मववेतलो}}{\Afootnote{K; \skt{वर्त्मस्था मृदवो}\ A; \emph{om.}\ H}}\skt{ }\edtext{\skt{मृदुरंकुरः}}{\Afootnote{K; \skt{ऽङ्कुराः । दा}\ A; \emph{om.}\ H}}\skt{ । }\pend
\pstart\versepar \skt{}\edtext{\skt{तद्रक्तजं}}{\Afootnote{K; \skt{हकण्डूरु}\ A; \emph{om.}\ H}}\skt{ }\edtext{\skt{शोणितार्शः}}{\Afootnote{K; \skt{जोपेतास्ते}\ A; \emph{om.}\ H}}\skt{ }\edtext{\edtext{\skt{छिन्नश्चाभिप्रवर्धते}}{\lemma{\skt{छिन्नश्}}\Afootnote{K; \skt{ऽर्शः}\ A; \emph{om.}\ H}}}{\lemma{\skt{चाभिप्रवर्धते}}\Afootnote{K; \skt{शोणितसंभवाः ॥}\ A; \emph{om.}\ H}}\skt{ ।}\pend

\pstart\versepar \ledsidenote{\passref{6.3.27}}\skt{अपाकः }\edtext{\skt{कठिनः}}{\Afootnote{A K; \skt{कठिनं}\ H}}\skt{ }\edtext{\skt{स्थूलोग्रन्थिवर्त्मानुजोरूजः}}{\Afootnote{K; \skt{स्थूलो ग्रन्थिर्वर्त्मभवोऽरुजः । सकण्डूः पिच्छिलः कोलप्रमाणो}\ A; \skt{स्थूलोग्रन्थिर्व्वत्मानुकोरूजः ।}\ H}}\skt{ । }\pend
\pstart\versepar \skt{}\edtext{\skt{लगणा}}{\Afootnote{K; \skt{लगणस्}\ A; \skt{लगण}\ H}}\skt{ }\edtext{\skt{सव्या}}{\Afootnote{K; \skt{तु सः ॥}\ A; \skt{नाम सव्या}\ H}}\skt{ }\edtext{\skt{निलम्वेत}}{\Afootnote{K; \emph{om.}\ A H}}\skt{ ।}\pend

\pstart\versepar \ledsidenote{\passref{6.3.28}}\edtext{\skt{गस्तीरोपगतं}}{\Afootnote{K; \skt{शूनं यद्वर्त्म बहुभिः सूक्ष्मैश्छिद्रैः समन्वितम् । बिसमन्तर्जल}\ A; \emph{om.}\ H}}\skt{ }\edtext{\skt{यथा}}{\Afootnote{K H; \skt{इव}\ A}}\skt{ । }\pend
\pstart\versepar \skt{}\edtext{\skt{विसवर्त्मतिजातीयात्विसप्रप्रख्योदकश्रवां}}{\Afootnote{K H; \skt{बिसवर्त्मेति तन्मतम् ॥}\ A}}\pend

\pstart\versepar \ledsidenote{\passref{6.3.29}}\edtext{\skt{प्रचालितानि}}{\Afootnote{K H; \emph{om.}\ A}}\skt{ }\edtext{\skt{वातेन}}{\Afootnote{K H; \emph{om.}\ A}}\skt{ }\edtext{\skt{पक्ष्माश्विसं ति}}{\Afootnote{\skt{दोषाः पक्ष्मा}\ A; \skt{पक्ष्मा}\ K; \skt{पक्ष्माश्ञ्नौअं ति}\ H}}\skt{ हि । }\pend
\pstart\versepar \skt{क्षि त्ति च ।}\edtext{}{\lemma{\skt{हि । क्षि त्ति च}}\Afootnote{\skt{शयगतास्तीक्ष्णाग्राणि खराणि च । निर्वर्तयन्ति पक्ष्माणि तैर्घुष्टं चाक्षि दूयते ॥}\ A; \skt{ण्यक्षि विशंति हि । घृष्यन्त्यक्षि मुहुस्तानि संरंतञ्जनयन्ति च ।}\ K; \skt{च ॥}\ H}}\pend

\pstart\versepar \ledsidenote{\passref{6.3.30}}\edtext{\skt{असिते}}{\Afootnote{H; \emph{om.}\ A; \skt{असित}\ K}}\skt{ }\edtext{\skt{सित}}{\Afootnote{H; \skt{उद्धृतैरुद्धृतैः}\ A; \skt{सित… श्}\ K}}\skt{ }\edtext{\skt{मूलकात्पतं }\edtext{\skt{न्ह्यपि}}{\lemma{\skt{अपि}}\Afootnote{H; \skt{द्वेषी}\ A; \emph{om.}\ K}}}{\lemma{\skt{मूलकात्पतं न्ह्य्}}\Afootnote{\skt{शान्तिः पक्ष्मभिश्चोपजायते । वातातपानल}\ A; \skt{च मूलकोषात्पतन्त्यपि ।}\ K; \skt{विन्दूश्च मूलकात्पतन्त्य्}\ H}}\skt{ । । }\pend
\pstart\versepar \skt{पक्ष्मकोपः }\edtext{\edtext{\skt{सविज्ञेयाव्याधिः}}{\Afootnote{K H; \skt{स}\ A}}\skt{ परमदारुण }\edtext{\skt{इति॥ {\symfont ❈}}}{\Afootnote{K; \skt{इति ॥}\ H}}\skt{ ॥}}{\lemma{\skt{विज्ञेयाव्याधिः… {\symfont ❈}}}\Afootnote{\emph{om.}\ A}}\pend

\pstart \ledsidenote{\passref{6.3.trailer}}\edtext{\skt{शालाक्ये}}{\Afootnote{K; \skt{इति सुश्रुतसंहितायामुत्तरतन्त्रे वर्त्मगतरोगविज्ञानीयो नाम तृतीयोऽध्यायः ॥३ ॥}\ A; \skt{शालाक्ये रोगविज्ञानीयन्तृतीयोऽध्याय ॥॥ {\symfont ❈} ॥ ॥}\ H}}\skt{ ३ ॥}\pend

\endnumbering
\beginnumbering
\eledsection*{Adhyāya 4}
\pstart \ledsidenote{\passref{6.4.1}}\edtext{\skt{अथातः}}{\Afootnote{A K; \skt{अथातो}\ H}}\skt{ }\edtext{\edtext{\skt{शुक्लगतरोगविज्ञानीयम्व्याख्यास्यामः}}{\lemma{\skt{शुक्लगतरोगविज्ञानीयम्}}\Afootnote{H; \skt{शुक्लगतरोगविज्ञानीयमध्यायं}\ A; \skt{शुक्लगतरोगविज्ञानीयं व्या व्याख्यास्यामः ॥ ॥}\ K}}}{\lemma{\skt{व्याख्यास्यामः}}\Afootnote{A H; \emph{om.}\ K}}\skt{ ॥}\pend

\pstart\versepar \ledsidenote{\passref{6.4.3}}\edtext{\skt{प्रस्तारिशुक्लक्षतजाधिमांसः}}{\Afootnote{K H; \skt{प्रस्तारिशुक्लक्षतजाधिमांस}\ A}}\skt{ }\edtext{\skt{स्नावर्मसम्ञ्ज्ञाः}}{\Afootnote{\skt{स्राय्वर्मसंज्ञाः}\ A; \skt{स्नावार्मसञ्ज्ञाः}\ K; \skt{स्नाय्वर्म्मसंज्ञाः}\ H}}\skt{ खलु पञ्चरोगाः । }\pend
\pstart\versepar \skt{}\edtext{\skt{स्याच्छुक्तिकोथार्जुनपिष्टकाख्यः}}{\Afootnote{H; \skt{स्युः शुक्तिका चार्जुनपिष्टकौ च}\ A; \skt{स्याच्छुक्तिकोथार्जनपिष्टकाख्यः}\ K}}\skt{ जालं सिराणां }\edtext{\skt{पिडकाञ्च याः}}{\Afootnote{K H; \skt{पिडकाश्च याः}\ A}}\skt{ स्युः ॥}\pend

\pstart\versepar \ledsidenote{\passref{6.4.4}}\skt{रोगा }\edtext{\skt{बलासग्रथितेन}}{\Afootnote{A; \skt{वलासग्रथितेन}\ K H}}\skt{ }\edtext{\skt{सार्धमेकादशैते}}{\lemma{\skt{एकादशैते}}\Afootnote{H; \skt{एकादशाक्ष्णोः}\ A; \skt{एकोदशैते}\ K}}\skt{ खलु शुक्लभागे ॥ }\pend
\pstart\versepar \skt{}\edtext{\skt{प्रस्तारि}}{\Afootnote{A K; \skt{प्रास्तारि}\ H}}\skt{ }\edtext{\skt{प्रथितमथार्मशुक्लभागे}}{\lemma{\skt{अथार्मशुक्लभागे}}\Afootnote{K H; \skt{इहार्म शुक्लभागे}\ A}}\skt{ विस्तीर्णन्तनु }\edtext{\skt{रूधिरप्रभं}}{\Afootnote{K H; \skt{रुधिरप्रभं}\ A}}\skt{ }\edtext{\skt{सलीलं}}{\Afootnote{K H; \skt{सनीलम् ॥}\ A}}\skt{ ।}\pend

\pstart\versepar \ledsidenote{\passref{6.4.5}}\skt{शुक्लाख्यं मृदु कथयन्ति शुक्लभागे }\edtext{\skt{सस्वेतं}}{\Afootnote{K H; \skt{सश्वेतं}\ A}}\skt{ }\edtext{\skt{सममभिवर्धते}}{\lemma{\skt{अभिवर्धते}}\Afootnote{K H; \skt{इह वर्धते}\ A}}\skt{ चिरान्तेत् ॥ }\pend
\pstart\versepar \skt{यन्मांसं प्रचयमुपैति शुक्लभागे }\edtext{\skt{पद्माभं}}{\Afootnote{A K; \skt{प}\ H\textsuperscript{ac}; \skt{पद्द्माभं}\ H\textsuperscript{pc}}}\skt{ }\edtext{\skt{तदुशन्ति}}{\Afootnote{K; \skt{तदुपदिशन्ति}\ A; \skt{द्नं मृदु तदुशन्ति}\ H\textsuperscript{ac}; \skt{नं मृदु तदुशन्ति}\ H\textsuperscript{pc}}}\skt{ लोहितार्म ।}\edtext{}{\lemma{\skt{चिरान्तेत् ॥ यन्मांसं}}\Afootnote{\skt{चिरेण । यन्मांसं}\ A; \skt{चिरान्तेत् ॥ य मांसं}\ K; \skt{चिरान्तते ॥ यन्मान्समम्}\ H}}\pend

\pstart\versepar \ledsidenote{\passref{6.4.6}}\skt{विस्तीर्ण्णं मृदु }\edtext{\skt{वहलं}}{\Afootnote{K H; \skt{बहलं}\ A}}\skt{ यकृत्प्रकाशं }\edtext{\edtext{\skt{श्यावम्पृथु}}{\lemma{\skt{श्यावम्}}\Afootnote{A K; \skt{श्यावा}\ H}}\skt{ च}}{\lemma{\skt{पृथु च}}\Afootnote{K; \skt{वा त}\ A; \skt{भं पृथ्य्च}\ H}}\skt{ }\edtext{\skt{तथापि}}{\Afootnote{K H; \skt{दधिक}\ A}}\skt{ मांसकार्म ॥ }\pend
\pstart\versepar \skt{स्या वा }\edtext{\skt{मस्छिरवहमांसवर्त्मभागे}}{\Afootnote{\skt{त् । शुक्ले यत्पिशितमुपैति वृद्धिमेतत्स्नाय्वर्मेत्यभिपठितं खरं प्रपाण्डु ॥}\ A; \skt{मस्छिरवहमांसवर्त्मभागे}\ K; \skt{स्नाय्वर्म्मस्छिरवद्धमांसवर्त्म्मभागे}\ H}}\skt{ }\edtext{\skt{प्रस्तोरि}}{\Afootnote{K H; \emph{om.}\ A}}\skt{ }\edtext{\skt{प्रभवति}}{\Afootnote{K H; \emph{om.}\ A}}\skt{ }\edtext{\skt{मारूताद्विशुष्कं}}{\Afootnote{K H; \emph{om.}\ A}}\skt{ ।}\edtext{}{\lemma{\skt{मांसकार्म ॥ स्या वा}}\Afootnote{K; \skt{मांसजार्म विद्या}\ A; \skt{मांसकार्म्म ॥}\ H}}\pend

\pstart\versepar \ledsidenote{\passref{6.4.7}}\skt{श्यावाः स्युः पिशितनिभाश्च }\edtext{\skt{बिन्दवो ये}}{\Afootnote{A; \skt{विन्दवो ये}\ K H}}\skt{ }\edtext{\skt{शुक्लाभाः}}{\Afootnote{K H; \skt{शुक्त्याभाः}\ A}}\skt{ }\edtext{\skt{सितनियताः स}}{\Afootnote{K; \skt{सितनयने स}\ A; \skt{सितनिमहाः स}\ H}}\skt{ }\edtext{\skt{शुक्लिसञ्ज्ञः}}{\Afootnote{K; \skt{शुक्तिसंज्ञः}\ A H}}\skt{ ॥ }\pend
\pstart\versepar \skt{एको यः }\edtext{\skt{शशरूधिरोपमश्च}}{\Afootnote{K H; \skt{शशरुधिरोपमस्तु}\ A}}\skt{ }\edtext{\skt{बिन्दुः}}{\Afootnote{A; \skt{विन्दुः}\ K H}}\skt{ शुक्लस्थो भवति }\edtext{\skt{तमर्जनम्वदन्ति}}{\lemma{\skt{अर्जनम्}}\Afootnote{K; \skt{अर्जुनं}\ A H}}\skt{ ।}\pend

\pstart\versepar \ledsidenote{\passref{6.4.8}}\skt{उत्सन्नः }\edtext{\skt{सलिलनिभो}}{\Afootnote{K H; \skt{सलिलनिभो}\ A}}\skt{ }\edtext{\skt{थ}}{\Afootnote{K H; \skt{ऽथ}\ A}}\skt{ पिष्टशुक्लो }\edtext{\skt{बिन्दुर्यः स}}{\Afootnote{\skt{बिन्दुर्यो}\ A; \skt{विन्दुर्यः स}\ K H}}\skt{ }\edtext{\skt{भवति}}{\Afootnote{K H; \skt{भवति स}\ A}}\skt{ पिष्टकः सुवृत्तः ॥ }\pend
\pstart\versepar \skt{जालाभः कठिनसिरो }\edtext{\skt{महां}}{\Afootnote{K; \skt{महा}\ A H}}\skt{ }\edtext{\skt{सरक्तः}}{\Afootnote{K; \skt{न्सरक्तः}\ A H}}\skt{ }\edtext{\skt{सन्तानः}}{\Afootnote{A K; \skt{सन्तान}\ H}}\skt{ स्मृत इह जालसंज्ञितस्तु ।}\pend

\pstart\versepar \ledsidenote{\passref{6.4.9}}\skt{शुक्लस्थाः }\edtext{\skt{सितपिडाकाः}}{\Afootnote{K; \skt{सितपिडकाः}\ A H}}\skt{ }\edtext{\skt{शिरावृताया यास्तां}}{\Afootnote{K H; \skt{सिरावृता यास्ता वि}\ A}}\skt{ }\edtext{\skt{यादसिअसमीपजाः}}{\Afootnote{K H; \skt{द्यादसितसमीपजाः}\ A}}\skt{ }\edtext{\skt{सिराजा}}{\Afootnote{K H; \skt{सिराजाः ।}\ A}}\skt{ ॥ }\pend
\pstart\versepar \skt{}\edtext{\skt{काशाभो}}{\Afootnote{K H; \skt{कांस्याभो}\ A}}\skt{ }\edtext{\edtext{\skt{मृदुरथवापि}}{\lemma{\skt{मृदुर्}}\Afootnote{K H; \emph{om.}\ A}}}{\lemma{\skt{अथवापि}}\Afootnote{K H; \emph{om.}\ A}}\skt{ }\edtext{\skt{विन्वकल्यौ}}{\Afootnote{K; \skt{भवति सितेऽम्बुबिन्दुतुल्यः}\ A; \skt{विन्दुकल्पो}\ H}}\skt{ }\edtext{\skt{विज्ञेयो}}{\Afootnote{K H; \skt{स ज्ञेयो}\ A}}\skt{ }\edtext{\skt{नयनसिते}}{\Afootnote{K; \skt{ऽमृदुररुजो}\ A; \skt{नयन्सिते}\ H}}\skt{ }\edtext{\skt{वलासकाख्य}}{\Afootnote{K; \skt{बलासकाख्यः ॥}\ A; \skt{वलासौप}\ H}}\skt{ }\edtext{\skt{इति}}{\Afootnote{K H; \emph{om.}\ A}}\skt{ ॥}\pend

\pstart \ledsidenote{\passref{6.4.trailer}}\edtext{\skt{॥ सालाक्य}}{\Afootnote{\skt{इति सुश्रुतसंहितायामुत्तरतन्त्रे शुक्लगतरोगविज्ञानीयो नाम चतुर्थोऽध्यायः ॥४ ॥}\ A; \skt{सालाक्य थ्व}\ K; \skt{शालाक्यस्थाने चतुर्थोऽध्यायः ॥}\ H}}\skt{ ४ ॥}\pend

\endnumbering
\beginnumbering
\eledsection*{Adhyāya 5}
\pstart \ledsidenote{\passref{6.5.1}}\edtext{\skt{अथातो}}{\Afootnote{K H; \skt{अथातः}\ A}}\skt{ }\edtext{\edtext{\skt{कृष्णगतरोगविञ्जानीयम्व्याख्यास्यामः}}{\lemma{\skt{कृष्णगतरोगविञ्जानीयम्}}\Afootnote{H; \skt{कृष्णगतरोगविज्ञानीयमध्यायं}\ A; \skt{कृष्णगतरोगविञ्जानीयम्व्या}\ K}}}{\lemma{\skt{व्याख्यास्यामः}}\Afootnote{A H; \emph{om.}\ K}}\pend

\pstart\versepar \ledsidenote{\passref{6.5.3}}\skt{यत्सव्रणं }\edtext{\skt{शुक्लमथाव्रणंअञ्च}}{\Afootnote{K H; \skt{शुक्ल(शुक्र)मथाव्रणं वा}\ A}}\skt{ }\edtext{\edtext{\skt{पाकाययश्चाप्यजा}}{\lemma{\skt{पाकाययश्}}\Afootnote{K; \skt{पाकात्ययश्}\ A; \skt{पाकात्ययंश्}\ H}}}{\lemma{\skt{चाप्यजा}}\Afootnote{K H; \skt{चाप्यजका}\ A}}\skt{ तथैव । }\pend
\pstart\versepar \skt{चत्वार एते }\edtext{\skt{नयनामयासु }\edtext{\skt{कृष्णप्रदेशे}}{\Afootnote{K; \skt{कृष्णप्रदेशं}\ H}}\skt{ नियता }\edtext{\skt{भवन्ति}}{\Afootnote{K H; \skt{ऽभिहिता विकाराः कृष्णाश्रयाः संग्रहतः पुरस्तात् ॥}\ A}}\skt{ ॥}}{\lemma{\skt{नयनामयासु… भ}}\Afootnote{\emph{om.}\ A}}\pend

\pstart\versepar \ledsidenote{\passref{6.5.4}}\skt{निमग्नरूपं हि भवेत्तु कृष्णे }\edtext{\skt{सूच्यैव}}{\Afootnote{K H; \skt{सूच्येव}\ A}}\skt{ }\edtext{\skt{विध्वश्च}}{\Afootnote{K; \skt{विद्धं प्र}\ A; \skt{विध्वे ।}\ H}}\skt{ }\edtext{\skt{विभाति}}{\Afootnote{K; \skt{तिभाति}\ A; \emph{om.}\ H}}\skt{ }\edtext{\skt{यद्वै}}{\Afootnote{A K; \emph{om.}\ H}}\skt{ । }\pend
\pstart\versepar \skt{स्रावं }\edtext{\skt{स्रवेदूष्णमतीव}}{\Afootnote{K H; \skt{स्रवेदुष्णमतीव}\ A}}\skt{ }\edtext{\skt{यत्र}}{\Afootnote{K H; \skt{रुक्च}\ A}}\skt{ तत्सव्रणं }\edtext{\edtext{\skt{शुक्रमुदाहरन्ति}}{\lemma{\skt{शुक्रम्}}\Afootnote{K H; \skt{शुक्र}\ A}}}{\lemma{\skt{उदाहरन्ति}}\Afootnote{K H; \skt{(शुक्ल)मुदाहरन्ति ॥}\ A}}\skt{ ॥}\pend

\pstart \ledsidenote{\passref{6.5.8cd}}\skt{वैहाय }\edtext{\skt{सात्रप्रभप्रकाशसल्य}}{\Afootnote{K; \skt{सात्रप्रतिमप्रकाश सल्य}\ H}}\skt{ }\edtext{\skt{व्रमात}}{\Afootnote{K; \skt{व्रणं साध्यत}\ H}}\skt{ }\edtext{\skt{शंशतंतत्}}{\Afootnote{K; \skt{मम्सतन्तत्}\ H}}\pend

\pstart\versepar \ledsidenote{\passref{6.5.6}}\skt{विच्छिन्नमध्यंम्पिशितावृतम्वा चलं }\edtext{\edtext{\skt{सिरासूक्ष्ममदृष्टिकृच्च}}{\lemma{\skt{सिरासूक्ष्मम्}}\Afootnote{K H; \skt{सिरासक्त}\ A}}}{\lemma{\skt{अदृष्टिकृच्च}}\Afootnote{K H; \skt{मदृष्टिकृच्च ।}\ A}}\skt{ । }\pend
\pstart\versepar \skt{द्वित्वग्गतं }\edtext{\skt{लोहितसन्ततश्च}}{\Afootnote{K H; \skt{लोहितमन्ततश्च}\ A}}\skt{ }\edtext{\edtext{\skt{चिरोत्थितश्वाचि}}{\lemma{\skt{चिरोत्थितश्}}\Afootnote{K; \skt{चिरोत्थितं}\ A H}}}{\lemma{\skt{वाचि}}\Afootnote{K; \skt{चापि}\ A; \skt{वापि}\ H}}\skt{ }\edtext{\skt{विवर्जनीत}}{\Afootnote{K H; \skt{विवर्जनीयम् ॥}\ A}}\skt{ ।}\pend

\pstart\versepar \ledsidenote{\passref{6.5.7}}\skt{उष्णाश्रुपातः पिडका च }\edtext{\skt{नेत्रे}}{\Afootnote{K H; \skt{कृष्णे}\ A}}\skt{ }\edtext{\skt{यस्मिन्भवेन्मुद्गनिभं श्च}}{\Afootnote{K; \skt{यस्मिन्भवेन्मुद्गनिभं च}\ A; \skt{यस्मिन्भवेन्मुद्गनिभ}\ H}}\skt{ }\edtext{\skt{शुक्रम्}}{\Afootnote{A K; \skt{ञ्शुक्रम् ॥}\ H}}\skt{ । }\pend
\pstart\versepar \skt{स्वेतः }\edtext{\skt{समाक्रामति सर्वतो हि दोषो हि }\edtext{\skt{यस्यासितमण्डलन्तु}}{\Afootnote{K; \skt{यस्यासितमशुलन्तु ।}\ H}}\skt{ ।}}{\lemma{\skt{स्वेतः… यस्यासि}}\Afootnote{\emph{om.}\ A}}\pend

\pstart\versepar \ledsidenote{\passref{6.5.10}}\edtext{\edtext{\skt{तमक्षिपाकातयमक्षिपाकं}}{\lemma{\skt{तमक्षिपाकातयम्}}\Afootnote{K H; \skt{तमक्षिपाकात्यय}\ A}}}{\lemma{\skt{अक्षिपाकं}}\Afootnote{K H; \skt{मक्षिकोप}\ A}}\skt{ }\edtext{\skt{सर्वात्मकं}}{\Afootnote{K H; \skt{समुत्थितं तीव्ररुजं}\ A}}\skt{ }\edtext{\edtext{\skt{वर्जयितव्यमाहुः}}{\lemma{\skt{वर्जयितव्यम्}}\Afootnote{K H; \skt{वदन्ति ।}\ A}}}{\lemma{\skt{आहुः}}\Afootnote{K H; \emph{om.}\ A}}\skt{ । }\pend
\pstart\versepar \skt{}\edtext{\skt{अजापूरीषप्रतिसोरूजावां}}{\Afootnote{K; \skt{अजापुरीषप्रतिमो रुजावा}\ A; \skt{अजापूरीषप्रतिसोरूजावा}\ H}}\skt{ }\edtext{\skt{सलोहितो}}{\Afootnote{K; \skt{न्सलोहितो}\ A H}}\skt{ }\edtext{\skt{लोहितपिच्छिलास्रः}}{\Afootnote{K H; \skt{लोहितपिच्छिलाश्रुः ।}\ A}}\skt{ । }\edtext{\skt{विगृह्यरुष्टं}}{\Afootnote{K H; \skt{विदार्य कृष्णं}\ A}}\skt{ }\edtext{\skt{प्रचयोभ्युपैति तञ्चाजकाजातमिति}}{\lemma{\skt{प्रचयोभ्युपैति तञ्}}\Afootnote{K; \skt{प्रचयोऽभ्युपैति तं}\ A H}}\skt{ }\edtext{\skt{व्यवस्येत्}}{\Afootnote{A K; \skt{व्यवस्येतिति ॥ {\symfont ❈} ॥}\ H}}\skt{ । ० ॥}\pend

\pstart \ledsidenote{\passref{6.5.trailer}}\edtext{\skt{शालाक्ये}}{\Afootnote{K; \skt{इति सुश्रुतसंहितायामुत्तरतन्त्रे कृष्णगतरोगविज्ञानीयो नाम पञ्चमोऽध्यायः ॥५ ॥}\ A; \skt{शालाक्यस्थाने पञ्चमोऽध्यायः ॥ {\symfont ❈} ॥ ॥}\ H}}\skt{ ५}\pend

\endnumbering
\beginnumbering
\eledsection*{Adhyāya 6}
\pstart \ledsidenote{\passref{6.6.1}}\skt{अथातः }\edtext{\skt{सर्व्वगतरोगविज्ञानीयं}}{\Afootnote{H; \skt{सर्वगतरोगविज्ञानीयमध्यायं}\ A; \skt{सर्व्वगतरोगविज्ञानीयं व्याव्याख्यास्यामः ॥}\ K}}\skt{ }\edtext{\skt{व्याख्यास्यामः}}{\Afootnote{A H; \emph{om.}\ K}}\skt{ ॥}\pend

\pstart\versepar \ledsidenote{\passref{6.6.3}}\skt{स्यन्दास्तु चत्वार इहोपदिष्टा स्तावन्त }\edtext{\skt{एव}}{\Afootnote{K H; \skt{एवे}\ A}}\skt{ }\edtext{\skt{प्रथिताधिमन्थाः}}{\Afootnote{K H; \skt{ह तथाधिमन्थाः ।}\ A}}\skt{ । }\pend
\pstart\versepar \skt{}\edtext{\skt{शोफाम्वितोशोफयुतै च}}{\Afootnote{K H; \skt{शोफान्वितोऽशोफयुतश्च}\ A}}\skt{ }\edtext{\skt{पाकावित्येवमेतेदश}}{\lemma{\skt{पाकाव्}}\Afootnote{A K; \skt{पाका}\ H}}\skt{ }\edtext{\skt{सम्प्रंदिष्टाः}}{\Afootnote{K H; \skt{संप्रदिष्टाः ॥}\ A}}\skt{ ।}\pend

\pstart\versepar \ledsidenote{\passref{6.6.4}}\edtext{\skt{हताधिमन्थोनिलपर्ययश्च}}{\Afootnote{K H; \skt{हताधिमन्थोऽनिलपर्ययश्च}\ A}}\skt{ }\edtext{\skt{शुष्काक्षिपाकोन्यत}}{\Afootnote{K H; \skt{शुष्काक्षिपाकोऽन्यत}\ A}}\skt{ एव वातः । }\pend
\pstart\versepar \skt{}\edtext{\skt{दृष्टिस्तथाम्लाध्युषिता}}{\Afootnote{K H; \skt{दृष्टिस्ततथाम्लाध्युषिता}\ A}}\skt{ }\edtext{\skt{सिराणां}}{\Afootnote{K H; \skt{सिराणा}\ A}}\skt{ }\edtext{\skt{उत्पातहर्षावपि}}{\Afootnote{K H; \skt{मुत्पातहर्षावपि}\ A}}\skt{ }\edtext{\skt{सर्व्वभागौ}}{\Afootnote{K H; \skt{सर्वभागाः ॥}\ A}}\skt{ ।}\pend

\pstart\versepar \ledsidenote{\passref{6.6.5}}\edtext{\skt{प्रायेन}}{\Afootnote{K; \skt{प्रायेण}\ A H}}\skt{ }\edtext{\skt{सर्व}}{\Afootnote{K H; \skt{सर्वे}\ A}}\skt{ }\edtext{\skt{नयनामयास्तु}}{\Afootnote{K H; \skt{नयनामयास्तु}\ A}}\skt{ }\edtext{\edtext{\skt{भवन्त्यभिष्यन्दनिमित्तमूलाः}}{\lemma{\skt{भवन्त्य्}}\Afootnote{K; \skt{भवन्}\ A H}}}{\lemma{\skt{अभिष्यन्दनिमित्तमूलाः}}\Afootnote{K; \skt{त्यभिष्यन्दनिमित्तमूलाः ।}\ A H}}\skt{ । }\pend
\pstart\versepar \skt{}\edtext{\edtext{\edtext{\skt{तस्मादभिष्यन्दमयोमानमुपकुर्यदामेताशु}}{\lemma{\skt{तस्मादभिष्यन्दमयोमानम्}}\Afootnote{K; \skt{तस्मादभिष्यन्दमुदीर्यमा}\ A; \skt{तस्मादभिष्यन्दमूयोमान म्}\ H}}}{\lemma{\skt{उपकुर्यदाम्}}\Afootnote{K; \skt{णमुपाच}\ A; \skt{उप}\ H}}}{\lemma{\skt{एताशु}}\Afootnote{K; \skt{रेदाशु}\ A; \skt{मेताशु}\ H}}\skt{ हिताय धीमान् ॥}\pend

\pstart\versepar \ledsidenote{\passref{6.6.6}}\edtext{\skt{निस्तोदनस्तम्भनरोमहर्ष}}{\Afootnote{K H; \skt{निस्तोदनं स्तम्भनरोमहर्ष}\ A}}\skt{ }\edtext{\skt{संहर्षपारूष्यशिरोभिघाताः}}{\Afootnote{K H; \skt{सङ्घर्षपारुष्यशिरोभितापाः ।}\ A}}\skt{ । }\pend
\pstart\versepar \skt{}\edtext{\skt{विशुष्कभावा}}{\Afootnote{K H; \skt{विशुष्कभावः}\ A}}\skt{ शिशिराश्रुता च वाताभिपन्ने नयने भवन्ति ॥}\pend

\pstart\versepar \ledsidenote{\passref{6.6.7}}\skt{दाहप्रपाकौ }\edtext{\skt{शिशिराभिनन्दौ}}{\Afootnote{K H; \skt{शिशिराभिनन्दा}\ A}}\skt{ }\edtext{\skt{धूमायनम्वाप्यसमुछ्रयश्च}}{\lemma{\skt{वाप्यसमुछ्रयश्च}}\Afootnote{K; \skt{बाष्पसमुच्छ्रयश्च ।}\ A; \skt{वापसमुछ्रयश्च ।}\ H}}\skt{ । }\pend
\pstart\versepar \skt{}\edtext{\skt{उष्णाश्रुतापीतकनेत्रता च}}{\Afootnote{A K; \skt{उष्णाश्रुता षीतकनेत्रता च}\ H}}\skt{ पित्ताभिपन्ने नयने भवन्ति ॥}\pend

\pstart\versepar \ledsidenote{\passref{6.6.9}}\skt{ताम्राश्रुता लोहितनेत्रता च राज्यः }\edtext{\skt{समन्तादपि}}{\Afootnote{K H; \skt{समन्तादति}\ A}}\skt{ लोहिताश्च । }\pend
\pstart\versepar \skt{पित्तस्य लिङ्गानि च }\edtext{\skt{कीर्तितानि}}{\Afootnote{K H; \skt{यानि तानि}\ A}}\skt{ रक्ताभिपन्ने नयने भवन्ति ॥}\pend

\pstart\versepar \ledsidenote{\passref{6.6.8}}\skt{उष्णाभिनन्दा }\edtext{\skt{गुरूताक्षिशोफाः}}{\Afootnote{K; \skt{गुरुताक्षिशोफः}\ A; \skt{गुरूताक्षिशोफः}\ H}}\skt{ । }\edtext{\skt{कण्डुपदेहावति}}{\Afootnote{K; \skt{कण्डूपदेहौ सित}\ A; \skt{कण्डुपदेह्यवति}\ H}}\skt{ }\edtext{\skt{शीतता च}}{\Afootnote{K; \skt{तातिशैत्यम् ।}\ A; \skt{शीतअ च ।}\ H}}\skt{ । }\pend
\pstart\versepar \skt{}\edtext{\skt{श्रावावः}}{\Afootnote{K; \skt{स्रावो मुहुः}\ A; \skt{श्रावः मुहुः}\ H}}\skt{ }\edtext{\skt{पिच्चिल}}{\Afootnote{K H; \skt{पिच्छिल}\ A}}\skt{ एव चापि }\edtext{\skt{कफाभिपन्ने}}{\Afootnote{A K; \skt{कफाभ्पन्ने}\ H}}\skt{ नयने भवन्ति ॥}\pend

\pstart\versepar \ledsidenote{\passref{6.6.10}}\edtext{\edtext{\edtext{\skt{वृद्धैरेतैरभिष्यन्दैर्नराणामक्रियवतां}}{\lemma{\skt{वृद्धैर्}}\Afootnote{K; \skt{वृद्धै}\ A H}}}{\lemma{\skt{एतैरभिष्यन्दैर्}}\Afootnote{K; \skt{रेतैरभिष्यन्दै}\ A; \skt{रेतैरभिष्यन्दैर्}\ H}}}{\lemma{\skt{नराणामक्रियवतां}}\Afootnote{K; \skt{र्नराणामक्रियावताम् ।}\ A; \skt{नराणामक्रियावतां ।}\ H}}\skt{ । }\pend
\pstart\versepar \skt{}\edtext{\skt{तावन्त्यस्त्वधिरोगाः}}{\Afootnote{K H; \skt{तावन्तस्त्वधिमन्थाः}\ A}}\skt{ }\edtext{\skt{स्यूर्नयने}}{\Afootnote{K H; \skt{स्युर्नयने}\ A}}\skt{ }\edtext{\skt{तीव्रवेदनाः}}{\Afootnote{A K; \skt{तीव्रवेदना ॥}\ H}}\skt{ ॥}\pend

\pstart\versepar \ledsidenote{\passref{6.6.11}}\edtext{\skt{उत्पाट्यते}}{\Afootnote{K H; \skt{उत्पाट्यत}\ A}}\skt{ }\edtext{\skt{भृशं च}}{\Afootnote{K; \skt{इवात्य}\ A; \skt{फलञ्च}\ H}}\skt{ }\edtext{\skt{कुर्यस्य}}{\Afootnote{K; \skt{र्थं नेत्रं}\ A; \skt{र्यस्य}\ H}}\skt{ }\edtext{\skt{निर्मर्थ्यते}}{\Afootnote{K H; \skt{निर्मथ्यते}\ A}}\skt{ तथा । }\pend
\pstart\versepar \skt{}\edtext{\skt{शिरसोर्धञ्च }\edtext{\skt{तम्विद्यादधिमन्थः}}{\lemma{\skt{विद्यादधिमन्थः}}\Afootnote{K; \skt{विद्यादधिमन्थं}\ A H}}}{\lemma{\skt{शिरसोर्धञ्च तम्}}\Afootnote{K H; \skt{शिरसोऽर्धं च तं}\ A}}\skt{ स्वलक्षणैः ॥}\pend

\pstart\versepar \ledsidenote{\passref{6.6.20}}\edtext{\skt{हन्यादृष्टिश्श्लैष्मिकः}}{\Afootnote{K H; \skt{हन्यादृष्टिं}\ A}}\skt{ }\edtext{\skt{सप्तरात्रात् । तथ्रक्तजाः}}{\Afootnote{K; \skt{सप्तरात्रात्कफोत्थोऽधीमन्थोऽसृक्संभवः}\ A; \skt{सप्तरात्रात्तथा रक्तजा}\ H}}\skt{ पञ्चरात्रात् । }\pend
\pstart\versepar \skt{षड्रात्राद्वआवातिका वै }\edtext{\skt{निहन्यात्मिथ्याचारात्पैत्तिकः}}{\Afootnote{K H; \skt{निहन्यान्मिथ्याचारात्पैत्तिकः}\ A}}\skt{ सद्य एव ॥}\edtext{}{\lemma{\skt{पञ्चरात्रात् । षड्रात्राद्वआवातिका वै}}\Afootnote{K H; \skt{पञ्चरात्रात् । षड्रात्राद्वौ मारुतोत्थो}\ A}}\pend

\pstart \ledsidenote{\passref{6.6.21ab}}\skt{कण्डूपदेहाश्रुयुतः }\edtext{\skt{पक्वोदुम्वरसन्निभः}}{\Afootnote{K H; \skt{पक्वोदुम्बरसन्निभः ।}\ A}}\skt{ ।}\pend

\pstart \ledsidenote{\passref{6.6.22}}\edtext{\skt{संरम्भी}}{\Afootnote{K H; \skt{जुष्टो मुहुः स्रवेच्चास्रमुष्णशीताम्बु पिच्छिलम् । संरम्भी}\ A}}\skt{ }\edtext{\skt{दद्यते यस्तु}}{\Afootnote{K H; \skt{पच्यते यश्च}\ A}}\skt{ }\edtext{\skt{सोक्षि}}{\Afootnote{K; \skt{नेत्र}\ A; \skt{सोऽक्षिक्षि}\ H}}\skt{ पाकः सशोफजः ॥}\pend

\pstart\versepar \ledsidenote{\passref{6.6.23}}\skt{शोफहीनानि लिङ्गानि नेत्रपाके त्वशोफजे । }\pend
\pstart\versepar \skt{}\edtext{\skt{उपेक्षणदक्षि}}{\Afootnote{K; \skt{उपेक्षणादक्षि}\ A H}}\skt{ }\edtext{\skt{यदाभिमन्थो}}{\Afootnote{K H; \skt{यदाधिमन्थो}\ A}}\skt{ वातात्मकः }\edtext{\skt{स्यन्दयति}}{\Afootnote{K H; \skt{सादयति}\ A}}\skt{ }\edtext{\skt{प्रसह्यः}}{\Afootnote{K; \skt{प्रसह्य ।}\ A H}}\skt{ । }\edtext{\edtext{\skt{रूजाभिरूग्राभिरसाध्य}}{\lemma{\skt{रूजाभिर्}}\Afootnote{K H; \skt{रुजाभिर्}\ A}}}{\lemma{\skt{ऊग्राभिर्}}\Afootnote{K H; \skt{उग्राभिर्}\ A}}\skt{ एष }\edtext{\skt{हताभिमन्थः}}{\Afootnote{K H; \skt{हताधिमन्थः}\ A}}\skt{ }\edtext{\skt{किल}}{\Afootnote{K H; \skt{खलु}\ A}}\skt{ नाम रोगः ॥}\pend

\pstart\versepar \ledsidenote{\passref{6.6.25}}\edtext{\skt{पक्ष्मद्वयेक्षिभ्रुवमाश्रितस्तु}}{\Afootnote{K; \skt{पक्ष्मद्वयाक्षिभ्रुवमाश्रितस्तु}\ A; \skt{पक्ष्मद्वयेक्षिब्रुवमाश्रितस्तु}\ H}}\skt{ }\edtext{\skt{यदानिलः}}{\Afootnote{K H; \skt{यन्नानिलः}\ A}}\skt{ सञ्चरति प्रदुष्टः }\pend
\pstart\versepar \skt{}\edtext{\skt{पर्यायतेश्चापि}}{\lemma{\skt{पर्यायतेश्}}\Afootnote{K H; \skt{पर्यायशश्}\ A}}\skt{ }\edtext{\skt{रूजाकरोति । तं}}{\Afootnote{K; \skt{रुजः करोति तं}\ A; \skt{रूजाः करोति तं}\ H}}\skt{ }\edtext{\skt{पर्ययं}}{\Afootnote{K H; \skt{वातप}\ A}}\skt{ }\edtext{\skt{वातकृतम्वदन्ति}}{\Afootnote{K H; \skt{र्यायमुदाहरन्ति ॥}\ A}}\skt{ ॥}\pend

\pstart\versepar \ledsidenote{\passref{6.6.26}}\edtext{\skt{कूनितं}}{\Afootnote{K; \skt{यत्कूणितं}\ A; \skt{कुकूणितं}\ H}}\skt{ }\edtext{\skt{दारूणरुक्ष्मवर्त्म}}{\Afootnote{K; \skt{दारुणरूक्षवर्त्म}\ A; \skt{दारूणपक्ष्मवर्त्म}\ H}}\skt{ }\edtext{\skt{सन्दिद्यते}}{\Afootnote{K H; \skt{विलोकने}\ A}}\skt{ चाविलदर्शनं यत् । }\pend
\pstart\versepar \skt{सुदारूणं }\edtext{\skt{यत्प्रतिबोधने च}}{\Afootnote{A; \skt{यत्प्रतिवोधने च}\ K H}}\skt{ }\edtext{\skt{शुष्काक्षिपाकोपहं}}{\Afootnote{K H; \skt{शुष्काक्षिपाकोपहतं}\ A}}\skt{ तदक्षि ॥}\edtext{}{\lemma{\skt{यत् । सुदारूणं}}\Afootnote{K; \skt{यत् । सुदारुणं}\ A; \skt{तत । सुदारूणं}\ H}}\pend

\pstart\versepar \ledsidenote{\passref{6.6.27}}\edtext{\skt{यस्यावटूकर्ण्णशिरोहनूषु}}{\Afootnote{K H; \skt{यस्यावटूकर्णशिरोहनुस्थो}\ A}}\skt{ मन्यागतो }\edtext{\edtext{\skt{वाप्यनिलोऽन्यथो}}{\lemma{\skt{वाप्यनिलो}}\Afootnote{\skt{वाप्यनिलो}\ A; \skt{वाप्यनिलो}\ K H}}\skt{ वा}}{\lemma{\skt{ऽन्यथो वा}}\Afootnote{\skt{ऽन्यतो वा ।}\ A; \skt{न्यथो वा ।}\ K; \skt{न्यतो वा ।}\ H}}\skt{ । }\pend
\pstart\versepar \skt{}\edtext{\skt{कुर्याद्रुजां}}{\Afootnote{K H; \skt{कुर्याद्रुजो}\ A}}\skt{ }\edtext{\skt{वै}}{\Afootnote{K H; \skt{ऽति}\ A}}\skt{ भ्रुवि लोचने }\edtext{\skt{वातमन्यतो}}{\Afootnote{K H; \skt{वा तमन्यतो}\ A}}\skt{ }\edtext{\skt{तमुदाहरन्ति}}{\Afootnote{K H; \skt{वातमुदाहरन्ति ॥}\ A}}\skt{ ॥}\pend

\pstart\versepar \ledsidenote{\passref{6.6.28}}\skt{अम्लेन भुक्तेन }\edtext{\skt{विदावा}}{\Afootnote{K; \skt{विदाहिना च}\ A; \skt{विदाहिना चा}\ H}}\skt{ }\edtext{\skt{पक्वायते}}{\Afootnote{K H; \skt{संछाद्यते}\ A}}\skt{ सर्वत एव नेत्रं । }\pend
\pstart\versepar \skt{शोफान्वितं }\edtext{\skt{लोहितकं}}{\Afootnote{K H; \skt{लोहितकैः}\ A}}\skt{ }\edtext{\skt{सदाहसेन्तदम्लाध्युषितम्वदन्ति}}{\lemma{\skt{सदाहसेन्तदम्लाध्युषितम्}}\Afootnote{K H; \skt{सनीलैरेतादृगम्लाध्युषितं}\ A}}\skt{ ॥}\pend

\pstart\versepar \ledsidenote{\passref{6.6.29}}\skt{अवेदना वापि सवेदना वा }\edtext{\skt{यस्याक्षिराज्योभिभवन्ति}}{\Afootnote{K; \skt{यस्याक्षिराज्यो हि भवन्ति}\ A; \skt{यस्याक्षिराज्योऽभिभवन्ति}\ H}}\skt{ ताम्राः । }\pend
\pstart\versepar \skt{मुहुर्विरज्यन्ति च ताः }\edtext{\skt{सतादृग्व्याधिः}}{\Afootnote{K; \skt{समन्ताद्व्याधिः}\ A; \skt{सता श्व्याधिः}\ H}}\skt{ }\edtext{\skt{शिरोत्पात}}{\Afootnote{K H; \skt{सिरोत्पात}\ A}}\skt{ इति प्रदिष्टः ॥}\pend

\pstart\versepar \ledsidenote{\passref{6.6.30}}\skt{मोहात्सिरोत्पात उपेक्षितस्तु जायेत रोगस्तु }\edtext{\skt{सिराप्रहर्षः}}{\Afootnote{A H; \skt{सिराप्र ।}\ K}}\skt{ । }\pend
\pstart\versepar \skt{}\edtext{\skt{तत्राश्रमच्छं}}{\Afootnote{K H; \skt{ताम्राच्छमस्रं}\ A}}\skt{ स्रवति }\edtext{\skt{प्रगाढं}}{\Afootnote{A K; \skt{प्रगाढंन्}\ H}}\skt{ तथा न }\edtext{\skt{शक्रोत्यभिवीक्षितुं}}{\Afootnote{K H; \skt{शक्नोत्यभिवीक्षितुं}\ A}}\skt{ }\edtext{\skt{चेति}}{\Afootnote{K H; \skt{च ॥}\ A}}\skt{ ॥ ॥}\pend

\pstart \ledsidenote{\passref{6.6.trailer}}\edtext{\skt{शालाक्ये}}{\Afootnote{K; \skt{इति सुश्रुतसंहितायामुत्तरतन्त्रे सर्वगतरोगविज्ञानीयो नाम षष्ठोऽध्यायः ॥६ ॥}\ A; \skt{शालाक्यस्थाने षष्ठोऽध्यायः ॥}\ H}}\skt{ ६}\pend

\endnumbering
\beginnumbering
\eledsection*{Adhyāya 7}
\pstart \ledsidenote{\passref{6.7.1}}\skt{अथातो दृष्टिगतरोगविज्ञानीयं }\edtext{\skt{व्याख्यास्यामः}}{\Afootnote{H; \skt{अध्यायं व्याख्यास्यामः ॥}\ A; \skt{व्याव्याव्याख्यास्यामः ।}\ K}}\skt{ ।}\pend

\pstart\versepar \ledsidenote{\passref{6.7.3}}\edtext{\skt{मसूरदलमात्रन्तु}}{\Afootnote{K H; \skt{मरूरदलमात्रां तु}\ A}}\skt{ }\edtext{\skt{पञ्चभूतप्रसादजां।}}{\lemma{\skt{पञ्चभूतप्रसादजां}}\Afootnote{A K; \skt{पञ्चभूतओरसाजं । ।}\ H}}\skt{ । }\pend
\pstart\versepar \skt{}\edtext{\skt{खद्योतविस्फुलिङ्गाभासिद्धां}}{\Afootnote{K H; \skt{खद्योतविस्फुलिङ्गाभामिद्धां}\ A}}\skt{ तेजोभिरव्ययैः ॥}\pend

\pstart\versepar \ledsidenote{\passref{6.7.4}}\edtext{\skt{आवृताम्पटलेनोक्ष्णोर्वाह्येन}}{\Afootnote{K H; \skt{आवृतां पटलेनाक्ष्णोर्बाह्येन}\ A}}\skt{ विवराकृतिं । । }\pend
\pstart\versepar \skt{शीतसात्म्यां नृणां }\edtext{\skt{हन्ति}}{\Afootnote{K H; \skt{दृष्टि}\ A}}\skt{ माहुर्नयनचिन्तकाः ।}\pend

\pstart\versepar \ledsidenote{\passref{6.7.5}}\edtext{\skt{रोगांस्तदाश्रयां}}{\Afootnote{K H; \skt{रोगांस्तदाश्रया}\ A}}\skt{ }\edtext{\skt{घोरां}}{\Afootnote{K H; \skt{न्घोरा}\ A}}\skt{ }\edtext{\skt{षट्वा}}{\Afootnote{K H; \skt{न्षट्च षट्च}\ A}}\skt{ }\edtext{\skt{प्रचक्षते}}{\Afootnote{K; \skt{प्रचक्ष्महे ।}\ A; \skt{प्रजायते । ।}\ H}}\skt{ । }\pend
\pstart\versepar \skt{पटलानुप्रविष्टस्य तिमिरस्य च लक्षणं ।}\pend

\pstart\versepar \ledsidenote{\passref{6.7.6}}\skt{सिराभिरभिसम्प्राप्य }\edtext{\skt{विगुलैभ्यन्तरे}}{\Afootnote{K; \skt{विगुणोऽभ्यन्तरे}\ A; \skt{विभुलभ्यन्तरे}\ H}}\skt{ भृशं । }\pend
\pstart\versepar \skt{प्रथमे पटले }\edtext{\skt{यस्य}}{\Afootnote{K H; \skt{दोषो यस्य}\ A}}\skt{ }\edtext{\skt{रोगादृष्टं}}{\Afootnote{K H; \skt{दृष्टौ}\ A}}\skt{ }\edtext{\skt{व्यवस्थिताः}}{\Afootnote{K; \skt{व्यवथितः ॥}\ A; \skt{व्यवस्थितः ॥}\ H}}\skt{ ॥}\pend

\pstart\versepar \ledsidenote{\passref{6.7.7}}\skt{अव्यक्तानि स रूपाणि }\edtext{\skt{कदाचिदथ}}{\Afootnote{K H; \skt{सर्वाण्येव प्र}\ A}}\skt{ पश्यति ॥ }\pend
\pstart\versepar \skt{}\edtext{\edtext{\skt{दृष्टिर्भृशम्विह्वलति}}{\lemma{\skt{दृष्टिर्}}\Afootnote{K H; \skt{दृष्टि}\ A}}}{\lemma{\skt{भृशम्विह्वलति}}\Afootnote{K H; \skt{र्भृशं विह्वलति}\ A}}\skt{ }\edtext{\skt{द्वितीयं}}{\Afootnote{A K; \skt{द्वितीय}\ H}}\skt{ पटलं गते ॥}\pend

\pstart\versepar \ledsidenote{\passref{6.7.8}}\edtext{\skt{मक्षिकां}}{\Afootnote{K H; \skt{मक्षिका}\ A}}\skt{ }\edtext{\edtext{\skt{मशकांश्चापि}}{\lemma{\skt{मशकांश्}}\Afootnote{K; \skt{मशका}\ A; \skt{मशकां}\ H}}}{\lemma{\skt{चापि}}\Afootnote{K H; \skt{न्केशाञ्}\ A}}\skt{ जालकानि च पश्यति । }\pend
\pstart\versepar \skt{}\edtext{\skt{मशुलानि}}{\Afootnote{K H; \skt{मण्डलानि}\ A}}\skt{ }\edtext{\skt{पताकांश्च}}{\Afootnote{A K; \skt{पताकानि}\ H}}\skt{ }\edtext{\skt{मरीच्यः}}{\Afootnote{K H; \skt{मरीचीः}\ A}}\skt{ कुण्डलानि च ।}\pend

\pstart\versepar \ledsidenote{\passref{6.7.9}}\edtext{\skt{परिपाकांश्च}}{\Afootnote{K; \skt{परिप्लवांश्च}\ A; \skt{परिपाकाश्च}\ H}}\skt{ }\edtext{\skt{विविधां}}{\Afootnote{K; \skt{विविधा}\ A H}}\skt{ }\edtext{\skt{वर्षमभ्रास्तमान्सि च}}{\Afootnote{K H; \skt{न्वर्षमभ्रं तमांसि च ।}\ A}}\skt{ । }\pend
\pstart\versepar \skt{}\edtext{\skt{दूरस्थान्यपि}}{\Afootnote{K H; \skt{दूरस्थान्यपि}\ A}}\skt{ रूपाणि मन्यते च समीपतः ।}\pend

\pstart\versepar \ledsidenote{\passref{6.7.10}}\edtext{\skt{समीपस्थानि}}{\Afootnote{A K; \skt{समीपस्थानि च}\ H}}\skt{ दूरे च दृष्टेर्दर्शनविभ्रमात् । }\pend
\pstart\versepar \skt{यत्नवानपि चात्यर्थं सूचीपाशं न पश्यति ।}\edtext{}{\lemma{\skt{दृष्टेर्दर्शनविभ्रमात् । यत्नवानपि}}\Afootnote{K H; \skt{दृष्टेर्गोचरविभ्रमात् । यत्नवानपि}\ A}}\pend

\pstart\versepar \ledsidenote{\passref{6.7.11}}\skt{ऊर्ध्वम्पश्यति }\edtext{\skt{नाधस्तां}}{\Afootnote{K H; \skt{गाधस्ता}\ A}}\skt{ }\edtext{\edtext{\skt{तृतीयम्पतलं}}{\lemma{\skt{तृतीयम्}}\Afootnote{K H; \skt{त्तृतीयं}\ A}}}{\lemma{\skt{पतलं}}\Afootnote{K H; \skt{पटलं}\ A}}\skt{ गते ॥ }\pend
\pstart\versepar \skt{}\edtext{\skt{महानपि च}}{\Afootnote{K H; \skt{महान्त्यपि च}\ A}}\skt{ रूपाणि }\edtext{\skt{च्छादितानीवचाम्वरैः}}{\Afootnote{K H; \skt{च्छादितानीव वाससा ॥}\ A}}\skt{ ।}\pend

\pstart\versepar \ledsidenote{\passref{6.7.12}}\skt{कर्णनासाक्षि }\edtext{\skt{हीनानि}}{\Afootnote{K; \skt{युक्तानि}\ A; \skt{कर्ण्णानि}\ H}}\skt{ }\edtext{\skt{विकृतानीवपश्यति}}{\Afootnote{K H; \skt{विपरीतानि वीक्षते ।}\ A}}\skt{ । }\pend
\pstart\versepar \skt{यथादोषञ्च }\edtext{\skt{रज्येतर्दृष्टिर्दोषे}}{\Afootnote{K; \skt{रज्येत दृष्टिर्दोषे}\ A H}}\skt{ }\edtext{\skt{वलीयसि}}{\Afootnote{K H; \skt{बलीयसि ॥}\ A}}\skt{ ।}\pend

\pstart\versepar \ledsidenote{\passref{6.7.13}}\edtext{\skt{अधस्थे च}}{\Afootnote{K; \skt{अधःस्थिते}\ A; \skt{अधः स्थे च}\ H}}\skt{ }\edtext{\skt{समीपस्थं}}{\Afootnote{A K; \skt{समीपस्थे}\ H}}\skt{ }\edtext{\skt{दूरस्थञ्चोपरिस्थितः}}{\lemma{\skt{चोपरिस्थितः}}\Afootnote{K H; \skt{चोपरिस्थिते ।}\ A}}\skt{ । }\pend
\pstart\versepar \skt{}\edtext{\skt{पार्श्वस्थे च}}{\Afootnote{K H; \skt{पार्श्वस्थिते}\ A}}\skt{ तथा दोषे }\edtext{\skt{पार्श्वस्थानि च}}{\Afootnote{K; \skt{पार्श्वस्थानि न}\ A H}}\skt{ पश्यति ।}\pend

\pstart\versepar \ledsidenote{\passref{6.7.14}}\edtext{\skt{समन्तात्तु}}{\Afootnote{K H; \skt{समन्ततः}\ A}}\skt{ स्थिते दोषे }\edtext{\skt{संकुलानि}}{\Afootnote{K H; \skt{सङ्कुलानी}\ A}}\skt{ }\edtext{\skt{प्रपश्यति}}{\Afootnote{K; \skt{व पश्यति ।}\ A; \skt{न पश्यति ।}\ H}}\skt{ । }\pend
\pstart\versepar \skt{}\edtext{\skt{दृष्टिमध्येस्थि ते}}{\Afootnote{K H; \skt{दृष्टिमध्यगते}\ A}}\skt{ दोषे }\edtext{\skt{सह ञ्च}}{\Afootnote{K H; \skt{स एकं}\ A}}\skt{ }\edtext{\skt{पश्यति}}{\Afootnote{K H; \skt{मन्यते द्विधा ॥}\ A}}\skt{ ।}\pend

\pstart\versepar \ledsidenote{\passref{6.7.15}}\edtext{\skt{द्विधाश्रिते}}{\Afootnote{K H; \skt{द्विधास्थिते}\ A}}\skt{ }\edtext{\skt{द्विधा}}{\Afootnote{K H; \skt{त्रिधा}\ A}}\skt{ }\edtext{\skt{पश्येद्वहुधा}}{\Afootnote{K; \skt{पश्येद्बहुधा}\ A; \skt{पश्यए वहुधा}\ H}}\skt{ चानवस्थिते । }\pend
\pstart\versepar \skt{दोषे दृष्टाश्रिते }\edtext{\skt{तियगेकं वै}}{\Afootnote{K; \skt{तिर्य शकं वै}\ H}}\skt{ सम्यते द्विधा । }\pend
\pstart\versepar \skt{}\edtext{\skt{तिमिराख्य स वै}}{\Afootnote{K H; \skt{तिमिराख्यः स वै}\ A}}\skt{ दोषश्च }\edtext{\skt{चतुर्थ}}{\Afootnote{K; \skt{तुर्थं}\ A; \skt{तुर्थ}\ H}}\skt{ }\edtext{\skt{पटल्ं}}{\Afootnote{K; \skt{पटलं}\ A; \skt{पतलं}\ H}}\skt{ }\edtext{\skt{गतंः}}{\Afootnote{K; \skt{गतः}\ A H}}\skt{ ।}\edtext{}{\lemma{\skt{दोषे… ति}}\Afootnote{\emph{om.}\ A}}\pend

\pstart\versepar \ledsidenote{\passref{6.7.16}}\skt{रुणद्धि सर्वतो }\edtext{\skt{दृष्टिलिङ्गनाशमतः}}{\Afootnote{K H; \skt{दृष्टिं लिङ्गनाशः स उ}\ A}}\skt{ }\edtext{\skt{परं}}{\Afootnote{K H; \skt{च्यते ।}\ A}}\skt{ । }\pend
\pstart\versepar \skt{}\edtext{\skt{अस्मिन्नपि न ते न}}{\Afootnote{K; \skt{तस्मिन्नपि त}\ A; \skt{अस्मिन्नपि त}\ H}}\skt{ मोभूते नातिरूढे }\edtext{\skt{महागद}}{\Afootnote{K; \skt{महागदे ॥}\ A; \skt{महागव ।}\ H}}\skt{ ।}\pend

\pstart\versepar \ledsidenote{\passref{6.7.17}}\skt{चन्द्रादित्यौ }\edtext{\skt{सनक्षत्रावन्तरिक्षे च}}{\Afootnote{K H; \skt{सनक्षत्रावन्तरीक्षे च}\ A}}\skt{ विद्युतः । }\pend
\pstart\versepar \skt{निर्मलानि च तेजांसि भ्राजिष्णूनि च पश्यति ।}\pend

\pstart\versepar \ledsidenote{\passref{6.7.18}}\skt{स एव }\edtext{\skt{लिंगनाशस्तु}}{\Afootnote{K H; \skt{लिङ्गनाशस्तु}\ A}}\skt{ }\edtext{\skt{नीलिकाकाचसंज्ञितः}}{\Afootnote{A K; \skt{नीलकाकाचसंज्ञितः ॥}\ H}}\skt{ । }\pend
\pstart\versepar \skt{}\edtext{\skt{वातेन}}{\Afootnote{K H; \skt{तत्र वातेन}\ A}}\skt{ }\edtext{\skt{चापि}}{\Afootnote{K H; \emph{om.}\ A}}\skt{ रूपाणि }\edtext{\skt{भ्रमन्तीवहि}}{\Afootnote{K H; \skt{भ्रमन्तीव स}\ A}}\skt{ पश्यति ।}\pend

\pstart\versepar \ledsidenote{\passref{6.7.19}}\edtext{\skt{आविलान्यरूणाभानि}}{\Afootnote{K; \skt{आविलान्यरुणाभानि}\ A; \skt{आविलान्यरूणाभानि}\ H}}\skt{ }\edtext{\skt{व्याविद्धानीव}}{\Afootnote{K H; \skt{व्याविद्धानि च}\ A}}\skt{ }\edtext{\skt{पार्श्वतः}}{\Afootnote{K H; \skt{मानवः ।}\ A}}\skt{ ॥ }\pend
\pstart\versepar \skt{}\edtext{\skt{पित्तेनादित्यखद्योतशक्रचापताशक्र}}{\Afootnote{K H; \skt{पित्तेनादित्यखद्योतशक्र}\ A}}\skt{ }\edtext{\skt{चापतडिड्गुणां}}{\Afootnote{K; \skt{चापतडिद्गुणान् ॥}\ A; \skt{चापतडिड्गुणान् ।}\ H}}\skt{ ।}\pend

\pstart\versepar \ledsidenote{\passref{6.7.20}}\edtext{\edtext{\skt{नृत्यन्तश्चैव}}{\lemma{\skt{नृत्यन्तश्}}\Afootnote{K H; \emph{om.}\ A}}}{\lemma{\skt{चैव}}\Afootnote{K H; \emph{om.}\ A}}\skt{ }\edtext{\skt{शिखिनः}}{\Afootnote{K H; \skt{शिखिब}\ A}}\skt{ }\edtext{\skt{सर्वं}}{\Afootnote{K; \skt{र्हविचित्राणि}\ A; \skt{स}\ H}}\skt{ }\edtext{\skt{कृष्तश्च}}{\Afootnote{K; \skt{नीलकृष्णानि}\ A; \skt{र्व्वश्च}\ H}}\skt{ पश्यति । }\pend
\pstart\versepar \skt{}\edtext{\edtext{\skt{कफेनाञ्भोजहारेन्द्रशंखाम्वुमिवतारकाः}}{\Afootnote{K; \skt{कफेन पश्येद्रूपाणि स्निग्धानि च सितानि च ॥}\ A; \skt{कफेनतम्भोजहारेन्द्रशंखाम्वुमिवतारकां ।}\ H}}\skt{ ।}}{\lemma{\skt{ञ्भोजहारेन्द्रशंखाम्वुमिवतारकाः ।}}\Afootnote{\emph{om.}\ A}}\pend

\pstart \ledsidenote{\passref{6.7.21cd}}\edtext{\skt{पश्यन्त}}{\Afootnote{K H; \skt{पश्येद}\ A}}\skt{ }\edtext{\skt{इव}}{\Afootnote{K H; \emph{om.}\ A}}\skt{ }\edtext{\skt{रूपाणि}}{\Afootnote{K H; \emph{om.}\ A}}\skt{ }\edtext{\skt{सिग्धानि च}}{\Afootnote{K H; \emph{om.}\ A}}\skt{ }\edtext{\skt{हिमानि च}}{\Afootnote{K H; \skt{सूक्ष्माण्यत्यर्थं व्यभ्रे चैवाभ्रसंप्लवम् ॥}\ A}}\skt{ ।}\pend

\pstart \ledsidenote{\passref{6.7.22cd}}\edtext{\skt{पश्येद्रक्तानि}}{\Afootnote{K H; \skt{तथा रक्तेन}\ A}}\skt{ }\edtext{\skt{रक्तेन}}{\Afootnote{K H; \skt{रक्तानि}\ A}}\skt{ तमांसि विविधानि च ॥}\pend

\pstart \ledsidenote{\passref{6.7.22ab}}\edtext{\edtext{\skt{ससितान्यकृष्टानि}}{\lemma{\skt{ससितान्य्}}\Afootnote{K; \skt{सलिलप्लावितानी}\ A; \skt{ससितान्य}\ H}}}{\lemma{\skt{अकृष्टानि}}\Afootnote{K; \skt{व}\ A; \skt{थकृष्टानि}\ H}}\skt{ }\edtext{\edtext{\skt{पीतान्यपि}}{\lemma{\skt{पीतान्य्}}\Afootnote{K; \skt{परिजा}\ A; \skt{पीता}\ H}}\skt{ च}}{\lemma{\skt{अपि च}}\Afootnote{K; \skt{ड्यानि}\ A; \skt{न्यपि च}\ H}}\skt{ मानवः ।}\pend

\pstart \ledsidenote{\passref{6.7.23cd}}\skt{सन्निपातेन चित्राणि }\edtext{\skt{विप्लतानीव}}{\Afootnote{K; \skt{विप्लुतानि च}\ A; \skt{विप्लुतानीव}\ H}}\skt{ पश्यति ॥}\pend

\pstart\versepar \ledsidenote{\passref{6.7.24}}\edtext{\skt{वहुधा वा}}{\Afootnote{K H; \skt{बहुधा वा}\ A}}\skt{ द्विधा वापि सर्वाण्येव समन्ततः । }\pend
\pstart\versepar \skt{हीनाधिकाङ्गान्यथवा }\edtext{\skt{ज्योतींष्यपि च}}{\Afootnote{K; \skt{ज्योतींष्यपि च}\ A; \skt{ज्योतीष्यपि च}\ H}}\skt{ }\edtext{\skt{भूयशः}}{\Afootnote{K H; \skt{पश्यति ॥}\ A}}\skt{ ।}\pend

\pstart\versepar \ledsidenote{\passref{6.7.25}}\skt{पित्तं कुर्यात्परिम्लायि }\edtext{\skt{मूर्च्छितं}}{\Afootnote{A K; \skt{मूर्च्छिततं}\ H}}\skt{ }\edtext{\skt{सूर्यतेजसा}}{\Afootnote{K H; \skt{रक्ततेजसा ।}\ A}}\skt{ । }\pend
\pstart\versepar \skt{}\edtext{\skt{पिता}}{\Afootnote{K H; \skt{पीता}\ A}}\skt{ }\edtext{\skt{दिशस्तथान्यो वा}}{\Afootnote{K H; \skt{दिशस्तथोद्यन्त}\ A}}\skt{ }\edtext{\skt{रवीनिव च}}{\Afootnote{K H; \skt{मादित्यमिव}\ A}}\skt{ पश्यति ।}\pend

\pstart\versepar \ledsidenote{\passref{6.7.26}}\edtext{\skt{विकीर्यमाणं}}{\Afootnote{K; \skt{विकीर्यमाणा}\ A H}}\skt{ }\edtext{\skt{खद्योतां}}{\Afootnote{K H; \skt{न्खद्योतैर्}\ A}}\skt{ }\edtext{\skt{क्षां}}{\Afootnote{K; \skt{वृक्षांस्}\ A; \skt{दक्षान्}\ H}}\skt{ }\edtext{\skt{तेःजांसि}}{\Afootnote{K; \skt{तेजोभिर्}\ A; \skt{तेजांसि}\ H}}\skt{ }\edtext{\skt{चोत्तरं}}{\Afootnote{K H; \skt{एव च ।}\ A}}\skt{ । }\pend
\pstart\versepar \skt{}\edtext{\skt{वक्षामि}}{\Afootnote{K; \skt{वक्ष्यामि}\ A; \skt{वक्ष्मि}\ H}}\skt{ षड्विधं }\edtext{\skt{रागैर्लिङ्गनाशमतः}}{\Afootnote{K H; \skt{रागैर्लिङ्गनाशमतः}\ A}}\skt{ परं ॥}\pend

\pstart\versepar \ledsidenote{\passref{6.7.27}}\edtext{\skt{रागोरूणे}}{\Afootnote{K; \skt{रागोऽरुणो}\ A; \skt{रोगारूणो}\ H}}\skt{ मारुतजः }\edtext{\skt{प्रदिष्टाम्ली च}}{\Afootnote{K H; \skt{प्रदि}\ A}}\skt{ }\edtext{\skt{नीलश्च}}{\Afootnote{K H; \skt{ष्टः}\ A}}\skt{ }\edtext{\skt{तथैव}}{\Afootnote{K H; \emph{om.}\ A}}\skt{ पित्तात् । }\pend
\pstart\versepar \skt{कफात्सितः शोणितजः }\edtext{\skt{सरक्तः}}{\Afootnote{K H; \skt{तु रक्तः}\ A}}\skt{ समस्त }\edtext{\skt{दोषप्रभवस्तु}}{\Afootnote{K H; \skt{दोषोऽथ वि}\ A}}\skt{ }\edtext{\skt{चित्रः}}{\Afootnote{K H; \skt{चित्ररूपः ॥}\ A}}\skt{ ।}\edtext{}{\lemma{\skt{पित्तात् । कफात्सितः}}\Afootnote{K H; \skt{पित्तात्परिम्लाय्यथवापि नीलः । कफात्सितः}\ A}}\pend

\pstart\versepar \ledsidenote{\passref{6.7.28}}\edtext{\skt{अरूणम्मण्दलं}}{\Afootnote{K H; \skt{रक्तजं मण्डलं}\ A}}\skt{ }\edtext{\skt{दृष्टास्थूलकाचोरूणस्तु सः}}{\Afootnote{K; \skt{दृष्टौ स्थूलकाचानलप्रभम् ।}\ A; \skt{दृष्टास्थूलकाचोरूणस्सः ।}\ H}}\skt{ । }\pend
\pstart\versepar \skt{परिम्लायिनि रोगे }\edtext{\skt{स्याम्लायीपञ्चमण्डलं}}{\Afootnote{K; \skt{स्यान्म्लाय्यानीलं च मण्डलम् ॥}\ A; \skt{स्यान्म्लायीपञ्चमण्डलं ॥}\ H}}\skt{ ।}\pend

\pstart\versepar \ledsidenote{\passref{6.7.29}}\edtext{\skt{दोषक्षयात्स्वयं}}{\Afootnote{K H; \skt{दोषक्षयात्कदाचित्स्यात्स्वयं}\ A}}\skt{ तत्र }\edtext{\skt{कदाचिन्स्यात्तु}}{\Afootnote{K H; \skt{च}\ A}}\skt{ दर्शनं । }\pend
\pstart\versepar \skt{}\edtext{\skt{अरूणं}}{\Afootnote{K H; \skt{अरुणं}\ A}}\skt{ मण्डलं }\edtext{\skt{वातात्चञ्चलम्परूषन्तथा}}{\lemma{\skt{परूषन्}}\Afootnote{K H; \skt{परुषं}\ A}}\skt{ ॥}\pend

\pstart\versepar \ledsidenote{\passref{6.7.30}}\edtext{\skt{पित्ततो}}{\Afootnote{K H; \skt{पित्ता}\ A}}\skt{ }\edtext{\skt{मण्डलं}}{\Afootnote{K H; \skt{न्मण्डलम्}\ A}}\skt{ }\edtext{\skt{नीलं}}{\Afootnote{K H; \skt{आनीलं}\ A}}\skt{ }\edtext{\skt{कांस्याभञ्वा}}{\Afootnote{K; \skt{कांस्याभं पी}\ A; \skt{कांस्याभम्वा}\ H}}\skt{ }\edtext{\skt{सपीतकं}}{\Afootnote{K H; \skt{तमेव वा ।}\ A}}\skt{ । }\pend
\pstart\versepar \skt{श्लेष्मणा }\edtext{\skt{वहुलं}}{\Afootnote{K; \skt{बहलं}\ A; \skt{वहं}\ H}}\skt{ }\edtext{\skt{स्निग्धां}}{\Afootnote{K; \skt{स्निग्धं}\ A H}}\skt{ }\edtext{\skt{शंख}}{\Afootnote{A K; \skt{शंखं}\ H}}\skt{ }\edtext{\skt{क्षीरेन्दुपाण्डुरं}}{\Afootnote{K H; \skt{कुन्देन्दुपाण्डुरम् ॥}\ A}}\skt{ ।}\pend

\pstart \ledsidenote{\passref{6.7.31ab}}\edtext{\skt{चलः}}{\Afootnote{K H; \skt{चल}\ A}}\skt{ }\edtext{\skt{पद्मपलाशस्थः}}{\Afootnote{K H; \skt{त्पद्मपलाशस्थः}\ A}}\skt{ शुक्लो }\edtext{\skt{विन्दुरिवाम्भसः}}{\lemma{\skt{विन्दुर्}}\Afootnote{K H; \skt{बिन्दुर्}\ A}}\skt{ ।}\pend

\pstart\versepar \ledsidenote{\passref{6.7.32}}\edtext{\skt{मृद्यमाने च}}{\Afootnote{K; \skt{मृद्यमा}\ A H}}\skt{ }\edtext{\skt{नयने च}}{\Afootnote{K; \skt{ने च}\ A H}}\skt{ नयने मण्डलं }\edtext{\skt{तद्विसर्पति}}{\Afootnote{K H; \skt{तद्विसर्पति ।}\ A}}\skt{ । }\pend
\pstart\versepar \skt{प्रवालपद्मपत्राभं मण्डलं शोणितात्मकं ।}\pend

\pstart\versepar \ledsidenote{\passref{6.7.33}}\skt{दृष्टिरागो भवेच्चित्रो लिंगनाशे त्रिदोषजे । }\pend
\pstart\versepar \skt{यथास्वं दोषलिंगानि }\edtext{\skt{सर्वेष्वेव}}{\Afootnote{K H; \skt{सर्वेष्वेव}\ A}}\skt{ भवन्ति हि ॥}\pend

\pstart \ledsidenote{\passref{6.7.34cd}}\skt{तथा नरः पित्तविदग्धदृष्टिः कफेन }\edtext{\skt{चान्यस्त्वथ}}{\Afootnote{K H; \skt{चान्यस्त्वथ}\ A}}\skt{ }\edtext{\skt{धूमदर्शी}}{\Afootnote{A K; \skt{धूमदर्शा ।}\ H}}\skt{ ।}\pend

\pstart \ledsidenote{\passref{6.7.35ab}}\edtext{\skt{यो}}{\Afootnote{K H; \skt{यो ह्रस्वजाड्यो}\ A}}\skt{ ह्रस्वजात्यो नकुलान्धता च गम्भीरसंज्ञा च तथैव दृष्टिः ।}\pend

\pstart \ledsidenote{\passref{6.7.34ab}}\edtext{\skt{षड्लिङ्गनाशो}}{\Afootnote{K H; \skt{षड्लिङ्गनाशाः}\ A}}\skt{ }\edtext{\skt{विषमे च}}{\Afootnote{K H\textsuperscript{pc}; \skt{षडिमे च}\ A; \skt{विषमे}\ H\textsuperscript{ac}}}\skt{ रोगा दृष्ट्याश्रयाः }\edtext{\skt{षडेव}}{\Afootnote{K; \skt{षट्च षडेव}\ A; \skt{षड्च षडेव}\ H}}\skt{ }\edtext{\skt{चैव}}{\Afootnote{K; \skt{च स्युः ।}\ A; \skt{यौ यि।}\ H}}\skt{ ॥}\pend

\pstart \ledsidenote{\passref{6.7.35cd}}\skt{पित्तेन दुष्टेन }\edtext{\skt{युता तु}}{\Afootnote{K H; \skt{गतेन}\ A}}\skt{ }\edtext{\skt{दृष्टिः}}{\Afootnote{K H; \skt{दृष्टिं}\ A}}\skt{ पीता }\edtext{\skt{भवेद्यस्य}}{\Afootnote{K H; \skt{भवेद्यस्य}\ A}}\skt{ नरस्य }\edtext{\skt{किञ्चित्}}{\Afootnote{K H; \skt{दृष्टिः ॥}\ A}}\skt{ ।}\pend

\pstart\versepar \ledsidenote{\passref{6.7.36}}\skt{पीतानि रूपाणि च }\edtext{\skt{मन्यते यः}}{\Afootnote{A K; \skt{मन्यते तु}\ H}}\skt{ }\edtext{\skt{सवै}}{\Afootnote{K H; \skt{स मा}\ A}}\skt{ }\edtext{\skt{ररः}}{\Afootnote{K H; \skt{नवः}\ A}}\skt{ पित्तविदग्धदृष्टिः ॥ }\pend
\pstart\versepar \skt{प्राप्ते }\edtext{\skt{तृतीयं}}{\Afootnote{A K; \skt{तृतीये}\ H}}\skt{ पटलन्तु दोषे दिवा न पश्येन्निशि }\edtext{\skt{वेक्षते}}{\Afootnote{K H; \skt{वीक्षते}\ A}}\skt{ }\edtext{\skt{सः}}{\Afootnote{K H; \skt{च ॥}\ A}}\skt{ ॥}\pend

\pstart \ledsidenote{\passref{6.7.37}}\edtext{\skt{तथा}}{\Afootnote{K H; \skt{(रात्रौ स शीतानुगृहीतदृष्टिः पित्ताल्पभावादपि तानि पश्येत्) । तथा}\ A}}\skt{ नरः }\edtext{\edtext{\edtext{\skt{श्लेष्मविदग्धदृष्टिस्तान्येव}}{\lemma{\skt{श्लेष्मविदग्धदृष्टिस्}}\Afootnote{K H; \skt{श्लेष्मविदग्धदृष्टि}\ A}}}{\lemma{\skt{तान्य्}}\Afootnote{K; \skt{स्ता}\ A; \skt{ता}\ H}}}{\lemma{\skt{एव}}\Afootnote{K; \skt{न्येव}\ A H}}\skt{ शुक्लानि }\edtext{\skt{हमन्यते}}{\Afootnote{K H; \skt{हि मन्यते}\ A}}\skt{ तु ।}\pend

\pstart\versepar \ledsidenote{\passref{6.7.38}}\skt{त्रिषु }\edtext{\skt{स्थितो यः}}{\Afootnote{K H; \skt{स्थितोल्पः}\ A}}\skt{ पटलेषु दोषो }\edtext{\edtext{\skt{नक्तान्ध्यमापादयति}}{\lemma{\skt{नक्तान्ध्यम्}}\Afootnote{K H; \skt{नक्तान्ध्य}\ A}}}{\lemma{\skt{आपादयति}}\Afootnote{K H; \skt{मापादयति}\ A}}\skt{ }\edtext{\skt{प्रसह्यं}}{\Afootnote{K H; \skt{प्रसह्य ।}\ A}}\skt{ । }\pend
\pstart\versepar \skt{दिवा स }\edtext{\skt{सूर्यानुगृहीतदृष्टिः}}{\Afootnote{K H; \skt{सूर्यानुगृहीतचक्षु}\ A}}\skt{ }\edtext{\skt{पश्येत्तु}}{\Afootnote{K H; \skt{रीक्षेत}\ A}}\skt{ रूपाणि कफाल्पभावात् ॥}\pend

\pstart\versepar \ledsidenote{\passref{6.7.39}}\edtext{\skt{शोकज्वरायासशिरोभितापैरभ्याहता}}{\Afootnote{A K; \skt{शोकज्वरायासशिरोऽभितापैरभ्याहता}\ H}}\skt{ यस्य नरस्य दृष्टिः ॥ }\pend
\pstart\versepar \skt{}\edtext{\skt{धूमार्न्स}}{\Afootnote{K H; \skt{सधूम}\ A}}\skt{ }\edtext{\skt{वैपश्यति}}{\Afootnote{K H; \skt{कान्पश्यति}\ A}}\skt{ सर्व्वभावांस्तन्धूमदर्शीति वदन्ति रोगं ॥}\pend

\pstart\versepar \ledsidenote{\passref{6.7.40}}\edtext{\skt{यो ह्रस्वजात्यो}}{\Afootnote{K H; \skt{स ह्रस्वजाड्यो}\ A}}\skt{ दिवसेषु }\edtext{\skt{कृच्छ्राध्ह्रस्वाणि}}{\Afootnote{K; \skt{कृच्छ्राद्ध्रस्वानि}\ A; \skt{कृच्छ्राध्ह्रस्वानि}\ H}}\skt{ रूपाणि च }\edtext{\skt{तेन}}{\Afootnote{K H; \skt{येन}\ A}}\skt{ पश्येत् । }\pend
\pstart\versepar \skt{विद्योतते }\edtext{\skt{यस्य}}{\Afootnote{K H; \skt{येन}\ A}}\skt{ नरस्य दृष्टिर्दोषाभिपन्ना नकुलस्य यद्वत् ।}\pend

\pstart\versepar \ledsidenote{\passref{6.7.41}}\skt{चित्राणि रूपाणि }\edtext{\skt{दिवासु}}{\Afootnote{K; \skt{दिवा स}\ A H}}\skt{ पश्येत्स वै विकारो नकुलान्ध्यसंज्ञः ॥ }\pend
\pstart\versepar \skt{}\edtext{\skt{दृष्टिर्विरूपा}}{\Afootnote{A K; \skt{दृष्टिर्व्विरूपाश्}\ H}}\skt{ }\edtext{\skt{श्वसनोपसृष्टा}}{\Afootnote{A K; \skt{च समापसृष्टा}\ H}}\skt{ }\edtext{\skt{संकुच्यतेभ्यन्तरतश्च}}{\Afootnote{K; \skt{सङ्कुच्यतेऽभ्यन्तरतश्च}\ A H}}\skt{ }\edtext{\skt{यानि}}{\Afootnote{K H; \skt{याति ॥}\ A}}\skt{ ।}\pend

\pstart\versepar \ledsidenote{\passref{6.7.42}}\edtext{\skt{रुजवगाढञ्च }\edtext{\skt{तमक्षिरोगं}}{\lemma{\skt{अक्षिरोगं}}\Afootnote{K; \skt{मक्षिरोगं}\ A H}}}{\lemma{\skt{रुजवगाढञ्च तम्}}\Afootnote{K; \skt{रुजावगाढा च त}\ A; \skt{रुअवगाढञ्च त}\ H}}\skt{ गम्भीरिकेति प्रवदन्ति तज्ज्ञाः ॥ }\pend
\pstart\versepar \skt{}\edtext{\skt{वाहौ}}{\Afootnote{K H; \skt{बाह्यौ}\ A}}\skt{ }\edtext{\skt{पुनर्वा}}{\Afootnote{K; \skt{पुनर्द्वा}\ A; \skt{पुनर्वा}\ H}}\skt{ }\edtext{\skt{विहयौ}}{\Afootnote{K H; \skt{विह सं}\ A}}\skt{ प्रदिष्टौ निमित्ततश्चाप्यनिमित्ततश्च ।}\pend

\pstart \ledsidenote{\passref{6.7.43}}\edtext{\edtext{\skt{निमित्तजास्तत्र}}{\lemma{\skt{निमित्तजास्}}\Afootnote{K; \skt{निमित्ततस्}\ A; \skt{निमित्तजा}\ H}}}{\lemma{\skt{तत्र}}\Afootnote{A K; \skt{स्तत्र}\ H}}\skt{ }\edtext{\edtext{\skt{मनोभिताज्ज्ञेयस्त्वभिष्यन्तनिदर्शनम्सः}}{\lemma{\skt{मनोभिताज्ज्ञेयस्}}\Afootnote{K; \skt{शिरोभितापाज्ज्ञेय}\ A; \skt{मनोभिताज्ज्ञेयस्}\ H}}}{\lemma{\skt{त्वभिष्यन्तनिदर्शनम्सः}}\Afootnote{K; \skt{स्त्वभिष्यन्दनिदर्शनैश्च ।}\ A; \skt{त्वभिष्यन्तनिदर्शनः सः ॥}\ H}}\skt{ ॥ सुरर्षिगन्धर्वमहोरगाणां सन्दर्शनेनापि च }\edtext{\skt{भास्करस्य}}{\Afootnote{K H; \skt{भास्वराणाम् ॥}\ A}}\skt{ ॥}\pend

\pstart\versepar \ledsidenote{\passref{6.7.44}}\edtext{\skt{हन्येत}}{\Afootnote{A K; \skt{हन्येन}\ H}}\skt{ दृष्टिर्मनुजस्य यस्य सलिङ्गनाशस्त्वनिमित्तसंज्ञः । }\pend
\pstart\versepar \skt{तत्राक्षि }\edtext{\skt{विस्पष्टं}}{\Afootnote{K; \skt{विस्पष्ट}\ A; \skt{विष्टस्य}\ H}}\skt{ }\edtext{\skt{मिवाभि}}{\Afootnote{K; \skt{मिवाव}\ A; \skt{त्सिभि}\ H}}\skt{ भाति वैदूर्यवर्णा विमला च }\edtext{\skt{दृष्टिः}}{\Afootnote{A K; \skt{दृष्टि ॥}\ H}}\skt{ ॥}\pend

\pstart\versepar \ledsidenote{\passref{6.7.46}}\skt{इत्येते नयनगता }\edtext{\skt{मये}}{\Afootnote{K H; \skt{मया}\ A}}\skt{ }\edtext{\skt{हरोगानिर्दिष्टाः}}{\Afootnote{K H; \skt{विकाराः संख्याताः}\ A}}\skt{ पृथगिह }\edtext{\skt{षट्वा}}{\Afootnote{K; \skt{षट्च}\ A H}}\skt{ सप्ततिश्च । }\pend
\pstart\versepar \skt{}\edtext{\skt{तेषां ये}}{\Afootnote{K; \skt{एतेषां}\ A; \skt{तेषां}\ H}}\skt{ }\edtext{\skt{पृथगपृथक्व}}{\Afootnote{K H; \skt{पृथगिह विस्त}\ A}}\skt{ }\edtext{\skt{स्पष्टार्थशृणुत}}{\Afootnote{K; \skt{रेण सर्वं वक्ष्येऽहं तदनु}\ A; \skt{लिष्टं स्पष्टार्थशृणु तं}\ H}}\skt{ चिकित्सितं }\edtext{\skt{यथावदिति ॥ {\symfont ❈}}}{\Afootnote{K; \skt{यथावत् ॥}\ A; \skt{यथावदिति ॥ ० ॥}\ H}}\skt{ ॥}\pend

\pstart \ledsidenote{\passref{6.7.trailer}}\edtext{\skt{शालाक्ये}}{\Afootnote{K; \skt{इति सुश्रुतसंहितायामुत्तरतन्त्रान्तर्गते शालाक्यतन्त्रे दृष्टिगतरोगविज्ञानीयो नाम सप्तमोऽध्यायः ॥७ ॥}\ A; \skt{शालाक्ये स्थाने सप्तमोऽध्यायः ॥ ॥}\ H}}\skt{ ७ ॥}\pend

\endnumbering
\beginnumbering
\eledsection*{Adhyāya 8}
\pstart \ledsidenote{\passref{6.8.1}}\edtext{\skt{अथातश्चिकित्साप्रविभागी}}{\lemma{\skt{चिकित्साप्रविभागी}}\Afootnote{\skt{चिकित्सितप्रविभागविज्ञानीयमध्यायं}\ A; \skt{चिकित्साप्रविभागी व्या व्याख्यास्यामः ।}\ K; \skt{चिकित्साप्रविभागविज्ञानीयं}\ H}}\skt{ }\edtext{\skt{व्याख्यास्यामः}}{\Afootnote{A H; \emph{om.}\ K}}\skt{ ।}\pend

\pstart\versepar \ledsidenote{\passref{6.8.3}}\edtext{\edtext{\skt{षट्सप्ततिर्येभिहिता}}{\lemma{\skt{षट्सप्ततिर्}}\Afootnote{K H; \skt{षट्सप्तति}\ A}}}{\lemma{\skt{येभिहिता}}\Afootnote{K; \skt{र्येऽभिहिता}\ A; \skt{येऽभिहिता}\ H}}\skt{ व्याधयो नामलक्षणैः । }\pend
\pstart\versepar \skt{}\edtext{\edtext{\edtext{\skt{तेषाञ्चिकिर्सितमिदं}}{\lemma{\skt{तेषाञ्}}\Afootnote{K H; \emph{om.}\ A}}}{\lemma{\skt{चिकिर्सितम्}}\Afootnote{K H; \skt{चिकित्सित}\ A}}}{\lemma{\skt{इदं}}\Afootnote{K H; \skt{मिदं}\ A}}\skt{ }\edtext{\edtext{\skt{समासाद्यासतः}}{\lemma{\skt{समासाद्यास्}}\Afootnote{K H; \skt{तेषां समासव्या}\ A}}}{\lemma{\skt{अतः}}\Afootnote{K H; \skt{सतः}\ A}}\skt{ }\edtext{\skt{शृणुः}}{\Afootnote{K; \skt{शृणु ॥}\ A H}}\skt{ ॥}\pend

\pstart\versepar \ledsidenote{\passref{6.8.4}}\edtext{\edtext{\skt{छेद्यास्तेषु}}{\lemma{\skt{छेद्यास्}}\Afootnote{K H; \skt{छेद्या}\ A}}}{\lemma{\skt{तेषु}}\Afootnote{K H; \skt{स्तेषु}\ A}}\skt{ }\edtext{\skt{दशैकञ्च}}{\Afootnote{K; \skt{दशैकश्च}\ A; \skt{दशलैश्च}\ H}}\skt{ नव }\edtext{\skt{लेख्या}}{\Afootnote{K H; \skt{लेख्याः}\ A}}\skt{ }\edtext{\skt{भवन्ति हि}}{\Afootnote{K H; \skt{प्रकीर्तिताः ।}\ A}}\skt{ ॥ }\pend
\pstart\versepar \skt{भेद्याः }\edtext{\skt{पञ्च तु}}{\Afootnote{K; \skt{पञ्च वि}\ A; \skt{पञ्चसु}\ H}}\skt{ }\edtext{\edtext{\skt{रोगाम्स्युर्व्यध्याः}}{\lemma{\skt{रोगाम्स्युर्}}\Afootnote{K; \skt{काराः स्यु}\ A; \skt{रोगाः स्युर्}\ H}}}{\lemma{\skt{व्यध्याः}}\Afootnote{K H; \skt{र्व्यध्याः}\ A}}\skt{ पञ्चदशैव तु ।}\pend

\pstart\versepar \ledsidenote{\passref{6.8.5}}\skt{द्वादशाशस्त्रकृत्याश्च }\edtext{\skt{यप्याः}}{\Afootnote{K H; \skt{याप्याः}\ A}}\skt{ सप्त }\edtext{\skt{भवन्ति च}}{\Afootnote{K H; \skt{भवन्ति हि ।}\ A}}\skt{ । }\pend
\pstart\versepar \skt{रोगा }\edtext{\skt{वर्जयितव्याः स्युर्दशपञ्च}}{\lemma{\skt{वर्जयितव्याः स्युर्}}\Afootnote{K; \skt{वर्जयितव्याः स्यु}\ A H}}\skt{ च जानता। }\edtext{\skt{असाधौ च}}{\Afootnote{K; \skt{असाध्यौ वा}\ A; \skt{असाधौ वा}\ H}}\skt{ भवेतां तु }\edtext{\skt{याप्यो}}{\Afootnote{K H; \skt{याप्यौ}\ A}}\skt{ }\edtext{\skt{वागन्तु}}{\Afootnote{K H; \skt{चागन्तु}\ A}}\skt{ संज्ञितौ ॥}\edtext{}{\lemma{\skt{दशपञ्च च}}\Afootnote{K; \skt{र्दश पञ्च च}\ A H}}\pend

\pstart\versepar \ledsidenote{\passref{6.8.6}}\edtext{\skt{अर्शान्वितं}}{\Afootnote{K H; \skt{अर्शोऽन्वितं}\ A}}\skt{ भवति वर्त्म तु }\edtext{\skt{यस्य}}{\Afootnote{K H; \skt{यत्त}\ A}}\skt{ }\edtext{\skt{चार्शः}}{\Afootnote{K H; \skt{थार्शः}\ A}}\skt{ । }\edtext{\edtext{\skt{शुक्लन्तथार्वुदमथो}}{\lemma{\skt{शुक्लन्}}\Afootnote{K H; \skt{शुष्कं}\ A}}}{\lemma{\skt{तथार्वुदमथो}}\Afootnote{K H; \skt{तथार्बुदमथो}\ A}}\skt{ पिडकाः सिराजाः । }\pend
\pstart\versepar \skt{जालं }\edtext{\skt{सिराजमथ}}{\lemma{\skt{अथ}}\Afootnote{K H; \skt{अपि}\ A}}\skt{ पञ्चविधं तथार्म छेद्या भवन्ति सह पर्वणि कामयेन ॥}\pend

\pstart\versepar \ledsidenote{\passref{6.8.7}}\skt{उत्सङ्गिनी }\edtext{\skt{वहलकर्दमवर्त्मनी च}}{\Afootnote{K H; \skt{बहलकर्दमवर्त्मनी च}\ A}}\skt{ श्यावञ्च यच्च पठितं त्विह }\edtext{\skt{वद्धवर्त्म}}{\Afootnote{K H; \skt{बद्धवर्त्म ।}\ A}}\skt{ । }\pend
\pstart\versepar \skt{क्लिष्टञ्च पोथकियुतं खलु }\edtext{\skt{यत्तु}}{\Afootnote{K H; \skt{यच्च}\ A}}\skt{ वर्त्म कुम्भीकिनी च सहशर्करया च }\edtext{\skt{लेख्याः}}{\Afootnote{A K; \skt{लेख्या ॥}\ H}}\skt{ ॥}\pend

\pstart\versepar \ledsidenote{\passref{6.8.8}}\edtext{\skt{श्लेष्मोपनाहलगणौ तु}}{\Afootnote{K H; \skt{श्लेष्मोपनाहलगणौ च}\ A}}\skt{ }\edtext{\skt{विसञ्च}}{\Afootnote{K H; \skt{बिसं च}\ A}}\skt{ }\edtext{\skt{नेद्या}}{\Afootnote{K; \skt{भेद्या}\ A H}}\skt{ ग्रन्थिश्च यः }\edtext{\skt{क्रिमिकृतो}}{\Afootnote{K; \skt{कृमिकृतो}\ A; \skt{क्रिमिकृतोञ्}\ H}}\skt{ }\edtext{\skt{ञ्जननामिका च}}{\Afootnote{A K; \skt{जननामिका च ।}\ H}}\skt{ ॥ }\pend
\pstart\versepar \skt{आदौ सिरा निगदितास्तु }\edtext{\skt{प्रयोगे}}{\Afootnote{K; \skt{ययोः प्रयोगे}\ A; \skt{सिराः प्रयोगे}\ H}}\skt{ }\edtext{\skt{पाकौ}}{\Afootnote{A K; \skt{पाकोऔ}\ H}}\skt{ }\edtext{\skt{चयो}}{\Afootnote{K H; \skt{च यौ}\ A}}\skt{ }\edtext{\skt{नयनः}}{\Afootnote{K; \skt{नयनयोः}\ A H}}\skt{ }\edtext{\skt{यवनोन्यतश्च}}{\lemma{\skt{यवनोन्यतश्}}\Afootnote{K; \skt{पवनोऽन्यतश्}\ A; \skt{ऽन्यतश्}\ H}}\skt{ ॥}\pend

\pstart\versepar \ledsidenote{\passref{6.8.9}}\skt{पूयालसानि }\edtext{\skt{मन्थसंज्ञा}}{\Afootnote{K; \skt{लविपर्ययमन्थसंज्ञाः}\ A; \skt{लविपर्ययमन्थसंज्ञा}\ H}}\skt{ }\edtext{\skt{स्यन्दाश्च}}{\Afootnote{K H; \skt{स्यन्दास्तु}\ A}}\skt{ }\edtext{\edtext{\skt{यान्त्युपशमन्ति}}{\lemma{\skt{यान्त्य्}}\Afootnote{K; \skt{यन्}\ A; \skt{यान्}\ H}}}{\lemma{\skt{उपशमन्ति}}\Afootnote{K; \skt{त्युपशमं हि}\ A; \skt{त्युपशमन्ति}\ H}}\skt{ }\edtext{\skt{सिराव्यधैस्ते}}{\Afootnote{K H; \skt{सिराव्यधेन ।}\ A}}\skt{ । }\pend
\pstart\versepar \skt{}\edtext{\skt{शुष्काक्षिपाककफपित्तविदग्धदृष्टिष्वम्लाख्यशुक्रसहितार्जुनपिष्टकेन।}}{\lemma{\skt{शुष्काक्षिपाककफपित्तविदग्धदृष्टिष्वम्लाख्यशुक्रसहितार्जुनपिष्टकेन}}\Afootnote{K H; \skt{शुष्काक्षिपाककफपित्तविदग्धदृष्टिष्वम्लाख्यशुक्रसहितार्जुनपिष्टकेषु ॥}\ A}}\pend

\pstart\versepar \ledsidenote{\passref{6.8.10}}\skt{अक्लिन्नवर्त्महुतभुग्ध्वजदर्शिशुक्ति प्रक्लिन्नवर्त्मसु तथैव }\edtext{\skt{वलाससंज्ञाः}}{\Afootnote{K H; \skt{बलाससंज्ञे ।}\ A}}\skt{ । }\pend
\pstart\versepar \skt{आगन्तुनाम य }\edtext{\skt{य्तेन च}}{\Afootnote{K; \skt{युगेन च}\ A; \skt{मुखेनच}\ H}}\skt{ दूषितायां }\edtext{\skt{दृष्ट्यां न}}{\Afootnote{K H; \skt{दृष्टौ न}\ A}}\skt{ शस्त्रपतनं प्रवदन्ति }\edtext{\skt{वैद्याः}}{\Afootnote{K H; \skt{तज्ज्ञाः ॥}\ A}}\skt{ ॥}\pend

\pstart\versepar \ledsidenote{\passref{6.8.11}}\edtext{\skt{सम्प्श्यतः}}{\Afootnote{K H; \skt{संपश्यतः}\ A}}\skt{ षडपि }\edtext{\skt{भितास्तु}}{\Afootnote{K; \skt{येऽभिहितास्तु}\ A; \skt{षेभितास्तु}\ H}}\skt{ }\edtext{\skt{कावास्ते}}{\Afootnote{K H; \skt{काचास्ते}\ A}}\skt{ }\edtext{\skt{पक्ष्मकोपसयितास्तु}}{\Afootnote{K H; \skt{पक्ष्मकोपसहितास्तु}\ A}}\skt{ }\edtext{\skt{भअदन्ति}}{\Afootnote{K; \skt{भवन्ति}\ A; \skt{वदन्ति}\ H}}\skt{ याप्याः ॥ }\pend
\pstart\versepar \skt{चत्वार एव }\edtext{\edtext{\skt{पवनप्रभवास्त्वसाध्या}}{\lemma{\skt{पवनप्रभवास्}}\Afootnote{K H; \skt{पवनप्रभवा}\ A}}\skt{ द्वौ}}{\lemma{\skt{त्वसाध्या द्वौ}}\Afootnote{K H; \skt{स्त्वसाध्या द्वौ}\ A}}\skt{ पित्तजौ कफनिमित्तज एक एव । }\edtext{\skt{चत्वार}}{\Afootnote{K H; \skt{अष्टार्ध}\ A}}\skt{ }\edtext{\skt{एव}}{\Afootnote{K H; \skt{का}\ A}}\skt{ }\edtext{\edtext{\edtext{\skt{रूधिरप्रभवोस्त्रिदोषास्तावन्त}}{\lemma{\skt{रूधिरप्रभवोस्}}\Afootnote{K H; \skt{रुधिरजाश्च गदा}\ A}}}{\lemma{\skt{त्रिदोषास्}}\Afootnote{K H; \skt{स्त्रिदोषा}\ A}}}{\lemma{\skt{तावन्त}}\Afootnote{K H; \skt{स्तावन्त}\ A}}\skt{ एव }\edtext{\skt{सुरुणौ च}}{\Afootnote{K H; \skt{गदिताव}\ A}}\skt{ }\edtext{\skt{तथापरौ }\edtext{\skt{द्वाविति}}{\lemma{\skt{इति}}\Afootnote{K; \emph{om.}\ A; \skt{विति ॥ ॥}\ H}}}{\lemma{\skt{तथापरौ द्वाव्}}\Afootnote{K; \skt{पि बाह्यजौ द्वौ ॥}\ A; \skt{तथापरौ द्वा}\ H}}\skt{ ॥}\pend

\pstart \ledsidenote{\passref{6.8.trailer}}\edtext{\skt{शालाक्यस्थाने}}{\Afootnote{K; \skt{इति सुश्रुतसंहितायामुत्तरतन्त्रान्तर्गते शालाक्यतन्त्रे चिकित्सितप्रविभागविज्ञानीयो नामाष्टमोऽध्यायः ॥८ ॥}\ A; \skt{शालाक्यस्थानेऽष्टमोऽध्यायः ॥ ॥}\ H}}\skt{ ॥ ० ॥}\pend

\endnumbering
\beginnumbering
\eledsection*{Adhyāya 9}
\pstart \ledsidenote{\passref{6.9.1}}\skt{अथातो वाताभिष्यन्दप्रतिषेधं व्याख्यास्यामः ॥}\pend

\pstart\versepar \ledsidenote{\passref{6.9.3}}\edtext{\skt{पुराणं}}{\Afootnote{K; \skt{पुराण}\ A H}}\skt{ सर्प्पिषास्निग्धौ स्यन्दाधीमन्थपीडितौ । }\pend
\pstart\versepar \skt{स्वेदयित्वा यथान्यायं }\edtext{\edtext{\skt{सिराव्यधमुपाचरेत्}}{\lemma{\skt{सिराव्यधम्}}\Afootnote{K H; \skt{सिरामोक्षे}\ A}}}{\lemma{\skt{उपाचरेत्}}\Afootnote{K H; \skt{ण योजयेत् ॥}\ A}}\skt{ ।}\pend

\pstart\versepar \ledsidenote{\passref{6.9.4}}\edtext{\skt{सम्पादयेद्बस्तिभिश्च}}{\Afootnote{\skt{संपादयेद्बस्तिभिस्तु}\ A; \skt{सम्पादयेद्वस्तिभिश्च}\ K H}}\skt{ }\edtext{\skt{सम्यक्स्नेहविरचितौ}}{\Afootnote{K H; \skt{सम्यक्स्नेहविरेचितौ ।}\ A}}\skt{ । }\pend
\pstart\versepar \skt{}\edtext{\skt{तर्पणं}}{\Afootnote{K H; \skt{तर्पणैः}\ A}}\skt{ पुटपाकैश्च }\edtext{\edtext{\skt{धूमैराश्च्योतनैस्तथा}}{\lemma{\skt{धूमैर्}}\Afootnote{K H; \skt{धूमै}\ A}}}{\lemma{\skt{आश्च्योतनैस्}}\Afootnote{K H; \skt{राश्च्योतनैस्}\ A}}\skt{ ।}\pend

\pstart\versepar \ledsidenote{\passref{6.9.5}}\edtext{\skt{नस्यैः}}{\Afootnote{K H; \skt{नस्य}\ A}}\skt{ स्नेहपरीषेकैः }\edtext{\skt{शिरोबस्तिभिरेव}}{\lemma{\skt{शिरोबस्तिभिर्}}\Afootnote{A; \skt{शिरोवस्तिभिर्}\ K H}}\skt{ च । }\pend
\pstart\versepar \skt{}\edtext{\skt{वातघ्नानूपजलजां}}{\Afootnote{K H; \skt{वातघ्ननूपजलज}\ A}}\skt{ }\edtext{\skt{मांसाम्लंक्वाथसेवनैः}}{\Afootnote{K H; \skt{मांसाम्लक्वाथसेचनैः ॥}\ A}}\skt{ ।}\pend

\pstart\versepar \ledsidenote{\passref{6.9.6}}\edtext{\skt{चतुः}}{\Afootnote{K H; \emph{om.}\ A}}\skt{ स्नेहैः }\edtext{\skt{सुखोष्णैश्च}}{\Afootnote{K H; \skt{चतुर्भिरुष्णैश्च}\ A}}\skt{ }\edtext{\skt{तत्पीताम्वरतापनैः}}{\lemma{\skt{वरतापनैः}}\Afootnote{K H; \skt{बरधारणैः ।}\ A}}\skt{ । }\pend
\pstart\versepar \skt{}\edtext{\skt{पयोभिर्वेशवारैश्च}}{\lemma{\skt{वेशवारैश्च}}\Afootnote{K H; \skt{वेसवारैश्च}\ A}}\skt{ }\edtext{\skt{पायसैः}}{\Afootnote{K; \skt{साल्वणैः}\ A; \skt{पायः}\ H}}\skt{ }\edtext{\skt{साल्वणैस्तथा}}{\lemma{\skt{साल्वणैस्}}\Afootnote{K; \skt{पायसैस्}\ A; \skt{शात्वलेस्}\ H}}\skt{ ।}\pend

\pstart\versepar \ledsidenote{\passref{6.9.7}}\edtext{\edtext{\skt{भिषक्संपादयेदेतामुपनाहेश्च}}{\lemma{\skt{भिषक्संपादयेदेताम्}}\Afootnote{K H; \skt{भिषक्संपादयेदेता}\ A}}}{\lemma{\skt{उपनाहेश्च}}\Afootnote{K H; \skt{वुपनाहैश्च}\ A}}\skt{ }\edtext{\skt{पूजितैः}}{\Afootnote{A K; \skt{पूजितः ।}\ H}}\skt{ । }\pend
\pstart\versepar \skt{}\edtext{\skt{ग्राम्यानूपोदकैः}}{\Afootnote{K; \skt{ग्राम्यानूपौदकरसैः}\ A; \skt{ग्राम्यानूपोदलैः}\ H}}\skt{ स्निग्धैः }\edtext{\skt{रसैः}}{\Afootnote{K H; \emph{om.}\ A}}\skt{ }\edtext{\skt{फलरसायुतैः}}{\Afootnote{K H; \skt{फलरसान्वितैः ॥}\ A}}\skt{ ।}\pend

\pstart\versepar \ledsidenote{\passref{6.9.8}}\skt{सुसंस्कृतैः पयोभिश्च }\edtext{\skt{तयोर्रहार}}{\lemma{\skt{रहार}}\Afootnote{K; \skt{आहार}\ A; \skt{आहारम्}\ H}}\skt{ इष्यते । }\pend
\pstart\versepar \skt{}\edtext{\skt{भुक्तोपरि}}{\Afootnote{K; \skt{तथा चोप}\ A; \skt{उक्तोप}\ H}}\skt{ }\edtext{\skt{पिवेयातां}}{\Afootnote{K; \skt{रि भ}\ A; \skt{रि यातां}\ H}}\skt{ }\edtext{\skt{घृतं तो च}}{\Afootnote{K; \skt{क्तस्य स}\ A; \skt{घृतं चा}\ H}}\skt{ }\edtext{\skt{पिपावनौ}}{\Afootnote{K H; \skt{र्पिःपानं(पाने) प्रशस्यते ॥}\ A}}\skt{ ।}\pend

\pstart \ledsidenote{\passref{6.9.11cd}}\skt{एरण्डपल्लवे मूले त्वचि वाथं पयः श्रितं ।}\pend

\pstart\versepar \ledsidenote{\passref{6.9.12}}\edtext{\skt{कण्ठकाराश्च}}{\Afootnote{K; \skt{कण्टकार्याश्च}\ A}}\skt{ }\edtext{\skt{मूलास्यां}}{\Afootnote{K; \skt{मूलेषु}\ A}}\skt{ }\edtext{\skt{सुखोष्टं}}{\Afootnote{K; \skt{सुखोष्णं}\ A}}\skt{ }\edtext{\skt{सेवनहितं}}{\Afootnote{K; \skt{सेचने हितम् ।}\ A}}\skt{ । }\pend
\pstart\versepar \skt{}\edtext{\edtext{\skt{सेन्धवोदीच्यमधुकपिप्पलीभिरयो}}{\lemma{\skt{सेन्धवोदीच्यमधुकपिप्पलीभिर्}}\Afootnote{K; \skt{सैन्धवोदीच्ययष्ट्याह्वपिप्पलीभिः}\ A; \skt{सैन्धवोदीच्यमधुकपिप्पलीभिर्}\ H}}\skt{ पि वा}}{\lemma{\skt{अयो पि वा}}\Afootnote{K; \skt{शृतं पयः ॥}\ A; \skt{अयो पि वा । हितमन्दोदकं सेके तथाश्च्योतनमेव च ॥}\ H}}\skt{ ।}\edtext{}{\lemma{\skt{… सेन्}}\Afootnote{\emph{om.}\ H}}\pend

\pstart\versepar \ledsidenote{\passref{6.9.13}}\skt{हितमर्धोदकं सेके }\edtext{\skt{तथाश्च्योतनमेव}}{\lemma{\skt{तथाश्च्योतनम्}}\Afootnote{K; \skt{तथाश्च्योतनम्}\ A}}\skt{ च । }\pend
\pstart\versepar \skt{}\edtext{\skt{ह्रीवेरचक्रशाङ्गष्ठोदुम्वरत्वक्षु}}{\Afootnote{K H\textsuperscript{pc}; \skt{ह्रीबेरवक्रमञ्जिष्ठोदुम्बरत्वक्षु}\ A; \skt{ह्रीवेरशाङ्गष्ठोदुम्वरत्वक्षु}\ H\textsuperscript{ac}}}\skt{ साधितं ।}\edtext{}{\lemma{\skt{… ह्री}}\Afootnote{\emph{om.}\ H\textsuperscript{ac}}}\edtext{}{\lemma{\skt{… ह्री}}\Afootnote{\emph{om.}\ H\textsuperscript{pc}}}\pend

\pstart\versepar \ledsidenote{\passref{6.9.14}}\edtext{\skt{साम्भसापयसाजेन}}{\Afootnote{K H; \skt{साम्भश्छागं पयो वापि}\ A}}\skt{ }\edtext{\skt{शूलाश्च्योतनमुत्तमं}}{\lemma{\skt{शूलाश्च्योतनम्}}\Afootnote{K\textsuperscript{pc} H; \skt{शूलाश्च्योतनम्}\ A; \skt{शूलाश्च्योतम्}\ K\textsuperscript{ac}}}\skt{ ॥ }\pend
\pstart\versepar \skt{मधुकं }\edtext{\skt{रजनीपथ्यादेवदारूश्च}}{\Afootnote{K H; \skt{रजनीं पथ्यां देवदारुं च}\ A}}\skt{ }\edtext{\skt{पीषयेत्}}{\Afootnote{K H; \skt{पेषयेत् ॥}\ A}}\skt{ ।}\pend

\pstart\versepar \ledsidenote{\passref{6.9.15}}\edtext{\skt{आजेनपयसा}}{\Afootnote{A K; \skt{अञ्जनं पयसा}\ H}}\skt{ }\edtext{\skt{श्रेष्ठमभिष्यने}}{\Afootnote{K H; \skt{श्रेष्ठमभिष्यन्दे}\ A}}\skt{ }\edtext{\skt{तदञ्जनं}}{\Afootnote{K H; \skt{तदञ्जनम् ।}\ A}}\skt{ ॥ }\pend
\pstart\versepar \skt{}\edtext{\skt{गैरिकः}}{\Afootnote{K H; \skt{गैरिकं}\ A}}\skt{ सैन्धवं }\edtext{\skt{कृष्णं}}{\Afootnote{K H; \skt{कृष्णां}\ A}}\skt{ }\edtext{\skt{नगरं च}}{\Afootnote{K H; \skt{नागरं च}\ A}}\skt{ यथोत्तरं ।}\pend

\pstart\versepar \ledsidenote{\passref{6.9.16}}\edtext{\skt{पिष्टं}}{\Afootnote{K H; \skt{द्विगुणं पिष्ट}\ A}}\skt{ }\edtext{\skt{द्विरंशतोभिर्वा}}{\Afootnote{K H; \skt{मद्भिस्तु}\ A}}\skt{ }\edtext{\edtext{\skt{काञ्जनमिष्यते}}{\lemma{\skt{काञ्जनम्}}\Afootnote{K; \skt{गुटिकाञ्जन}\ A; \skt{गुटिकाञ्जनम्}\ H}}}{\lemma{\skt{इष्यते}}\Afootnote{K H; \skt{मिष्यते ।}\ A}}\skt{ ॥ }\pend
\pstart\versepar \skt{}\edtext{\edtext{\skt{स्नेहाञ्जनञ्चात्र}}{\lemma{\skt{स्नेहाञ्जनञ्}}\Afootnote{A K; \skt{स्नेहाञ्जन}\ H}}}{\lemma{\skt{चात्र}}\Afootnote{K; \skt{हितं}\ A; \skt{त्र}\ H}}\skt{ }\edtext{\skt{हितं}}{\Afootnote{K H; \skt{चात्र}\ A}}\skt{ वक्ष्यते }\edtext{\skt{तद्यथाविधिः}}{\Afootnote{K H; \skt{तद्यथाविधि ॥}\ A}}\skt{ ।}\pend

\pstart\versepar \ledsidenote{\passref{6.9.17}}\edtext{\skt{रोगो यश्चान्यतोवातो}}{\lemma{\skt{रोगो यश्}}\Afootnote{A K; \skt{रोगायश्}\ H}}\skt{ यश्च }\edtext{\skt{मारूतपर्ययः}}{\Afootnote{K H; \skt{मारुतपर्ययः ।}\ A}}\skt{ । }\pend
\pstart\versepar \skt{अनेनैव विधानेन }\edtext{\skt{भिषक्तान्यपि}}{\Afootnote{K; \skt{भिषक्तावपि}\ A; \skt{भिषकानपि}\ H}}\skt{ साधयेत् ॥}\edtext{}{\lemma{\skt{चान्यतोवातो यश्च}}\Afootnote{A K; \skt{चान्यतावातो यश्च}\ H}}\pend

\pstart \ledsidenote{\passref{6.9.18}}\edtext{\skt{पूर्वभक्तहितं}}{\Afootnote{K; \skt{पूर्वभक्तं हितं}\ A; \skt{पूर्व्वभक्ते हितं}\ H}}\skt{ }\edtext{\skt{सर्पिः}}{\Afootnote{A K; \skt{सर्प्पि}\ H}}\skt{ }\edtext{\skt{क्षीरञ्चात्राथ}}{\Afootnote{K; \skt{क्षीरं वाप्यथ}\ A; \skt{क्षीरञ्जनाथ}\ H}}\skt{ }\edtext{\skt{भोजयेत् । वृक्षादन्यां}}{\Afootnote{K; \skt{भोजने । वृक्षादन्यां}\ A; \skt{भोजयेत् ॥ वृक्षादन्या}\ H}}\skt{ }\edtext{\skt{कपित्थेच}}{\Afootnote{A K; \skt{न्कपित्थेच}\ H}}\skt{ पञ्चमूले महत्यपि।}\pend

\pstart\versepar \ledsidenote{\passref{6.9.19}}\edtext{\skt{सक्षीरकर्कटरसे}}{\Afootnote{K; \skt{सक्षीरं कर्कटरसे}\ A; \skt{सक्षीरे कर्कटरसे}\ H\textsuperscript{pc}}}\skt{ }\edtext{\skt{सिद्धञ्चा}}{\Afootnote{K H\textsuperscript{pc}; \skt{सिद्धं चात्र}\ A}}\skt{ }\edtext{\skt{पिवेत्घृतं}}{\Afootnote{K H\textsuperscript{pc}; \skt{घृतं पिबेत् । सिद्धं}\ A}}\skt{ । }\pend
\pstart\versepar \skt{}\edtext{\skt{पानम्वा}}{\Afootnote{K H; \skt{वा}\ A}}\skt{ }\edtext{\edtext{\skt{हितमत्राहुः}}{\lemma{\skt{हितम्}}\Afootnote{K H; \skt{हित}\ A}}}{\lemma{\skt{अत्राहुः}}\Afootnote{K H; \skt{मत्राहुः}\ A}}\skt{ }\edtext{\skt{पक्वाश्चार्त्रगलाग्निकैः}}{\Afootnote{K; \skt{पत्तूरार्तगलाग्निकैः ॥}\ A; \skt{पक्वार्तमलाग्निकैः ॥}\ H}}\skt{ ।}\edtext{}{\lemma{\skt{… पा}}\Afootnote{\emph{om.}\ H\textsuperscript{ac}}}\pend

\pstart\versepar \ledsidenote{\passref{6.9.20}}\skt{सक्षीरं }\edtext{\skt{मेषशृंगाश्च}}{\Afootnote{K H; \skt{मेषशृङ्ग्या वा}\ A}}\skt{ }\edtext{\skt{सर्पिर्वीरतरेपि च}}{\Afootnote{K; \skt{सर्पिर्वीरतरेण वा ।}\ A; \skt{सर्प्पिर्वीरतरेऽपि च ।}\ H}}\skt{ । }\pend
\pstart\versepar \skt{सैन्धवं }\edtext{\skt{दारूशुण्ठी च}}{\Afootnote{K; \skt{दारु शुण्ठी च}\ A; \skt{दारूस्वर्ण्णञ्च}\ H}}\skt{ मातुलुङ्गरसो घृतं ।}\pend

\pstart\versepar \ledsidenote{\passref{6.9.21}}\edtext{\skt{स्तन्योदकावं}}{\Afootnote{K; \skt{स्तन्योदकाभ्यां}\ A H}}\skt{ कर्त्तव्यं शुष्कपाके }\edtext{\skt{तदञ्जनं}}{\Afootnote{K H; \skt{तदञ्जनम् ।}\ A}}\skt{ । }\pend
\pstart\versepar \skt{पूजितं सर्पिषश्चात्र पानमक्ष्णोश्च तर्पणं ।}\pend

\pstart\versepar \ledsidenote{\passref{6.9.22}}\skt{घृतेन जीवनीयेन नस्यं तैलेन }\edtext{\skt{चानुना}}{\Afootnote{K H; \skt{चाणुना ।}\ A}}\skt{ । }\pend
\pstart\versepar \skt{}\edtext{\skt{परिषेकोहितश्चात्र}}{\lemma{\skt{परिषेकोहितश्}}\Afootnote{K H; \skt{परिषेके हितं}\ A}}\skt{ पयः शीतं ससैन्धवं ।}\pend

\pstart\versepar \ledsidenote{\passref{6.9.23}}\edtext{\skt{रजनादारूसिद्धम्वा}}{\Afootnote{K; \skt{रजनीदारुसिद्धं वा}\ A; \skt{रजनीदारूसिद्धम्वा}\ H}}\skt{ }\edtext{\skt{सेन्धवेन}}{\Afootnote{K; \skt{सैन्धवेन}\ A; \skt{सनिन्धवेन}\ H}}\skt{ }\edtext{\skt{सामायुतं}}{\Afootnote{K; \skt{समायुतम्}\ A H}}\skt{ । }\pend
\pstart\versepar \skt{}\edtext{\skt{सर्पिर्युतम्स्तन्य}}{\lemma{\skt{सर्पिर्युतम्}}\Afootnote{A K; \skt{सर्प्पियुतं}\ H}}\skt{ }\edtext{\skt{धृष्टमञ्जने च}}{\Afootnote{K; \skt{घृष्टमञ्जनं वा}\ A; \skt{घृष्टमञ्जने च}\ H}}\skt{ महौषधं ।}\pend

\pstart\versepar \ledsidenote{\passref{6.9.24}}\skt{वसा }\edtext{\skt{वानूपजलजाः}}{\Afootnote{K H; \skt{वानूपजलजा}\ A}}\skt{ }\edtext{\skt{सैन्धवेन}}{\Afootnote{A K; \skt{सन्धवेन}\ H}}\skt{ समायुताः । }\pend
\pstart\versepar \skt{}\edtext{\skt{नागरोन्मिश्रिताः}}{\Afootnote{K H; \skt{नागरोन्मिश्रिता}\ A}}\skt{ किञ्चिच्छुष्कपाके }\edtext{\skt{तदञ्जनं}}{\Afootnote{K H; \skt{तदञ्जनम् ॥}\ A}}\skt{ ।}\pend

\pstart\versepar \ledsidenote{\passref{6.9.25}}\skt{पवनप्रभवा रोगा ये }\edtext{\skt{केचिद्दृष्टिनाशनाः}}{\Afootnote{K H; \skt{केचिद्दृष्टिनाशनाः ।}\ A}}\skt{ । }\pend
\pstart\versepar \skt{}\edtext{\skt{वीजेनानेन}}{\Afootnote{K; \skt{बीजेनानेन}\ A; \skt{जीवेनानेन}\ H}}\skt{ }\edtext{\skt{मतिमान्स्तेषु}}{\Afootnote{K; \skt{मतिमान्तेषु}\ A H}}\skt{ }\edtext{\skt{कर्मप्रयोजयेतिति}}{\Afootnote{K H; \skt{कर्म प्रयोजयेत् ॥}\ A}}\skt{ ॥}\pend

\pstart \ledsidenote{\passref{6.9.trailer}}\edtext{\skt{शालाक्ये}}{\Afootnote{K; \skt{इति सुश्रुतसंहितायामुत्त्रतन्त्रान्तर्गते शालाक्यतन्त्रे वाताभिष्यन्दप्रतिषेधो नाम नवमोऽध्यायः ॥९ ॥}\ A; \skt{शालाक्यस्थाने नवमोऽध्यायः ॥ ॥ ॥}\ H}}\skt{ ॥ ॥ ० ॥}\pend

\endnumbering
\beginnumbering
\eledsection*{Adhyāya 10}
\pstart \ledsidenote{\passref{6.10.1}}\skt{अथातः पित्ताभिष्यन्दप्रतिषेधं }\edtext{\skt{व्याख्यास्यामः}}{\Afootnote{A H; \skt{व्या व्याख्यास्यामः ।}\ K}}\skt{ ।}\pend

\pstart\versepar \ledsidenote{\passref{6.10.3}}\edtext{\skt{पित्तस्यान्दे}}{\Afootnote{K; \skt{पित्तस्यन्दे}\ A H}}\skt{ पैत्तिके }\edtext{\skt{चाधिमन्थे}}{\Afootnote{A K; \skt{चाधिमन्थ}\ H}}\skt{ }\edtext{\skt{सिराव्यधः}}{\Afootnote{K H; \skt{रक्तास्रावः}\ A}}\skt{ }\edtext{\skt{स्रंसनञ्चाध्युदारं}}{\lemma{\skt{चाध्युदारं}}\Afootnote{K; \skt{चापि कार्यम् ।}\ A; \skt{जादारं ।}\ H}}\skt{ । }\pend
\pstart\versepar \skt{}\edtext{\skt{अक्ष्णो}}{\Afootnote{K H; \skt{अक्ष्णोः}\ A}}\skt{ }\edtext{\skt{सृ}}{\Afootnote{K H; \skt{सेका}\ A}}\skt{ }\edtext{\skt{पुटपाकाञ्जनानि च}}{\Afootnote{K H; \skt{लेपनस्याञ्जनानि}\ A}}\skt{ क्रमः }\edtext{\skt{सर्वः }\edtext{\skt{पित्तवीसर्पहा च}}{\Afootnote{K; \skt{पित्तवीसर्प्पहाश्च ॥ ॥}\ H}}\skt{ ॥}}{\lemma{\skt{च… पित्तवीसर्प}}\Afootnote{\emph{om.}\ A}}\pend

\pstart \ledsidenote{\passref{6.10.4ab}}\edtext{\skt{गुन्द्रं}}{\Afootnote{K H; \skt{गुन्द्रां}\ A}}\skt{ शालिं शैवलं }\edtext{\edtext{\skt{दर्भमिक्षुर्नलं}}{\lemma{\skt{दर्भमिक्षुर्}}\Afootnote{K; \skt{शैलभेदं दार्वीमेला}\ A; \skt{दर्भमिक्षुन्}\ H}}}{\lemma{\skt{नलं}}\Afootnote{K; \skt{मुत्पलं}\ A; \skt{ननलं}\ H}}\skt{ }\edtext{\skt{लोध्रं}}{\Afootnote{K H; \skt{रोध्र}\ A}}\skt{ }\edtext{\edtext{\edtext{\skt{वेतसम्पद्मकञ्च}}{\lemma{\skt{वेतसम्}}\Afootnote{K H; \skt{मभ्रम् ।}\ A}}}{\lemma{\skt{पद्मकञ्}}\Afootnote{K H; \emph{om.}\ A}}}{\lemma{\skt{च}}\Afootnote{K H; \emph{om.}\ A}}\skt{ ।}\pend

\pstart \ledsidenote{\passref{6.10.5ab}}\edtext{\skt{द्राक्षाक्षौद्रं}}{\Afootnote{K; \skt{द्राक्षां क्षौद्रं}\ A; \skt{द्राक्षान्क्षौद्रञ्}\ H}}\skt{ }\edtext{\skt{चन्दनञ्चोत्पलञ्च}}{\lemma{\skt{चोत्पलञ्च}}\Afootnote{K; \skt{यष्टिकाह्वं योषित्क्षीरं रात्र्यनन्}\ A; \skt{चात्पलञ्च}\ H}}\skt{ }\edtext{\skt{स्तियाः }\edtext{\skt{स्तन्यरसमिक्षाः}}{\Afootnote{K; \skt{स्तन्यं रसमिक्षाः}\ H}}\skt{ सिता }\edtext{\skt{च}}{\Afootnote{K H; \skt{च पिष्ट्वा ।}\ A}}\skt{ ।}}{\lemma{\skt{च… सिता}}\Afootnote{\emph{om.}\ A}}\pend

\pstart \ledsidenote{\passref{6.10.4cd}}\edtext{\edtext{\skt{पद्मम्पत्रं}}{\lemma{\skt{पद्मम्}}\Afootnote{K; \skt{पद्मा}\ A; \skt{पत्रं}\ H}}}{\lemma{\skt{पत्रं}}\Afootnote{K; \skt{त्पत्रं}\ A; \skt{पद्मं}\ H}}\skt{ }\edtext{\skt{यष्ठि}}{\Afootnote{K H; \skt{शर्क}\ A}}\skt{ }\edtext{\skt{साकं}}{\Afootnote{K H; \skt{रा द}\ A}}\skt{ }\edtext{\skt{हरिद्रां}}{\Afootnote{K H; \skt{र्भमिक्षुं}\ A}}\skt{ }\edtext{\skt{तथानन्ता चापि }\edtext{\skt{सम्भृत्यसर्पिः}}{\Afootnote{K; \skt{सम्भृत्य सर्प्पिः सिद्धं सेकं तर्प्पणं चाञ्जने च ।}\ H}}\skt{ ।}}{\lemma{\skt{तथानन्ता… सम्भृत्यसर्पिः}}\Afootnote{\emph{om.}\ A}}\pend

\pstart \ledsidenote{\passref{6.10.5cd}}\edtext{\skt{सिद्धं}}{\Afootnote{K; \skt{सर्पिःसिद्धं}\ A; \skt{वातघ्नानूपजलजां मांसाम्लंक्वाथ}\ H}}\skt{ }\edtext{\skt{सेकतर्पणो}}{\Afootnote{K; \skt{तर्पणे}\ A; \skt{सेवनैः ॥}\ H}}\skt{ }\edtext{\skt{चाञ्जने च}}{\Afootnote{K; \skt{सेकनस्ये}\ A; \emph{om.}\ H}}\skt{ }\edtext{\skt{हितं}}{\Afootnote{K; \skt{शस्तं}\ A; \emph{om.}\ H}}\skt{ }\edtext{\skt{क्षीरं चा मे ते}}{\Afootnote{K; \skt{क्षीरं सिद्धमे}\ A; \emph{om.}\ H}}\skt{ }\edtext{\skt{धूमकै}}{\Afootnote{K; \skt{तेषु चाजम् ॥}\ A; \emph{om.}\ H}}\skt{ ।}\pend

\pstart\versepar \ledsidenote{\passref{6.10.6}}\skt{योज्यो वर्गो व्यस्त एषो }\edtext{\skt{द्विधा वा च}}{\Afootnote{K; \skt{ऽन्यथा वा स}\ A}}\skt{ सचापि नस्येथावत् । }\pend
\pstart\versepar \skt{क्रियाः }\edtext{\skt{सर्वाः}}{\Afootnote{A K; \skt{सर्व्वा}\ H}}\skt{ }\edtext{\skt{पित्तहराम्च}}{\Afootnote{K; \skt{पित्तहर्यः प्र}\ A; \skt{पित्तहराश्च}\ H}}\skt{ }\edtext{\skt{शस्ताः}}{\Afootnote{K H; \skt{शस्ता}\ A}}\skt{ }\edtext{\skt{त्र्यहादूर्ध्वं}}{\Afootnote{K H; \skt{स्त्र्यहाच्चोर्ध्वं}\ A}}\skt{ क्षीरसर्पिश्च नस्यं ।}\edtext{}{\lemma{\skt{… क्रि}}\Afootnote{\emph{om.}\ H}}\edtext{}{\lemma{\skt{सचापि नस्येथावत् । क्रियाः}}\Afootnote{K; \skt{म्यङ्नस्येऽष्टार्धसंख्येऽपि नित्यम् । क्रियाः}\ A}}\pend

\pstart\versepar \ledsidenote{\passref{6.10.7}}\skt{पालाशं स्याच्छोणितं }\edtext{\edtext{\skt{चाञ्जनार्थम्भवेत्कार्यं}}{\lemma{\skt{चाञ्जनार्थम्}}\Afootnote{K H; \skt{चाञ्जनार्थे}\ A}}}{\lemma{\skt{भवेत्कार्यं}}\Afootnote{K H; \skt{शल्लक्या वा}\ A}}\skt{ शर्करा क्षौद्रयुक्तं । }\pend
\pstart\versepar \skt{}\edtext{\skt{रसक्रियां}}{\Afootnote{A K; \skt{रसक्रिया}\ H}}\skt{ }\edtext{\skt{शर्कराक्षौद्रयुक्तं}}{\Afootnote{K; \skt{शर्कराक्षौद्रयुक्तां}\ A; \skt{शर्लराक्षौद्रयुक्तं}\ H}}\skt{ पालिन्द्याम्वा मधुके }\edtext{\skt{चापि}}{\Afootnote{K H; \skt{वापि}\ A}}\skt{ कुर्यात् ।}\pend

\pstart \ledsidenote{\passref{6.10.8}}\edtext{\skt{मुस्तं}}{\Afootnote{K H; \skt{मुस्ता}\ A}}\skt{ }\edtext{\skt{फेनः}}{\Afootnote{A K; \skt{फेनं}\ H}}\skt{ }\edtext{\skt{सागरं}}{\Afootnote{K H; \skt{सागर}\ A}}\skt{ }\edtext{\skt{सोत्पलञ्च}}{\Afootnote{K H; \skt{स्योत्पलं च}\ A}}\skt{ }\edtext{\skt{क्रिमिघ्नैलाधात्रिवीजाद्रसञ्च}}{\Afootnote{K H; \skt{कृमिघ्नैलाधात्रिबीजाद्रसश्च । तालीशैलागैरिकोशीरशङ्खैरेव ं युञ्ज्यादञ्जनं स्तन्यपिष्टैः ॥}\ A}}\skt{ ।}\pend

\pstart\versepar \ledsidenote{\passref{6.10.9}}\edtext{\skt{कार्यं}}{\Afootnote{K H; \emph{om.}\ A}}\skt{ }\edtext{\skt{चूर्णाञ्जनार्थं}}{\Afootnote{K H; \skt{चूर्णं कुर्यादञ्जनार्थे}\ A}}\skt{ रसो वा स्तन्योपेतो धातकीस्यन्दनाभ्यां । }\pend
\pstart\versepar \skt{}\edtext{\skt{योषित्स्तन्य}}{\Afootnote{K; \skt{योषित्स्तन्यं}\ A; \skt{योषित्स्तन्यो}\ H}}\skt{ शातकुम्भं विघृष्टं }\edtext{\skt{सिताक्षौद्रे}}{\Afootnote{K H; \skt{क्षौद्रो}\ A}}\skt{ }\edtext{\skt{किंशुकाद्वापि}}{\Afootnote{K H; \skt{पेतं कैशुकं चापि}\ A}}\skt{ }\edtext{\skt{पुष्पां}}{\Afootnote{K H; \skt{पुष्पम् ॥}\ A}}\skt{ ।}\pend

\pstart\versepar \ledsidenote{\passref{6.10.10}}\edtext{\skt{लोध्रं}}{\Afootnote{K H; \skt{रोध्रं}\ A}}\skt{ }\edtext{\skt{द्राक्षां}}{\Afootnote{A K; \skt{द्राक्षा}\ H}}\skt{ }\edtext{\skt{शर्कराञ्चोत्पलञ्स्त्रियाः}}{\Afootnote{K; \skt{शर्करामुत्पलं च नार्याः}\ A; \skt{शर्क्कराञ्त्पलञ्स्त्रियाः}\ H}}\skt{ }\edtext{\skt{स्तन्ये}}{\Afootnote{K H; \skt{क्षीरे}\ A}}\skt{ }\edtext{\skt{षष्टिसाकं}}{\Afootnote{K H; \skt{यष्टिकाह्वं}\ A}}\skt{ }\edtext{\skt{कृतञ्च}}{\Afootnote{K H; \skt{वचां च ।}\ A}}\skt{ । }\pend
\pstart\versepar \skt{पिष्ट्वा क्षीरे }\edtext{\skt{वर्णकस्य}}{\Afootnote{A K; \skt{वर्णलस्य}\ H}}\skt{ त्वचं च }\edtext{\edtext{\skt{वैजलोन्मिश्रञ्चन्दनोदूम्वरञ्च}}{\lemma{\skt{वैजलोन्मिश्रञ्}}\Afootnote{K H; \skt{तोयोन्मिश्रे}\ A}}}{\lemma{\skt{चन्दनोदूम्वरञ्}}\Afootnote{K H; \skt{चन्दनोदुम्बरे}\ A}}\skt{ ।}\pend

\pstart\versepar \ledsidenote{\passref{6.10.12}}\edtext{\edtext{\skt{क्षौमावर्धम्पथ्यमाश्च्योतने}}{\lemma{\skt{क्षौमावर्धम्}}\Afootnote{K H; \skt{क्षौमाबद्धं}\ A}}\skt{ वा}}{\lemma{\skt{पथ्यमाश्च्योतने वा}}\Afootnote{K; \skt{पथ्यमाश्च्योतने वा}\ A; \skt{पथ्यमाश्च्योतने वा}\ H}}\skt{ । }\edtext{\skt{घृतभृष्टं}}{\Afootnote{K H; \skt{सर्पिर्घृष्टं}\ A}}\skt{ }\edtext{\skt{यष्टिसाह्वं}}{\Afootnote{K H; \skt{यष्टिकाह्वं}\ A}}\skt{ }\edtext{\skt{सलोध्रं}}{\Afootnote{K H; \skt{सरोध्रम् ।}\ A}}\skt{ । }\pend
\pstart\versepar \skt{}\edtext{\skt{तोयोन्मिश्रं}}{\Afootnote{K; \skt{तोयोन्मिश्राः}\ A; \skt{तोयोमिश्रं}\ H}}\skt{ }\edtext{\skt{कार्शमरीधात्रिपथ्यास्तथैवाः}}{\Afootnote{K; \skt{काश्मरीधात्रिपथ्यास्तद्वच्चाहुः}\ A; \skt{शर्क्करीधात्रिपथ्यास्तथैवाः}\ H}}\skt{ }\edtext{\skt{कट्फलञ्चाम्वुनैव}}{\lemma{\skt{चाम्वुनैव}}\Afootnote{K H; \skt{चाम्बुनैव ॥}\ A}}\skt{ ।}\pend

\pstart \ledsidenote{\passref{6.10.13}}\edtext{\skt{एषा}}{\Afootnote{K H; \skt{एषो}\ A}}\skt{ }\edtext{\skt{क्रियाम्लाख्युर्षितसशुक्तौ}}{\Afootnote{K H; \skt{ऽम्लाख्येऽनुक्रमश्चापि शुक्तौ}\ A}}\skt{ }\edtext{\skt{कार्या}}{\Afootnote{K H; \skt{कार्यः}\ A}}\skt{ }\edtext{\skt{सर्वाव्यपवर्ज्यः}}{\Afootnote{K; \skt{सर्वः स्यात्सिरामोक्षवर्ज्यः ॥}\ A; \skt{सर्व्वः स्थाप्यमोक्षं}\ H}}\skt{ }\edtext{\skt{सिराणां}}{\Afootnote{K H; \emph{om.}\ A}}\skt{ ।}\pend

\pstart\versepar \ledsidenote{\passref{6.10.14}}\skt{सर्पिः पेयं त्रैफलं तैल्वकम्वा पेयं वा }\edtext{\skt{केवलं}}{\Afootnote{K H; \skt{स्यात्केवलं}\ A}}\skt{ यत्पुराणं । }\pend
\pstart\versepar \skt{दोषे }\edtext{\skt{कृतेस्तु}}{\Afootnote{K; \emph{om.}\ A; \skt{तेशु}\ H}}\skt{ }\edtext{\edtext{\edtext{\skt{क्रिकयामधस्ताच्छीतैर्द्रव्यैरञ्जनान्यस्य}}{\lemma{\skt{क्रिकयामधस्ताच्छीतैर्}}\Afootnote{\skt{ऽधस्ताच्छुक्तिकायामपास्ते शीतै}\ A; \skt{क्रिकयामधस्ताच्छीतैर्}\ K; \skt{त्तिकायामधस्ताच्छीतैर्}\ H}}}{\lemma{\skt{द्रव्यैर्}}\Afootnote{K H; \skt{र्द्रव्यै}\ A}}}{\lemma{\skt{अञ्जनान्यस्य}}\Afootnote{K H; \skt{रञ्जनं कार्यमा}\ A}}\skt{ }\edtext{\skt{कुर्यात्}}{\Afootnote{K H; \skt{शु ॥}\ A}}\skt{ ।}\pend

\pstart\versepar \ledsidenote{\passref{6.10.15}}\edtext{\skt{वैदूर्यम्वा}}{\Afootnote{K; \skt{वैदूर्यं}\ A; \skt{वैदूर्यम्व}\ H}}\skt{ }\edtext{\skt{स्फाटिकं}}{\Afootnote{K; \skt{यत्स्फाटिकं}\ A; \skt{रस्फाटिकं}\ H}}\skt{ }\edtext{\skt{विद्रुमञ्च}}{\Afootnote{K H; \skt{वैद्रुमं च}\ A}}\skt{ }\edtext{\skt{मुक्तां}}{\Afootnote{K H; \skt{मौक्तं}\ A}}\skt{ }\edtext{\skt{शंखं}}{\Afootnote{K H; \skt{शाङ्खं}\ A}}\skt{ }\edtext{\skt{चाञ्जनं}}{\Afootnote{K H; \skt{राजतं}\ A}}\skt{ चाञ्चनञ्च । }\pend
\pstart\versepar \skt{स }\edtext{\skt{चन्दनं}}{\Afootnote{K H; \skt{र्णं सूक्ष्मं}\ A}}\skt{ शर्कराक्षौद्रयुक्तं शुक्तिं }\edtext{\skt{हन्यादञ्जे}}{\Afootnote{K H; \skt{हन्यादञ्जनं}\ A}}\skt{ }\edtext{\skt{नैतदाशुः}}{\Afootnote{K; \skt{चैतदाशु ॥}\ A; \skt{नैतदाशु ।}\ H}}\skt{ ।}\edtext{}{\lemma{\skt{चाञ्चनञ्च । स}}\Afootnote{K H; \skt{शातकुम्भम् । चू}\ A}}\pend

\pstart\versepar \ledsidenote{\passref{6.10.16ab}}\edtext{\skt{घृतं}}{\Afootnote{K H; \skt{युञ्}\ A}}\skt{ }\edtext{\skt{पिवेद्धूमदर्शी}}{\Afootnote{K H; \skt{ज्यात्सर्पिर्धूमदर्शी}\ A}}\skt{ नरस्तु }\pend
\pstart\versepar \skt{}\edtext{\skt{विधिं}}{\Afootnote{K H; \skt{शेषं}\ A}}\skt{ }\edtext{\skt{कुर्यात्पित्तहरञ्च}}{\Afootnote{K H; \skt{कुर्याद्रक्त}\ A}}\skt{ }\edtext{\edtext{\skt{सर्वमिति}}{\lemma{\skt{सर्वम्}}\Afootnote{K H; \emph{om.}\ A}}}{\lemma{\skt{इति}}\Afootnote{K H; \skt{पित्ते विधानम् ।}\ A}}\skt{ ॥}\pend

\pstart\versepar \ledsidenote{\passref{6.10.16.add1}}\skt{औपद्रवं सन्धिगतं वर्त्मशुक्लगतं तथा । }\pend
\pstart\versepar \skt{}\edtext{\skt{कृष्णगतं}}{\Afootnote{K; \skt{कृष्णागतं}\ H}}\skt{ सर्वगतं दृष्टिरोगमथापरं ॥}\pend

\pstart\versepar \ledsidenote{\passref{6.10.16.add2}}\skt{चिकित्साप्रविभाञ्च }\edtext{\skt{वाताभिष्यन्दपित्तकैः}}{\Afootnote{\skt{वाताभिष्यन्दपित्तकै}\ K H}}\skt{ । }\pend
\pstart\versepar \skt{शालाक्यकाशिराजेन प्रोक्तं वै }\edtext{\skt{प्रथमो दश}}{\Afootnote{\skt{प्रथमोदश}\ K H}}\skt{ ॥}\pend

\pstart \ledsidenote{\passref{6.10.trailer}}\edtext{\skt{शालाक्ये ॥ {\symfont ❈}}}{\Afootnote{K; \skt{इति श्रीसुश्रुतसंहितायामुत्तरतन्त्रान्तर्गते शालाक्यतन्त्रे पित्ताभिष्यन्दप्रतिषेधो नाम दशमोऽध्यायः ॥१० ॥}\ A; \skt{शालाक्यस्थाने दशमोऽध्यायः ॥ ॥ ॥}\ H}}\skt{ ॥}\pend

\endnumbering
\beginnumbering
\eledsection*{Adhyāya 11}
\pstart \ledsidenote{\passref{6.11.1}}\skt{अथातः }\edtext{\skt{श्लेष्माभिष्यन्दप्रतिषेधम्व्याख्यास्यामः}}{\lemma{\skt{व्याख्यास्यामः}}\Afootnote{A H; \skt{व्या व्याख्यास्यामः ।}\ K}}\skt{ ।}\pend

\pstart\versepar \ledsidenote{\passref{6.11.3}}\edtext{\skt{स्यन्दाधिमन्थैकफजौ}}{\Afootnote{K H; \skt{स्यन्दाधिमन्थौ कफजौ}\ A}}\skt{ }\edtext{\edtext{\skt{जयेत्शिरोणामथ}}{\lemma{\skt{जयेत्शिरोणाम्}}\Afootnote{K H; \skt{प्रवृद्धौ जयेत्सिराणा}\ A}}}{\lemma{\skt{अथ}}\Afootnote{K H; \skt{मथ}\ A}}\skt{ मोक्षणेन । }\pend
\pstart\versepar \skt{}\edtext{\skt{स्वेदावपीडाञ्जनस्य}}{\Afootnote{K; \skt{स्वेदावपीडाञ्जनधूमसे}\ A; \skt{स्वेदावपीडाञ्जनधूम}\ H}}\skt{ }\edtext{\edtext{\skt{धूमैतथै}}{\Afootnote{K; \skt{कप्रलेप}\ A}}\skt{ }\edtext{\skt{तीक्ष्णैः}}{\Afootnote{K; \skt{योगैः}\ A}}\skt{ }\edtext{\skt{कवडग्रहैश्च}}{\Afootnote{K; \skt{कवलग्रहैश्च ॥}\ A}}\skt{ ।}}{\lemma{\skt{स्य… च}}\Afootnote{\emph{om.}\ H}}\pend

\pstart\versepar \ledsidenote{\passref{6.11.4}}\edtext{\edtext{\edtext{\skt{उष्णैस्तथाश्च्योतनसम्विधानैस्तथैव}}{\lemma{\skt{उष्णैस्}}\Afootnote{K; \skt{रूक्षैस्}\ A}}}{\lemma{\skt{तथाश्च्योतनसम्विधानैस्}}\Afootnote{K; \skt{तथाश्च्योतनसंविधानैस्}\ A}}\skt{ }\edtext{\skt{रुक्षैः}}{\Afootnote{K; \skt{रूक्सैः}\ A; \skt{र्ताक्षैः}\ H}}}{\lemma{\skt{… तथैव}}\Afootnote{\emph{om.}\ H}}\skt{ पुटपाकयोगैः । }\pend
\pstart\versepar \skt{}\edtext{\skt{त्र्यहां}}{\Afootnote{K; \skt{त्र्यहा}\ A H}}\skt{ }\edtext{\skt{त्र्यहा}}{\Afootnote{K; \skt{त्त्र्यहा}\ A; \skt{त्र्यह्वा}\ H}}\skt{ }\edtext{\edtext{\skt{च्चाप्यपतर्पणान्ते}}{\lemma{\skt{च्चाप्य्}}\Afootnote{K; \skt{च्चाप्य्}\ A; \skt{प्य्}\ H}}}{\lemma{\skt{अपतर्पणान्ते}}\Afootnote{A K; \skt{अवतर्प्पणान्ते}\ H}}\skt{ }\edtext{\skt{तौ}}{\Afootnote{K H; \skt{प्रातस्तयोस्तिक्तघृतं}\ A}}\skt{ }\edtext{\edtext{\skt{प्रातरन्नं}}{\lemma{\skt{प्रातर्}}\Afootnote{K; \skt{प्रशस्तम् ॥}\ A; \skt{प्रात}\ H}}}{\lemma{\skt{अन्नं}}\Afootnote{K; \emph{om.}\ A; \skt{रन्नं}\ H}}\skt{ }\edtext{\edtext{\skt{हितमाप्नुयातां}}{\lemma{\skt{हितम्}}\Afootnote{K H; \emph{om.}\ A}}}{\lemma{\skt{आप्नुयातां}}\Afootnote{K; \emph{om.}\ A; \skt{आप्नुयां ।}\ H}}\skt{ ।}\pend

\pstart\versepar \ledsidenote{\passref{6.11.5}}\edtext{\skt{फणिज्झकोस्फुतकपित्थविल्वेष}}{\Afootnote{K; \skt{तदन्नपानं च समाचरेद्धि यच्छ्लेष्मणो नैव करोति वृद्धिम् । कुटन्नटास्फोटफणिज्ज्ञबिल्व}\ A; \skt{फलिकोस्फुतकपित्थविल्व}\ H}}\skt{ }\edtext{\skt{धत्तूरपीलूसुरसार्कभागोः}}{\Afootnote{K; \skt{पत्तूरपिल्वर्ककपित्थभङ्गैः ॥}\ A; \skt{धत्तूरपिलूसुरसार्क्कभङ्गैः ॥}\ H}}\skt{ । }\pend
\pstart\versepar \skt{}\edtext{\edtext{\skt{स्वेदम्विदध्यादथवात्रलेपं}}{\lemma{\skt{स्वेदम्}}\Afootnote{K H; \emph{om.}\ A}}}{\lemma{\skt{विदध्यादथवात्रलेपं}}\Afootnote{K H; \emph{om.}\ A}}\skt{ }\edtext{\skt{वर्हिष्ठशुण्ठीसुरदारुकुष्ठैः}}{\Afootnote{K H; \emph{om.}\ A}}\skt{ ।}\pend

\pstart \ledsidenote{\passref{6.11.6}}\edtext{\skt{सिन्धूप्रसुतं}}{\Afootnote{K; \skt{स्वेदं विदध्यातथवानुलेपं बर्हिष्ठशुण्ठीसुरकाष्ठकुष्ठैः । सिन्धूत्थहिङ्गु}\ A; \skt{सिन्धूप्रसुत}\ H}}\skt{ }\edtext{\skt{त्रिफलामधुकप्रपौण्डरीकाञ्जनतुत्थताम्रैः}}{\Afootnote{K H; \skt{त्रिफलामधूकप्रपौण्डरीकाञ्जनतुत्थताम्रैः ॥}\ A}}\skt{ ।}\pend

\pstart\versepar \ledsidenote{\passref{6.11.7}}\edtext{\skt{वर्त्म्याञ्जने }\edtext{\skt{स्युर्जलसम्प्रपिष्ठा}}{\lemma{\skt{जलसम्प्रपिष्ठा}}\Afootnote{K H; \skt{जनवर्तयः स्युः}\ A}}}{\lemma{\skt{वर्त्म्याञ्जने स्युर्}}\Afootnote{K; \skt{पिष्टैर्जलेनाञ्}\ A; \skt{ञ्जने स्युर्}\ H}}\skt{ पथ्या हरिद्रा }\edtext{\skt{मरिचाञ्जनैर्वा}}{\lemma{\skt{मरिचाञ्}}\Afootnote{K H; \skt{मधुकाञ्}\ A}}\skt{ }\pend
\pstart\versepar \skt{}\edtext{\skt{त्रीण्युपणानि}}{\Afootnote{K H; \skt{त्रीण्यूषणानि}\ A}}\skt{ }\edtext{\skt{त्रिफला}}{\Afootnote{A K; \skt{त्रिफलां}\ H}}\skt{ }\edtext{\edtext{\skt{हरिद्राम्विडंगसारञ्च}}{\lemma{\skt{हरिद्राम्}}\Afootnote{K H; \skt{हरिद्रा}\ A}}}{\lemma{\skt{विडंगसारञ्च}}\Afootnote{K H; \skt{विडङ्गसारश्च}\ A}}\skt{ समानि च स्युः ।}\pend

\pstart\versepar \ledsidenote{\passref{6.11.8}}\edtext{\skt{वर्हिष्ठकुष्ठेसुरदारूणशंखपाठानखव्योषमनः}}{\Afootnote{K; \skt{बर्हिष्ठकुष्ठामरकाष्ठशङ्खपाठामलव्योषमनः}\ A; \skt{वर्हिष्ठकुष्ठेसुरदारूणं खं पाठानखव्योष}\ H}}\skt{ }\edtext{\skt{शिला च}}{\Afootnote{K; \skt{शिलाश्च ।}\ A; \skt{ला च ।}\ H}}\skt{ । }\pend
\pstart\versepar \skt{}\edtext{\skt{मूर्व्वाश्च}}{\Afootnote{K; \skt{पिष्ट्वाम्बु}\ A; \skt{मूर्व्वा च}\ H}}\skt{ }\edtext{\skt{जाश्या}}{\Afootnote{K; \skt{ना वा}\ A; \skt{जाञ्ज}\ H}}\skt{ }\edtext{\skt{मूकुलानि}}{\Afootnote{K H; \skt{कुसुमानि}\ A}}\skt{ }\edtext{\skt{चापि}}{\Afootnote{K H; \skt{जाति}\ A}}\skt{ }\edtext{\skt{पिष्ठ्वाम्वुनैतानि}}{\Afootnote{K; \skt{करञ्जशोभाञ्जनजानि}\ A; \skt{पिष्ठाबुनैतानि}\ H}}\skt{ }\edtext{\skt{भिषग्विदध्यात्}}{\Afootnote{K H; \skt{युञ्ज्यात् ॥}\ A}}\skt{ ।}\pend

\pstart\versepar \ledsidenote{\passref{6.11.9}}\skt{फलं }\edtext{\skt{प्रकीर्यात्वथवापिशिग्राः}}{\Afootnote{K H\textsuperscript{pc}; \skt{प्रकीर्यादथवापि शिग्रोः}\ A; \skt{प्रकीर्यात्वथवाशिग्राः}\ H\textsuperscript{ac}}}\skt{ पुष्पञ्च तुल्यं }\edtext{\skt{वृहदर्शनस्य}}{\Afootnote{K H; \skt{बृहतीद्वयस्य ।}\ A}}\skt{ । }\pend
\pstart\versepar \skt{}\edtext{\skt{रसाञ्जनञ्चन्दनसैन्धवञ्च}}{\lemma{\skt{चन्दनसैन्धवञ्च}}\Afootnote{K; \skt{सैन्धवचन्दनं च}\ A; \skt{चन्दनसैन्धवे च}\ H}}\skt{ मनःशिलाले लशुनञ्च तुल्यं ।}\pend

\pstart\versepar \ledsidenote{\passref{6.11.10}}\edtext{\skt{पिष्ठ्वाथ वर्त्ती}}{\Afootnote{K H; \skt{पिष्ट्वाञ्जनार्थे कफजेषु धीमान्वर्ती}\ A}}\skt{ }\edtext{\skt{विदधीत}}{\Afootnote{K H; \skt{र्विदध्यान्न}\ A}}\skt{ }\edtext{\skt{सम्यग्न}}{\Afootnote{K H; \skt{यनाम}\ A}}\skt{ }\edtext{\skt{देतु}}{\Afootnote{K H; \skt{येषु ।}\ A}}\skt{ }\edtext{\skt{धीमां}}{\Afootnote{K; \skt{धीमा}\ H}}\skt{ }\edtext{\skt{कफजेञ्जनार्थं}}{\Afootnote{K; \skt{न्कफजेञ्जनार्थं ।}\ H}}\skt{ । }\pend
\pstart\versepar \skt{रोगे }\edtext{\skt{वलासग्रथितेन्तुजान्ः तु}}{\Afootnote{K; \skt{बलासग्रथितेऽञ्जनं ज्ञैः}\ A; \skt{वलासग्रथितेञ्जनन्तु}\ H}}\skt{ }\edtext{\skt{कुर्यादिदं}}{\Afootnote{K H; \skt{कर्त}\ A}}\skt{ }\edtext{\skt{स्यात्प्रविशाष्यदेहः}}{\Afootnote{K H; \skt{व्यमेतत्सुविशुद्धकाये ॥}\ A}}\skt{ ।}\edtext{}{\lemma{\skt{धीमां… रो}}\Afootnote{\emph{om.}\ A}}\pend

\pstart \ledsidenote{\passref{6.11.11}}\edtext{\skt{नीलां}}{\Afootnote{K H; \skt{नीला}\ A}}\skt{ }\edtext{\skt{यवानाजपयोन्नुपीतां}}{\Afootnote{K; \skt{न्यवान्गव्यपयोऽनुपीता}\ A; \skt{यवानाजपयोन्नपीतां}\ H}}\skt{ }\edtext{\skt{संशाध्य}}{\Afootnote{K H; \skt{न्शलाकि}\ A}}\skt{ }\edtext{\skt{संशोष्य}}{\Afootnote{K H; \skt{नः शुष्क}\ A}}\skt{ }\edtext{\skt{ततो}}{\Afootnote{K H; \skt{त}\ A}}\skt{ }\edtext{\skt{नूदह्य}}{\Afootnote{K; \skt{नून्विदह्य । तथार्जकास्फोतकपित्थबिल्वनिर्गुण्डिजातीकुसुमानि चैव ॥}\ A; \skt{ऽनूदह्य ।}\ H}}\skt{ ।}\pend

\pstart\versepar \ledsidenote{\passref{6.11.12}}\edtext{\skt{तक्षारकल्येन च}}{\Afootnote{K H; \skt{तत्क्षारवत्सैन्धव}\ A}}\skt{ }\edtext{\skt{भस्मधीमां}}{\Afootnote{K; \skt{तुत्थरोचनं}\ A; \skt{भस्मधीमा}\ H}}\skt{ }\edtext{\edtext{\skt{पक्वम्विदध्यादथवेशूनाड्यां}}{\lemma{\skt{पक्वम्}}\Afootnote{A K; \skt{न्पक्वम्}\ H}}}{\lemma{\skt{विदध्यादथवेशूनाड्यां}}\Afootnote{K H; \skt{विदध्यादथ लोहनाड्याम् ।}\ A}}\skt{ । }\pend
\pstart\versepar \skt{}\edtext{\skt{एतद्वलासग्रथितेञ्जनं}}{\Afootnote{K H; \skt{एतद्बलासग्रथितेऽञ्जनं}\ A}}\skt{ }\edtext{\edtext{\skt{स्यात्कल्पस्तत्थैवैषफलिञ्जकेषु}}{\lemma{\skt{स्यात्कल्पस्}}\Afootnote{K; \skt{स्यादेषोऽनुकल्पस्}\ A; \skt{स्यात्कल्पस्त}\ H}}}{\lemma{\skt{तत्थैवैषफलिञ्जकेषु}}\Afootnote{K; \skt{तु फणिज्ज्ञकादौ ॥}\ A; \skt{त्थैवैषफलिञ्जकेषु ।}\ H}}\skt{ ।}\pend

\pstart\versepar \ledsidenote{\passref{6.11.13}}\skt{महौषधं }\edtext{\skt{मागधिकाश्च}}{\Afootnote{K H; \skt{मागधिकां च}\ A}}\skt{ }\edtext{\skt{मुख्यां}}{\Afootnote{K; \skt{मुस्तां}\ A; \skt{मुख्या}\ H}}\skt{ }\edtext{\skt{ससन्धवं}}{\Afootnote{K; \skt{ससैन्धवं}\ A; \skt{संसिन्धवं}\ H}}\skt{ }\edtext{\skt{यत्मरिचञ्च}}{\Afootnote{K H; \skt{यन्मरिचं च}\ A}}\skt{ }\edtext{\skt{शुष्कं}}{\Afootnote{K H; \skt{शुक्लम् ।}\ A}}\skt{ । }\pend
\pstart\versepar \skt{}\edtext{\skt{तत्मातुलुङ्गस्य}}{\Afootnote{K H; \skt{तन्मातुलुङ्गस्व}\ A}}\skt{ रसेन पिष्टं नेत्राञ्जनं पिष्टकमाशु हन्यात् ।}\pend

\pstart\versepar \ledsidenote{\passref{6.11.14}}\edtext{\edtext{\skt{फलेर्वृहत्या}}{\lemma{\skt{फलेर्}}\Afootnote{K H; \skt{फले}\ A}}}{\lemma{\skt{वृहत्या}}\Afootnote{K H; \skt{बृहत्या}\ A}}\skt{ }\edtext{\skt{मगधोद्भवानां}}{\Afootnote{K; \skt{मगधोद्भवाना}\ A; \skt{मगधा भवानां}\ H}}\skt{ }\edtext{\skt{कल्कं}}{\Afootnote{K H; \skt{निधाय कल्कं}\ A}}\skt{ }\edtext{\skt{निरवात्वात्फलपाकसीन्ध्य}}{\Afootnote{K H; \skt{फलपाककाले ।}\ A}}\skt{ }\pend
\pstart\versepar \skt{}\edtext{\skt{श्रोतोजयुक्तं}}{\Afootnote{K; \skt{स्रोतोजयुक्तं}\ A H}}\skt{ }\edtext{\skt{तत}}{\Afootnote{K H; \skt{च त}\ A}}\skt{ }\edtext{\skt{उद्धृतं}}{\Afootnote{K H; \skt{दुद्धृतं}\ A}}\skt{ }\edtext{\skt{स्यात्तद्वात्र}}{\Afootnote{K H; \skt{स्यात्तद्वत्तु}\ A}}\skt{ }\edtext{\skt{पथ्यं}}{\Afootnote{K; \skt{पिष्टे}\ A; \skt{पथ्याम्}\ H}}\skt{ }\edtext{\skt{विधिरेष}}{\Afootnote{K H; \skt{विधिरेष}\ A}}\skt{ चापि ॥}\pend

\pstart\versepar \ledsidenote{\passref{6.11.15}}\edtext{\skt{वार्ताकिशिग्रुमृगज}}{\Afootnote{K; \skt{वार्ताकशिग्र्विन्द्रसु}\ A; \skt{वार्ताकिशिग्रुमृगजो}\ H}}\skt{ }\edtext{\skt{पटोलकिराततिक्तामलकीफलेषु}}{\Afootnote{K; \skt{रापटोलकिराततिक्तामलकीफलेषु ।}\ A; \skt{पटोलकिरागतिकामलकीफलेषु ॥}\ H}}\skt{ ॥ }\pend
\pstart\versepar \skt{}\edtext{\skt{कासीससामुद्ररसाञ्जनानि}}{\Afootnote{A K; \skt{कासीससामुद्रसाञ्जनानि}\ H}}\skt{ }\edtext{\edtext{\skt{जात्यास्तथा}}{\lemma{\skt{जात्यास्}}\Afootnote{A K; \skt{जात्यो}\ H}}}{\lemma{\skt{तथा}}\Afootnote{A K; \skt{थ}\ H\textsuperscript{ac}; \skt{थवा}\ H\textsuperscript{pc}}}\skt{ }\edtext{\edtext{\skt{क्षारकमेव}}{\lemma{\skt{क्षारकम्}}\Afootnote{A K; \emph{om.}\ H\textsuperscript{ac}; \skt{क्षीरिकम्}\ H\textsuperscript{pc}}}}{\lemma{\skt{एव}}\Afootnote{K H\textsuperscript{pc}; \skt{(कोरकम्)एव}\ A; \skt{क्षीरिकमेव}\ H\textsuperscript{ac}}}\skt{ चापि ।}\pend

\pstart\versepar \ledsidenote{\passref{6.11.16}}\edtext{\edtext{\skt{पक्लिन्नवर्त्मन्युपदिश्येतेतु}}{\lemma{\skt{पक्लिन्नवर्त्मन्य्}}\Afootnote{K; \skt{प्रक्लिन्नवर्त्म}\ A; \skt{प्रक्लिन्नवर्त्मन्य्}\ H}}}{\lemma{\skt{उपदिश्येतेतु}}\Afootnote{K H; \skt{न्युपदिश्यते तु}\ A}}\skt{ योगाञ्जनं }\edtext{\skt{तन्मधुना}}{\Afootnote{K H; \skt{तन्मधुना}\ A}}\skt{ }\edtext{\skt{विघृष्टं}}{\Afootnote{K; \skt{वघृष्टम् ।}\ A; \skt{विघृष्टा ॥}\ H}}\skt{ । }\pend
\pstart\versepar \skt{}\edtext{\skt{नादेयमग्र्यलवलञ्च}}{\Afootnote{K; \skt{नादेयमग्र्यं मरिचं च}\ A; \skt{नादेयमग्र्यलवलं च}\ H}}\skt{ }\edtext{\skt{शुक्लं}}{\Afootnote{A K; \skt{शुक्ल}\ H}}\skt{ }\edtext{\skt{नेपालजाताञ्च}}{\Afootnote{K; \skt{नेपालजाता च}\ A; \skt{नेपालजाता}\ H}}\skt{ }\edtext{\edtext{\skt{समासतस्तु}}{\lemma{\skt{समासतस्}}\Afootnote{K; \skt{समप्रमाणा ॥}\ A; \skt{तु समास तस्}\ H}}}{\lemma{\skt{तु}}\Afootnote{K H; \emph{om.}\ A}}\skt{ ।}\pend

\pstart\versepar \ledsidenote{\passref{6.11.17}}\edtext{\skt{समातुलुंगाद्रस}}{\Afootnote{K H; \skt{समातुलुङ्गद्रव}\ A}}\skt{ एषयोगः }\edtext{\skt{कण्डुन्निहन्यात्सकृदञ्जनेन}}{\lemma{\skt{कण्डुन्}}\Afootnote{K H; \skt{कण्डूं}\ A}}\skt{ ॥ }\pend
\pstart\versepar \skt{सशृङ्गवेरं }\edtext{\skt{सुरदारू}}{\Afootnote{K H; \skt{सुरदारु}\ A}}\skt{ }\edtext{\skt{मुख्यं}}{\Afootnote{K H; \skt{मुस्तं}\ A}}\skt{ }\edtext{\skt{सिन्धुप्रसूतं}}{\Afootnote{K H; \skt{सिन्धुप्रभूतं(प्रसूतं)}\ A}}\skt{ }\edtext{\skt{मुकुलञ्च}}{\Afootnote{K; \skt{मुकुलानि}\ A H}}\skt{ जात्याः ।}\pend

\pstart\versepar \ledsidenote{\passref{6.11.18}}\edtext{\skt{सुरासुपिष्टं}}{\Afootnote{K; \skt{सुराप्रपिष्टं}\ A; \skt{सुरापिष्ट}\ H}}\skt{ }\edtext{\skt{स्विदमञ्जनन्तु}}{\Afootnote{K; \skt{त्विदमञ्जनं हि}\ A; \skt{न्स्विदमञ्जनन्तु}\ H}}\skt{ }\edtext{\skt{कण्ड्वाञ्च}}{\Afootnote{A K; \skt{कण्ड्वा}\ H}}\skt{ }\edtext{\skt{शोफे च}}{\Afootnote{A K; \skt{ञ्शोफे च}\ H}}\skt{ हितं वदन्ति । }\pend
\pstart\versepar \skt{}\edtext{\edtext{\skt{स्यन्दाधिमन्थक्रममाचरेच्च}}{\lemma{\skt{स्यन्दाधिमन्थक्रमम्}}\Afootnote{K H; \skt{स्यन्दाधिमन्थक्रम}\ A}}}{\lemma{\skt{आचरेच्च}}\Afootnote{K H; \skt{माचरेच्च}\ A}}\skt{ सर्वेषु चैतेषु सदा }\edtext{\skt{प्रमत्त}}{\Afootnote{K H; \skt{प्रमत्तः ॥}\ A}}\skt{ }\edtext{\skt{इति}}{\Afootnote{K H; \skt{(विशेषतो नावनमेव कार्यं संसर्जनं चापि यथोपदिष्टम्)}\ A}}\skt{ ॥}\pend

\pstart \ledsidenote{\passref{6.11.trailer}}\edtext{\skt{शालाक्ये}}{\Afootnote{\skt{इति सुश्रुतसंहितायामुत्तरतन्त्रान्तर्गते शालाक्यतन्त्रे कफाभिष्यन्दप्रतिषेधो नामैकादशोऽध्यायः ॥११ ॥}\ A; \skt{शालाक्ये लृ ॥ ० ॥}\ K; \skt{शालाक्यस्थाने एकादशमोध्यायः ॥ ॥}\ H}}\skt{ ॥ ० ॥}\pend

\endnumbering
\beginnumbering
\eledsection*{Adhyāya 12}
\pstart \ledsidenote{\passref{6.12.1}}\skt{अथातो }\edtext{\skt{रक्ताभिष्यन्दप्रतिषेढं}}{\Afootnote{K H; \skt{रक्ताभिष्यन्दप्रतिषेधं}\ A}}\skt{ }\edtext{\skt{व्याख्यास्यामः}}{\Afootnote{A H; \skt{व्या व्याख्यास्यामः ॥}\ K}}\skt{ ॥}\pend

\pstart\versepar \ledsidenote{\passref{6.12.3}}\edtext{\skt{मन्थस्यन्दौसिरोत्पातं}}{\Afootnote{K H; \skt{मन्थं स्यन्दं सिरोत्पातं}\ A}}\skt{ सिराहर्षञ्च रक्तजं । }\pend
\pstart\versepar \skt{एकेनैव विधानेन चिकित्सेच्चतुरो }\edtext{\skt{गदां}}{\Afootnote{K; \skt{गदान् ॥}\ A H}}\skt{ ।}\pend

\pstart\versepar \ledsidenote{\passref{6.12.4}}\edtext{\edtext{\skt{व्याध्यर्ताश्चतुरोप्यतां}}{\lemma{\skt{व्याध्यर्ताश्}}\Afootnote{K; \skt{व्याध्यार्तांश्}\ A; \skt{व्याध्यर्त्ता ॥ ॥}\ H}}}{\lemma{\skt{चतुरोप्यतां}}\Afootnote{K; \skt{चतुरोऽप्येता}\ A; \skt{चतुरो प्येतां}\ H}}\skt{ }\edtext{\skt{स्निग्धां}}{\Afootnote{K; \skt{न्स्निग्धान्}\ A; \skt{स्निग्धा}\ H}}\skt{ }\edtext{\skt{कोम्भेन}}{\Afootnote{K H; \skt{कौम्भेन}\ A}}\skt{ सर्प्पिषा । }\pend
\pstart\versepar \skt{}\edtext{\skt{रसैरूदारैरथवा}}{\lemma{\skt{रसैरूदारैर्}}\Afootnote{K H; \skt{रसैरुदारैर्}\ A}}\skt{ सिरामोक्षेण योजयेत् ।}\pend

\pstart\versepar \ledsidenote{\passref{6.12.5}}\edtext{\skt{विरिक्तानाम्प्रगाढञ्च}}{\lemma{\skt{प्रगाढञ्च}}\Afootnote{K H; \skt{प्रकामं च}\ A}}\skt{ }\edtext{\skt{सिरान्येषाम्विशोधयेत्}}{\lemma{\skt{सिरान्येषाम्}}\Afootnote{K H; \skt{शिरांस्येषां}\ A}}\skt{ । }\pend
\pstart\versepar \skt{वैरेचनिकसिद्धेन सितायुक्तेन }\edtext{\skt{सर्पिषा}}{\Afootnote{K H; \skt{सर्पिषा ॥ (मज्ज्ञा वा तद्विमिश्रेण मेदसा तच्छृतेन वा । )}\ A}}\skt{ ।}\pend

\pstart\versepar \ledsidenote{\passref{6.12.6}}\skt{ततः }\edtext{\skt{प्रदेहाम्परिषेचनानि}}{\lemma{\skt{प्रदेहाम्}}\Afootnote{K H; \skt{प्रदेहाः}\ A}}\skt{ }\edtext{\skt{तस्यानि}}{\Afootnote{K H; \skt{नस्यानि}\ A}}\skt{ }\edtext{\skt{धूमानि}}{\Afootnote{K H; \skt{धूमाश्च}\ A}}\skt{ यथा स्वमेव । }\pend
\pstart\versepar \skt{}\edtext{\skt{आश्च्योतनाभ्यञ्जन}}{\Afootnote{K H; \skt{आश्च्योतनाभ्यञ्जन}\ A}}\skt{ }\edtext{\skt{तर्पशानि}}{\Afootnote{K H; \skt{तर्पणानि}\ A}}\skt{ }\edtext{\skt{स्विग्धांश्च}}{\Afootnote{K H; \skt{स्निग्धाश्च}\ A}}\skt{ }\edtext{\skt{कार्यान्पुटपाकयोगां}}{\Afootnote{K H; \skt{कार्याः पुटपाकयोगाः ॥}\ A}}\skt{ ।}\pend

\pstart\versepar \ledsidenote{\passref{6.12.8}}\edtext{\skt{रुजायाञ्चाप्यर्भृशं}}{\lemma{\skt{चाप्यर्भृशं}}\Afootnote{K; \skt{चाप्यतिभृशं}\ A; \skt{चाप्यथफलं}\ H}}\skt{ }\edtext{\skt{स्वेदास्तु}}{\Afootnote{K H; \skt{स्वेदाश्च}\ A}}\skt{ मृदवो }\edtext{\skt{हिताः}}{\Afootnote{A K; \skt{हितं ।}\ H}}\skt{ । }\pend
\pstart\versepar \skt{}\edtext{\skt{अक्ष्णोरपि}}{\lemma{\skt{अपि}}\Afootnote{K H; \skt{समन्}\ A}}\skt{ }\edtext{\skt{समन्ताच्च}}{\Afootnote{K H; \skt{ततः कार्यं}\ A}}\skt{ }\edtext{\skt{पातनन्तु}}{\Afootnote{K H; \skt{पातनं च}\ A}}\skt{ जलौकसां ।}\pend

\pstart\versepar \ledsidenote{\passref{6.12.9}}\skt{घृतस्य महती मात्रा पीता }\edtext{\skt{चार्ति}}{\Afootnote{K; \skt{चार्तिं}\ A H}}\skt{ नियच्छति । }\pend
\pstart\versepar \skt{पिताभिष्यन्दशमनो }\edtext{\edtext{\skt{गणम्चाष्प्युपयादितः}}{\lemma{\skt{गणम्}}\Afootnote{K H; \skt{विधिश्}\ A}}}{\lemma{\skt{चाष्प्युपयादितः}}\Afootnote{K H; \skt{चाप्युपपादितः ॥}\ A}}\skt{ ।}\pend

\pstart\versepar \ledsidenote{\passref{6.12.10}}\edtext{\skt{कसेरुमधुकानाञ्च}}{\Afootnote{K H; \skt{कशेरुमधुकाभ्यां वा}\ A}}\skt{ }\edtext{\skt{चूर्णमम्वरसम्वृतं}}{\Afootnote{K H; \skt{चूर्णमम्बरसंवृतम् ।}\ A}}\skt{ । }\pend
\pstart\versepar \skt{}\edtext{\skt{स्तन्यमप्स्वान्तरिक्षासु}}{\Afootnote{K H; \skt{न्यस्तमप्स्वान्तरिक्षासु}\ A}}\skt{ }\edtext{\skt{हितमाश्च्योतनं}}{\Afootnote{K H; \skt{हितमाश्च्योतनं}\ A}}\skt{ }\edtext{\skt{सदा}}{\Afootnote{K; \skt{भवेत् ॥}\ A; \skt{तथा ॥}\ H}}\skt{ ।}\pend

\pstart\versepar \ledsidenote{\passref{6.12.11}}\edtext{\skt{श्रीपर्णीपाटलीधात्री}}{\Afootnote{K H; \skt{पाटल्यर्जुनश्रीपर्णीधातकीधात्रि}\ A}}\skt{ धातकीविल्वकार्जुनात् । }\pend
\pstart\versepar \skt{पुष्पाण्यथ }\edtext{\skt{वृहत्यौ च}}{\Afootnote{K H; \skt{बृहत्योश्च}\ A}}\skt{ }\edtext{\skt{विम्वीलोटाश्च}}{\Afootnote{K H; \skt{बिम्बीलोटाच्च}\ A}}\skt{ तुल्यशः ।}\edtext{}{\lemma{\skt{धातकीविल्वकार्जुनात् । पुष्पाण्यथ}}\Afootnote{K H; \skt{बिल्वतः । पुष्पाण्यथ}\ A}}\pend

\pstart\versepar \ledsidenote{\passref{6.12.12}}\edtext{\skt{मञ्जिष्ठा}}{\Afootnote{K H; \skt{समञ्जिष्ठा}\ A}}\skt{ }\edtext{\skt{चापि}}{\Afootnote{K H; \skt{नि}\ A}}\skt{ मधुना }\edtext{\skt{पिष्ट्वा}}{\Afootnote{K H; \skt{पिष्टा}\ A}}\skt{ }\edtext{\skt{चक्षुरसेन वा}}{\Afootnote{K H; \skt{नीक्षुरसेन वा ।}\ A}}\skt{ । }\pend
\pstart\versepar \skt{}\edtext{\edtext{\skt{रुधिरस्यन्दशान्त्यर्थमेतदञ्जनमिष्यते}}{\lemma{\skt{रुधिरस्यन्दशान्त्यर्थम्}}\Afootnote{\skt{रक्ताभिष्यन्दशान्त्यर्थम्}\ A; \skt{रुधिरस्यन्दशान्त्यर्थ}\ K H}}}{\lemma{\skt{एतदञ्जनम्}}\Afootnote{A; \skt{मेतदञ्जनम्}\ K H}}\skt{ ।}\pend

\pstart\versepar \ledsidenote{\passref{6.12.13}}\edtext{\edtext{\skt{चन्दनन्तुगरम्पत्रं}}{\lemma{\skt{चन्दनन्}}\Afootnote{A K; \skt{चन्दन}\ H}}}{\lemma{\skt{तुगरम्}}\Afootnote{K; \skt{कुमुदं}\ A; \skt{ञ्कुमुददम्}\ H}}\skt{ शिलाजतु सकुङ्कुमं । }\pend
\pstart\versepar \skt{}\edtext{\skt{रजस्ताम्रमयं तुत्थं}}{\Afootnote{K; \skt{अयस्ताम्ररजस्तुत्थं}\ A; \skt{रजस्ताम्रयतुथ}\ H\textsuperscript{ac}; \skt{रजस्ताम्रमयतुथ}\ H\textsuperscript{pc}}}\skt{ }\edtext{\skt{निम्वकाशी}}{\Afootnote{K H; \skt{निम्बनिर्या}\ A}}\skt{ }\edtext{\skt{समञ्जनं}}{\Afootnote{K H; \skt{समञ्जनम् ॥}\ A}}\skt{ ।}\pend

\pstart\versepar \ledsidenote{\passref{6.12.14}}\skt{त्रपु }\edtext{\skt{कासं}}{\Afootnote{K H; \skt{कांस्य}\ A}}\skt{ }\edtext{\skt{मलश्चापि}}{\lemma{\skt{मलश्}}\Afootnote{K; \skt{मलं}\ A; \skt{मलौ}\ H}}\skt{ पिष्ट्वा पुष्परसेन तु । }\pend
\pstart\versepar \skt{}\edtext{\skt{विद्वुदायां}}{\Afootnote{K H; \skt{विपुला याः}\ A}}\skt{ कृता वर्त्यः }\edtext{\edtext{\skt{पूजितश्चाञ्जने}}{\lemma{\skt{पूजितश्}}\Afootnote{K; \skt{पूजिताश्}\ A H}}\skt{ तु}}{\lemma{\skt{चाञ्जने तु}}\Afootnote{K H; \skt{चाञ्जने स}\ A}}\skt{ }\edtext{\skt{ताः}}{\Afootnote{K H; \skt{दा ॥}\ A}}\skt{ ।}\pend

\pstart\versepar \ledsidenote{\passref{6.12.15}}\edtext{\skt{सर्पिः}}{\Afootnote{K H; \skt{स्यादञ्जनं घृतं}\ A}}\skt{ क्षौद्रं }\edtext{\skt{चाञ्जनं}}{\Afootnote{K; \emph{om.}\ A; \skt{चाञ्जन}\ H}}\skt{ }\edtext{\skt{स्यात्सिरोत्पातस्य}}{\Afootnote{K H; \skt{सिरोत्पातस्य}\ A}}\skt{ भेषजं । }\pend
\pstart\versepar \skt{तद्वत्सैन्धवकासीसं }\edtext{\skt{स्तन्ययुक्तं तु}}{\Afootnote{K H; \skt{स्तन्यघृष्टं च}\ A}}\skt{ पूजितं ॥}\pend

\pstart\versepar \ledsidenote{\passref{6.12.16}}\skt{मधुना }\edtext{\skt{शंखनेपाली}}{\Afootnote{K H; \skt{शङ्खनैपाली}\ A}}\skt{ तुत्थदार्व्यः ससैन्धवा । }\pend
\pstart\versepar \skt{रसः शिरीष }\edtext{\skt{पुष्पाद्वा}}{\Afootnote{K H; \skt{पुष्पाच्च}\ A}}\skt{ सुरामरिचमाक्षिकैः ।}\pend

\pstart\versepar \ledsidenote{\passref{6.12.17}}\edtext{\skt{युक्ताथ}}{\Afootnote{K; \skt{युक्तं तु}\ A; \skt{युक्ताऽथ}\ H}}\skt{ मधुना वापि }\edtext{\skt{गैरिको}}{\Afootnote{K H; \skt{गैरिकं}\ A}}\skt{ }\edtext{\skt{हितसञ्जनं}}{\Afootnote{K H; \skt{हितमञ्जने ।}\ A}}\skt{ ॥ }\pend
\pstart\versepar \skt{}\edtext{\skt{शिराहर्षञ्जनं}}{\Afootnote{K H; \skt{सिराहर्षेऽञ्जनं}\ A}}\skt{ }\edtext{\skt{कुर्यात्फाणितं}}{\Afootnote{K H; \skt{कुर्यात्फणितं}\ A}}\skt{ }\edtext{\skt{मधुसैन्धवं}}{\Afootnote{K H; \skt{मधुसंयुतम् ॥}\ A}}\skt{ ।}\pend

\pstart\versepar \ledsidenote{\passref{6.12.18}}\skt{मधुना तार्क्ष्यजं }\edtext{\skt{शैलम्वा}}{\Afootnote{K H; \skt{वापि}\ A}}\skt{ }\edtext{\skt{कासीसम्वा}}{\Afootnote{K H; \skt{कासीअं वा}\ A}}\skt{ }\edtext{\skt{ससैन्धवं}}{\Afootnote{K H; \skt{ससैधवम् ।}\ A}}\skt{ । }\pend
\pstart\versepar \skt{}\edtext{\skt{वेतमाम्लं}}{\Afootnote{K; \skt{वेत्राम्लस्त}\ A; \emph{om.}\ H}}\skt{ }\edtext{\skt{स्तन्ययुक्तं}}{\Afootnote{K; \skt{न्यसंयुक्तं}\ A; \emph{om.}\ H}}\skt{ }\edtext{\skt{कासीसम्वा}}{\Afootnote{K; \skt{फाणितं च}\ A; \emph{om.}\ H}}\skt{ }\edtext{\skt{ससैन्धवं}}{\Afootnote{A K; \emph{om.}\ H}}\skt{ ।}\pend

\pstart\versepar \ledsidenote{\passref{6.12.19}}\edtext{\skt{पैत्तिकान्तु}}{\Afootnote{K; \skt{पैत्तं}\ A; \skt{पैत्तिकं तु}\ H}}\skt{ }\edtext{\skt{विधिं}}{\Afootnote{K H; \skt{विधि}\ A}}\skt{ }\edtext{\skt{कुर्याद्कृत्स्नं}}{\Afootnote{K; \skt{मशेषेण कुर्या}\ A; \skt{कृत्स्नं कुर्या}\ H}}\skt{ }\edtext{\skt{अर्जुनशान्तये}}{\Afootnote{K; \skt{दर्जुनशान्तये}\ A H}}\skt{ ॥ }\pend
\pstart\versepar \skt{}\edtext{\skt{सितेक्षु}}{\Afootnote{K H; \skt{इक्षुक्षौ}\ A}}\skt{ }\edtext{\skt{मधुकक्षेद्र}}{\Afootnote{K H; \skt{द्रसितास्तन्य}\ A}}\skt{ }\edtext{\skt{दार्व्वीस्तन्यैः}}{\Afootnote{K H; \skt{दार्वीमधु}\ A}}\skt{ }\edtext{\skt{ससन्धवैः}}{\Afootnote{K H; \skt{कसैन्धवैः ॥}\ A}}\skt{ ।}\pend

\pstart\versepar \ledsidenote{\passref{6.12.20}}\skt{सेकाञ्जनञ्चात्र }\edtext{\skt{हितमम्लैराश्च्योतनैस्तथा}}{\lemma{\skt{हितमम्लैराश्च्योतनैस्}}\Afootnote{K; \skt{हितमम्लैराश्च्योतनं}\ A; \skt{हितमम्लैराश्च्योतनं}\ H}}\skt{ । }\pend
\pstart\versepar \skt{}\edtext{\skt{सितामधुककट्वाङ्गमस्तुक्षौद्राम्लसैन्धवैः}}{\Afootnote{K H; \skt{सितामधुककट्वङ्गमस्तुक्षौद्राम्लसैन्धवैः ॥}\ A}}\skt{ ।}\pend

\pstart\versepar \ledsidenote{\passref{6.12.21}}\edtext{\skt{वीजपूरकको लाम्लदाभिमाम्लैश्च}}{\Afootnote{K; \skt{बीजपूरककोलाम्लदाडिमाम्लैश्च}\ A; \skt{वीजपूरककोम्लदाडिमाम्लैश्च}\ H\textsuperscript{ac}; \skt{वीजपूरककोलाम्लदाडिमाम्लैश्च}\ H\textsuperscript{pc}}}\skt{ युक्तितः । }\pend
\pstart\versepar \skt{एकशो वा द्विशो वापि योजितम्वा त्रिभिस्त्रिभिः ।}\pend

\pstart\versepar \ledsidenote{\passref{6.12.22}}\skt{स्फटिकं }\edtext{\skt{कुंकुमं}}{\Afootnote{K; \skt{विद्रुमं}\ A; \skt{कुमं}\ H}}\skt{ }\edtext{\skt{शंखा}}{\Afootnote{K H; \skt{शङ्खो}\ A}}\skt{ मधुकं }\edtext{\skt{मधु}}{\Afootnote{K H; \skt{कधु}\ A}}\skt{ }\edtext{\skt{चापि हि}}{\Afootnote{K H; \skt{चैव हि ।}\ A}}\skt{ । }\pend
\pstart\versepar \skt{शंखक्षौद्रसितायुक्तः }\edtext{\skt{सामुद्र}}{\Afootnote{K H; \skt{सामुद्रः}\ A}}\skt{ फेन एव }\edtext{\skt{च}}{\Afootnote{K H; \skt{वा ॥}\ A}}\skt{ ।}\pend

\pstart\versepar \ledsidenote{\passref{6.12.23}}\edtext{\skt{द्वावेतौ}}{\Afootnote{K H; \skt{द्वाविमौ}\ A}}\skt{ }\edtext{\skt{विहितौ}}{\Afootnote{A K; \skt{विहि}\ H\textsuperscript{ac}; \skt{विहिता}\ H\textsuperscript{pc}}}\skt{ }\edtext{\skt{योगावञ्जनार्जुननाशनौ}}{\Afootnote{K H; \skt{योगावञ्जनेऽर्जुननाशनौ ।}\ A}}\skt{ ॥ }\pend
\pstart\versepar \skt{}\edtext{\skt{सैन्धवक्षौद्रकतकं}}{\Afootnote{K; \skt{सैन्धवक्षौद्रकतकाः}\ A; \skt{सन्धैवक्षौद्रकतकाः}\ H}}\skt{ सक्षौद्रम्वा रसाञ्जनं ।}\pend

\pstart\versepar \ledsidenote{\passref{6.12.24}}\skt{कासीसम्मधुना }\edtext{\skt{चापि}}{\Afootnote{K H; \skt{वापि}\ A}}\skt{ योज्यमत्राञ्जने सदा । }\pend
\pstart\versepar \skt{लोहचूर्ण्णानि सर्वाणि धातवो लवणानि च ।}\pend

\pstart\versepar \ledsidenote{\passref{6.12.25}}\skt{रत्नानि दन्ताः }\edtext{\skt{शृङ्गाणि}}{\Afootnote{A K; \skt{शृङ्गानि}\ H}}\skt{ }\edtext{\edtext{\skt{गणाश्चाप्युपपादितः}}{\lemma{\skt{गणाश्}}\Afootnote{K; \skt{गणश्}\ A H}}}{\lemma{\skt{उपपादितः}}\Afootnote{K H; \skt{अवसादनः ।}\ A}}\skt{ । }\pend
\pstart\versepar \skt{}\edtext{\skt{कुक्कुटाशुकपालानि}}{\Afootnote{K H; \skt{कुक्कुटाण्डकपालानि}\ A}}\skt{ }\edtext{\skt{लसुनं}}{\Afootnote{K; \skt{लशुनं}\ A; \skt{लसूनं}\ H}}\skt{ }\edtext{\skt{कटुक}}{\Afootnote{A K; \skt{कटु}\ H}}\skt{ }\edtext{\skt{त्रिकं}}{\Afootnote{K H; \skt{त्रयम् ॥}\ A}}\skt{ ।}\pend

\pstart\versepar \ledsidenote{\passref{6.12.26}}\edtext{\skt{करञ्जवीजमेला च}}{\Afootnote{K H; \skt{करञ्जबीजमेला च}\ A}}\skt{ }\edtext{\skt{लेखनीयाञ्जनं}}{\Afootnote{K H; \skt{लेख्याञ्जनमिदं}\ A}}\skt{ स्मृतं ॥ }\pend
\pstart\versepar \skt{}\edtext{\skt{रक्तप्रसादनवता}}{\Afootnote{K H; \emph{om.}\ A}}\skt{ }\edtext{\skt{पूटपाके न}}{\Afootnote{K\textsuperscript{pc}; \skt{पुटपाकाव}\ A; \skt{पूटपाके}\ K\textsuperscript{ac}; \skt{पूटपाका च}\ H}}\skt{ }\edtext{\skt{साधिता}}{\Afootnote{K H; \skt{सानेन रक्तविस्रावणादिना ॥}\ A}}\skt{ ।}\pend

\pstart\versepar \ledsidenote{\passref{6.12.27}}\skt{सम्पादितस्य }\edtext{\skt{विध्ना}}{\Afootnote{K H; \skt{विधिना}\ A}}\skt{ कृत्स्नेन स्यन्दघातिना । }\pend
\pstart\versepar \skt{}\edtext{\skt{अनेनापाचरेच्चुक्रमव्रणं}}{\Afootnote{K; \skt{अनेनापहरेच्छुक्रमव्रणं}\ A; \skt{अनेनापचरेच्चुक्रमव्रणं}\ H}}\skt{ कुशलो भिषक् ।}\pend

\pstart\versepar \ledsidenote{\passref{6.12.28}}\edtext{\skt{उत्तानम्वा}}{\Afootnote{K H; \skt{उत्तानम}\ A}}\skt{ वगाढम्वा }\edtext{\skt{कर्षयेच्चापि}}{\Afootnote{K H; \skt{कर्कशं वापि}\ A}}\skt{ सव्रणं । }\pend
\pstart\versepar \skt{}\edtext{\skt{शिरीषवीजमरिचपिप्पलीसैन्धवैरपि}}{\lemma{\skt{शिरीषवीजमरिचपिप्पलीसैन्धवैर्}}\Afootnote{K H; \skt{शिरीषबीजमरिचपिप्पलीसैन्धवैर्}\ A}}\skt{ ।}\pend

\pstart \ledsidenote{\passref{6.12.29}}\edtext{\skt{शुक्रप्रघर्षणं}}{\Afootnote{K H; \skt{शुक्रस्य घर्षणं}\ A}}\skt{ कार्य }\edtext{\skt{मधुना}}{\Afootnote{K H; \skt{मथवा}\ A}}\skt{ }\edtext{\skt{सैन्धवे न तु}}{\Afootnote{K H; \skt{सैन्धवेन तु । कुर्यात्ताम्ररजःशङ्खशिलामरिचसैन्धवैः ॥}\ A}}\skt{ ।}\pend

\pstart \ledsidenote{\passref{6.12.30}}\edtext{\skt{कुर्यादञ्जनयोगा वा}}{\Afootnote{K H; \skt{अन्त्याद्द्विगुणितैरेभिरञ्जनं शुक्रनाशनम् । कुर्यादञ्जनयोगौ वा}\ A}}\skt{ सम्यक्श्लोकार्धिकाविमौ ।}\pend

\pstart\versepar \ledsidenote{\passref{6.12.31}}\edtext{\skt{शंखकोलास्तिकनक}}{\Afootnote{K; \skt{शङ्खकोलास्थिकतक}\ A; \skt{शंखकानोस्तिकनक}\ H}}\skt{ }\edtext{\skt{द्राक्षामधुकमाक्षिकैः}}{\Afootnote{K H; \skt{द्राक्सामधुकमाक्षिकैः ।}\ A}}\skt{ । }\pend
\pstart\versepar \skt{}\edtext{\skt{सूरादन्त्यार्णवमलैः}}{\Afootnote{K H; \skt{क्षौद्रदन्तार्णवमल}\ A}}\skt{ }\edtext{\skt{शिरीषकुसुमेन}}{\Afootnote{K H; \skt{शिरीषकुसुमैर}\ A}}\skt{ }\edtext{\skt{च}}{\Afootnote{K H; \skt{पि ॥}\ A}}\skt{ ।}\pend

\pstart\versepar \ledsidenote{\passref{6.12.32}}\edtext{\skt{क्षाराञ्जनं वा}}{\Afootnote{A K; \skt{क्षाराजनं वा}\ H}}\skt{ }\edtext{\skt{वितरेद्वलासग्रथितापहं}}{\Afootnote{K; \skt{वितरेद्बलासग्रथितापहम् ।}\ A; \skt{वितरेद्वल्-असग्रथितापहं ।}\ H}}\skt{ । }\pend
\pstart\versepar \skt{}\edtext{\skt{मुद्गाम्वा}}{\Afootnote{K H; \skt{मुद्गान्वा}\ A}}\skt{ }\edtext{\skt{निस्तुषां}}{\Afootnote{K H; \skt{निस्तुषा}\ A}}\skt{ }\edtext{\skt{पिष्ट्वा}}{\Afootnote{K H; \skt{न्भृष्टा}\ A}}\skt{ }\edtext{\skt{शंखक्षौद्रसितायुतां}}{\Afootnote{K H; \skt{न्शङ्खक्षौद्रसितायुतान् ॥}\ A}}\skt{ ।}\pend

\pstart\versepar \ledsidenote{\passref{6.12.33}}\edtext{\skt{मधुकसारं}}{\Afootnote{K H; \skt{मधूकसारं}\ A}}\skt{ मधुना }\edtext{\skt{योजयेद्चाञ्जने}}{\Afootnote{K H; \skt{योजयेच्चाञ्जने}\ A}}\skt{ सदा । }\pend
\pstart\versepar \skt{}\edtext{\skt{विभीतकास्थिमज्जा वा}}{\Afootnote{K H; \skt{बिभीतकास्थिमज्जा वा}\ A}}\skt{ }\edtext{\skt{सक्षौद्रा}}{\Afootnote{K; \skt{सक्षौद्रः}\ A H}}\skt{ }\edtext{\skt{शुक्रशान्तये}}{\Afootnote{K H; \skt{शुक्रनाशनः ॥}\ A}}\skt{ ॥}\pend

\pstart \ledsidenote{\passref{6.12.34}}\edtext{\skt{शंखास्य्न्त}}{\Afootnote{K; \skt{शङ्खशुक्ति}\ A; \skt{शंखान्त}\ H}}\skt{ }\edtext{\skt{मधुदाक्षा}}{\Afootnote{K H; \skt{मधुद्राक्षा}\ A}}\skt{ }\edtext{\skt{मयूतकतकानि वा}}{\Afootnote{K H; \skt{मधुकं कतकानि च । द्वित्वग्गते सशूले वा वातघ्नं तर्पणं हितम् ॥}\ A}}\skt{ ।}\pend

\pstart\versepar \ledsidenote{\passref{6.12.35}}\edtext{\skt{वंशजारूष्करौ}}{\Afootnote{K H; \skt{वंशजारुष्करौ}\ A}}\skt{ तालं }\edtext{\skt{नालिकेरञ्च}}{\Afootnote{K H; \skt{नारिकेलं च}\ A}}\skt{ तद्दहेत् । }\pend
\pstart\versepar \skt{क्षारवन्श्रावयेच्चूर्ण्णं }\edtext{\skt{करभातु}}{\Afootnote{K; \emph{om.}\ A; \skt{करभाम्तु}\ H}}\skt{ }\edtext{\skt{भावयेत्}}{\Afootnote{K H; \skt{भावयेत्करभास्थिजम् ॥}\ A}}\skt{ ।}\edtext{}{\lemma{\skt{तद्दहेत् । क्षारवन्श्रावयेच्चूर्ण्णं}}\Afootnote{K\textsuperscript{pc}; \skt{दाहयेत् । विस्राव्य क्सारवच्चूर्णं}\ A; \skt{द्दहेत् । क्षारवन्श्रावयेच्चूर्ण्णं}\ K\textsuperscript{ac}; \skt{तद्दहेत् ॥ क्षारवन्चयेच्चूर्ण्ण}\ H}}\pend

\pstart \ledsidenote{\passref{6.12.36}}\edtext{\skt{वहुशोञ्जनयेतत्स्याच्छुक्रवैवर्ण्णनाशनं}}{\Afootnote{K H; \skt{बहुशोऽञ्जनमेतत्स्याच्छुक्रवैवर्ण्यनाशनम् । अजकां पार्श्वतो विद्ध्वा सूच्या विस्राव्य चोदकम् ॥}\ A}}\skt{ ।}\pend

\pstart\versepar \ledsidenote{\passref{6.12.38}}\edtext{\skt{सशोफश्चाप्यशोफञ्च द्वौ}}{\lemma{\skt{चाप्यशोफञ्च द्वौ}}\Afootnote{K H; \skt{चाप्यशोफश्च द्वौ}\ A}}\skt{ पाकौ यौ प्रकीर्तितौ । }\pend
\pstart\versepar \skt{}\edtext{\skt{शारीरा}}{\Afootnote{K H; \skt{स्ने}\ A}}\skt{ }\edtext{\edtext{\skt{हितास्तत्र}}{\lemma{\skt{हितास्}}\Afootnote{K; \skt{हस्वेदोपपन्नस्य}\ A; \skt{हितास्ते}\ H}}}{\lemma{\skt{तत्र}}\Afootnote{A K; \skt{त्र}\ H}}\skt{ }\edtext{\skt{विद}}{\Afootnote{K H; \skt{विद्ध्वा}\ A}}\skt{ सिराम्भिषक् ।}\pend

\pstart\versepar \ledsidenote{\passref{6.12.39}}\edtext{\skt{सेकाश्च्योतननस्यानि}}{\Afootnote{K H; \skt{सेकाश्च्योतननस्यानि}\ A}}\skt{ पुटपाकाञ्श्च योजयेत् । }\pend
\pstart\versepar \skt{}\edtext{\skt{सर्वतश्चाप्यशुद्धस्य}}{\lemma{\skt{चाप्यशुद्धस्य}}\Afootnote{K H; \skt{चापि शुद्धस्य}\ A}}\skt{ }\edtext{\skt{कर्तव्यमिदमजनं}}{\lemma{\skt{इदमजनं}}\Afootnote{K; \skt{इदमञ्जनम्}\ A H}}\skt{ ।}\edtext{}{\lemma{\skt{योजयेत् । सर्वतश्}}\Afootnote{K H; \skt{कारयेत् । सर्वतश्}\ A}}\pend

\pstart\versepar \ledsidenote{\passref{6.12.40}}\edtext{\skt{ताम्रपात्रस्थितम्मांसं}}{\lemma{\skt{मांसं}}\Afootnote{K H; \skt{मासं}\ A}}\skt{ सर्पिः सैन्धवसंयुतं । }\pend
\pstart\versepar \skt{मैरेयम्वा }\edtext{\skt{विधातव्यंदध्युत्तरकमेव}}{\lemma{\skt{विधातव्यंदध्युत्तरकम्}}\Afootnote{K; \skt{पि दध्येवं दध्युत्तरकम्}\ A; \skt{विधातव्यदध्युत्तरकम्}\ H}}\skt{ वा ।}\pend

\pstart\versepar \ledsidenote{\passref{6.12.41}}\edtext{\skt{घृतां}}{\Afootnote{K; \skt{घृतं}\ A H}}\skt{ }\edtext{\skt{कांसमलोपेतां}}{\Afootnote{K; \skt{कांस्यमलोपेतं}\ A; \skt{कांसमलोपेतं}\ H}}\skt{ }\edtext{\skt{स्तन्यम्वापिं}}{\lemma{\skt{वापिं}}\Afootnote{K H; \skt{वापि}\ A}}\skt{ }\edtext{\skt{ससन्धवं}}{\Afootnote{K H; \skt{ससैन्धवम् ।}\ A}}\skt{ । }\pend
\pstart\versepar \skt{मधूकसारम्मधुना तुल्यांशं }\edtext{\skt{गैरिकेन}}{\Afootnote{K H; \skt{गैरिकेण}\ A}}\skt{ वा ।}\pend

\pstart\versepar \ledsidenote{\passref{6.12.42}}\skt{सर्पिः सैन्धवताम्राणि योषित्स्तन्ययुतानि वा । }\pend
\pstart\versepar \skt{}\edtext{\skt{दाडिमैरावताश्मन्तकोलाम्लमधुसैन्धवैः}}{\Afootnote{K H; \skt{दाडिमारेवताश्मन्तकोलाम्लैश्च ससैन्धवाम् ।}\ A}}\skt{ । रसक्रियाम्वा }\edtext{\skt{विपचेत्सम्यक्पाकजिघांसया}}{\Afootnote{K H; \skt{वितरेत्सम्यक्पाकजिघांसया ॥}\ A}}\skt{ ।}\pend

\pstart\versepar \ledsidenote{\passref{6.12.43}}\edtext{\skt{माससैन्धवसंयुक्तं}}{\Afootnote{K; \skt{मासं सैन्धवसंयुक्तं}\ A; \skt{मांसं सैन्धवसंयुक्तं}\ H}}\skt{ स्थितं सर्पिषि नागरं ॥ }\pend
\pstart\versepar \skt{}\edtext{\skt{आश्च्योतनाञ्जने}}{\Afootnote{K H; \skt{आश्च्योतनाञ्जनं}\ A}}\skt{ }\edtext{\skt{योज्यामवलाक्षीरसंयुतं}}{\Afootnote{K H; \skt{योज्यमबलाक्षीरसंयुतम् ॥}\ A}}\skt{ ॥}\pend

\pstart\versepar \ledsidenote{\passref{6.12.44}}\edtext{\skt{पुष्पं}}{\Afootnote{K H; \skt{जात्याः पुष्पं}\ A}}\skt{ }\edtext{\skt{मालत्याः}}{\Afootnote{K H; \emph{om.}\ A}}\skt{ सैन्धवं शृङ्गवेरं }\edtext{\skt{सारो}}{\Afootnote{K; \skt{कृष्णा}\ A; \skt{सा}\ H}}\skt{ }\edtext{\skt{वैडङ्गःः}}{\Afootnote{K; \skt{बीजं की}\ A; \emph{om.}\ H}}\skt{ }\edtext{\skt{पिप्पलीतण्डूअश्च}}{\Afootnote{K; \skt{टशत्रोश्च}\ A; \skt{रवैरं ः पिप्पलीतण्डूअश्च ।}\ H}}\skt{ । }\pend
\pstart\versepar \skt{}\edtext{\skt{एतत्पिष्ट्वात्क्षौद्रयुक्त्याञ्जनार्थं}}{\Afootnote{K H; \skt{सारम् । एतत्पिष्टं नेत्रपाकेऽञ्जनार्थं}\ A}}\skt{ }\edtext{\skt{सम्य}}{\Afootnote{K H; \skt{क्षौद्रो}\ A}}\skt{ }\edtext{\skt{युक्तं}}{\Afootnote{K H; \skt{पेतं}\ A}}\skt{ }\edtext{\skt{निर्विशंकः}}{\Afootnote{K; \skt{निर्विशङ्कं}\ A H}}\skt{ }\edtext{\skt{प्रदद्यात्}}{\Afootnote{K; \skt{प्रयोज्यम् ॥}\ A; \skt{विदद्यात् ॥}\ H}}\skt{ ॥}\pend

\pstart\versepar \ledsidenote{\passref{6.12.45}}\skt{पूयालसे शोणितमोक्षणं }\edtext{\skt{एव}}{\Afootnote{K H; \skt{च}\ A}}\skt{ }\edtext{\skt{विधिस्तथैवाप्युपनाहनश्च}}{\Afootnote{K H; \skt{हितं तथैवाप्युपनाहनं च ।}\ A}}\skt{ । }\pend
\pstart\versepar \skt{कृत्स्नो }\edtext{\skt{विधिश्चेक्षणपाकहन्ता}}{\Afootnote{K H; \skt{विधिश्चेक्षणपाकघाती}\ A}}\skt{ }\edtext{\skt{यथाभिधानं}}{\Afootnote{K H; \skt{यथाविधानं}\ A}}\skt{ भिषजा }\edtext{\skt{प्रयोज्यां}}{\Afootnote{K; \skt{प्रयोज्यः ॥}\ A H}}\skt{ ।}\pend

\pstart\versepar \ledsidenote{\passref{6.12.46}}\edtext{\skt{कासीससिन्धुप्रभवाजनेश्च}}{\Afootnote{K H; \skt{कासीससिन्धुप्रभवार्द्रकैस्तु}\ A}}\skt{ हितं भवेदञ्जनमेव चात्र । }\pend
\pstart\versepar \skt{}\edtext{\edtext{\edtext{\skt{क्षौद्रान्द्वितैरेभिरथैव}}{\lemma{\skt{क्षौद्रान्द्वितैर्}}\Afootnote{K; \skt{क्षौद्रान्वितैर्}\ A; \skt{क्षौद्रान्वितै}\ H}}}{\lemma{\skt{एभिर्}}\Afootnote{A K; \skt{रेभिर्}\ H}}\skt{ च}}{\lemma{\skt{अथैव च}}\Afootnote{K; \skt{अथोप}\ A; \skt{अथो पि}\ H}}\skt{ }\edtext{\skt{स्यादन्यत्र}}{\Afootnote{K; \skt{युञ्ज्यादन्यत्तु}\ A; \skt{यस्यादन्यत्र}\ H}}\skt{ ताम्रायसचूर्ण्णयुक्तैः ।}\pend

\pstart\versepar \ledsidenote{\passref{6.12.47}}\edtext{\skt{स्नेहादिना}}{\Afootnote{K H; \skt{स्नेहादिभिः}\ A}}\skt{ }\edtext{\skt{सम्यगुपास्य}}{\Afootnote{K H; \skt{सम्यगपास्य}\ A}}\skt{ दोषां }\edtext{\skt{तृप्तिं}}{\Afootnote{K H; \skt{स्तृप्तिं}\ A}}\skt{ }\edtext{\skt{विधिनेन}}{\Afootnote{K H; \skt{विधायाथ}\ A}}\skt{ }\edtext{\edtext{\skt{यथास्वमेव}}{\lemma{\skt{यथास्वम्}}\Afootnote{A K; \skt{यथास्व}\ H}}}{\lemma{\skt{एव}}\Afootnote{A K; \skt{व ।}\ H}}\skt{ । }\pend
\pstart\versepar \skt{}\edtext{\edtext{\skt{अक्लिन्नवर्त्मानमुपक्रमेत}}{\lemma{\skt{अक्लिन्नवर्त्मानम्}}\Afootnote{K H; \skt{प्रक्लिन्नवर्त्मानम्}\ A}}}{\lemma{\skt{उपक्रमेत}}\Afootnote{A K; \skt{उपल्रमेत}\ H}}\skt{ }\edtext{\skt{सेकाञ्जनाश्च्योतननस्यधूमैः}}{\Afootnote{K H; \skt{सेकाञ्जनाश्च्योतननस्यधूमैः ॥}\ A}}\skt{ ।}\pend

\pstart\versepar \ledsidenote{\passref{6.12.48}}\skt{मुस्ताहरिद्रामधुकप्रियंगु }\edtext{\skt{सिद्धार्थलोध्रोसितशारिवाभिः}}{\Afootnote{K; \skt{सिद्धार्थरोध्रोत्पलसारिवाभिः ।}\ A; \skt{सिद्धार्थरोध्रोसितशारिवाभिः ।}\ H}}\skt{ । }\pend
\pstart\versepar \skt{}\edtext{\skt{क्षुण्णाभिराश्च्योतनमेव}}{\lemma{\skt{आश्च्योतनम्}}\Afootnote{K H; \skt{आश्च्योतनम्}\ A}}\skt{ }\edtext{\edtext{\skt{कार्यमत्राञ्जनं}}{\lemma{\skt{कार्यम्}}\Afootnote{K H; \skt{कार्य}\ A}}}{\lemma{\skt{अत्राञ्जनं}}\Afootnote{K H; \skt{मत्राञ्जनं}\ A}}\skt{ }\edtext{\skt{चाञ्जनमाक्षिक}}{\Afootnote{K H; \skt{चाञ्जनमाक्षिकं}\ A}}\skt{ स्यात् ॥}\pend

\pstart\versepar \ledsidenote{\passref{6.12.49}}\edtext{\skt{पत्रं}}{\Afootnote{A K; \skt{पत्र}\ H}}\skt{ }\edtext{\skt{फलं}}{\Afootnote{A K; \skt{मूलञ्}\ H}}\skt{ }\edtext{\skt{चामलकस्य}}{\Afootnote{A K; \skt{चामलकसुअ}\ H}}\skt{ }\edtext{\skt{पक्वा}}{\Afootnote{K H; \skt{पक्त्वा}\ A}}\skt{ }\edtext{\skt{क्रियाम्विदध्यादथवाञ्जनार्थं}}{\lemma{\skt{विदध्यादथवाञ्जनार्थं}}\Afootnote{K; \skt{विदध्यादथवाञ्जनार्थे ।}\ A; \skt{विदध्यादथवाञ्जनार्थं ।}\ H}}\skt{ । }\pend
\pstart\versepar \skt{}\edtext{\skt{वंशस्य चा}}{\Afootnote{K H; \skt{वंशस्य}\ A}}\skt{ }\edtext{\skt{मूलेन}}{\Afootnote{K H; \skt{भूलेन}\ A}}\skt{ }\edtext{\skt{सक्रियां तु}}{\Afootnote{K H; \skt{रसक्रियां वा}\ A}}\skt{ वर्तीकृतां ताम्रकपालपक्वां ।}\pend

\pstart\versepar \ledsidenote{\passref{6.12.50}}\skt{रसक्रियां वा }\edtext{\skt{विफलाविपक्वां}}{\Afootnote{K; \skt{त्रिफलाविपक्वां}\ A; \skt{विफलाविपक्वं}\ H}}\skt{ पलाशपुष्पैः }\edtext{\skt{करवीरजैर्वा}}{\Afootnote{K H; \skt{स्वरमञ्जरेर्वा ।}\ A}}\skt{ । }\pend
\pstart\versepar \skt{}\edtext{\skt{पिष्ट्वान्च}}{\Afootnote{K; \skt{पिष्ट्वा छ}\ A; \skt{पिष्ट्वाल}\ H}}\skt{ }\edtext{\skt{जानां}}{\Afootnote{K H; \skt{गल्याः}\ A}}\skt{ पयसामलम्वा कांसस्य }\edtext{\skt{दग्धा}}{\Afootnote{K H; \skt{दग्ध्वा}\ A}}\skt{ सह तान्तवेन ।}\pend

\pstart \ledsidenote{\passref{6.12.51}}\edtext{\skt{प्रत्यञ्जनतत्मरिचैरूपेतं}}{\Afootnote{K; \skt{प्रत्यञ्जनं तन्मरिचैरुपेतं}\ A; \skt{प्रत्यञ्जनतत्मरिचैरूपेतं}\ H}}\skt{ }\edtext{\skt{चूर्ण्णताम्रस्य}}{\Afootnote{K H; \skt{चूर्णेन ताम्रस्य}\ A}}\skt{ }\edtext{\skt{सहाभियुञ्ज्यात्}}{\Afootnote{K H; \skt{सहोपयोज्यम् । समुद्रफेनं लवणोत्तमं च शङ्खोऽथ मुद्गो मरिचं च शुक्लम् ॥}\ A}}\skt{ ।}\pend

\pstart \ledsidenote{\passref{6.12.52}}\edtext{\edtext{\skt{प्रक्लिन्नवर्त्मान्यरेत}}{\lemma{\skt{प्रक्लिन्नवर्त्मान्य्}}\Afootnote{K; \skt{चूर्णाञ्जनं जाड्यमथापि कण्डूमक्लिन्नवर्त्मान्युपहन्ति शीघ्रम् । प्रक्लिन्नवर्त्मन्य्}\ A; \skt{प्रक्लिन्नवर्त्मा}\ H}}}{\lemma{\skt{अरेत}}\Afootnote{K; \skt{अपि चैत}\ A; \skt{न्यरेत}\ H}}\skt{ एव योगाः प्रयोज्याः }\edtext{\skt{प्रसमीक्ष्य}}{\Afootnote{K; \skt{च समीक्ष्य}\ A; \skt{समीक्ष्य}\ H}}\skt{ }\edtext{\skt{दोषानिति}}{\Afootnote{K; \skt{दोषम् ॥}\ A; \skt{दोस्̤आनिति ॥ ॥}\ H}}\skt{ ॥ ॥}\pend

\pstart \ledsidenote{\passref{6.12.trailer}}\edtext{\skt{शालाक्ये}}{\Afootnote{K; \skt{इति सुश्रुतसंहितायामुत्तरतन्त्रान्तर्गते शालाक्यतन्त्रे रक्ताभिष्यन्दप्रतिषेधो नाम द्वादशोऽध्यायः ॥१२ ॥}\ A; \skt{शालाक्यस्थाने रक्ताभिष्यन्दप्रतिषेधो द्वादशमोऽध्यायः ॥ ॥}\ H}}\skt{ २ ॥ ० ॥}\pend

\endnumbering
\beginnumbering
\eledsection*{Adhyāya 13}
\pstart \ledsidenote{\passref{6.13.1}}\skt{अथातो }\edtext{\skt{लेख्यरोगप्रतिषेधमध्यायं}}{\lemma{\skt{अध्यायं}}\Afootnote{K; \emph{om.}\ A H}}\skt{ }\edtext{\skt{व्याख्यास्यामः}}{\Afootnote{A H; \skt{व्या व्याख्यास्यामः ॥}\ K}}\skt{ ॥}\pend

\pstart\versepar \ledsidenote{\passref{6.13.3}}\skt{नव ये }\edtext{\skt{लेख्यसाध्यास्तु}}{\Afootnote{K H; \skt{ऽभिहिता लेख्याः}\ A}}\skt{ }\edtext{\skt{सामान्यस्तेष्वसम्विधिः}}{\lemma{\skt{असम्विधिः}}\Afootnote{K H; \skt{अयं विधिः ।}\ A}}\skt{ । }\pend
\pstart\versepar \skt{स्निग्धवान्तविरिक्तस्य }\edtext{\skt{निवातातपसंसदि}}{\Afootnote{K H; \skt{निवातातपसद्मनि ॥}\ A}}\skt{ ।}\pend

\pstart \ledsidenote{\passref{6.13.4}}\edtext{\skt{आप्तैर्दृढं}}{\lemma{\skt{आप्तैर्}}\Afootnote{A K; \skt{आप्तै}\ H}}\skt{ गृहीतस्य }\edtext{\skt{वेश्मन्युतानशायिनः}}{\Afootnote{K; \skt{वेश्मन्युत्तानशायिनः ।)}\ A; \skt{वेश्मन्युत्तानशायिनः ।}\ H}}\skt{ । सुखोदकप्रतप्तेन वाससा सुसमाहितः । स्वेदयेद्वर्त्मनिर्भुज्य }\edtext{\skt{वामाङ्गुष्ठाङ्गुलिस्थितं}}{\Afootnote{K H; \skt{वामाङ्गुष्ठाङ्गुलिस्ह्तितम् ॥}\ A}}\skt{ ।}\pend

\pstart \ledsidenote{\passref{6.13.5.1}}\edtext{\skt{वर्त्व्यावेश्च त}}{\Afootnote{K; \skt{वर्त्मनावेश्च तं}\ H}}\skt{ प्लोतं }\edtext{\skt{सस्यस्न्नं}}{\Afootnote{K; \skt{सस्यस्थिन्न}\ H}}\skt{ }\edtext{\skt{ततो}}{\Afootnote{K; \skt{लतो}\ H}}\skt{ लिखेत् ।}\pend

\pstart \ledsidenote{\passref{6.13.6}}\edtext{\skt{शस्त्रेण}}{\Afootnote{K H; \skt{ततः प्रमृज्य प्लोतेन वर्त्म शस्त्रपदाङ्कितम् । लिखेच्छस्त्रेण}\ A}}\skt{ }\edtext{\skt{पत्रैरथवा}}{\lemma{\skt{अथवा}}\Afootnote{K H; \skt{वा}\ A}}\skt{ ततो रक्ते स्थिते पुनः ।}\pend

\pstart \ledsidenote{\passref{6.13.7}}\edtext{\skt{स्वेद्यसैन्धवचूर्ण्णेन}}{\Afootnote{K H; \skt{स्विन्नं मनोह्वाकासीसव्योषार्द्राञ्जनसैन्धवैः । श्लक्ष्णपिष्टैः समाक्षीकैः}\ A}}\skt{ }\edtext{\skt{प्रतिसार्यसुखाम्वुना}}{\Afootnote{K H; \skt{प्रतिसार्योष्णवारिणा ॥}\ A}}\skt{ ।}\pend

\pstart \ledsidenote{\passref{6.13.8}}\skt{प्रक्षाल्य }\edtext{\skt{सिक्तं}}{\Afootnote{K H; \emph{om.}\ A}}\skt{ हविषा }\edtext{\skt{व्रणवत्न्समुपाचरेत्}}{\Afootnote{K H; \skt{सिक्तं व्रणवत्समुपाचरेत् । स्वेदावपीडप्रभृतींस्त्र्यहादूर्ध्वं प्रयोजयेत् ॥}\ A}}\skt{ ।}\pend

\pstart \ledsidenote{\passref{6.13.9}}\edtext{\skt{व्यासतस्ते}}{\Afootnote{K H; \skt{व्यासतस्ते}\ A}}\skt{ }\edtext{\skt{तदखिलं}}{\Afootnote{K H; \skt{समुद्दिष्टं}\ A}}\skt{ विधानं लेख्य }\edtext{\skt{वर्त्मणि}}{\Afootnote{K; \skt{कर्मणि । असृगास्रावरहितं कण्डूशोफविवर्जितम् ॥}\ A; \skt{वर्त्मनि ॥}\ H}}\skt{ ।}\pend

\pstart\versepar \ledsidenote{\passref{6.13.10}}\edtext{\skt{नखस्नायुनिभं}}{\Afootnote{K; \skt{समं नखनिभं}\ A; \skt{नखस्रायुनिभं}\ H}}\skt{ वर्त्म लिखितं सम्यगिष्यते । }\pend
\pstart\versepar \skt{}\edtext{\skt{रक्तमक्षि}}{\Afootnote{K H; \skt{रक्तमक्षि}\ A}}\skt{ स्रवे }\edtext{\skt{स्कन्नं}}{\Afootnote{K H; \skt{त्स्कन्नं}\ A}}\skt{ }\edtext{\skt{तिमिरं}}{\Afootnote{K; \skt{क्षताच्छस्त्रकृताद्ध्रुवम् ॥}\ A; \skt{निमिरं}\ H}}\skt{ }\edtext{\skt{व्याध्यनिर्जयः}}{\Afootnote{K; \emph{om.}\ A; \skt{व्याधिनिर्जयः ॥}\ H}}\skt{ ।}\pend

\pstart \ledsidenote{\passref{6.13.11}}\edtext{\edtext{\skt{वर्त्मञ्च}}{\lemma{\skt{वर्त्मञ्}}\Afootnote{K; \skt{रागशोफपरिस्रावास्तिमिरं व्याध्यनिर्जयः । वर्त्म}\ A; \skt{वर्त्म}\ H}}}{\lemma{\skt{च}}\Afootnote{K H; \skt{श्यावं}\ A}}\skt{ गुरु }\edtext{\skt{शूनञ्च}}{\Afootnote{K H; \skt{स्तब्धं}\ A}}\skt{ कण्डूहर्षोपदेहवत् ।}\pend

\pstart \ledsidenote{\passref{6.13.12}}\edtext{\skt{दुर्विरिक्तं}}{\Afootnote{K; \skt{नेत्रपाकमुदीर्णं वा कुर्वीताप्र}\ A; \skt{दुर्व्विरिक्त}\ H}}\skt{ }\edtext{\skt{विजानीयादतत्सिन्नं}}{\Afootnote{K; \skt{तिकारिणः । एतद्दुर्लिखितं ज्ञेयं स्नेहयित्वा}\ A; \skt{विजानीयादेतत्सिन्नं}\ H}}\skt{ पुनर्लिखेत् ।}\pend

\pstart\versepar \ledsidenote{\passref{6.13.13}}\edtext{\skt{व्यावर्ततेवर्त्म}}{\Afootnote{K H; \skt{व्यावर्तते यदा}\ A}}\skt{ }\edtext{\skt{यदा}}{\Afootnote{K H; \skt{वर्त्म}\ A}}\skt{ }\edtext{\skt{पक्ष्ं}}{\Afootnote{K; \skt{पक्ष्म}\ A; \skt{पक्ष}\ H}}\skt{ चापि }\edtext{\skt{विमुञ्चति}}{\Afootnote{K H; \skt{विमुह्यति ।}\ A}}\skt{ । }\pend
\pstart\versepar \skt{}\edtext{\skt{वद्धशुरूक्तव }\edtext{\skt{वश्चैव}}{\lemma{\skt{चैव}}\Afootnote{K H; \skt{लं}\ A}}}{\lemma{\skt{वद्धशुरूक्तव वश्}}\Afootnote{K; \skt{स्यात्सरुक्स्रावबहु}\ A; \skt{वद्धशुरूक्वं}\ H}}\skt{ }\edtext{\skt{तदतिस्रावितं}}{\Afootnote{\skt{तदतिस्रावितं}\ A K H}}\skt{ विदुः ।}\pend

\pstart \ledsidenote{\passref{6.13.14ab}}\edtext{\skt{स्नेहस्वेदादिरिष्टं}}{\Afootnote{K; \skt{स्नेहस्वेदादिरिष्टः}\ A}}\skt{ }\edtext{\skt{स्यात्क्रमस्तत्रानिलापहः}}{\Afootnote{K; \skt{स्यात्क्रमस्तत्रानिलापहः ।}\ A}}\skt{ ।}\pend

\pstart \ledsidenote{\passref{6.13.14cd}}\edtext{\skt{वर्त्माववद्धक्लिष्टं च}}{\Afootnote{K H; \skt{वर्त्मावबन्धं क्लिष्टं च}\ A}}\skt{ }\edtext{\skt{वहलं}}{\Afootnote{K H; \skt{बहलं}\ A}}\skt{ यच्च कीर्तितं ।}\pend

\pstart\versepar \ledsidenote{\passref{6.13.15}}\edtext{\skt{पोत्थकीश्चापि}}{\lemma{\skt{पोत्थकीश्}}\Afootnote{A K; \skt{पौत्थकी}\ H}}\skt{ }\edtext{\skt{विलिखेत्प्रच्छयित्वाग्रतः}}{\Afootnote{\skt{अवलिखेत्प्रच्छयित्वाग्रतः}\ A; \skt{अविलिखेत्प्रच्छयित्वाग्रतः}\ K; \skt{विलिखेत्प्रच्छयित्वा यतः}\ H}}\skt{ शनैः । }\pend
\pstart\versepar \skt{समं }\edtext{\skt{लिखेत}}{\Afootnote{K; \skt{लिखेत्तु}\ A; \skt{लिखेद्धि}\ H}}\skt{ }\edtext{\skt{मनिमाञ्श्यावकर्दमवर्त्मनी}}{\Afootnote{K; \skt{मेधावी श्यावकर्दमवर्त्मनी ॥}\ A; \skt{मनिमाञ्श्यावकर्द्दमवर्त्मनीं ।}\ H}}\skt{ ।}\pend

\pstart\versepar \ledsidenote{\passref{6.13.16}}\edtext{\skt{कुभीकिनींशर्कराञ्च}}{\Afootnote{K; \skt{कुम्भीकिनीं शर्करां च}\ A; \skt{कुभीकिंनीशर्क्कराञ्च}\ H}}\skt{ }\edtext{\skt{तथैवोत्सञ्गिनीम}}{\Afootnote{K; \skt{तथैवोत्सङ्गिनीमपि}\ A H}}\skt{ }\pend
\pstart\versepar \skt{}\edtext{\skt{धीमं}}{\Afootnote{K; \emph{om.}\ A; \skt{धीम}\ H}}\skt{ }\edtext{\skt{विकल्प्य}}{\Afootnote{K; \skt{कल्प}\ A; \skt{न्विकल्प्य}\ H}}\skt{ }\edtext{\skt{शस्त्रेण}}{\Afootnote{K H; \skt{यित्वा तु शस्त्रेण}\ A}}\skt{ }\edtext{\skt{निखेत्पश्चादतन्द्रितः}}{\Afootnote{K; \skt{लिखेत्पश्चादतन्द्रितः ॥}\ A; \skt{निखेत्पश्चादतन्द्रितः ।}\ H}}\skt{ ।}\pend

\pstart\versepar \ledsidenote{\passref{6.13.17}}\edtext{\edtext{\skt{भवेयर्वर्त्मसु}}{\lemma{\skt{भवेयर्}}\Afootnote{K; \skt{भवेयु}\ A; \skt{भेवेयर्}\ H}}\skt{ च याः}}{\lemma{\skt{वर्त्मसु च याः}}\Afootnote{K; \skt{र्वर्त्मसु च याः}\ A; \skt{वर्त्मसु च या}\ H}}\skt{ पिडकाः कठिना भृशं । }\pend
\pstart\versepar \skt{}\edtext{\skt{ह्रस्वास्ताम्राश्च ताः}}{\Afootnote{\skt{ह्रस्वास्ताम्राश्च ताः}\ A H K}}\skt{ पक्वा }\edtext{\skt{भिन्द्याद्भित्वा}}{\Afootnote{\skt{भिन्द्याद्भिन्ना}\ A; \skt{भिन्द्याद्भित्वा}\ K H}}\skt{ }\edtext{\skt{लिखेदपि}}{\Afootnote{K H; \skt{लिखेदपि ॥}\ A}}\skt{ ।}\pend

\pstart\versepar \ledsidenote{\passref{6.13.18}}\edtext{\skt{तरुण्यो}}{\Afootnote{K H; \skt{तरुणीश्}\ A}}\skt{ }\edtext{\skt{वाल्पसंरम्भाः}}{\Afootnote{K H; \skt{चाल्पसंरम्भा}\ A}}\skt{ पिडका }\edtext{\skt{वाह्यवर्त्मजाः}}{\Afootnote{K H; \skt{बाह्यवर्त्मजाः ।}\ A}}\skt{ । }\pend
\pstart\versepar \skt{}\edtext{\skt{विदित्वैताम्प्रशमयेत्स्वेदालेपनशोधनैः}}{\lemma{\skt{विदित्वैताम्}}\Afootnote{K H; \skt{विदित्वैताः}\ A}}\skt{ ।}\pend

\pstart \ledsidenote{\passref{6.13.18-n1}}\skt{Tहे Nएपलेसे वेर्सिओन्हस्नो छप्तेर्ब्रेअखेरे. Nएपलेसे छप्तेर् १३ चोन्तिनुएस्तो स्नेहपूर्वकम्, इ.ए., वुल्गते ६.१४.११अब्}\pend

\endnumbering
\beginnumbering
\eledsection*{Adhyāya 14}
\pstart\versepar \ledsidenote{\passref{6.14.3}}\skt{स्वेदयित्वा }\edtext{\skt{विसग्रन्थिपक्वं}}{\Afootnote{K; \skt{बिसग्रन्थिं छिद्राण्य्}\ A; \skt{विसग्रन्थिम्पक्वन्}\ H}}\skt{ }\edtext{\skt{तस्य}}{\Afootnote{K H; \skt{अस्य}\ A}}\skt{ निराशयं । }\pend
\pstart\versepar \skt{}\edtext{\skt{भित्वाच्छिद्राणि}}{\Afootnote{K H; \skt{पक्वं भित्त्वा तु}\ A}}\skt{ शस्त्रेण }\edtext{\skt{सैन्धवेनावचूर्णयेत्}}{\Afootnote{K H; \skt{सिन्धवेनावचूर्णयेत् ॥}\ A}}\skt{ ।}\pend

\pstart \ledsidenote{\passref{6.14.4}}\edtext{\skt{कासीसमार्गधीमुश्चीनेपाल्येलायुतेन तु}}{\Afootnote{K; \skt{कासीसमागधीपुष्पनेपाल्येलायुतेन तु । ततः क्षौद्रघृतं दत्त्वा सम्यग्बन्धमथाचरेत् ॥}\ A; \skt{कासीसमागधाचूर्णेन नेपाल्येलातथैव च ॥}\ H}}\skt{ ।}\pend

\pstart\versepar \ledsidenote{\passref{6.14.5}}\edtext{\skt{रोचनक्षारतुत्थानि}}{\Afootnote{K H; \skt{रोचनाक्षारतुत्थानि}\ A}}\skt{ पिप्पल्यः क्षौद्रमेव च । }\pend
\pstart\versepar \skt{}\edtext{\edtext{\skt{प्रतिसारणमेकैकं}}{\lemma{\skt{प्रतिसारणम्}}\Afootnote{K H; \skt{प्रतिसारण}\ A}}}{\lemma{\skt{एकैकं}}\Afootnote{K H; \skt{मेकैकं}\ A}}\skt{ भिन्ने लगण इष्यते ।}\pend

\pstart \ledsidenote{\passref{6.14.6}}\edtext{\edtext{\skt{महान्त्यथ}}{\lemma{\skt{महान्त्य्}}\Afootnote{K H; \skt{महत्य्}\ A}}\skt{ च}}{\lemma{\skt{अथ च}}\Afootnote{K H; \skt{अपि च}\ A}}\skt{ युञ्जीत }\edtext{\skt{क्षाराग्नि}}{\Afootnote{K H; \skt{क्षाराग्नी}\ A}}\skt{ }\edtext{\skt{विधिकोविदः}}{\Afootnote{K H; \skt{विधिकोविदः । स्विन्नां भिन्नां विनिष्पीड्य भिषगञ्जननामिकाम् ॥}\ A}}\skt{ ।}\pend

\pstart \ledsidenote{\passref{6.14.7}}\edtext{\skt{रसाञ्जनमधुभ्यान्तु}}{\Afootnote{K H; \skt{शिलैलानतसिन्धूत्थैः सक्षौद्रैः प्रतिसारयेत् । रसाञ्जनमधुभ्यां तु}\ A}}\skt{ }\edtext{\skt{भिन्ना}}{\Afootnote{K H; \emph{om.}\ A}}\skt{ }\edtext{\skt{चाञ्जनेनामिकां}}{\Afootnote{K; \skt{भित्त्वा वा शस्त्रकर्मवित् ॥}\ A; \skt{मञ्जनेनामिकां ॥}\ H}}\skt{ ।}\pend

\pstart\versepar \ledsidenote{\passref{6.14.8}}\edtext{\skt{प्रतिसायाञ्जने}}{\Afootnote{K H; \skt{प्रतिसार्याञ्जनैर्}\ A}}\skt{ }\edtext{\skt{युञ्ज्यान्क्रियां}}{\Afootnote{K H; \skt{युञ्ज्या दुष्णैर्}\ A}}\skt{ }\edtext{\skt{लवणवर्जितां}}{\Afootnote{K H; \skt{दीपशिखोद्भवैः ।}\ A}}\skt{ ॥ }\pend
\pstart\versepar \skt{सम्यक्स्विन्ने }\edtext{\skt{क्रिमिग्रन्थौ}}{\Afootnote{K H; \skt{कृमिग्रन्थौ}\ A}}\skt{ भिन्ने स्यात्प्रतिसारणं ।}\pend

\pstart\versepar \ledsidenote{\passref{6.14.9}}\skt{त्रिफलातुथकासीस सैन्धवैः }\edtext{\skt{सरसाञ्जनैः}}{\Afootnote{K H; \skt{च र}\ A}}\skt{ । }\pend
\pstart\versepar \skt{}\edtext{\skt{रसक्रियां}}{\Afootnote{K\textsuperscript{pc} H; \skt{स}\ A}}\skt{ }\edtext{\skt{क्रिमिग्रन्थै}}{\Afootnote{K\textsuperscript{pc} H; \skt{क्रि}\ A}}\skt{ भिन्न स्यात्प्रतिसारणं }\pend
\pstart\versepar \skt{}\edtext{\skt{भित्तोपनाहं}}{\Afootnote{K H; \skt{भित्त्वोपनाहं}\ A}}\skt{ }\edtext{\skt{कफजं}}{\Afootnote{A K; \skt{कजं}\ H}}\skt{ पिप्पलीमधुसैन्धवैः ।}\edtext{}{\lemma{\skt{रसक्रियां… भि}}\Afootnote{\emph{om.}\ K\textsuperscript{ac}}}\edtext{}{\lemma{\skt{मिग्रन्थै… प्रतिसा}}\Afootnote{\emph{om.}\ A}}\pend

\pstart\versepar \ledsidenote{\passref{6.14.10}}\edtext{\skt{विलिखेल्मंडलाग्रेण}}{\Afootnote{K H; \skt{लेखयेन्मण्डलाग्रेण}\ A}}\skt{ }\edtext{\skt{प्रच्छयेद्वा}}{\Afootnote{K H; \skt{समन्तात्प्रच्छये}\ A}}\skt{ }\edtext{\skt{समन्ततः}}{\Afootnote{K H; \skt{दपि ।}\ A}}\skt{ । }\pend
\pstart\versepar \skt{संस्नेह्य पत्रभङ्गैश्च स्वेदयित्वा }\edtext{\skt{यथावलं}}{\Afootnote{K H; \skt{यथासुखम् ॥}\ A}}\skt{ ।}\pend

\pstart\versepar \ledsidenote{\passref{6.14.11}}\skt{आपाका }\edtext{\skt{विधिनोक्तेन}}{\Afootnote{K; \skt{द्विधिनोक्तेन}\ H}}\skt{ पञ्चभेद्यानुपाचरेत् । }\pend
\pstart\versepar \skt{}\edtext{\skt{सर्वेष्वतेषु}}{\lemma{\skt{अतेषु}}\Afootnote{K H; \skt{तेषु}\ A}}\skt{ }\edtext{\skt{विहितं}}{\Afootnote{A K; \skt{च हितं}\ H}}\skt{ विधानं }\edtext{\skt{स्नेहपूर्वकमिति}}{\lemma{\skt{इति}}\Afootnote{K H; \skt{संपक्वे प्रयतो भूत्वा कुर्वीत व्रणरोपणम् ॥}\ A}}\skt{ ॥}\edtext{}{\lemma{\skt{… स}}\Afootnote{\emph{om.}\ A}}\pend

\pstart \ledsidenote{\passref{6.14.trailer}}\edtext{\skt{शालाक्ये १३ ॥ {\symfont ❈}}}{\Afootnote{K; \skt{इति सुश्रुतसंहितायामुत्तरतन्त्रान्तर्गते शालाक्यतन्त्रे भेद्यरोगप्रतिषेधो नाम चतुर्दशोऽध्यायः ॥१४ ॥}\ A; \skt{शालाक्यस्थाने लेह्यरोगप्रतिषेधस्त्रयोदशमोऽध्यायः ॥ ॥}\ H}}\skt{ ॥}\pend

\endnumbering
\beginnumbering
\eledsection*{Adhyāya 16}
\pstart \ledsidenote{\passref{6.16.1}}\skt{अथातः }\edtext{\skt{पक्ष्मगतरोगप्रतिषेधं}}{\Afootnote{K H; \skt{पक्ष्मकोपप्रतिषेधं}\ A}}\skt{ }\edtext{\skt{व्याख्यास्यामः}}{\Afootnote{A H; \skt{व्या व्याख्यास्यामः ॥}\ K}}\skt{ ॥}\pend

\pstart\versepar \ledsidenote{\passref{6.16.3}}\skt{याप्यस्तु यो }\edtext{\skt{वर्त्मभवो विः}}{\Afootnote{K; \skt{वर्त्मभवो वि}\ A H}}\skt{ }\edtext{\skt{पक्ष्मप्रकोपोभिहितः}}{\Afootnote{K; \skt{कारः पक्ष्मप्रकोपोऽभिहितः}\ A H}}\skt{ पुरस्तात् । }\pend
\pstart\versepar \skt{तत्रोपदिष्तस्य }\edtext{\skt{नरस्य}}{\Afootnote{A K; \skt{नर॥…॥स्य}\ H}}\skt{ }\edtext{\skt{चर्म}}{\Afootnote{A K; \skt{वर्त्म्म}\ H}}\skt{ }\edtext{\skt{वर्त्मोपरिष्टादनुतिर्यगेव}}{\Afootnote{K H; \skt{वर्त्मोपरिष्टादनुतिर्यगेव ॥}\ A}}\skt{ ।}\edtext{}{\lemma{\skt{पुरस्तात् । तत्रोपदिष्तस्य}}\Afootnote{K H; \skt{पुरस्तात् । तत्रोपविष्टस्य}\ A}}\pend

\pstart\versepar \ledsidenote{\passref{6.16.5}}\edtext{\skt{उद्भृत्यशस्त्रेण}}{\Afootnote{K H; \skt{उत्कृत्य शस्त्रेण}\ A}}\skt{ }\edtext{\skt{यवप्रमाणम्वालेन}}{\lemma{\skt{वालेन}}\Afootnote{K H; \skt{बालेन}\ A}}\skt{ }\edtext{\skt{सीव्येद्भिषगप्रमत्तः}}{\Afootnote{K H; \skt{सीव्योद्भिषगप्रमत्तः ।}\ A}}\skt{ । }\pend
\pstart\versepar \skt{}\edtext{\skt{दत्वा तु}}{\Afootnote{K H; \skt{दत्त्वा च}\ A}}\skt{ }\edtext{\skt{सर्पिधुचापिशेषं}}{\Afootnote{K; \skt{सर्पिर्मधुनावशेषं}\ A; \skt{सर्प्पिम्र्मधुचापिशेषं}\ H}}\skt{ }\edtext{\skt{कुर्यात्क्रमं}}{\Afootnote{K; \skt{कुर्याद्विधानं}\ A; \skt{कुर्यात्क्रमं यो}\ H}}\skt{ }\edtext{\skt{योभिहितो}}{\Afootnote{K; \skt{विहितं}\ A; \skt{ऽभिहितो}\ H}}\skt{ }\edtext{\skt{व्रणेषु}}{\Afootnote{K H; \skt{व्रणे यत् ॥}\ A}}\skt{ ।}\pend

\pstart \ledsidenote{\passref{6.16.4ab}}\skt{भ्रुवोरधस्तात्परिमुच्य भागौ पक्ष्माश्रितं }\edtext{\edtext{\skt{चैकमतस्तु}}{\lemma{\skt{चैकमतस्}}\Afootnote{K; \skt{चैकमतो}\ A; \skt{चैकतमस्}\ H}}}{\lemma{\skt{तु}}\Afootnote{K H; \skt{ऽव}\ A}}\skt{ }\edtext{\skt{च्छिन्द्यात्}}{\Afootnote{K H; \skt{कृन्तेत् ।}\ A}}\skt{ ।}\pend

\pstart\versepar \ledsidenote{\passref{6.16.6}}\skt{ललाटदेशे च }\edtext{\skt{निवध्य}}{\Afootnote{K H; \skt{निबद्ध}\ A}}\skt{ }\edtext{\skt{पट्टां}}{\Afootnote{K H; \skt{पट्टं}\ A}}\skt{ }\edtext{\edtext{\skt{प्राक्सीवितन्तत्र}}{\lemma{\skt{प्राक्सीवितन्}}\Afootnote{K H; \skt{प्राक्स्यूतम्}\ A}}}{\lemma{\skt{तत्र}}\Afootnote{K H; \skt{अत्राप्य्}\ A}}\skt{ }\edtext{\skt{परञ्च}}{\Afootnote{K H; \skt{अपरं च}\ A}}\skt{ वन्धेत् । }\pend
\pstart\versepar \skt{स्थैर्यं गते }\edtext{\skt{चाप्यथ}}{\Afootnote{K H; \skt{चाप्यथ}\ A}}\skt{ शस्त्रमार्गे }\edtext{\edtext{\skt{वालाम्विमुञ्चेत्कुशलाभिक्ष्य}}{\lemma{\skt{वालाम्}}\Afootnote{K H; \skt{बाला}\ A}}}{\lemma{\skt{विमुञ्चेत्कुशलाभिक्ष्य}}\Afootnote{K H; \skt{न्विमुञ्चेत्कुशलोऽभिवीक्ष्य ॥}\ A}}\skt{ ।}\edtext{}{\lemma{\skt{वन्धेत् । स्थैर्यं}}\Afootnote{K H; \skt{बद्ध्वा । स्थैर्यं}\ A}}\pend

\pstart\versepar \ledsidenote{\passref{6.16.7}}\skt{एवन्न चेच्छाम्यति तस्य वर्त्म निर्भुज्य }\edtext{\skt{दोषोपचितां}}{\Afootnote{K H; \skt{दोषोपहतां}\ A}}\skt{ }\edtext{\skt{वलीन्तु}}{\Afootnote{K H; \skt{वलिं च ।}\ A}}\skt{ । }\pend
\pstart\versepar \skt{}\edtext{\skt{ततोग्निना वा}}{\Afootnote{K; \skt{ततोऽग्निना वा}\ A H}}\skt{ }\edtext{\skt{प्रतिसारयीतक्षारणवातद्विमीक्ष्यधीरः}}{\Afootnote{K H; \skt{प्रतिसारयेत्तां क्षारेण वा सम्यगवेक्ष्य धीरः ॥}\ A}}\skt{ ।}\pend

\pstart\versepar \ledsidenote{\passref{6.16.8}}\skt{छित्त्वा समं }\edtext{\edtext{\skt{चाप्युपपक्ष्ममालां}}{\lemma{\skt{चाप्य्}}\Afootnote{K H; \skt{वा}\ A}}}{\lemma{\skt{उपपक्ष्ममालां}}\Afootnote{K H; \skt{प्युपपक्ष्ममालां}\ A}}\skt{ सम्यग्गृहीत्वा }\edtext{\edtext{\skt{वडिशैस्त्रिभिस्तु}}{\lemma{\skt{वडिशैस्}}\Afootnote{K H; \skt{बडिशै}\ A}}}{\lemma{\skt{त्रिभिस्तु}}\Afootnote{K H; \skt{स्त्रिभिस्तु ।}\ A}}\skt{ । }\pend
\pstart\versepar \skt{पथ्याफलेन }\edtext{\skt{प्रतिसारयीत}}{\Afootnote{K H; \skt{प्रतिसारयेत्तु}\ A}}\skt{ घृष्टेन वा }\edtext{\skt{तौरुवके न}}{\Afootnote{K H; \skt{तौवरकेण}\ A}}\skt{ सम्यक् ।}\pend

\pstart\versepar \ledsidenote{\passref{6.16.9}}\skt{चत्वार एते }\edtext{\skt{विधियो}}{\Afootnote{K; \skt{विधयो}\ A H}}\skt{ }\edtext{\skt{निहन्ति}}{\Afootnote{K; \skt{विहन्तुं}\ A; \skt{निहन्तुः}\ H}}\skt{ पक्ष्मोपरोधम्पृथगेव }\edtext{\skt{युक्ताः}}{\Afootnote{K H; \skt{शस्ताः ।}\ A}}\skt{ । }\pend
\pstart\versepar \skt{}\edtext{\edtext{\skt{विरेचनाश्च्योतननस्य}}{\lemma{\skt{विरेचनाश्}}\Afootnote{A K; \skt{विरेचना}\ H}}}{\lemma{\skt{च्योतननस्य}}\Afootnote{K; \skt{च्योतन}\ A; \skt{श्च्योतननस्य}\ H}}\skt{ }\edtext{\edtext{\skt{धूमलेपांजनस्नेहरसाक्रियाश्चेति}}{\lemma{\skt{धूमलेपांजनस्नेहरसाक्रियाश्}}\Afootnote{K H; \skt{धूमनस्यलेपाञ्जनस्नेहरसक्रियाश्}\ A}}}{\lemma{\skt{चेति}}\Afootnote{K H; \skt{च ॥}\ A}}\skt{ ॥}\pend

\pstart \ledsidenote{\passref{6.16.trailer}}\edtext{\skt{शालाक्ये}}{\Afootnote{K; \skt{इति सुश्रुतसंहितायामुत्तरतन्त्रान्तर्गते शालाक्यतन्त्रे पक्ष्मगतरोगप्रतिषेधो नाम षोडशोऽध्यायः ॥१६ ॥}\ A; \skt{शालाक्ये पक्ष्मगतरोगप्रतिषेधम्पञ्चदशमोऽध्यायः ॥ ॥}\ H}}\skt{ १५ ॥ ० ॥}\pend

\endnumbering
\beginnumbering
\eledsection*{Adhyāya 17}
\pstart \ledsidenote{\passref{6.17.1}}\skt{अथातो }\edtext{\skt{दृष्टिगतरोगप्रतिषेधं}}{\Afootnote{A K; \skt{दृष्टिगतरोगविशेषविज्ञानीयम्}\ H}}\skt{ }\edtext{\skt{व्याख्यास्यामः}}{\Afootnote{A H; \emph{om.}\ K}}\skt{ ॥}\pend

\pstart\versepar \ledsidenote{\passref{6.17.3}}\skt{त्रयः }\edtext{\edtext{\edtext{\skt{साध्यास्त्रयोऽसाध्याः}}{\lemma{\skt{साध्यास्}}\Afootnote{K H; \skt{साध्या}\ A}}}{\lemma{\skt{त्रयो}}\Afootnote{K H; \skt{स्त्रयो}\ A}}}{\lemma{\skt{ऽसाध्याः}}\Afootnote{K H; \skt{ऽसाध्या याप्याः}\ A}}\skt{ षट्च }\edtext{\skt{याप्या}}{\Afootnote{K H; \emph{om.}\ A}}\skt{ }\edtext{\skt{भवन्ति ह}}{\Afootnote{K H; \skt{भवन्ति हि ।}\ A}}\skt{ । }\pend
\pstart\versepar \skt{}\edtext{\skt{रोगा}}{\Afootnote{H; \skt{रोग}\ K}}\skt{ दृष्टिगता }\edtext{\skt{नॄणां}}{\Afootnote{\skt{नॄणां}\ K H}}\skt{ तेषु }\edtext{\skt{साध्यास्त्रयस्त्रये}}{\lemma{\skt{साध्यास्}}\Afootnote{K; \skt{सार्ध्यस्}\ H}}\skt{ ॥ }\pend
\pstart\versepar \skt{तत्रैकस्य प्रतीकारः कीर्तितो धूमदर्शिनः ।}\edtext{}{\lemma{\skt{रोगा… त}}\Afootnote{\emph{om.}\ A}}\pend

\pstart\versepar \ledsidenote{\passref{6.17.4}}\edtext{\skt{दृष्ट्याम्पित्तविदग्धायां}}{\lemma{\skt{दृष्ट्याम्}}\Afootnote{H; \skt{दृष्टौ}\ A; \skt{दृष्ट्या}\ K}}\skt{ विदग्धायां कफेन च । }\pend
\pstart\versepar \skt{पित्तश्लेष्महरं कुर्याद्विधिं }\edtext{\skt{शस्त्रक्षताद्विना}}{\Afootnote{K H; \skt{शस्त्रक्षतादृते ॥}\ A}}\skt{ ॥}\pend

\pstart\versepar \ledsidenote{\passref{6.17.5}}\skt{नस्यसेकाञ्जनालेपपुटपाकैः सतर्पणैः । }\pend
\pstart\versepar \skt{आद्ये तु }\edtext{\skt{त्रैफलं}}{\Afootnote{A K; \skt{त्रैफला}\ H}}\skt{ पेयं सर्पिस्त्रैवृतमुत्तरे ॥}\pend

\pstart\versepar \ledsidenote{\passref{6.17.6}}\skt{तैल्वकञ्चोभयोः पथ्यं केवलं जीर्णमेव वा । }\pend
\pstart\versepar \skt{}\edtext{\skt{गैरिकः}}{\Afootnote{K H; \skt{गैरिकं}\ A}}\skt{ }\edtext{\skt{सैन्धवः}}{\Afootnote{K H; \skt{सैन्धवं}\ A}}\skt{ कृष्णा गोदन्तस्य }\edtext{\skt{मषी}}{\Afootnote{A; \skt{मसीन्}\ K H}}\skt{ तथा ॥}\pend

\pstart\versepar \ledsidenote{\passref{6.17.7}}\skt{गोमांसं मरिचं }\edtext{\skt{बीजं}}{\Afootnote{A; \skt{वीजं}\ K H}}\skt{ }\edtext{\skt{शिरीषस्य}}{\Afootnote{A H; \skt{सिरीषस्य}\ K}}\skt{ मनःशिला । }\pend
\pstart\versepar \skt{वृन्तं }\edtext{\skt{कपित्थान्मधुना}}{\Afootnote{A; \skt{कपित्थात्मधुना}\ K H}}\skt{ स्वयंगुप्ताफलानि }\edtext{\skt{वा}}{\Afootnote{K H; \skt{च ॥}\ A}}\skt{ ॥}\pend

\pstart\versepar \ledsidenote{\passref{6.17.8}}\edtext{\edtext{\skt{एतेऽञ्जने}}{\lemma{\skt{एते}}\Afootnote{H; \emph{om.}\ A; \skt{ए}\ K}}\skt{ तु}}{\lemma{\skt{ऽञ्जने तु}}\Afootnote{H; \emph{om.}\ A; \skt{ते ञ्जने तु}\ K}}\skt{ }\edtext{\skt{चत्वारो}}{\Afootnote{K H; \skt{चत्वार}\ A}}\skt{ }\edtext{\skt{योगाः }\edtext{\skt{स्युरुभयोर्हिताः}}{\lemma{\skt{हिताः}}\Afootnote{K H; \skt{अञ्जने हिताः ।}\ A}}}{\lemma{\skt{योगाः स्युर्}}\Afootnote{H; \skt{एते योगाः स्युर्}\ A; \skt{योगा स्युर्}\ K}}\skt{ । }\pend
\pstart\versepar \skt{}\edtext{\skt{कुप्यकाशोकशालाम्रप्रियङ्गुनलिनोत्पलैः}}{\Afootnote{K; \skt{कुब्जकाशोकशालाम्रप्रियङ्गुनलिनोत्पलैः ॥}\ A; \skt{… प्यकाशोकशालाम्रप्रियङ्गुनलिनोत्पलैः ।}\ H}}\skt{ ॥}\pend

\pstart\versepar \ledsidenote{\passref{6.17.9}}\edtext{\edtext{\skt{पिष्टैर्हरेण्वामलकपथ्यापिप्पलिसंयुतैः}}{\lemma{\skt{पिष्टैर्}}\Afootnote{K H; \skt{पुष्पैर्}\ A}}}{\lemma{\skt{हरेण्वामलकपथ्यापिप्पलिसंयुतैः}}\Afootnote{K H; \skt{हरेणुकृष्णाह्वापथ्यामलकसंयुतैः ।}\ A}}\skt{ । }\pend
\pstart\versepar \skt{}\edtext{\edtext{\skt{कारयेदञ्जनं}}{\Afootnote{K; \skt{कारयेदञ्जन}\ H}}\skt{ वैद्यः }\edtext{\skt{सर्पिःक्षौद्रसमायुतं}}{\Afootnote{K; \skt{सर्पिर्मधुयुतैश्चूर्णैर्वेणुनाड्यामवस्थितैः ॥}\ A; \skt{सर्प्पिःक्षैद्रसमायुतः ॥}\ H}}}{\lemma{\skt{कारयेद्… स}}\Afootnote{\emph{om.}\ A}}\skt{ ॥}\pend

\pstart\versepar \ledsidenote{\passref{6.17.10}}\edtext{\skt{अञ्जयेद्द्वावपि}}{\lemma{\skt{अञ्जयेद्द्वाव्}}\Afootnote{A; \skt{अञ्जयेद्वाव्}\ K H}}\skt{ }\edtext{\skt{भिषक्पित्तस्लेष्मविभावितौ}}{\Afootnote{\skt{भिषक्पित्तश्लेष्मविभावितौ ।}\ A; \skt{भिषक्पित्तस्मेष्मविभावितौ ।}\ K; \skt{भिषक्पित्तश्लेष्मविभिवितौ ।}\ H}}\skt{ । }\pend
\pstart\versepar \skt{}\edtext{\edtext{\skt{हरेणवोऽम्रजम्बूभ्यां}}{\lemma{\skt{हरेणवो}}\Afootnote{K H; \emph{om.}\ A}}}{\lemma{\skt{ऽम्रजम्बूभ्यां}}\Afootnote{H; \skt{आम्रजम्बूद्भवं}\ A; \skt{म्रजम्बूभ्यां}\ K}}\skt{ }\edtext{\edtext{\skt{पुष्पस्वरसमेव}}{\lemma{\skt{पुष्पस्वरसम्}}\Afootnote{K H; \skt{पुष्पं तद्र}\ A}}}{\lemma{\skt{एव}}\Afootnote{K H; \skt{सेन}\ A}}\skt{ }\edtext{\skt{च}}{\Afootnote{K H; \skt{हरेणुकाम् ॥}\ A}}\skt{ ॥}\pend

\pstart\versepar \ledsidenote{\passref{6.17.11}}\edtext{\skt{घृतक्षौद्रविपक्वन्तु}}{\Afootnote{K; \skt{पिष्ट्वा क्षौद्राज्यसंयु}\ A; \skt{घृतक्षौद्रविपक्वन्}\ H}}\skt{ }\edtext{\edtext{\skt{तत्प्रयोज्यमथाञ्जनं}}{\lemma{\skt{तत्प्रयोज्यम्}}\Afootnote{K H; \skt{क्तं प्रयोज्यम्}\ A}}}{\lemma{\skt{अथाञ्जनं}}\Afootnote{K; \skt{अथवाञ्जनम्}\ A H}}\skt{ । }\pend
\pstart\versepar \skt{}\edtext{\skt{नलिनोत्पलकिञ्जल्कगैरिका}}{\Afootnote{K H; \skt{नलिनोत्पलकिञ्जल्कगैरिकैर्}\ A}}\skt{ }\edtext{\skt{गोशकृद्रसः}}{\Afootnote{H; \skt{गोशकृद्रसैः ॥}\ A; \skt{गोसकृद्रसः ।}\ K}}\skt{ ॥}\pend

\pstart\versepar \ledsidenote{\passref{6.17.12}}\skt{गुडिकाञ्जनमेतद्वा दिनरात्र्यन्धयोर्हितम् । }\pend
\pstart\versepar \skt{}\edtext{\skt{रसाञ्जनघृतक्षौद्रतालीशं}}{\Afootnote{\skt{रसाञ्जनरसक्षौद्रतालीश}\ A; \skt{रसाञ्जनघृतक्षौद्रतालीसं}\ K; \skt{रसाञ्जनमधुक्षौद्रतालीस}\ H}}\skt{ }\edtext{\skt{स्वर्णगैरिकैः}}{\Afootnote{K H; \skt{स्वर्णगैरिकम् ॥}\ A}}\skt{ ॥}\pend

\pstart\versepar \ledsidenote{\passref{6.17.13}}\skt{गोशकृद्रससंयुक्तं पित्तोपहतदृष्टये । }\pend
\pstart\versepar \skt{}\edtext{\edtext{\skt{शीतसौवीरकम्वापि}}{\lemma{\skt{शीतसौवीरकम्}}\Afootnote{\skt{शीतं सौवीरकं}\ A; \skt{सीतसौवीरकम्}\ K; \skt{शीतसौवीरकञ्च}\ H}}}{\lemma{\skt{वापि}}\Afootnote{A K; \skt{गपि}\ H}}\skt{ पिष्ट्वाथ रसभावितं ॥}\pend

\pstart\versepar \ledsidenote{\passref{6.17.14}}\skt{कूर्मपित्तेन }\edtext{\skt{मतिमान्भावयेद्रौहितेन वा}}{\Afootnote{A H; \skt{मतिमां भावयेद्रौहितेन वा ।}\ K}}\skt{ । }\pend
\pstart\versepar \skt{}\edtext{\skt{चूर्णाञ्जनेन तु}}{\Afootnote{K H; \skt{चूर्णाञ्जनमिदं}\ A}}\skt{ }\edtext{\skt{सदा}}{\Afootnote{K H; \skt{नित्यं}\ A}}\skt{ }\edtext{\skt{प्रयोज्यः}}{\Afootnote{K H; \skt{प्रयोज्यं}\ A}}\skt{ पित्तशान्तये ॥}\pend

\pstart\versepar \ledsidenote{\passref{6.17.15}}\edtext{\skt{कार्श्मरीपुष्पमधुकदार्वीलोध्ररसाञ्जनैः}}{\Afootnote{\skt{काश्मरीपुष्पमधुकदार्वीरोध्ररसाञ्जनैः ।}\ A; \skt{कार्स्मरीपुष्पमधुकदार्वीलोध्ररसाञ्जनैः ।}\ K; \skt{काश्मरीपुष्पमधुकदार्व्वीलोध्ररसाञ्जनैः ।}\ H}}\skt{ । }\pend
\pstart\versepar \skt{}\edtext{\skt{सक्षौद्रमञ्जनं}}{\lemma{\skt{अञ्जनं}}\Afootnote{A H; \skt{अञ्जन}\ K}}\skt{ }\edtext{\skt{तद्वा}}{\Afootnote{H; \skt{तद्व}\ A; \skt{तंद्वा}\ K}}\skt{ }\edtext{\edtext{\skt{हितमत्राञ्जने}}{\lemma{\skt{हितम्}}\Afootnote{K; \skt{द्धितम्}\ A; \skt{र्हितम्}\ H}}}{\lemma{\skt{अत्राञ्जने}}\Afootnote{K H; \skt{अत्रामये}\ A}}\skt{ सदा ॥}\pend

\pstart\versepar \ledsidenote{\passref{6.17.27}}\edtext{\edtext{\skt{नदीजशङ्खत्रिकटून्यथाञ्जनं}}{\lemma{\skt{नदीजशङ्खत्रिकटून्य्}}\Afootnote{\skt{नदीजशिम्बीत्रिकटून्य्}\ A; \skt{नदीजसांखत्रिकटून्}\ K; \skt{नदीजशखत्रिकटुं}\ H}}}{\lemma{\skt{अथाञ्जनं}}\Afootnote{A; \skt{तथाञ्जनं}\ K H}}\skt{ }\pend
\pstart\versepar \skt{मनःशिला द्वे }\edtext{\skt{रजने}}{\Afootnote{K H; \skt{च निशे}\ A}}\skt{ }\edtext{\skt{यकृद्रसः}}{\Afootnote{\skt{य(श)कृद्गवाम् ।}\ A; \skt{चयद्रसः ।}\ K; \skt{समञ्जय}\ H}}\skt{ । }\pend
\pstart\versepar \skt{}\edtext{\skt{सचन्दनेयं}}{\Afootnote{A K; \skt{त् । सचन्दनेयं}\ H}}\skt{ }\edtext{\skt{गुडिकाथवाञ्जने}}{\Afootnote{K; \skt{गुटिकाथवाञ्जनं}\ A; \skt{गुटिकाथवाञ्जने}\ H}}\skt{ }\pend
\pstart\versepar \skt{}\edtext{\skt{प्रशस्यते वै}}{\Afootnote{A H; \skt{प्रसस्यते वै}\ K}}\skt{ }\edtext{\edtext{\skt{दिवसेष्वपश्यतां}}{\lemma{\skt{दिवसेष्व्}}\Afootnote{A K; \skt{दिवसे}\ H}}}{\lemma{\skt{अपश्यतां}}\Afootnote{A; \skt{अपस्यतां ॥}\ K; \skt{ष्वपश्यतां ॥}\ H}}\skt{ ॥}\pend

\pstart\versepar \ledsidenote{\passref{6.17.16}}\skt{स्रोतोजं सैन्धवं }\edtext{\skt{कृष्णां}}{\Afootnote{K H; \emph{om.}\ A}}\skt{ }\edtext{\skt{हरेणुञ्चापि}}{\lemma{\skt{हरेणुञ्}}\Afootnote{K H; \skt{रेणुकां}\ A}}\skt{ पेषयेत् । }\pend
\pstart\versepar \skt{अजामूत्रेण ता वर्त्यः }\edtext{\skt{क्षणदान्धाञ्जने}}{\Afootnote{K H; \skt{क्षणदान्ध्याञ्जने}\ A}}\skt{ }\edtext{\skt{हिताः}}{\Afootnote{A H; \skt{हिता ॥}\ K}}\skt{ ॥}\edtext{}{\lemma{\skt{पेषयेत् । अजामूत्रेण ता}}\Afootnote{A K; \skt{पीषयेत् । अजामूत्रेण ता}\ H}}\pend

\pstart\versepar \ledsidenote{\passref{6.17.17}}\edtext{\skt{कालानुसारिवां}}{\Afootnote{A K; \skt{कालानुशारिवां}\ H}}\skt{ कृष्णां नागरं मधुकन्तथा । }\pend
\pstart\versepar \skt{}\edtext{\skt{तालीशपत्रं}}{\Afootnote{A; \skt{तालीसपत्रं}\ K H}}\skt{ }\edtext{\skt{रजने}}{\Afootnote{K; \skt{क्षणदे}\ A; \skt{रुजने}\ H}}\skt{ }\edtext{\edtext{\skt{सङ्खञ्चैव}}{\lemma{\skt{सङ्खञ्}}\Afootnote{K; \skt{गाङ्गे}\ A; \skt{शङ्खं}\ H}}}{\lemma{\skt{चैव}}\Afootnote{K H; \skt{यं च}\ A}}\skt{ }\edtext{\skt{यकृद्रसः}}{\Afootnote{\skt{यकृद्रसे ॥}\ A K; \skt{शकृद्रसे ॥}\ H}}\skt{ ॥}\pend

\pstart\versepar \ledsidenote{\passref{6.17.18}}\edtext{\skt{पिष्ट्वा}}{\Afootnote{K H; \skt{कृतास्ता वर्तयः पिष्टाश्}\ A}}\skt{ }\edtext{\edtext{\skt{वर्त्योऽथवा}}{\lemma{\skt{वर्त्यो}}\Afootnote{\emph{om.}\ A; \skt{वर्त्यो}\ K H}}}{\lemma{\skt{ऽथवा}}\Afootnote{\emph{om.}\ A; \skt{थवा}\ K H}}\skt{ }\edtext{\skt{तत्र}}{\Afootnote{K H; \emph{om.}\ A}}\skt{ }\edtext{\skt{छायाशुष्का}}{\Afootnote{H; \skt{छायाशुष्काः}\ A; \skt{छायासुष्का}\ K}}\skt{ }\edtext{\skt{रुजापहा}}{\Afootnote{K H; \skt{सुखावहाः ।}\ A}}\skt{ । }\pend
\pstart\versepar \skt{}\edtext{\skt{मनःशिलाभयाव्योषं}}{\Afootnote{K; \skt{मनःशिलाभयाव्योष}\ A H}}\skt{ }\edtext{\skt{तथा}}{\Afootnote{K H; \skt{बला}\ A}}\skt{ }\edtext{\skt{कालानुसारिवां}}{\Afootnote{\skt{कालानुसारिवाः ॥}\ A; \skt{कालानुशारिवां ॥}\ K H}}\skt{ ॥}\pend

\pstart\versepar \ledsidenote{\passref{6.17.19}}\edtext{\skt{समुद्रफेनश्च}}{\Afootnote{K H; \skt{स}\ A}}\skt{ }\edtext{\skt{हिता}}{\Afootnote{K H; \skt{फेना}\ A}}\skt{ }\edtext{\edtext{\skt{वर्त्यश्छागपयोन्विताः}}{\lemma{\skt{वर्त्यश्}}\Afootnote{K H; \skt{वर्तयः}\ A}}}{\lemma{\skt{छागपयोन्विताः}}\Afootnote{K H; \skt{पिष्टाश्छागक्षीरसमन्विताः ।}\ A}}\skt{ । }\pend
\pstart\versepar \skt{}\edtext{\skt{गोमूत्रपित्तमदिरायकृद्धात्रीरसैः}}{\Afootnote{K H; \skt{गोमूत्रपित्तमदिरायकृद्धात्रीरसे}\ A}}\skt{ पचेत् ॥}\pend

\pstart\versepar \ledsidenote{\passref{6.17.20}}\edtext{\skt{क्षौद्राञ्जनं}}{\Afootnote{K H; \skt{क्षुद्राञ्जनं}\ A}}\skt{ रसे }\edtext{\edtext{\skt{वान्यद्यकृतस्त्रैफलेन}}{\lemma{\skt{वान्यद्यकृतस्}}\Afootnote{K H; \skt{नान्यद्यक्द्त}\ A}}\skt{ वा}}{\lemma{\skt{त्रैफलेन वा}}\Afootnote{K H; \skt{स्त्रैफलेऽपि वा ।}\ A}}\skt{ । }\pend
\pstart\versepar \skt{}\edtext{\skt{गोमूत्राज्यार्णवमलं}}{\Afootnote{K; \skt{गोमूत्राज्यार्णवमल}\ A; \skt{गोमूत्राज्यार्णवमलः}\ H}}\skt{ पिप्पलीक्षौद्रकट्फलैः ॥}\pend

\pstart\versepar \ledsidenote{\passref{6.17.21}}\skt{सैन्धवोपहितं युञ्ज्यान्निहितं वेणुगह्वरे । }\pend
\pstart\versepar \skt{}\edtext{\skt{मेषाद्यकृद्घृतञ्चाजं}}{\lemma{\skt{मेषाद्यकृद्घृतञ्}}\Afootnote{K H; \skt{मेदो यकृद्घृतं}\ A}}\skt{ }\edtext{\skt{पिप्पलीसैधवं}}{\Afootnote{K H; \skt{पिप्पल्यः सैन्धवं}\ A}}\skt{ }\edtext{\skt{मधु}}{\Afootnote{A H; \skt{मधुः ॥}\ K}}\skt{ ॥}\pend

\pstart\versepar \ledsidenote{\passref{6.17.22}}\edtext{\skt{रसमामलकानाञ्च}}{\lemma{\skt{आमलकानाञ्च}}\Afootnote{K H; \skt{आमलकाच्चापि}\ A}}\skt{ }\edtext{\skt{पक्वा}}{\Afootnote{K; \skt{पक्वं}\ A; \skt{पक्वां}\ H}}\skt{ }\edtext{\skt{सम्यङ्निधापयेत्}}{\lemma{\skt{सम्यङ्}}\Afootnote{A K; \skt{सम्य}\ H}}\skt{ । }\pend
\pstart\versepar \skt{कोशके }\edtext{\skt{खादिरे}}{\Afootnote{\skt{खदिरनिर्माणे}\ A; \skt{खादिरे य}\ K; \skt{खदिरे}\ H}}\skt{ }\edtext{\skt{तद्वा}}{\Afootnote{H; \skt{तद्व}\ A; \skt{द्वा}\ K}}\skt{ }\edtext{\skt{हितं}}{\Afootnote{K H; \emph{om.}\ A}}\skt{ }\edtext{\skt{क्षौद्राञ्जनं}}{\Afootnote{K H; \skt{त्क्षुद्राञ्जनं}\ A}}\skt{ }\edtext{\skt{भवेत्}}{\Afootnote{K H; \skt{हितम् ॥}\ A}}\skt{ ॥}\edtext{}{\lemma{\skt{निधापयेत् । कोशके}}\Afootnote{\skt{निधापयेत् । कोशे}\ A; \skt{निधापयेत् । कोसके}\ K; \skt{ग्निधापयेत् । कोशिके}\ H}}\pend

\pstart\versepar \ledsidenote{\passref{6.17.23}}\edtext{\skt{मागध्यजास्थिमज्जैलायकृद्युक्ता}}{\Afootnote{K; \skt{हरेणुमगधाजास्थिमज्जैलायकृदन्}\ A; \skt{मागध्यजास्थिमज्जेलायकृद्युक्ता}\ H}}\skt{ }\edtext{\skt{हरेणवः}}{\Afootnote{K H; \skt{वितम् ।}\ A}}\skt{ ॥ }\pend
\pstart\versepar \skt{यकृद्रसेनाञ्जनम्वा श्लेष्मोपहतदृष्टये ॥}\pend

\pstart\versepar \ledsidenote{\passref{6.17.24}}\edtext{\skt{पचेत}}{\Afootnote{K H; \skt{विपाच्य}\ A}}\skt{ }\edtext{\skt{गोधायकृदन्त्रकल्पितं}}{\Afootnote{K H; \skt{गोधायकृदर्धपाटितं}\ A}}\skt{ }\pend
\pstart\versepar \skt{सुपूरितं मागधिकाभिरग्निना । }\pend
\pstart\versepar \skt{निषेवितं }\edtext{\skt{यद्यकृदञ्जनेन वा}}{\Afootnote{K H; \skt{तद्यकृदञ्जनेन}\ A}}\skt{ }\pend
\pstart\versepar \skt{निहन्ति }\edtext{\edtext{\skt{नक्तान्ध्यमसंशयं}}{\lemma{\skt{नक्तान्ध्यम्}}\Afootnote{A K; \skt{अनक्तान्ध्यम्}\ H}}}{\lemma{\skt{असंशयं}}\Afootnote{A H; \skt{असंसयं}\ K}}\skt{ खलु ॥}\pend

\pstart\versepar \ledsidenote{\passref{6.17.26ab}}\skt{प्लीहा }\edtext{\skt{यकृच्चाप्युपभक्षयेदुभे}}{\lemma{\skt{उपभक्षयेदुभे}}\Afootnote{\skt{उपभाक्षिते उभे}\ A}}\skt{ }\pend
\pstart\versepar \skt{प्रकल्प्य शूल्ये घृततैलसंयुते । }\pend
\pstart\versepar \skt{भवन्ति याप्याः खलु ये षडामया }\pend
\pstart\versepar \skt{हरेद्धि तत्र क्षतजं सिराव्यधैः ॥}\edtext{}{\lemma{\skt{भवन्ति… सिराव्यधैः}}\Afootnote{\emph{om.}\ A}}\pend

\pstart\versepar \ledsidenote{\passref{6.17.28cd}}\skt{विरेचयेच्चापि पुराणसर्पिषा }\pend
\pstart\versepar \skt{विरेचनाङ्गोपहितेन }\edtext{\skt{भूतये}}{\Afootnote{\skt{सर्वदा ॥}\ A}}\skt{ ॥}\pend

\pstart\versepar \ledsidenote{\passref{6.17.29}}\edtext{\skt{पयोविमिश्रं}}{\Afootnote{A H; \skt{पयोविमिस्रं}\ K}}\skt{ पवनोद्भवे हितं }\pend
\pstart\versepar \skt{वदन्ति पञ्चाङ्गुलतैलमेव तु । }\pend
\pstart\versepar \skt{}\edtext{\skt{पिबेद्घृतं}}{\Afootnote{K H; \skt{भवेद्घृतं}\ A}}\skt{ त्रैफलमेव }\edtext{\skt{शोधनं}}{\Afootnote{A H; \skt{सोधनं}\ K}}\skt{ }\pend
\pstart\versepar \skt{}\edtext{\skt{विशेषतः}}{\Afootnote{A H; \skt{विसेषतः}\ K}}\skt{ }\edtext{\skt{शोणितपित्तरोगयोः}}{\Afootnote{A H; \skt{सोणितपित्तरोगयोः ॥}\ K}}\skt{ ॥}\pend

\pstart\versepar \ledsidenote{\passref{6.17.30}}\skt{त्रिवृद्विरेकः कफजे }\edtext{\skt{प्रशस्यते}}{\Afootnote{A H; \skt{प्रसस्यते}\ K}}\skt{ }\pend
\pstart\versepar \skt{त्रिदोषजे }\edtext{\skt{तैलसुगन्धि}}{\Afootnote{K H; \skt{तैलमुशन्ति}\ A}}\skt{ }\edtext{\skt{तत्कृतं}}{\Afootnote{A; \skt{यत्कृतं ।}\ K H}}\skt{ । }\pend
\pstart\versepar \skt{पुराणसर्पिस्तिमिरेषु }\edtext{\skt{शस्यते}}{\Afootnote{\skt{सर्वशो}\ A H; \skt{सस्यते}\ K}}\skt{ }\pend
\pstart\versepar \skt{हितं }\edtext{\skt{भवेदायसभाजने}}{\Afootnote{K H; \skt{भवेदायसभाजन}\ A}}\skt{ स्थितं ॥}\pend

\pstart\versepar \ledsidenote{\passref{6.17.31}}\skt{हितञ्च विद्यात्त्रिफलाघृतं सदा }\pend
\pstart\versepar \skt{कृतञ्च }\edtext{\skt{तन्मेषविषाणसञ्ज्ञया}}{\Afootnote{K H; \skt{यन्मेषविषाणनामभिः । सदाव}\ A}}\skt{ । }\pend
\pstart\versepar \skt{}\edtext{\skt{लिहेत्सपाणं}}{\Afootnote{K; \skt{लिह्या}\ A; \skt{लिह्यात्सपाणं}\ H}}\skt{ }\edtext{\skt{त्रिफलां}}{\Afootnote{K H; \skt{त्त्रिफलां}\ A}}\skt{ सुचूर्णितां }\pend
\pstart\versepar \skt{}\edtext{\edtext{\skt{घृतप्रगाढान्तिमिरातुरो}}{\lemma{\skt{घृतप्रगाढान्}}\Afootnote{A K; \skt{घृतम्प्रगाढान्}\ H}}}{\lemma{\skt{तिमिरातुरो}}\Afootnote{K H; \skt{तिमिरे}\ A}}\skt{ }\edtext{\skt{नरः}}{\Afootnote{K H; \skt{ऽथ पित्तजे ॥}\ A}}\skt{ ॥}\pend

\pstart\versepar \ledsidenote{\passref{6.17.32}}\edtext{\skt{कफाभिभूतः}}{\Afootnote{K H; \skt{समीरजे तैलयुतां कफा}\ A}}\skt{ खलु तैलसंयुतां }\pend
\pstart\versepar \skt{}\edtext{\skt{मधुप्रगाढामथवापि}}{\lemma{\skt{अथवापि}}\Afootnote{K H; \skt{विदधीत}\ A}}\skt{ योजयेत् । }\pend
\pstart\versepar \skt{गवां शकृत्क्वाथविपक्वमुत्तमं }\pend
\pstart\versepar \skt{हितन्तु तैलन्तिमिरेषु }\edtext{\skt{नस्यतः}}{\Afootnote{K H; \skt{नावनम् ॥}\ A}}\skt{ ॥}\edtext{}{\lemma{\skt{भिभूतः… तैलसंयुतां}}\Afootnote{\emph{om.}\ A}}\edtext{}{\lemma{\skt{योजयेत् । गवां}}\Afootnote{K H; \skt{युक्तितः । गवां}\ A}}\pend

\pstart\versepar \ledsidenote{\passref{6.17.33}}\edtext{\skt{घृतं}}{\Afootnote{K H; \skt{हितं}\ A}}\skt{ }\edtext{\skt{हितं}}{\Afootnote{K H; \skt{घृतं}\ A}}\skt{ }\edtext{\skt{केवलमेव}}{\lemma{\skt{केवलम्}}\Afootnote{K H; \skt{केवल}\ A}}\skt{ }\edtext{\skt{पैत्तिके}}{\Afootnote{K H; \skt{पैत्तिके ह्य्}\ A}}\skt{ }\pend
\pstart\versepar \skt{}\edtext{\skt{तथा तु}}{\Afootnote{K; \skt{अजाविकं यन्मधुरैर्विपाचितम् । तैलं स्थिरादौ मधुरे च यद्गणे तथाणु}\ A; \skt{तथा नु}\ H}}\skt{ तैलं }\edtext{\skt{पवनासृगुत्थिते}}{\Afootnote{K H; \skt{पवनासृगुत्थयोः ॥}\ A}}\skt{ ।}\pend

\pstart\versepar \ledsidenote{\passref{6.17.34}}\edtext{\skt{तथाश्वगन्धातिबलाबलाश्रितं}}{\Afootnote{\skt{सहाश्वगन्धातिबलावरीशृतं हितं च}\ A; \skt{तथास्वगन्धातिवलावलाश्रितं}\ K; \skt{तथाश्वगन्धातिवलावलं श्रितं}\ H}}\skt{ }\pend
\pstart\versepar \skt{}\edtext{\skt{नस्ये तु}}{\Afootnote{K H; \skt{नस्ये}\ A}}\skt{ }\edtext{\skt{वाते}}{\Afootnote{K H; \emph{om.}\ A}}\skt{ त्रिवृतं प्रयोजयेत् ॥ }\pend
\pstart\versepar \skt{जलोद्भवानूपजमांससंस्कृताद् }\pend
\pstart\versepar \skt{घृतम्विधेयं पयसो यदुत्थितं ॥}\edtext{}{\lemma{\skt{प्रयोजयेत् ॥ जलोद्भवानूपजमांससंस्कृताद्घृतम्}}\Afootnote{H; \skt{यदीरितम् । जलोद्भवानूपजमांससंस्कृताद्घृतं}\ A; \skt{प्रयोजयेत् ॥ जलोद्भ…ऊपजमांससंस्कृतां घृतम्}\ K}}\pend

\pstart\versepar \ledsidenote{\passref{6.17.35}}\edtext{\skt{ससैन्धवक्रव्यभुगेणमांसजो}}{\Afootnote{\skt{ससैन्धवः क्रव्यभुगेणमांसयोर्}\ A; \skt{ससैन्धवक्रव्यभुजेनमास}\ K; \skt{ससैन्धवक्रव्यभुजेनमासजो}\ H}}\skt{ }\pend
\pstart\versepar \skt{}\edtext{\skt{युतः}}{\Afootnote{H; \skt{हितः}\ A; \emph{om.}\ K}}\skt{ }\edtext{\skt{पुटाख्यो}}{\Afootnote{H; \skt{ससर्पिः समधुः पुटा}\ A; \skt{जा युतः पुटाख्यो}\ K}}\skt{ }\edtext{\skt{मधुना}}{\Afootnote{K H; \emph{om.}\ A}}\skt{ }\edtext{\skt{घृतेन च}}{\Afootnote{K H; \skt{ह्वयः ।}\ A}}\skt{ । }\pend
\pstart\versepar \skt{}\edtext{\skt{वसाश्वगृध्रोरगताम्रचूडजा}}{\Afootnote{H; \skt{वसाथ गृध्रोरगताम्रचूडजा}\ A; \skt{वसास्वगृध्रोरगताम्रचूडजा}\ K}}\skt{ }\pend
\pstart\versepar \skt{}\edtext{\skt{मधूकयुक्ता च}}{\Afootnote{K H; \skt{सदा प्रशस्ता म}\ A}}\skt{ }\edtext{\skt{हिता}}{\Afootnote{K H; \skt{धुकान्}\ A}}\skt{ }\edtext{\skt{सदाञ्जने}}{\Afootnote{K H; \skt{विताञ्जने ॥}\ A}}\skt{ ॥}\pend

\pstart\versepar \ledsidenote{\passref{6.17.36}}\edtext{\skt{तथाञ्जनं}}{\Afootnote{K H; \skt{प्रत्यञ्जनं}\ A}}\skt{ }\edtext{\edtext{\skt{स्रोतसिजन्निषेवितं}}{\lemma{\skt{स्रोतसिजन्}}\Afootnote{K; \skt{स्रोतसि य}\ A; \skt{सोतसिजं}\ H}}}{\lemma{\skt{निषेवितं}}\Afootnote{K H; \skt{त्समुत्थितं}\ A}}\skt{ }\pend
\pstart\versepar \skt{}\edtext{\skt{क्रमाद्रसक्षीरघृतेषु}}{\Afootnote{A K; \skt{क्रमाद्रसः क्षीरघृतेषु}\ H}}\skt{ }\edtext{\skt{युज्यते}}{\Afootnote{K H; \skt{भावितम् ।}\ A}}\skt{ । }\pend
\pstart\versepar \skt{स्थितं }\edtext{\skt{दशाहत्रयमेतदञ्जनं}}{\lemma{\skt{एतदञ्जनं}}\Afootnote{A K; \skt{एव चाञ्जनं}\ H}}\skt{ }\pend
\pstart\versepar \skt{कृष्णोरगास्ये कुशसंप्रवेष्टिते ॥}\pend

\pstart\versepar \ledsidenote{\passref{6.17.37}}\edtext{\skt{समागधीक्षारकसैन्धवान्वितं}}{\Afootnote{K H; \skt{तन्मालतीकोरकसैन्धवायुतं}\ A}}\skt{ }\pend
\pstart\versepar \skt{}\edtext{\skt{तदञ्जनं}}{\Afootnote{H; \skt{सदाञ्जनं}\ A; \skt{तदाञ्जनं}\ K}}\skt{ }\edtext{\skt{स्यात्तिमिरेऽथ}}{\lemma{\skt{ऽथ}}\Afootnote{A; \skt{थ}\ K H}}\skt{ रागिणि । }\pend
\pstart\versepar \skt{}\edtext{\edtext{\skt{सुभावितञ्चापि}}{\lemma{\skt{सुभावितञ्}}\Afootnote{K H; \skt{सुभावितं वा}\ A}}}{\lemma{\skt{चापि}}\Afootnote{K H; \skt{पयसा दिनत्रयं काचाप}\ A}}\skt{ }\edtext{\skt{पयस्तु}}{\Afootnote{K H; \skt{हं शास्त्रविदः प्र}\ A}}\skt{ यद्भवेत् }\pend
\pstart\versepar \skt{सर्पस्य }\edtext{\skt{कृष्णस्य}}{\Afootnote{K H; \emph{om.}\ A}}\skt{ }\edtext{\skt{मुखेन}}{\Afootnote{K H; \emph{om.}\ A}}\skt{ }\edtext{\skt{भूरिणा}}{\Afootnote{K H; \emph{om.}\ A}}\skt{ ॥}\edtext{}{\lemma{\skt{यद्भवेत्सर्पस्य}}\Afootnote{K H; \skt{चक्षते ॥}\ A}}\pend

\pstart\versepar \ledsidenote{\passref{6.17.38}}\edtext{\skt{घृतं}}{\Afootnote{K H; \skt{हविर्}\ A}}\skt{ }\edtext{\skt{हितन्तत्प्रभवन्तु}}{\lemma{\skt{तत्प्रभवन्तु}}\Afootnote{K H; \skt{क्षीरभवं तु}\ A}}\skt{ पैत्तिके }\pend
\pstart\versepar \skt{वदन्ति }\edtext{\skt{नस्यं}}{\Afootnote{K H; \skt{नस्ये}\ A}}\skt{ }\edtext{\skt{मधुरौषधश्रितं}}{\Afootnote{H; \skt{मधुरौषधैः}\ A; \skt{मधुरौषधं स्रितं ।}\ K}}\skt{ । }\pend
\pstart\versepar \skt{}\edtext{\skt{हितं च}}{\Afootnote{K H; \skt{कृतम् ।}\ A}}\skt{ }\edtext{\edtext{\skt{तत्तर्पणमत्र}}{\lemma{\skt{तत्तर्पणम्}}\Afootnote{K H; \skt{तत्तर्पणे चैव हितं प्रयोजितं सजाङ्}\ A}}}{\lemma{\skt{अत्र}}\Afootnote{K H; \skt{गल}\ A}}\skt{ }\edtext{\skt{सम्युतः}}{\Afootnote{K H; \skt{स्तेषु च यः}\ A}}\skt{ }\pend
\pstart\versepar \skt{पुटाह्वयो }\edtext{\skt{जाङ्गलमांसजश्च}}{\Afootnote{K; \emph{om.}\ A; \skt{जाङ्गलजौमांसजाश्च}\ H\textsuperscript{ac}; \skt{जाङ्गलमांसजाश्च}\ H\textsuperscript{pc}}}\skt{ }\edtext{\skt{यः}}{\Afootnote{K H; \emph{om.}\ A}}\skt{ ॥}\pend

\pstart\versepar \ledsidenote{\passref{6.17.39}}\edtext{\skt{रसाञ्जनक्षौद्रयुता}}{\Afootnote{K; \skt{रसाञ्जनक्षौद्रसिता}\ A; \skt{रसाञ्जनं क्षौद्रसिता}\ H}}\skt{ }\edtext{\skt{मनःशिला}}{\Afootnote{K H; \skt{मनःशिलाः}\ A}}\skt{ }\pend
\pstart\versepar \skt{}\edtext{\edtext{\skt{द्रवाञ्जनञ्चात्र}}{\lemma{\skt{द्रवाञ्जनञ्}}\Afootnote{K H; \skt{क्षुद्राञ्जनं}\ A}}}{\lemma{\skt{चात्र}}\Afootnote{K H; \skt{तन्}\ A}}\skt{ }\edtext{\skt{मधूकसंयुतम्}}{\Afootnote{K H; \skt{मधुकेन संयुतम् ।}\ A}}\skt{ । }\pend
\pstart\versepar \skt{समाञ्जनं वा कनकाकरोद्भवं }\pend
\pstart\versepar \skt{}\edtext{\skt{सुपेषितं}}{\Afootnote{H; \skt{सुचूर्णितं}\ A; \skt{सुपेसितं}\ K}}\skt{ }\edtext{\edtext{\skt{तुत्थमुशन्ति}}{\lemma{\skt{तुत्थम्}}\Afootnote{K H; \skt{श्रेष्ठम्}\ A}}}{\lemma{\skt{उशन्ति}}\Afootnote{A H; \skt{उसन्ति}\ K}}\skt{ तद्विदः ॥}\pend

\pstart\versepar \ledsidenote{\passref{6.17.40}}\edtext{\skt{समेषशृङ्गाञ्जनभागसम्मितः}}{\Afootnote{H; \skt{भिल्लोटगन्धोदकसेकसेचितं प्रत्यञ्जने चात्र हितं तु तुत्थकम् । समेषशृङ्गाञ्जनभागसंमि}\ A; \skt{समेषसृङ्गाञ्जनभागसम्मितः}\ K}}\skt{ }\pend
\pstart\versepar \skt{}\edtext{\edtext{\skt{शङ्खोऽञ्जनात्काचमलं}}{\lemma{\skt{शङ्खो}}\Afootnote{H; \skt{तं}\ A; \skt{सङ्खो}\ K}}}{\lemma{\skt{ऽञ्जनात्काचमलं}}\Afootnote{K H; \skt{जलोद्भवं काचमलं}\ A}}\skt{ व्यपोहति ।}\pend

\pstart\versepar \ledsidenote{\passref{6.17.41}}\edtext{\skt{पलाशरोहीतमधूकजा}}{\Afootnote{A H; \skt{पलासरोहीतमधूकजा}\ K}}\skt{ }\edtext{\skt{रसाः}}{\Afootnote{A K; \skt{रसा}\ H}}\skt{ }\pend
\pstart\versepar \skt{}\edtext{\skt{प्रयोजिता वा}}{\Afootnote{K H; \skt{क्षौद्रेण युक्ता}\ A}}\skt{ }\edtext{\skt{मदिराग्रपेषिताः}}{\Afootnote{\skt{मदिराग्रमिश्रिताः ।}\ A; \skt{मदिग्रपेसिमा ।}\ K\textsuperscript{ac}; \skt{मदिराग्रपेसिमा ।}\ K\textsuperscript{pc}; \skt{मदिराग्रपेषिमाः ।}\ H}}\skt{ ॥ }\pend
\pstart\versepar \skt{}\edtext{\skt{उशीरलोध्रत्रिफलाप्रियङ्गुभिः}}{\Afootnote{A H; \skt{उसीरलोध्रत्रिफलाप्रियङ्गुभिः}\ K}}\skt{ }\pend
\pstart\versepar \skt{}\edtext{\skt{पचेत}}{\Afootnote{K\textsuperscript{pc} H; \skt{पचेत्तु}\ A; \skt{पचे}\ K\textsuperscript{ac}}}\skt{ नस्यं कफरोगशान्तये ।}\pend

\pstart\versepar \ledsidenote{\passref{6.17.42}}\skt{विडङ्गपाठाकिणिहीङ्गुदीत्वचः }\pend
\pstart\versepar \skt{}\edtext{\edtext{\skt{प्रयोजयेद्धूममुशीरमेव}}{\lemma{\skt{उशीरम्}}\Afootnote{K H; \skt{उशीर}\ A}}\skt{ च}}{\lemma{\skt{एव च}}\Afootnote{K H; \skt{संयुताः ।}\ A}}\skt{ ॥ }\pend
\pstart\versepar \skt{}\edtext{\skt{वनस्पतिक्वाथविभावितं}}{\Afootnote{K H; \skt{वनस्पतिक्वाथविपाचितं}\ A}}\skt{ घृतं }\pend
\pstart\versepar \skt{हितं }\edtext{\skt{हरिद्रानलदे च}}{\Afootnote{A K; \skt{हरिद्रानलदे व}\ H}}\skt{ }\edtext{\skt{तर्पणे}}{\Afootnote{K H; \skt{तर्पणम् ॥}\ A}}\skt{ ।}\pend

\pstart\versepar \ledsidenote{\passref{6.17.43}}\edtext{\skt{समागधासैन्धवमाक्षिकाढ्य}}{\Afootnote{K; \skt{समागधो माक्षिकसैन्}\ A; \skt{समागर्धासैन्धवमाक्षिकाढ्य}\ H}}\skt{ }\edtext{\skt{वा}}{\Afootnote{K H; \skt{धवाढ्यः}\ A}}\skt{ }\pend
\pstart\versepar \skt{}\edtext{\skt{सजाङ्गलः}}{\Afootnote{A H; \skt{सजाङ्गल}\ K}}\skt{ }\edtext{\skt{स्यात्पुटपाकमेव}}{\lemma{\skt{स्यात्पुटपाकम्}}\Afootnote{K H; \skt{स्यात्पुटपाक}\ A}}\skt{ च ॥ }\pend
\pstart\versepar \skt{मनःशिलात्र्यूषणशङ्खमाक्षिकैः }\pend
\pstart\versepar \skt{ससिन्धुकासीसरसाञ्जनैः }\edtext{\skt{क्रिया}}{\Afootnote{K H; \skt{क्रियाः ॥}\ A}}\skt{ ।}\pend

\pstart\versepar \ledsidenote{\passref{6.17.44}}\edtext{\skt{हितं च}}{\Afootnote{\skt{हिते च}\ A K H}}\skt{ }\edtext{\skt{कासीसरसाञ्जनाञ्जनं}}{\Afootnote{K H; \skt{कासीसरसाञ्जने तथा}\ A}}\skt{ }\pend
\pstart\versepar \skt{वदन्ति }\edtext{\edtext{\skt{पथ्यागुडनागरैर्युतं}}{\lemma{\skt{पथ्यागुडनागरैर्}}\Afootnote{K H; \skt{पथ्ये गुडनागरैर्}\ A}}}{\lemma{\skt{युतं}}\Afootnote{K H; \skt{युते ।}\ A}}\skt{ ॥ }\pend
\pstart\versepar \skt{यदञ्जनम्वा }\edtext{\skt{बहुशो}}{\Afootnote{A H; \skt{वहुशो}\ K}}\skt{ }\edtext{\skt{निषेवितं}}{\Afootnote{K H; \skt{निषेचितं}\ A}}\skt{ }\pend
\pstart\versepar \skt{}\edtext{\skt{समूत्रवर्गे}}{\Afootnote{A K; \skt{समूत्रवर्ग्ने}\ H}}\skt{ त्रिफलोदके }\edtext{\skt{श्रितं}}{\Afootnote{H; \skt{शृते ॥}\ A; \skt{स्रितं ॥}\ K}}\skt{ ।}\pend

\pstart \ledsidenote{\passref{6.17.44-app}}\skt{Tहे उननिमोउस्Nएपलेसे रेअदिन्घिते इस्हर्द्तो चोन्स्त्रुए बुत्मय्बे ओरिगिनल्.}\pend

\pstart\versepar \ledsidenote{\passref{6.17.45}}\edtext{\skt{निशाचरास्ये}}{\Afootnote{H; \skt{निशाचरास्थि}\ A; \skt{निसाचरास्ये}\ K}}\skt{ }\edtext{\skt{प्रणिधाय}}{\Afootnote{K H; \skt{स्थितमेत}\ A}}\skt{ }\edtext{\skt{चाञ्जनं}}{\Afootnote{K H; \skt{दञ्जनं}\ A}}\skt{ }\pend
\pstart\versepar \skt{}\edtext{\skt{न्यसेद्द्विमासं}}{\Afootnote{\skt{क्षिपेच्च मासं}\ A; \skt{न्यसेद्विमासं}\ K H}}\skt{ }\edtext{\skt{सलिलोत्थिते}}{\Afootnote{K; \skt{सलिलेऽस्थिरे}\ A; \skt{सलिलोत्थितं}\ H}}\skt{ पुनः ॥ }\pend
\pstart\versepar \skt{}\edtext{\skt{मधूकशिग्रोः}}{\Afootnote{H; \skt{मेषस्य पुष्पैर्मधुकेन}\ A; \skt{मधूकसिग्रोः}\ K}}\skt{ }\edtext{\skt{कुसुमैः}}{\Afootnote{K H; \emph{om.}\ A}}\skt{ }\edtext{\skt{समायुतं}}{\Afootnote{K H; \skt{संयुतं}\ A}}\skt{ }\pend
\pstart\versepar \skt{तदञ्जनं सर्वकृते प्रयोजयेत् ।}\pend

\pstart\versepar \ledsidenote{\passref{6.17.46}}\edtext{\skt{क्रियास्तु}}{\Afootnote{K H; \skt{क्रियाश्च}\ A}}\skt{ }\edtext{\skt{सर्वास्त्वथवापि}}{\lemma{\skt{त्वथवापि}}\Afootnote{K H; \skt{क्षत}\ A}}\skt{ }\edtext{\skt{रक्तजे}}{\Afootnote{K H; \skt{जोद्भवे}\ A}}\skt{ }\pend
\pstart\versepar \skt{}\edtext{\skt{हितः}}{\Afootnote{A H; \skt{हित}\ K}}\skt{ क्रमो }\edtext{\skt{म्लायिनि}}{\Afootnote{K H; \skt{परिम्लायिनि}\ A}}\skt{ }\edtext{\skt{वापि}}{\Afootnote{K H; \skt{चापि}\ A}}\skt{ }\edtext{\skt{पित्तहा}}{\Afootnote{K H; \skt{पित्तहृ}\ A}}\skt{ । }\pend
\pstart\versepar \skt{दोषान्वितस्यापि मुखप्रलेपनं }\pend
\pstart\versepar \skt{}\edtext{\skt{कुर्याच्च}}{\Afootnote{K; \skt{कुद्याच्च}\ H}}\skt{ सर्वेषु विधिं समीक्ष्य तु । }\pend
\pstart\versepar \skt{}\edtext{\skt{क्रमो}}{\Afootnote{K H; \skt{त् । क्रमो}\ A}}\skt{ हितः }\edtext{\skt{स्यन्दहरः}}{\Afootnote{A; \skt{स्यन्दहितः}\ K H}}\skt{ }\edtext{\skt{सुयोजितः}}{\Afootnote{K H; \skt{प्रयोजितः}\ A}}\skt{ }\pend
\pstart\versepar \skt{}\edtext{\edtext{\skt{सर्वेष्वथैतेषु}}{\lemma{\skt{सर्वेष्व्}}\Afootnote{K H; \skt{समी}\ A}}}{\lemma{\skt{अथैतेषु}}\Afootnote{K H; \skt{क्ष्य दोषेषु}\ A}}\skt{ }\edtext{\skt{यथास्वमेव}}{\lemma{\skt{एव}}\Afootnote{K H; \skt{एव च ॥}\ A}}\skt{ ॥}\edtext{}{\lemma{\skt{दोषान्वितस्यापि… क्र}}\Afootnote{\emph{om.}\ A}}\pend

\pstart \ledsidenote{\passref{6.17.46-app}}\skt{Tहिस्हेमिस्तिछप्पेअर्स्तो बे पर्त्ल्योमित्तेदन्द्पर्त्ल्य्रेअर्रन्गेदिन्थे वुल्गते. Tहे एमेन्दतिओनोf स्यन्दहितः तो स्यन्दहरः इस्बसेदोन्थे वुल्गते रेअदिन्ग्. Tहे Nएपलेसे MSS मय्बे एxहिबितिन्ग्दित्तोग्रफ्य्.}\pend

\pstart\versepar \ledsidenote{\passref{6.17.47}}\edtext{\skt{दोषोच्छ्रये}}{\Afootnote{K\textsuperscript{pc} H; \skt{दोषोदये}\ A}}\skt{ }\edtext{\skt{नैव च}}{\Afootnote{A K\textsuperscript{pc}; \skt{णैव च}\ H}}\skt{ }\edtext{\skt{विप्लुते}}{\Afootnote{K\textsuperscript{pc} H; \skt{विप्लुतिं}\ A}}\skt{ }\edtext{\skt{गदे}}{\Afootnote{K\textsuperscript{pc}; \skt{गते}\ A H}}\skt{ }\pend
\pstart\versepar \skt{द्रव्याणि नस्यादिषु योजयेद्भिषक् । }\pend
\pstart\versepar \skt{पुनश्च }\edtext{\skt{कल्पेऽञ्जनविस्तरः}}{\Afootnote{K H; \skt{कल्पेऽञ्जनविस्तरः}\ A}}\skt{ शुभः }\pend
\pstart\versepar \skt{प्रवक्ष्यते }\edtext{\skt{तच्च}}{\Afootnote{K H; \skt{ऽन्य}\ A}}\skt{ }\edtext{\skt{समीक्ष्य}}{\Afootnote{K; \skt{स्तमपीह}\ A; \skt{महीह}\ H}}\skt{ योजयेत् ॥}\edtext{}{\lemma{\skt{… पु}}\Afootnote{\emph{om.}\ K\textsuperscript{ac}}}\edtext{}{\lemma{\skt{योजयेद्भिषक् । पुनश्च}}\Afootnote{K\textsuperscript{pc} H; \skt{योजयेद्बुधः । पुनश्च}\ A}}\pend

\pstart\versepar \ledsidenote{\passref{6.17.48}}\skt{घृतम्पुराणं त्रिफलां }\edtext{\skt{शतावरीन्}}{\Afootnote{A H; \skt{सतावरीन्}\ K}}\skt{ }\pend
\pstart\versepar \skt{}\edtext{\skt{तथैव}}{\Afootnote{K H; \skt{पटोल}\ A}}\skt{ }\edtext{\skt{मुद्गामलकं}}{\Afootnote{A K; \skt{मुद्गामलका}\ H}}\skt{ }\edtext{\skt{यवानपि}}{\Afootnote{A K; \skt{न्यवानपि ।}\ H}}\skt{ । }\pend
\pstart\versepar \skt{}\edtext{\skt{निषेवमाणस्य}}{\Afootnote{A K; \skt{निषेव्यमाणस्य}\ H}}\skt{ }\edtext{\skt{हि}}{\Afootnote{K H; \skt{नरस्य}\ A}}\skt{ यत्नतो भयं }\pend
\pstart\versepar \skt{}\edtext{\skt{न}}{\Afootnote{K H; \skt{सुघोररात्तिमिरान्न}\ A}}\skt{ }\edtext{\skt{विद्यते वै}}{\Afootnote{K H; \skt{विद्यते ॥}\ A}}\skt{ }\edtext{\skt{तिमिरात्सुघोरतः}}{\Afootnote{K H; \emph{om.}\ A}}\skt{ ॥}\pend

\pstart\versepar \ledsidenote{\passref{6.17.49}}\edtext{\skt{शतावरीपायस}}{\Afootnote{A H; \skt{सतावरीपायस}\ K}}\skt{ एव }\edtext{\skt{केवलम्}}{\Afootnote{K H; \skt{केवलस्}\ A}}\skt{ }\pend
\pstart\versepar \skt{तथा कृतो वामलकेषु पायसः । }\pend
\pstart\versepar \skt{प्रभूतसर्पिस्त्रिफलोदकोत्तरो }\pend
\pstart\versepar \skt{यवौदनो वा तिमिरम्व्यपोहति ॥}\pend

\pstart\versepar \ledsidenote{\passref{6.17.52}}\edtext{\skt{तिमिरे}}{\Afootnote{K H; \skt{विवर्जयेत्सिरामोक्षं तिमिरे}\ A}}\skt{ }\edtext{\skt{रागिणि}}{\Afootnote{K H; \emph{om.}\ A}}\skt{ }\edtext{\skt{सिरा न}}{\Afootnote{K H; \skt{राग}\ A}}\skt{ }\edtext{\skt{व्यध्या }\edtext{\skt{स्युर्विजानता}}{\lemma{\skt{विजानता}}\Afootnote{K H; \skt{मागते ।}\ A}}}{\lemma{\skt{व्यध्या स्युर्}}\Afootnote{K; \emph{om.}\ A; \skt{व्यध्याः स्युर्}\ H}}\skt{ । }\pend
\pstart\versepar \skt{यन्त्रेणोत्पीडितो }\edtext{\skt{दोषस्तथा}}{\lemma{\skt{तथा}}\Afootnote{K H; \skt{नि}\ A}}\skt{ }\edtext{\edtext{\skt{हन्त्याशु}}{\lemma{\skt{हन्त्य्}}\Afootnote{K H; \skt{हन्या}\ A}}}{\lemma{\skt{आशु}}\Afootnote{H; \skt{दाशु}\ A; \skt{आसु}\ K}}\skt{ }\edtext{\skt{दर्शनं}}{\Afootnote{A H; \skt{दर्सनं ॥}\ K}}\skt{ ॥}\pend

\pstart\versepar \ledsidenote{\passref{6.17.53}}\edtext{\skt{अरागं}}{\Afootnote{K H; \skt{अरागि}\ A}}\skt{ तिमिरं साध्यमाद्यं }\edtext{\edtext{\skt{पटलमाश्रितं}}{\lemma{\skt{पटलम्}}\Afootnote{H; \skt{पटल}\ A; \skt{पटालम्}\ K}}}{\lemma{\skt{आश्रितं}}\Afootnote{H; \skt{माश्रितम् ।}\ A; \skt{आस्रितं ।}\ K}}\skt{ । }\pend
\pstart\versepar \skt{}\edtext{\skt{रागि}}{\Afootnote{K H; \skt{कृच्छ्रं}\ A}}\skt{ द्वितीये }\edtext{\skt{कृच्छ्रेण}}{\Afootnote{K H; \skt{रागि स्या}\ A}}\skt{ }\edtext{\skt{तृतीये}}{\Afootnote{K H; \skt{त्तृतीये}\ A}}\skt{ }\edtext{\skt{याप्यमुच्यते}}{\lemma{\skt{याप्यम्}}\Afootnote{A K; \skt{याप्य}\ H}}\skt{ ॥}\pend

\pstart\versepar \ledsidenote{\passref{6.17.55}}\skt{श्लैष्मिके }\edtext{\skt{लिङ्गनाशे तु}}{\Afootnote{A H; \skt{लिङ्गनासे तु}\ K}}\skt{ कर्म वक्ष्यामि सिद्धये । }\pend
\pstart\versepar \skt{}\edtext{\skt{भवेद्वृद्धेन्द्वातपत्रमुक्तावर्ताकृतिः}}{\Afootnote{K; \skt{न चेदर्धेन्दुघर्माम्बुबिन्दुमुक्ताकृतिः}\ A; \skt{भवेद्वृद्धेद्धातपत्रमुक्तावर्त्ताकृतिः}\ H}}\skt{ }\edtext{\skt{सितः}}{\Afootnote{K H; \skt{स्थिरः ॥}\ A}}\skt{ ॥}\pend

\pstart\versepar \ledsidenote{\passref{6.17.56}}\skt{विषमो वा तनुर्मध्ये राजिमान्वा }\edtext{\skt{बहुप्रभः}}{\Afootnote{A H; \skt{वहुप्रदः ।}\ K}}\skt{ । }\pend
\pstart\versepar \skt{दृष्टिस्थो लक्ष्यते दोषः सरुजो वा सलोहितः ॥}\pend

\pstart\versepar \ledsidenote{\passref{6.17.57}}\skt{स्निग्धस्विन्नस्य तस्याथ काले नात्युष्णशीतले । }\pend
\pstart\versepar \skt{यन्त्रितस्योपविष्टस्य स्वान्नासां पश्यतः समं ॥}\pend

\pstart\versepar \ledsidenote{\passref{6.17.58}}\edtext{\edtext{\skt{मतिमाञ्छुक्लभागौ}}{\lemma{\skt{मतिमाञ्}}\Afootnote{K H; \skt{मतिमा}\ A}}\skt{ द्वौ}}{\lemma{\skt{छुक्लभागौ द्वौ}}\Afootnote{K H; \skt{न्शुक्लभागौ द्वौ}\ A}}\skt{ कृष्णान्मुक्त्वा }\edtext{\skt{ह्यपाङ्गतः}}{\lemma{\skt{अपाङ्गतः}}\Afootnote{K H; \skt{अपाङ्गतः ।}\ A}}\skt{ । }\pend
\pstart\versepar \skt{}\edtext{\skt{निपीड्य}}{\Afootnote{K H; \skt{उन्मील्य}\ A}}\skt{ नयने }\edtext{\skt{सम्यक्छिद्रे}}{\Afootnote{K H; \skt{सम्यक्सिरा}\ A}}\skt{ }\edtext{\skt{दैवकृते}}{\Afootnote{K H; \skt{जालविवर्जिते ॥}\ A}}\skt{ }\edtext{\skt{ततः}}{\Afootnote{K H; \emph{om.}\ A}}\skt{ ॥}\pend

\pstart \ledsidenote{\passref{6.17.59}}\edtext{\skt{शलाकया}}{\Afootnote{K H; \skt{नाधो नोर्ध्वं न पार्श्वाभ्यां छिद्रे दैवकृते ततः । शलाकया}\ A}}\skt{ }\edtext{\skt{ताम्रमय्या}}{\Afootnote{K H; \skt{प्रयत्ने}\ A}}\skt{ }\edtext{\skt{लौह्या वा}}{\Afootnote{K H; \skt{न विश्वस्तं}\ A}}\skt{ }\edtext{\skt{यववक्रया}}{\Afootnote{A; \skt{यववत्क्रया ।}\ K; \skt{यववक्त्रया ।}\ H}}\skt{ ।}\pend

\pstart\versepar \ledsidenote{\passref{6.17.60}}\skt{मध्यप्रदेशिन्यङ्गुष्ठस्थिरहस्तगृहीतया ॥५३॥ }\pend
\pstart\versepar \skt{}\edtext{\skt{अदक्षिणं}}{\Afootnote{K H; \emph{om.}\ A}}\skt{ दक्षिणेन }\edtext{\edtext{\skt{हस्तेनेतरमन्यथा}}{\lemma{\skt{हस्तेनेतरम्}}\Afootnote{K H; \skt{भिषक्सव्यं विध्येत्सव्येन चेतर}\ A}}}{\lemma{\skt{अन्यथा}}\Afootnote{K H; \skt{त् ॥}\ A}}\skt{ ।}\pend

\pstart\versepar \ledsidenote{\passref{6.17.61}}\edtext{\skt{वारिबिन्द्वागमः}}{\Afootnote{A H; \skt{वारिविन्द्वागमः}\ K}}\skt{ }\edtext{\skt{सम्यक्छब्दश्च}}{\Afootnote{K H; \skt{सम्यग्भवेच्छ}\ A}}\skt{ }\edtext{\skt{व्यधने}}{\Afootnote{K H; \skt{ब्देस्तथा}\ A}}\skt{ भवेत् ॥५४॥ }\pend
\pstart\versepar \skt{संसिच्य }\edtext{\skt{विद्धमात्रन्तु}}{\Afootnote{A K; \skt{विद्धमानन्तु}\ H}}\skt{ योषित्स्तन्येन कोविदः ।}\edtext{}{\lemma{\skt{भवेत् ॥५४॥ संसिच्य}}\Afootnote{K H; \skt{व्यधे । संसिच्य}\ A}}\pend

\pstart\versepar \ledsidenote{\passref{6.17.63}}\skt{शलाकाग्रेण तु ततो }\edtext{\skt{निर्लिखेद्दृष्टिमण्डलं}}{\Afootnote{A K; \skt{निर्द्दिशेद्दृष्टिमण्डलं ।}\ H}}\skt{ ॥५५॥ }\pend
\pstart\versepar \skt{}\edtext{\skt{अबाधमानः }\edtext{\skt{शनकैर्नासाम्प्रति}}{\lemma{\skt{शनकैर्}}\Afootnote{H; \skt{सनकैर्}\ K}}\skt{ }\edtext{\skt{नुदंस्ततः}}{\lemma{\skt{ततः}}\Afootnote{K H; \skt{ध्यतो योऽन्यपार्श्वेऽक्ष्णस्तं रुद्ध्वा नासिकापुटम् ॥}\ A}}\skt{ ।}}{\lemma{\skt{अबाधमानः… नुदंस्}}\Afootnote{\emph{om.}\ A}}\pend

\pstart \ledsidenote{\passref{6.17.64a}}\edtext{\skt{उच्छिङ्घनाच्चापहरेद्दृष्टिमण्डलगं}}{\Afootnote{K; \skt{उच्छिङ्घनेन हर्तव्यो दृष्टिमण्डलगः}\ A; \skt{उच्छिङ्घनाच्चाये हरेद्दृष्टिमण्डलगङ्}\ H}}\skt{ }\edtext{\skt{कफं}}{\Afootnote{K H; \skt{कफः ।}\ A}}\skt{ ॥५६॥}\pend

\pstart\versepar \ledsidenote{\passref{6.17.62}}\edtext{\skt{स्त्याने}}{\Afootnote{K H; \skt{स्थिरे}\ A}}\skt{ दोषे चले }\edtext{\skt{वापि}}{\Afootnote{A H; \skt{चापि}\ K}}\skt{ स्वेदयेदक्षि }\edtext{\skt{बाह्यतः}}{\Afootnote{A; \skt{वाह्यतः ॥}\ K H}}\skt{ । }\pend
\pstart\versepar \skt{}\edtext{\edtext{\skt{सम्यक्सूचीमवस्थाप्य}}{\lemma{\skt{सम्यक्सूचीम्}}\Afootnote{K; \skt{सम्यक्शला}\ A; \skt{सम्यक्सचीम्}\ H}}}{\lemma{\skt{अवस्थाप्य}}\Afootnote{K H; \skt{कां संस्थाप्य}\ A}}\skt{ भङ्गैरनिलनाशनैः ॥५७॥}\pend

\pstart \ledsidenote{\passref{6.17.65}}\edtext{\skt{एवं}}{\Afootnote{K H; \skt{तदासौ लिखिता सम्यग्ज्ञेया या चापि निर्व्यथा । (एवं}\ A}}\skt{ }\edtext{\skt{त्वशक्ये}}{\Afootnote{A; \skt{त्वसक्ये}\ K; \skt{त्वशक्यो}\ H}}\skt{ }\edtext{\skt{निर्हन्तुं}}{\Afootnote{K H; \skt{निर्हर्तुं}\ A}}\skt{ दोषे }\edtext{\skt{प्रत्यागतेऽपि}}{\Afootnote{K H; \skt{प्रत्यागतेऽपि}\ A}}\skt{ वा ।}\pend

\pstart \ledsidenote{\passref{6.17.66a}}\edtext{\skt{स्नेहाद्यैरुपपन्नस्य}}{\lemma{\skt{स्नेहाद्यैर्}}\Afootnote{K H; \skt{स्नेहाद्यैर्}\ A}}\skt{ }\edtext{\skt{विदधीत}}{\Afootnote{K H; \emph{om.}\ A}}\skt{ }\edtext{\edtext{\skt{पुनर्व्यधं}}{\lemma{\skt{पुनर्}}\Afootnote{K H; \emph{om.}\ A}}}{\lemma{\skt{व्यधं}}\Afootnote{K H; \skt{व्यधो भूयो विधीयते ।)}\ A}}\skt{ ॥५८॥}\pend

\pstart \ledsidenote{\passref{6.17.64b}}\edtext{\skt{अथाभ्रमुक्ते च}}{\Afootnote{\skt{निर}\ A; \skt{अथाग्रमुक्ते च}\ K; \skt{अथाग्रमुक्तै}\ H}}\skt{ }\edtext{\edtext{\skt{हरिर्यथा}}{\lemma{\skt{हरिर्}}\Afootnote{\skt{भ्र इव घर्मांशु}\ A; \skt{हरि}\ K; \skt{व च हरिर्}\ H}}}{\lemma{\skt{यथा}}\Afootnote{H; \skt{र्यदा}\ A; \skt{यदा}\ K}}\skt{ }\edtext{\skt{दृष्टिः}}{\Afootnote{A K; \skt{दृष्टि}\ H}}\skt{ }\edtext{\skt{प्रकाशते}}{\Afootnote{A H; \skt{प्रकासते ।}\ K}}\skt{ ।}\pend

\pstart \ledsidenote{\passref{6.17.66b}}\skt{ततो दृष्टेषु रूपेषु }\edtext{\skt{शलाकामाहरेच्छनैः}}{\lemma{\skt{शलाकाम्}}\Afootnote{A H; \skt{सलाकाम्}\ K}}\skt{ ॥५९॥}\pend

\pstart \ledsidenote{\passref{6.17.64b-app}}\skt{अथाभ्रमुक्ते इस चोन्जेच्तुरे ब्य्Aन्द्रेय्Kलेबनोव्बसेदोन्थे वुल्गतेऽस्रेअदिन्गोf निरभ्र इव.}\pend

\pstart\versepar \ledsidenote{\passref{6.17.67}}\edtext{\skt{सर्पिषाभ्यज्य}}{\Afootnote{K H; \skt{घृतेनाभ्यज्य}\ A}}\skt{ }\edtext{\skt{नयने}}{\Afootnote{K H; \skt{नयनं}\ A}}\skt{ }\edtext{\skt{ततः}}{\Afootnote{K H; \skt{वस्त्र}\ A}}\skt{ पट्टेन वेष्टयेत् । }\pend
\pstart\versepar \skt{ततो गृहे निराबाधे }\edtext{\skt{शयीतोत्तानमेव}}{\lemma{\skt{शयीतोत्तानम्}}\Afootnote{\skt{शयीतोत्तान}\ A H; \skt{सयीतोत्तानम्}\ K}}\skt{ }\edtext{\skt{तु}}{\Afootnote{K; \skt{च ॥}\ A H}}\skt{ ॥६०॥}\pend

\pstart\versepar \ledsidenote{\passref{6.17.68}}\edtext{\skt{उद्गारकासक्षयथुष्ठीवनोत्कम्पनानि च}}{\Afootnote{K; \skt{उद्गारकासक्षवथुष्ठीवनोत्कम्पनानि च ।}\ A H}}\skt{ । }\pend
\pstart\versepar \skt{तत्कालन्नाचरेदूर्ध्वं }\edtext{\skt{यन्त्रणा}}{\Afootnote{A H; \skt{यन्त्रण}\ K}}\skt{ स्नेहपीतवत् ॥६१॥}\pend

\pstart\versepar \ledsidenote{\passref{6.17.69}}\skt{त्र्यहात्त्र्यहाच्च धावेत कषायैरनिलापहैः । }\pend
\pstart\versepar \skt{वायोर्भयात्त्र्यहादूर्ध्वं स्वेदयेदक्षि }\edtext{\skt{बाह्यतः}}{\Afootnote{\skt{पूर्ववत् ॥}\ A; \skt{वाह्यतः ॥}\ K H}}\skt{ ॥६२॥}\pend

\pstart\versepar \ledsidenote{\passref{6.17.70}}\skt{दशाहमेवं संयम्य हितं दृष्टिप्रसादनं । }\pend
\pstart\versepar \skt{पश्चात्कर्म च सेवेत }\edtext{\skt{लघ्वन्नञ्चैव}}{\lemma{\skt{चैव}}\Afootnote{K H; \skt{चापि}\ A}}\skt{ मात्रया ॥६३॥}\pend

\pstart \ledsidenote{\passref{6.17.72}}\skt{पूर्यते }\edtext{\skt{शोणितस्याक्षि}}{\Afootnote{K H; \skt{शोणितेनाक्षि}\ A}}\skt{ }\edtext{\skt{सिरामोक्षाद्विलोचने}}{\Afootnote{H; \skt{सिरावेधाद्विसर्पता । तत्र स्त्रीस्तन्ययष्ट्याह्वपक्वं सेके हितं घृतम् ॥}\ A; \skt{सिरामोक्षाद्विरोचने ।}\ K}}\skt{ ।}\pend

\pstart\versepar \ledsidenote{\passref{6.17.82}}\edtext{\skt{शलाकात्कर्कशाच्छूलं}}{\Afootnote{K; \skt{शलाका कर्कशा शूलं}\ A; \skt{शलाकात्कर्क्कशात्थूलं}\ H}}\skt{ }\edtext{\skt{कृशाद्दोषपरिप्लवं}}{\Afootnote{K; \skt{खरा दोषपरिप्लुतिम् ।}\ A; \skt{कृशोद्दोषपरिप्लवं ॥}\ H}}\skt{ ॥६४॥ }\pend
\pstart\versepar \skt{व्रणं विशालं }\edtext{\edtext{\skt{स्थूलाग्रान्तीक्ष्णं}}{\lemma{\skt{स्थूलाग्रान्}}\Afootnote{K H; \skt{स्थूलाग्रा}\ A}}}{\lemma{\skt{तीक्ष्णं}}\Afootnote{K H; \skt{तीक्ष्णा}\ A}}\skt{ हिंस्यादनेकधा ।}\pend

\pstart\versepar \ledsidenote{\passref{6.17.83}}\skt{जलास्रावं }\edtext{\skt{सुविषमात्क्रियासङ्गकरी}}{\Afootnote{K H; \skt{तु विषमा क्रियासङ्गमथा}\ A}}\skt{ स्थिरा ॥६५॥ }\pend
\pstart\versepar \skt{करोति }\edtext{\skt{वर्जितां}}{\Afootnote{K; \skt{वर्जिता}\ A; \skt{वर्ज्जिनां}\ H}}\skt{ }\edtext{\skt{दोषैस्तस्मादेभिर्हितां}}{\lemma{\skt{हितां}}\Afootnote{\skt{हिता}\ A K H}}\skt{ भवेत् ।}\pend

\pstart\versepar \ledsidenote{\passref{6.17.84}}\skt{अष्टाङ्गुलायता मध्ये सूत्रेण परिवेष्टिता ॥६६॥ }\pend
\pstart\versepar \skt{}\edtext{\skt{अङ्गुष्ठपर्वतुल्या तु}}{\Afootnote{K; \skt{अङ्गुष्ठपर्वसमिता}\ A; \skt{अङ्गुष्ठपर्व्वतु …}\ H}}\skt{ वक्त्रयोर्मुकुलाकृतिः ।}\pend

\pstart\versepar \ledsidenote{\passref{6.17.85}}\edtext{\skt{तारारी}}{\Afootnote{K H; \skt{ताम्रायसी}\ A}}\skt{ }\edtext{\skt{शातकुम्भी च}}{\Afootnote{\skt{शातकुम्भी}\ A; \skt{सातकुम्भी च}\ K; \skt{सातकोम्भी च}\ H}}\skt{ शलाका }\edtext{\skt{स्यादकुत्सिता}}{\Afootnote{K; \skt{स्यादनिन्दिता ।}\ A; \skt{स्यादकुम्भिता ।}\ H}}\skt{ ॥६७॥ }\pend
\pstart\versepar \skt{}\edtext{\skt{रागशोफार्बुदं}}{\Afootnote{\skt{रागः शोफोऽर्बुदं}\ A; \skt{रागशोफार्वुदं}\ K H}}\skt{ चोषो बुद्बुदं }\edtext{\skt{स्{\symfont `}ऊकराक्षिता}}{\Afootnote{\skt{शूकराक्षिता ॥}\ A; \skt{सूकराक्षिता}\ K H}}\skt{ ।}\pend

\pstart\versepar \ledsidenote{\passref{6.17.86}}\skt{अधिमन्थादयश्चान्ये }\edtext{\skt{रोगाः }\edtext{\skt{स्युर्व्यधदोषतः}}{\lemma{\skt{व्यधदोषतः}}\Afootnote{K H; \skt{व्यधदोषजाः ।}\ A}}}{\lemma{\skt{रोगाः स्युर्}}\Afootnote{A H; \skt{रोगा स्युर्}\ K}}\skt{ ॥६८॥ }\pend
\pstart\versepar \skt{अहिताचारतो वापि यथास्वं तानुपाचरेत् ।}\pend

\pstart\versepar \ledsidenote{\passref{6.17.87}}\edtext{\edtext{\skt{रुजास्वक्षिषु}}{\lemma{\skt{रुजास्व्}}\Afootnote{K; \skt{रुजायाम्}\ A; \skt{राजास्व्}\ H}}}{\lemma{\skt{अक्षिषु}}\Afootnote{\skt{अक्षि}\ A; \skt{अक्षि तु}\ K H}}\skt{ }\edtext{\skt{रागेषु}}{\Afootnote{K H; \skt{रागे वा योगा}\ A}}\skt{ }\edtext{\skt{भूयो}}{\Afootnote{K H; \skt{न्भूयो}\ A}}\skt{ }\edtext{\skt{योगान्निबोध}}{\lemma{\skt{योगान्}}\Afootnote{K H; \emph{om.}\ A}}\skt{ मे ॥६९॥ }\pend
\pstart\versepar \skt{}\edtext{\skt{गैरिकः}}{\Afootnote{K H; \skt{गैरिकं}\ A}}\skt{ }\edtext{\skt{शारिवा}}{\Afootnote{K H; \skt{सारिवा}\ A}}\skt{ दूर्वा यवपिष्टं }\edtext{\skt{घृतन्तथा}}{\lemma{\skt{तथा}}\Afootnote{K H; \skt{पयः ॥}\ A}}\skt{ ।}\pend

\pstart\versepar \ledsidenote{\passref{6.17.88}}\edtext{\skt{मुखालेपः}}{\Afootnote{K H; \skt{सुखालेपः}\ A}}\skt{ }\edtext{\skt{प्रयोज्योऽयं}}{\Afootnote{\skt{प्रयोज्योऽयं}\ A; \skt{प्रयोज्यो यं}\ K H}}\skt{ वेदनारागशान्तये ॥७०॥ }\pend
\pstart\versepar \skt{मृदु भृष्टैस्तिलैर्वापि सिद्धार्थकसमायुतैः ।}\pend

\pstart\versepar \ledsidenote{\passref{6.17.89}}\edtext{\skt{मातुलुङ्गरसोपेतः}}{\Afootnote{K H; \skt{मातुलुङ्गरसोपेतैः}\ A}}\skt{ }\edtext{\skt{सद्यः}}{\Afootnote{K H; \skt{सुखालेप}\ A}}\skt{ }\edtext{\skt{शर्मार्थिने}}{\Afootnote{K H; \skt{स्तदर्थ}\ A}}\skt{ }\edtext{\skt{हितः}}{\Afootnote{K H; \emph{om.}\ A}}\skt{ ॥७१॥ }\pend
\pstart\versepar \skt{}\edtext{\skt{पयस्याशारिवापत्रमञ्जिष्ठामधुकैरपि}}{\lemma{\skt{पयस्याशारिवापत्रमञ्जिष्ठामधुकैर्}}\Afootnote{\skt{कृत् । पयस्याशारिवापत्रमञ्जिष्ठामधुकैर्}\ A; \skt{पयस्यासारिवापत्रमञ्जिष्ठामधुकैर्}\ K H}}\skt{ ।}\pend

\pstart\versepar \ledsidenote{\passref{6.17.90}}\edtext{\edtext{\skt{अजाक्षीरार्दितैर्लेपः}}{\lemma{\skt{अजाक्षीरार्दितैर्}}\Afootnote{\skt{अजाक्षीरान्वितैर्}\ A; \skt{अजाक्षीराद्धितैर्}\ K; \skt{अजक्षीरार्द्दितैर्}\ H}}}{\lemma{\skt{लेपः}}\Afootnote{A H; \skt{लेपिः}\ K}}\skt{ सुखोष्णः पथ्य उच्यते ॥७२॥ }\pend
\pstart\versepar \skt{}\edtext{\skt{दारुपद्मकशुण्ठीभिरेवमेव}}{\lemma{\skt{दारुपद्मकशुण्ठीभिर्}}\Afootnote{A H; \skt{दारुपद्मकसुण्ठीभिर्}\ K}}\skt{ }\edtext{\skt{कृतोऽथवा}}{\Afootnote{K H; \skt{कृतोऽपि वा ॥}\ A}}\skt{ ।}\pend

\pstart\versepar \ledsidenote{\passref{6.17.91}}\edtext{\skt{द्राक्षामधुकलोध्रैर्वा}}{\Afootnote{K H; \skt{द्राक्षामधुककुष्ठैर्वा}\ A}}\skt{ }\edtext{\skt{स्यादेवं}}{\Afootnote{K H; \skt{तद्व}\ A}}\skt{ }\edtext{\skt{सैधवायुतैः}}{\Afootnote{\skt{त्सैन्धवसंयुतैः ।}\ A; \skt{सैन्धवायुतैः ।}\ K H}}\skt{ ॥७३॥ }\pend
\pstart\versepar \skt{}\edtext{\skt{लोध्रसैन्धवमृद्वीकामधुकैर्वाप्यजापयः}}{\lemma{\skt{लोध्रसैन्धवमृद्वीकामधुकैर्}}\Afootnote{\skt{रोध्रसैन्धवमृद्वीकामधुकैर्}\ A; \skt{लोध्रं सैन्धवमृद्वीकामधुकैर्}\ K; \skt{लोध्रसैन्धवमृद्वीकामधूकैर्}\ H}}\skt{ ।}\pend

\pstart\versepar \ledsidenote{\passref{6.17.92}}\edtext{\skt{श्रितं}}{\Afootnote{\skt{शृतं}\ A; \skt{स्रितं}\ K; \skt{स्रिते}\ H}}\skt{ सेके प्रयोक्तव्यं }\edtext{\skt{रुजारागविनाशनं}}{\Afootnote{H; \skt{रुजारागनिवारणम् ।}\ A; \skt{रुजारागविनासनं ।}\ K}}\skt{ ॥७४॥ }\pend
\pstart\versepar \skt{}\edtext{\skt{मधुकोत्पलकुष्ठैर्वा}}{\Afootnote{A K; \skt{मधूकोत्पलकुष्ठैर्व्वा}\ H}}\skt{ }\edtext{\skt{द्राक्षालाक्षासितायुतैः}}{\Afootnote{K H; \skt{द्राक्षालाक्षासितायुतैः । ससैन्धवैः शृतं क्षीरं रुजारागनिबर्हणम् ॥}\ A}}\skt{ ।}\pend

\pstart\versepar \ledsidenote{\passref{6.17.93}}\edtext{\skt{शतावरीपृथक्पर्णीमुस्तामधुकपद्मकैः}}{\Afootnote{H; \skt{शतावरीपृथक्पर्णीमुस्तामलकपद्मकैः ।}\ A; \skt{सतावरीपृथक्पर्ण्णीमुस्तामधुकपद्मकैः ।}\ K}}\skt{ ॥७५॥ }\pend
\pstart\versepar \skt{}\edtext{\skt{ससैन्धवैः}}{\Afootnote{K H; \skt{साजक्षीरैः}\ A}}\skt{ }\edtext{\skt{श्रितं}}{\Afootnote{\skt{शृतं}\ A; \skt{स्रितं}\ K H}}\skt{ चापि }\edtext{\edtext{\skt{सुखोष्णमवचारयेत्}}{\lemma{\skt{अवचारयेत्}}\Afootnote{K; \skt{अवतारयेत् ॥}\ H}}\skt{ ।}}{\lemma{\skt{चापि… अवचारयेत्}}\Afootnote{\emph{om.}\ A}}\pend

\pstart\versepar \ledsidenote{\passref{6.17.94}}\skt{वातघ्नसिद्धे पयसि सिद्धं }\edtext{\skt{सर्पिश्चतुर्गुणे}}{\lemma{\skt{चतुर्गुणे}}\Afootnote{A H; \skt{चतुर्गुणेः ।}\ K}}\skt{ ॥७६॥ }\pend
\pstart\versepar \skt{काकोल्यादिप्रतीवापं }\edtext{\skt{सर्वकर्मसु}}{\Afootnote{K H; \skt{तद्युञ्ज्यात्सर्वकर्मसु ॥}\ A}}\skt{ }\edtext{\skt{चादिशेत्}}{\Afootnote{K H; \emph{om.}\ A}}\skt{ ॥}\pend

\pstart\versepar \ledsidenote{\passref{6.17.95}}\edtext{\skt{शाम्यत्येवं}}{\lemma{\skt{शाम्यत्य्}}\Afootnote{A H; \skt{साम्यत्य्}\ K}}\skt{ न चेच्छूलं स्निग्धस्विन्नस्य मोक्षयेत् ॥७७॥ }\pend
\pstart\versepar \skt{ततः सिरां }\edtext{\skt{दहेच्चापि}}{\Afootnote{K H; \skt{दहेद्वापि}\ A}}\skt{ }\edtext{\skt{मतिमान्कीर्त्तितं}}{\Afootnote{A H; \skt{मतिमां कीर्त्तितं}\ K}}\skt{ यथा ।}\pend

\pstart\versepar \ledsidenote{\passref{6.17.96}}\edtext{\edtext{\skt{दृष्टेरतः}}{\lemma{\skt{दृष्टेर्}}\Afootnote{A K; \skt{दृष्टैर्}\ H}}}{\lemma{\skt{अतः}}\Afootnote{A H; \emph{om.}\ K}}\skt{ }\edtext{\edtext{\skt{प्रसादार्थमञ्जने}}{\lemma{\skt{प्रसादार्थम्}}\Afootnote{A H; \skt{प्रसादे}\ K}}}{\lemma{\skt{अञ्जने}}\Afootnote{A H; \skt{जने}\ K}}\skt{ शृणु मे शुभे ॥७८॥ }\pend
\pstart\versepar \skt{}\edtext{\skt{मेषशृङ्ग्यास्तु}}{\Afootnote{K H; \skt{मेषशृङ्गस्य}\ A}}\skt{ पुष्पाणि शिरीषधवयोरपि ।}\pend

\pstart\versepar \ledsidenote{\passref{6.17.97}}\edtext{\edtext{\skt{जात्याश्चैव}}{\lemma{\skt{जात्याश्}}\Afootnote{K H; \skt{सुमनायाश्}\ A}}\skt{ हि}}{\lemma{\skt{चैव हि}}\Afootnote{K H; \skt{च}\ A}}\skt{ }\edtext{\skt{तुल्यानि}}{\Afootnote{\skt{पुष्पाणि}\ A; \skt{तुल्याणि}\ K H}}\skt{ }\edtext{\skt{मुक्तावैडूर्यमेव}}{\lemma{\skt{मुक्तावैडूर्यम्}}\Afootnote{K H; \skt{मुक्ता वैदूर्यम्}\ A}}\skt{ च ॥७९॥ }\pend
\pstart\versepar \skt{अजाक्षीरेण }\edtext{\skt{सम्पिष्य}}{\Afootnote{A K\textsuperscript{pc} H; \emph{om.}\ K\textsuperscript{ac}}}\skt{ ताम्रे सप्ताहमावपेत् ।}\pend

\pstart\versepar \ledsidenote{\passref{6.17.98}}\edtext{\skt{वर्त्तीस्तास्तु}}{\Afootnote{K H; \skt{प्रविधाय च तद्वर्तीर्}\ A}}\skt{ }\edtext{\skt{ततः}}{\Afootnote{K H; \emph{om.}\ A}}\skt{ }\edtext{\skt{कृत्वा}}{\Afootnote{K H; \emph{om.}\ A}}\skt{ }\edtext{\skt{योजयेदञ्जने}}{\Afootnote{K H; \skt{योजयेच्चाञ्जने}\ A}}\skt{ भिषक् ॥८०॥ }\pend
\pstart\versepar \skt{स्रोतोजं विद्रुमं }\edtext{\skt{फेनः}}{\Afootnote{K H; \skt{फेनं}\ A}}\skt{ सागरस्य }\edtext{\skt{मनःशिला}}{\Afootnote{K H; \skt{मनःशिलाम् ॥}\ A}}\skt{ ।}\pend

\pstart\versepar \ledsidenote{\passref{6.17.99}}\edtext{\skt{मरिचानि च ता}}{\Afootnote{K H; \skt{मरिचानि च त}\ A}}\skt{ }\edtext{\skt{वर्तीः}}{\Afootnote{K; \skt{द्वर्तीः}\ A; \skt{वर्ती}\ H}}\skt{ }\edtext{\skt{कारयेद्वापि}}{\Afootnote{K H; \skt{कारयेच्चापि}\ A}}\skt{ पूर्ववत् । }\pend
\pstart\versepar \skt{दृष्टेः }\edtext{\edtext{\skt{स्थैर्यार्थमेतास्तु}}{\lemma{\skt{स्थैर्यार्थम्}}\Afootnote{A K; \skt{स्थैर्यम्}\ H}}}{\lemma{\skt{एतास्तु}}\Afootnote{K; \skt{एतत्तु}\ A; \skt{एतास्थ}\ H}}\skt{ विदध्यादञ्जने }\edtext{\skt{हिताः}}{\Afootnote{K; \skt{हितम् ॥}\ A; \skt{हिता ॥}\ H}}\skt{ ॥८१॥}\edtext{}{\lemma{\skt{पूर्ववत् । दृष्टेः}}\Afootnote{H; \skt{पूर्ववत् । दृष्टि}\ A; \skt{पूर्ववत् । दृष्टैः}\ K}}\pend

\pstart\versepar \ledsidenote{\passref{6.17.100}}\skt{भूयो वक्ष्यामि मुख्यानि }\edtext{\skt{विस्तरेणाञ्जनानि तु}}{\Afootnote{K; \skt{विस्तरेणाञ्जनानि च ।}\ A H}}\skt{ । }\pend
\pstart\versepar \skt{कल्पे नानाप्रकाराणि }\edtext{\edtext{\skt{तान्यपीह}}{\lemma{\skt{तान्य्}}\Afootnote{K H; \skt{ता}\ A}}}{\lemma{\skt{अपीह}}\Afootnote{K H; \skt{न्यपीह}\ A}}\skt{ }\edtext{\skt{प्रयोजयेत् ॥ ८२ ॥ इति}}{\Afootnote{H; \skt{प्रयोजयेत् ॥}\ A; \skt{प्रयोजयेत् ॥ दिति ॥ सालाक्ये ॥ १० ६ ॥}\ K}}\skt{ ॥}\pend

\endnumbering
\beginnumbering
\eledsection*{Adhyāya 18}
\pstart \ledsidenote{\passref{6.18.1}}\skt{अथातः क्रियाकल्पं }\edtext{\skt{व्याख्यास्यामः}}{\Afootnote{A H; \skt{व्याव्याख्यास्यामः ॥}\ K}}\skt{ ॥}\pend

\pstart\versepar \ledsidenote{\passref{6.18.3}}\skt{सर्वशास्त्रार्थतत्वज्ञस्तपोदृष्टिरुदारधीः । }\pend
\pstart\versepar \skt{}\edtext{\skt{वैश्वामित्रिं}}{\Afootnote{K H; \skt{वैश्वामित्रं}\ A}}\skt{ शशासाथ शिष्यं }\edtext{\skt{काषिपतिर्मुनिः}}{\lemma{\skt{काषिपतिर्}}\Afootnote{K; \skt{काशिपतिर्}\ A H}}\skt{ ।}\pend

\pstart\versepar \ledsidenote{\passref{6.18.4}}\skt{तर्पणम्पुटपाकश्च }\edtext{\skt{सेका}}{\Afootnote{K H; \skt{सेक}\ A}}\skt{ }\edtext{\skt{आश्च्योतनाञ्जने}}{\Afootnote{A K; \skt{श्च्योतनमञ्जनं ।}\ H}}\skt{ । }\pend
\pstart\versepar \skt{तत्र तत्रोपदिष्टानि तेषां व्यासं निबोध मे ।}\pend

\pstart\versepar \ledsidenote{\passref{6.18.5}}\skt{संशुद्धदेहशिरसो }\edtext{\skt{जीर्ण्णान्नस्येच्छतः}}{\Afootnote{K; \skt{जीर्णान्नस्य}\ A; \skt{जीर्ण्णान्नस्येप्सतः}\ H}}\skt{ }\edtext{\skt{शुभं}}{\Afootnote{K H; \skt{शुभे}\ A}}\skt{ । }\pend
\pstart\versepar \skt{}\edtext{\skt{पूर्वाह्णे}}{\Afootnote{K H; \skt{दिने । पूर्वाह्णे}\ A}}\skt{ }\edtext{\skt{चापराह्णे वा}}{\Afootnote{K H; \skt{वापराह्णे वा}\ A}}\skt{ }\edtext{\skt{कार्यमक्ष्णोस्तु}}{\lemma{\skt{अक्ष्णोस्तु}}\Afootnote{A K; \skt{अक्ष्णोश्च}\ H}}\skt{ }\edtext{\skt{तर्पनं}}{\Afootnote{K; \skt{तर्पणम्}\ A H}}\skt{ ।}\pend

\pstart\versepar \ledsidenote{\passref{6.18.6}}\skt{वातातपरजोहीने वेश्मन्युत्तानशायिनः । }\pend
\pstart\versepar \skt{}\edtext{\skt{आधारौ}}{\Afootnote{A K; \skt{आधारो}\ H}}\skt{ माषचूर्ण्णेन क्लिन्नेन परिमण्डलौ ।}\pend

\pstart\versepar \ledsidenote{\passref{6.18.7}}\edtext{\skt{नेत्रकोशस्य}}{\Afootnote{K H; \skt{समौ दृढावसंबाधौ कर्तव्यौ नेत्रकोश}\ A}}\skt{ }\edtext{\skt{संवेश्य}}{\Afootnote{K H; \emph{om.}\ A}}\skt{ }\edtext{\skt{त्वसम्बाधौ}}{\Afootnote{K H; \skt{योः ।}\ A}}\skt{ }\edtext{\skt{दृढौ}}{\Afootnote{K; \emph{om.}\ A; \skt{समौ}\ H}}\skt{ }\edtext{\skt{समौ}}{\Afootnote{K; \emph{om.}\ A; \skt{दृढौ ।}\ H}}\skt{ । }\pend
\pstart\versepar \skt{पूरयेद्घृतमण्डस्य }\edtext{\skt{निषर्ण्णस्य}}{\Afootnote{K H; \skt{विलीनस्य}\ A}}\skt{ सुखोदके ।}\pend

\pstart\versepar \ledsidenote{\passref{6.18.8}}\edtext{\skt{आपक्ष्माग्रान्ततः}}{\Afootnote{K H; \skt{आपक्ष्माग्रात्ततः}\ A}}\skt{ }\edtext{\skt{स्थाप्य}}{\Afootnote{K; \skt{स्थाप्यं}\ A H}}\skt{ पञ्च }\edtext{\skt{षड्वाक्छतानि तु}}{\Afootnote{K H; \skt{तद्वाक्शतानि तु ।}\ A}}\skt{ । }\pend
\pstart\versepar \skt{स्वस्थे कफे }\edtext{\skt{षट्पित्ते}}{\Afootnote{A K\textsuperscript{pc} H; \skt{ष पित्ते}\ K\textsuperscript{ac}}}\skt{ }\edtext{\skt{ष्टौ}}{\Afootnote{K; \skt{ऽष्टौ}\ A H}}\skt{ दश वाते }\edtext{\skt{तदुत्तमं}}{\Afootnote{\skt{तदुत्तमम्}\ A K H}}\skt{ ।}\pend

\pstart\versepar \ledsidenote{\passref{6.18.9}}\skt{रोगस्थानविशेषेण केचित्कालं }\edtext{\skt{प्रवक्षते}}{\Afootnote{K; \skt{प्रचक्षते ।}\ A; \skt{प्रवक्ष्यते ।}\ H}}\skt{ । }\pend
\pstart\versepar \skt{यथाक्रमोपदिष्टेषु }\edtext{\edtext{\skt{त्रीण्येकं}}{\lemma{\skt{त्रीण्य्}}\Afootnote{K H; \skt{त्री}\ A}}}{\lemma{\skt{एकं}}\Afootnote{K H; \skt{ण्येकं}\ A}}\skt{ पञ्च सप्त }\edtext{\skt{वा}}{\Afootnote{K H; \skt{च ॥}\ A}}\skt{ ।}\pend

\pstart\versepar \ledsidenote{\passref{6.18.10}}\edtext{\edtext{\skt{दशादृष्ट्यामथाष्टौ}}{\lemma{\skt{दशादृष्ट्याम्}}\Afootnote{K; \skt{दश दृष्ट्याम्}\ A H}}\skt{ वा}}{\lemma{\skt{अथाष्टौ वा}}\Afootnote{K H; \skt{अथाष्टौ च}\ A}}\skt{ वाक्छतानि }\edtext{\skt{विशारदाः}}{\Afootnote{K H; \skt{विभावये}\ A}}\skt{ । }\pend
\pstart\versepar \skt{}\edtext{\skt{पूर्ण्णे}}{\Afootnote{K H; \skt{त् । ततश्}\ A}}\skt{ चापांगतः स्नेहं }\edtext{\skt{संस्राव्याक्षिविशोधयेत्}}{\Afootnote{K H; \skt{स्रावयित्वाक्षि शोधयेत् ॥}\ A}}\skt{ ।}\pend

\pstart\versepar \ledsidenote{\passref{6.18.11}}\skt{स्विन्नेन यवपिष्टेन स्नेहवीर्येरितं ततः । }\pend
\pstart\versepar \skt{यथास्वं धूमपानेन कफमस्य विशोधयेत् ।}\pend

\pstart\versepar \ledsidenote{\passref{6.18.12}}\skt{एकाहं वा त्र्यहम्वापि }\edtext{\skt{यथादोषं}}{\Afootnote{K; \skt{पञ्चाहं चे}\ A H}}\skt{ }\edtext{\skt{विशेषतः}}{\Afootnote{K; \skt{ष्यते परम्}\ A H}}\skt{ । }\pend
\pstart\versepar \skt{तर्पणे तृप्तिलिङ्गानि }\edtext{\skt{नेत्रस्येमानि}}{\Afootnote{A K; \skt{नेत्रस्येशनि}\ H}}\skt{ लक्षयेत् ।}\pend

\pstart\versepar \ledsidenote{\passref{6.18.13}}\skt{सुखस्वप्नावबोधत्वं }\edtext{\skt{वैशद्यं}}{\Afootnote{A K; \skt{…}\ H}}\skt{ वर्ण्णपाटवं । }\pend
\pstart\versepar \skt{}\edtext{\skt{निर्वृत्तिर्व्याधिसंशान्तिः}}{\lemma{\skt{व्याधिसंशान्तिः}}\Afootnote{K; \skt{व्याधिविध्वंसः}\ A; \skt{व्व्याधिशान्तिश्च}\ H}}\skt{ क्रियालाघवमेव च ।}\pend

\pstart\versepar \ledsidenote{\passref{6.18.14}}\skt{गुर्वाविलमतिस्निग्धमश्रुकण्डूपदेहवत् । }\pend
\pstart\versepar \skt{ज्ञेयन्दोषसमुत्क्लिष्टं नेत्रमत्यर्थतर्पितं ।}\pend

\pstart\versepar \ledsidenote{\passref{6.18.15}}\edtext{\skt{रूक्षमाविलमस्रावमसहं}}{\lemma{\skt{अस्रावम्}}\Afootnote{K H; \skt{अस्राढ्यम्}\ A}}\skt{ रूपदर्शने । }\pend
\pstart\versepar \skt{व्याधिवृद्धिश्च }\edtext{\skt{तज्ञेयं}}{\Afootnote{K; \skt{तज्ज्ञेयं}\ A; \skt{तज्जेयं}\ H}}\skt{ }\edtext{\edtext{\skt{हीनतर्पणमक्षितः}}{\lemma{\skt{हीनतर्पणम्}}\Afootnote{K H; \skt{हीनतर्पित}\ A}}}{\lemma{\skt{अक्षितः}}\Afootnote{K H; \skt{मक्षि च ॥}\ A}}\skt{ ।}\pend

\pstart \ledsidenote{\passref{6.18.16}}\edtext{\skt{धूमनस्याञ्जनैः}}{\Afootnote{K H; \skt{अजयोर्दोषबाहुल्यात्प्रयतेत चिकित्सिते । धूमनस्याञ्जनैः}\ A}}\skt{ }\edtext{\edtext{\skt{सेकैरवपीडैश्च}}{\lemma{\skt{सेकैर्}}\Afootnote{K H; \skt{सेकै}\ A}}\skt{ }\edtext{\skt{तञ्जयेत्}}{\lemma{\skt{जयेत्}}\Afootnote{K H; \skt{गवित् ॥}\ A}}}{\lemma{\skt{अवपीडैश्च तञ्}}\Afootnote{K H; \skt{रूक्षैः स्निग्धैश्च यो}\ A}}\skt{ ।}\pend

\pstart\versepar \ledsidenote{\passref{6.18.17}}\edtext{\skt{ताम्यते}}{\Afootnote{K; \skt{ताम्यत्य}\ A; \skt{ताम्येद्वि}\ H}}\skt{ }\edtext{\skt{परिशुष्कं}}{\Afootnote{K H; \skt{तिविशुष्कं}\ A}}\skt{ }\edtext{\skt{यद्रूक्षं}}{\Afootnote{K H; \skt{यद्रूक्षं}\ A}}\skt{ }\edtext{\skt{यच्चापि}}{\Afootnote{K H; \skt{यच्चाति}\ A}}\skt{ दारुणं । }\pend
\pstart\versepar \skt{शीर्ण्णपक्ष्माविलं जिह्मं रोगक्लिष्टञ्च यद्भृशं ।}\pend

\pstart\versepar \ledsidenote{\passref{6.18.18}}\edtext{\skt{तदक्षि}}{\Afootnote{K H; \skt{तदक्षि}\ A}}\skt{ }\edtext{\skt{तर्पणादेव}}{\Afootnote{K H; \skt{तर्पणादेव}\ A}}\skt{ }\edtext{\edtext{\edtext{\skt{रोगैरेतैर्विमुच्यते}}{\lemma{\skt{रोगैर्}}\Afootnote{K H; \emph{om.}\ A}}}{\lemma{\skt{एतैर्}}\Afootnote{K H; \emph{om.}\ A}}}{\lemma{\skt{विमुच्यते}}\Afootnote{K H; \skt{लभेतो}\ A}}\skt{ । }\pend
\pstart\versepar \skt{}\edtext{\skt{दुर्दिनात्युष्णशीतेषु}}{\Afootnote{K; \skt{र्जामसंशयम् । दुर्दिनात्युष्णशीतेषु}\ A; \skt{दुर्द्दिने नात्युष्णशीते}\ H}}\skt{ }\edtext{\skt{चिन्तायासश्रमेषु}}{\Afootnote{K; \skt{चिन्तायासभ्रमेषु}\ A; \skt{चिन्तायासस्रमेषु}\ H}}\skt{ च ।}\pend

\pstart\versepar \ledsidenote{\passref{6.18.19}}\skt{अशान्तोपद्रवे }\edtext{\skt{चाक्षितर्पनं}}{\Afootnote{K; \skt{चाक्षिण तर्पणं}\ A; \skt{चाक्षितर्प्पणं}\ H}}\skt{ }\edtext{\skt{नोपदिश्यते}}{\Afootnote{K H; \skt{न प्रशस्यते ।}\ A}}\skt{ । }\pend
\pstart\versepar \skt{}\edtext{\skt{पुटपाकश्च}}{\lemma{\skt{च}}\Afootnote{K H; \skt{तथैतेषु नस्यं येषु च}\ A}}\skt{ }\edtext{\skt{नस्ये च}}{\Afootnote{K; \skt{गर्हितम् ॥}\ A; \skt{नस्येन}\ H}}\skt{ }\edtext{\skt{प्रतिषिद्धास्तु ये}}{\Afootnote{K; \emph{om.}\ A; \skt{प्रतिषिद्धाश्च ये}\ H}}\skt{ }\edtext{\skt{गदाः}}{\Afootnote{K H; \emph{om.}\ A}}\skt{ ।}\pend

\pstart \ledsidenote{\passref{6.18.20}}\edtext{\skt{तस्मात्प्रशान्तदोषेषु}}{\Afootnote{K H; \skt{तर्पणार्हा न ये प्रोक्ताः स्नेहपानाक्षमाश्च ये । ततः प्रशान्तदोषेषु}\ A}}\skt{ पुटपाकक्षमेषु च ।}\pend

\pstart\versepar \ledsidenote{\passref{6.18.21}}\edtext{\skt{पुटपाकः}}{\Afootnote{A K\textsuperscript{pc} H; \skt{पुटपाक}\ K\textsuperscript{ac}}}\skt{ प्रयोक्तव्यो नेत्रेषु भिषजा भवेत् । }\pend
\pstart\versepar \skt{}\edtext{\skt{स्नेहनोलेखनश्चापि}}{\lemma{\skt{चापि}}\Afootnote{K H; \skt{यश्}\ A}}\skt{ }\edtext{\skt{तथा}}{\Afootnote{K H; \skt{च}\ A}}\skt{ }\edtext{\skt{रोपण}}{\Afootnote{K; \skt{रोपणी}\ A; \skt{रोपणम्}\ H}}\skt{ }\edtext{\skt{एव च}}{\Afootnote{K H; \skt{यश्च स त्रिधा ॥}\ A}}\skt{ ।}\edtext{}{\lemma{\skt{भवेत् । स्नेहनोलेखनश्}}\Afootnote{K; \skt{भवेत् । स्नेहनो लेखनी}\ A; \skt{भवेत् । स्नेहनोल्लेखनश्}\ H}}\pend

\pstart\versepar \ledsidenote{\passref{6.18.22}}\edtext{\skt{त्रिविधः}}{\Afootnote{K H; \emph{om.}\ A}}\skt{ }\edtext{\skt{सहितो}}{\Afootnote{K H; \skt{हितः}\ A}}\skt{ }\edtext{\skt{रूक्षो}}{\Afootnote{K; \skt{स्निग्धोऽतिरूक्षस्य}\ A; \skt{रूक्षे}\ H}}\skt{ }\edtext{\skt{स्निग्धः}}{\Afootnote{K H; \skt{स्निग्ध}\ A}}\skt{ }\edtext{\skt{स्निग्धे च}}{\Afootnote{K H; \skt{स्यापि च}\ A}}\skt{ लेखनः । }\pend
\pstart\versepar \skt{}\edtext{\edtext{\skt{दृष्टेर्बलार्थमितरः}}{\lemma{\skt{बलार्थम्}}\Afootnote{K; \skt{बलार्थ}\ A H}}}{\lemma{\skt{इतरः}}\Afootnote{K; \skt{मपरः}\ A; \skt{तिमिरः}\ H}}\skt{ }\edtext{\skt{पित्तासृग्व्रणवातहा}}{\Afootnote{K; \skt{पित्तासृग्व्रणवातनुत् ॥}\ A; \skt{पित्तासृग्ग्रहवातहा ॥}\ H}}\skt{ ।}\pend

\pstart\versepar \ledsidenote{\passref{6.18.23}}\edtext{\skt{स्नेहमांसवसामज्जामेदःस्वाद्वौषधीकृतः}}{\Afootnote{K H; \skt{स्नेहमांसवसामज्जमेदःस्वाद्वौषधैः कृतः ।}\ A}}\skt{ ॥ }\pend
\pstart\versepar \skt{}\edtext{\skt{स्नेहार्थः}}{\Afootnote{K H; \skt{स्नेहतः}\ A}}\skt{ पुटपाकः }\edtext{\skt{स्याद्धार्यो द्वे}}{\Afootnote{K; \skt{तु धार्यो द्वे}\ A; \skt{स्याद्वायोर्द्वो}\ H}}\skt{ }\edtext{\skt{वाक्छते तु}}{\Afootnote{A K; \skt{वाक्षते तु}\ H}}\skt{ सः ॥}\pend

\pstart \ledsidenote{\passref{6.18.24}}\skt{जांगलानां }\edtext{\skt{यकृत्मान्सै}}{\Afootnote{K H; \skt{यकृन्मांसै}\ A}}\skt{ }\edtext{\edtext{\skt{लेखनैश्चापि}}{\lemma{\skt{लेखनैश्}}\Afootnote{K; \skt{र्लेखन}\ A; \skt{क्लेखनैश्}\ H}}}{\lemma{\skt{चापि}}\Afootnote{K H\textsuperscript{pc}; \skt{द्रव्य}\ A; \skt{चा}\ H\textsuperscript{ac}}}\skt{ }\edtext{\skt{सम्भृतः}}{\Afootnote{K H; \skt{संभृतैः । कृष्णलोहरजस्ताम्रशङ्खविद्रुमसिन्धुजैः ॥}\ A}}\pend

\pstart \ledsidenote{\passref{6.18.25}}\edtext{\skt{लेहनो}}{\Afootnote{K H; \skt{समुद्रफेनकासीसस्रोतोजदधिमस्तुभिः । लेहनो}\ A}}\skt{ वाक्छतं तस्य }\edtext{\skt{हितं}}{\Afootnote{K H; \skt{परं}\ A}}\skt{ }\edtext{\edtext{\skt{धारणमिष्यते}}{\lemma{\skt{धारणम्}}\Afootnote{K H; \skt{धारण}\ A}}}{\lemma{\skt{इष्यते}}\Afootnote{K H; \skt{मुच्यते ॥}\ A}}\skt{ ।}\pend

\pstart\versepar \ledsidenote{\passref{6.18.26}}\edtext{\skt{स्तन्यजाङ्गलमध्वाज्यजीवनद्रव्यविपाचितः}}{\Afootnote{K; \skt{स्तन्यजाङ्गलमध्वाज्यतिक्तद्रव्यविपाचितः ।}\ A; \skt{स्तन्यजाङ्गलमध्वाज्यस्तिक्तद्रव्यविपाचितः ।}\ H}}\skt{ । }\pend
\pstart\versepar \skt{लेखनात्त्रिगुणं }\edtext{\skt{धार्य}}{\Afootnote{K; \skt{धार्यः}\ A H}}\skt{ पुटपाकस्तु }\edtext{\skt{रोपनः}}{\Afootnote{K; \skt{रोपणः ॥}\ A H}}\skt{ ।}\pend

\pstart \ledsidenote{\passref{6.18.27}}\edtext{\skt{स्वेदधूमौ च}}{\Afootnote{K H; \emph{om.}\ A}}\skt{ }\edtext{\skt{वितरेद्द्वयोर्हित्वा}}{\lemma{\skt{वितरेद्द्वयोर्}}\Afootnote{K H; \skt{वितरेत्तर्पणोक्तं तु धूमं}\ A}}\skt{ तु }\edtext{\skt{रोपणं}}{\Afootnote{K H; \skt{रोपणम् । स्नेहस्वेदौ द्वयोः कार्यौ कार्यो नैव च रोपणे ॥}\ A}}\skt{ ।}\pend

\pstart\versepar \ledsidenote{\passref{6.18.28}}\skt{एकाहं वा }\edtext{\edtext{\skt{द्व्यहम्वापि}}{\lemma{\skt{द्व्यहम्}}\Afootnote{A K; \skt{अहञ्}\ H}}}{\lemma{\skt{वापि}}\Afootnote{A K; \skt{चैषां}\ H}}\skt{ }\edtext{\skt{त्र्यह}}{\Afootnote{K; \skt{त्र्यहं}\ A H}}\skt{ }\edtext{\edtext{\skt{वाप्यवचारणं}}{\lemma{\skt{वाप्य्}}\Afootnote{K H; \skt{वा}\ A}}}{\lemma{\skt{अवचारणं}}\Afootnote{K H; \skt{प्यवचारणम् ।}\ A}}\skt{ । }\pend
\pstart\versepar \skt{यन्त्रणा तु क्रियाकाला }\edtext{\skt{द्विगुणं}}{\Afootnote{K H; \skt{द्द्विगुणं}\ A}}\skt{ कालमिष्यते ।}\pend

\pstart\versepar \ledsidenote{\passref{6.18.29}}\edtext{\edtext{\edtext{\edtext{\skt{तेजांस्यनिलमाकाशमादर्शम्भास्वराणि}}{\lemma{\skt{तेजांस्य्}}\Afootnote{K H; \skt{तेजां}\ A}}}{\lemma{\skt{अनिलम्}}\Afootnote{K H; \skt{स्यनिल}\ A}}}{\lemma{\skt{आकाशम्}}\Afootnote{K H; \skt{माकाश}\ A}}}{\lemma{\skt{आदर्शम्}}\Afootnote{K H; \skt{मादर्शं}\ A}}\skt{ च । }\pend
\pstart\versepar \skt{न }\edtext{\skt{पश्येत्तर्पिताक्षश्च}}{\Afootnote{K H; \skt{क्षेत तर्पिते नेत्रे}\ A}}\skt{ }\edtext{\skt{पुटपाककृतो}}{\Afootnote{K; \skt{पुटपाककृते}\ A; \skt{पुटाख्यो यस्य वाक्यतः ।}\ H}}\skt{ }\edtext{\skt{पि}}{\Afootnote{K; \skt{तथा ॥}\ A; \emph{om.}\ H}}\skt{ }\edtext{\skt{वा}}{\Afootnote{K; \emph{om.}\ A H}}\skt{ ।}\edtext{}{\lemma{\skt{भास्वराणि च । न}}\Afootnote{K H; \skt{भास्वराणि च । ने}\ A}}\pend

\pstart\versepar \ledsidenote{\passref{6.18.30}}\skt{मिथ्योपचारादनयोर्यो }\edtext{\skt{दोषः}}{\Afootnote{K H\textsuperscript{pc}; \skt{व्याधि}\ A; \skt{दोष}\ H\textsuperscript{ac}}}\skt{ }\edtext{\skt{समुदीर्यते}}{\Afootnote{K H; \skt{रुपजायते ।}\ A}}\skt{ । }\pend
\pstart\versepar \skt{}\edtext{\edtext{\skt{अञ्जनाश्च्योतनास्वेदैर्यथास्वं}}{\lemma{\skt{अञ्जनाश्च्योतनास्वेदैर्}}\Afootnote{K; \skt{अञ्जनाश्च्योतनस्वेदै}\ A; \skt{अञ्जनाश्च्योतनाश्चैदैर्}\ H}}\skt{ }\edtext{\skt{तमुपाचरेत्}}{\lemma{\skt{उपाचरेत्}}\Afootnote{K H; \skt{मुपाचरेत् ॥}\ A}}}{\lemma{\skt{यथास्वं तम्}}\Afootnote{K H; \skt{र्यथास्वं त}\ A}}\skt{ ।}\pend

\pstart \ledsidenote{\passref{6.18.33}}\edtext{\skt{अतीसारेषु}}{\Afootnote{K H; \skt{अत ऊर्ध्वं प्र}\ A}}\skt{ वक्ष्यामि }\edtext{\skt{पुटपाकप्रसाधनं}}{\Afootnote{K H; \skt{पुटपाकप्रसाधनम् । द्वौ बिल्वमात्रौ श्लक्ष्णस्य पिण्डौ मांसस्य पेषितौ ॥}\ A}}\skt{ ।}\pend

\pstart\versepar \ledsidenote{\passref{6.18.34}}\skt{वनस्पत्युरुवूकानां पाटल्याः ककुभस्य वा । }\pend
\pstart\versepar \skt{}\edtext{\skt{काष्ठैः }\edtext{\skt{स्युर्गोमयैर्वापि}}{\lemma{\skt{गोमयैर्}}\Afootnote{K; \skt{गोमये}\ H}}}{\lemma{\skt{काष्ठैः स्युर्}}\Afootnote{K\textsuperscript{pc}; \skt{काष्ठै स्युर्}\ K\textsuperscript{ac}; \skt{काष्ठे स्याङ्}\ H}}\skt{ }\edtext{\edtext{\skt{पक्तिर्वेष्टनकर्मसु}}{\lemma{\skt{पक्तिर्}}\Afootnote{K; \skt{पक्त}\ H}}}{\lemma{\skt{वेष्टनकर्मसु}}\Afootnote{K; \skt{व्येषु न कर्म्मसु ॥}\ H}}\skt{ ।}\edtext{}{\lemma{\skt{… वेष्टनकर्मसु}}\Afootnote{\emph{om.}\ A}}\pend

\pstart \ledsidenote{\passref{6.18.35}}\edtext{\skt{काश्मरीकुमुदैरण्डपद्मिनीकदलच्छदैः}}{\Afootnote{K; \skt{(काश्मरीकुमुदैरण्डपद्मिनीकदलीभवैः । ) मृदावलिप्तमङ्गारैः खादिरैरवकूलयेत् ॥}\ A; \skt{काश्मीरकुमुदैरण्डपद्मिनकदलच्छदैः ।}\ H}}\skt{ ।}\pend

\pstart\versepar \ledsidenote{\passref{6.18.44}}\skt{यथादोषोपयुक्तन्तु नातिप्रबलमोजसा । }\pend
\pstart\versepar \skt{}\edtext{\edtext{\skt{रोगमाश्च्योतनं}}{\lemma{\skt{रोगम्}}\Afootnote{K H; \skt{रोग}\ A}}}{\lemma{\skt{आश्च्योतनं}}\Afootnote{K H; \skt{माश्च्योतनं}\ A}}\skt{ हन्ति सेको }\edtext{\skt{पि}}{\Afootnote{K H; \skt{तु}\ A}}\skt{ }\edtext{\skt{बलवत्तरं}}{\Afootnote{A K; \skt{बलवत्तरः ।}\ H}}\skt{ ।}\pend

\pstart\versepar \ledsidenote{\passref{6.18.45}}\edtext{\skt{तौ तु}}{\Afootnote{K H; \skt{तौ}\ A}}\skt{ }\edtext{\skt{त्रिधा}}{\Afootnote{K H; \skt{त्रिधै}\ A}}\skt{ }\edtext{\skt{प्रयुञ्ज्येत}}{\Afootnote{K; \skt{वोपयुज्येते}\ A; \skt{प्रयुञ्जीत}\ H}}\skt{ }\edtext{\skt{रोगेषु}}{\Afootnote{A K; \skt{रोगे}\ H}}\skt{ पुटपाकवत् । }\pend
\pstart\versepar \skt{लेखने सप्त }\edtext{\skt{वाष्टौ वा}}{\Afootnote{K H; \skt{चाष्टौ वा}\ A}}\skt{ बिन्दवः स्नैहिके दश ।}\pend

\pstart\versepar \ledsidenote{\passref{6.18.46}}\edtext{\skt{आश्च्योतने}}{\Afootnote{K H; \skt{आश्च्योतने}\ A}}\skt{ प्रयोक्तव्या द्वादशैव तु रोपणे । }\pend
\pstart\versepar \skt{सेकस्य }\edtext{\skt{कालो}}{\Afootnote{K H; \emph{om.}\ A}}\skt{ द्विगुणः }\edtext{\skt{पुटपाकात्परो}}{\Afootnote{K H; \skt{कालः पुटपाकात्परो}\ A}}\skt{ मतः ।}\pend

\pstart\versepar \ledsidenote{\passref{6.18.47}}\skt{अथ वा }\edtext{\edtext{\skt{कार्यनिर्वृत्तेरपयोगो}}{\lemma{\skt{कार्यनिर्वृत्तेर्}}\Afootnote{A K; \skt{कार्यनिर्व्वृत्ते}\ H}}}{\lemma{\skt{अपयोगो}}\Afootnote{K; \skt{उपयोगो}\ A; \skt{रूपयोगो}\ H}}\skt{ यथाक्रमं । }\pend
\pstart\versepar \skt{}\edtext{\skt{पूर्वापराह्णमध्याह्नरुजाकालेषु}}{\Afootnote{K H; \skt{पूर्वापराह्णे मध्याह्ने रुजाकालेषु}\ A}}\skt{ चोभयोः ।}\pend

\pstart\versepar \ledsidenote{\passref{6.18.48}}\edtext{\skt{योगायोगां स्नेहसेके}}{\Afootnote{K; \skt{योगायोगात्स्नेहसेके}\ A; \skt{योगायोगात्स्नेहसेका}\ H}}\skt{ }\edtext{\edtext{\skt{तर्पणोक्ताम्प्रवक्ष्यते}}{\lemma{\skt{तर्पणोक्ताम्}}\Afootnote{K H; \skt{तर्पणोक्ता}\ A}}}{\lemma{\skt{प्रवक्ष्यते}}\Afootnote{K H; \skt{न्प्रचक्षते ।}\ A}}\skt{ । }\pend
\pstart\versepar \skt{}\edtext{\skt{रोगां}}{\Afootnote{A K H\textsuperscript{pc}; \skt{रोगा}\ H\textsuperscript{ac}}}\skt{ }\edtext{\skt{शिरसि}}{\Afootnote{K; \skt{छिरसि}\ A H\textsuperscript{pc}; \skt{न्ञ्छिरसि}\ H\textsuperscript{ac}}}\skt{ }\edtext{\skt{सम्भूतां}}{\Afootnote{K H; \skt{संभूता}\ A}}\skt{ }\edtext{\skt{हित्वातिप्रबलां गुणां}}{\Afootnote{K; \skt{न्हत्वातिप्रबलान्गुणान् ॥}\ A; \skt{हत्वातिप्रबलान्गुणान् ।}\ H}}\skt{ ।}\pend

\pstart\versepar \ledsidenote{\passref{6.18.49}}\skt{करोति }\edtext{\skt{शिरसो}}{\Afootnote{A K; \skt{शिरसा}\ H}}\skt{ }\edtext{\skt{बस्तिर्य}}{\Afootnote{K H; \skt{बस्तिरु}\ A}}\skt{ }\edtext{\skt{उक्तो}}{\Afootnote{K H; \skt{क्ता ये}\ A}}\skt{ }\edtext{\skt{मूर्ध्नि}}{\Afootnote{K H; \skt{मूर्ध}\ A}}\skt{ तैलिकाः । }\pend
\pstart\versepar \skt{}\edtext{\skt{शुद्धदेहस्यापराह्णे}}{\Afootnote{K H; \skt{शुद्धदेहस्य सायाह्ने}\ A}}\skt{ }\edtext{\skt{यथाव्याधिशिरस्य}}{\Afootnote{K H; \skt{यथाव्याध्यशितस्य}\ A}}\skt{ तु ।}\pend

\pstart\versepar \ledsidenote{\passref{6.18.50}}\edtext{\skt{दृढं}}{\Afootnote{K H; \skt{ऋज्वासीनस्य बध्नीयाद्बस्तिकोशं ततो दृढम् ।}\ A}}\skt{ केशान्तमावध्य बस्तिं }\edtext{\skt{दत्वा}}{\Afootnote{K; \skt{दद्या}\ H}}\skt{ }\edtext{\skt{विचक्षणः}}{\Afootnote{K; \skt{द्विचक्षणः ॥}\ H}}\skt{ । }\pend
\pstart\versepar \skt{}\edtext{\skt{यथाव्याधिः}}{\Afootnote{K; \skt{यथाव्याधि}\ A H}}\skt{ }\edtext{\skt{श्रितस्नेहपूर्ण्णं सं}}{\Afootnote{K\textsuperscript{pc}; \skt{शृतस्नेहपूर्णं सं}\ A; \skt{श्रितः स्नेहः पूर्ण्णं सं}\ K\textsuperscript{ac}; \skt{श्रितं स्नेहं पूर्ण्णं स}\ H}}\skt{ }\edtext{\skt{यस्य}}{\Afootnote{K H; \skt{यम्य}\ A}}\skt{ धारयेत् ।}\edtext{}{\lemma{\skt{केशान्तम्… य}}\Afootnote{\emph{om.}\ A}}\pend

\pstart\versepar \ledsidenote{\passref{6.18.51}}\skt{तर्पणोक्तं दशगुणं यथादोषं विधाणवित् । }\pend
\pstart\versepar \skt{व्यक्तरूपेशु }\edtext{\skt{रोगेषु}}{\Afootnote{K H; \skt{दोषेषु}\ A}}\skt{ }\edtext{\skt{शुद्धदेहस्य}}{\Afootnote{K H; \skt{शुद्धकायस्य}\ A}}\skt{ केवले ।}\edtext{}{\lemma{\skt{विधाणवित् । व्यक्तरूपेशु}}\Afootnote{K; \skt{विधानवित्व्यक्तरूपेषु}\ A H}}\pend

\pstart\versepar \ledsidenote{\passref{6.18.52}}\skt{नेत्र एव स्थिते दोषे प्राप्तमञ्जनमाचरेत् । }\pend
\pstart\versepar \skt{लेखनं }\edtext{\edtext{\skt{रोपनम्वापि}}{\lemma{\skt{रोपनम्}}\Afootnote{K; \skt{रोपणं}\ A H}}}{\lemma{\skt{वापि}}\Afootnote{K H; \skt{चापि}\ A}}\skt{ प्रसादनमथापि वा ।}\pend

\pstart\versepar \ledsidenote{\passref{6.18.53}}\skt{तत्र पञ्च रसा }\edtext{\skt{वर्ज्याः}}{\Afootnote{K; \skt{न्व्यस्ता}\ A H}}\skt{ }\edtext{\skt{मधुरैकविवर्जितां}}{\Afootnote{K; \skt{नाद्यैकरसवर्जिता}\ A H}}\skt{ । }\pend
\pstart\versepar \skt{}\edtext{\skt{पञ्चधा}}{\Afootnote{K; \skt{न्पञ्चधा}\ A H}}\skt{ लेखनं }\edtext{\skt{युञ्ज्यद्यथादोषमतन्द्रितः}}{\lemma{\skt{युञ्ज्यद्यथादोषम्}}\Afootnote{K H; \skt{युञ्ज्याद्यथादोषम्}\ A}}\skt{ ।}\pend

\pstart\versepar \ledsidenote{\passref{6.18.54}}\edtext{\skt{नेत्रवर्त्मसिराकोशस्रोतःशृङ्गाटकाश्रितं}}{\Afootnote{A K; \skt{नेत्रवर्त्म्मसिराकोशस्रोतःशृङ्गाटकाश्रितं ॥}\ H}}\skt{ । }\pend
\pstart\versepar \skt{मुखनासाक्षिभिर्दोषमोजसा }\edtext{\skt{स्रावयेत्तु तं}}{\Afootnote{K; \skt{स्रावयेत्तु तत् ॥}\ A; \skt{सावयोत्तु तत् ।}\ H}}\skt{ ।}\pend

\pstart\versepar \ledsidenote{\passref{6.18.55}}\skt{कषायन्तिक्तकं वापि सस्नेहं रोपणं }\edtext{\skt{हितं}}{\Afootnote{K H; \skt{मतम् ।}\ A}}\skt{ । }\pend
\pstart\versepar \skt{}\edtext{\skt{तत्स्नेहशैत्याद्वर्ण्णं}}{\Afootnote{K; \skt{तत्स्नेहशैत्याद्वर्ण्यं}\ A; \skt{तत्स्नेहशैत्याद्वर्ण्णं}\ H}}\skt{ स्याद्दृष्टेश्च बलवर्धनं ।}\pend

\pstart\versepar \ledsidenote{\passref{6.18.56}}\skt{मधुरं }\edtext{\skt{स्नेहसम्पनमञ्जनन्तु}}{\lemma{\skt{स्नेहसम्पनम्}}\Afootnote{K H; \skt{स्नेहसंपन्नम्}\ A}}\skt{ प्रसादनं । }\pend
\pstart\versepar \skt{}\edtext{\edtext{\skt{दृष्टेर्दोषप्रसादार्थं}}{\lemma{\skt{दृष्टेर्}}\Afootnote{K H; \skt{दृष्टि}\ A}}}{\lemma{\skt{दोषप्रसादार्थं}}\Afootnote{A K; \skt{अतः प्रसादार्थं}\ H}}\skt{ }\edtext{\skt{स्नेहनार्थं तु}}{\Afootnote{K H; \skt{स्नेहनार्थं च}\ A}}\skt{ }\edtext{\skt{यद्धितं}}{\Afootnote{K; \skt{तद्धितम् ॥}\ A; \skt{तद्धितं ।}\ H}}\skt{ ।}\pend

\pstart\versepar \ledsidenote{\passref{6.18.57}}\skt{यथादोषं प्रयोज्यानि तानि रोगविशारदैः । }\pend
\pstart\versepar \skt{}\edtext{\skt{पूर्वापराह्णक्षणदा वाञ्जने }\edtext{\skt{काल}}{\Afootnote{K; \skt{कालम्}\ H}}\skt{ इष्यते ।}}{\lemma{\skt{पूर्वापराह्णक्षणदा… इष्यते}}\Afootnote{\emph{om.}\ A}}\pend

\pstart\versepar \ledsidenote{\passref{6.18.58}}\edtext{\skt{त्रिधैवाञ्जनसंयोगो }\edtext{\skt{गुडिकारसचूर्ण्णतः}}{\Afootnote{K H; \skt{गुटिकारसचूर्णानि त्रिविधान्यञ्जनानि तु ।}\ A}}}{\lemma{\skt{… गु}}\Afootnote{\emph{om.}\ A}}\skt{ । }\pend
\pstart\versepar \skt{यथापूर्वं बलं }\edtext{\skt{तेषु}}{\Afootnote{K H; \skt{तेषां}\ A}}\skt{ }\edtext{\skt{श्रेष्ठमाहुर्मनीषिणः}}{\lemma{\skt{मनीषिणः}}\Afootnote{A K\textsuperscript{pc} H; \skt{मनीषिणि ः ।}\ K\textsuperscript{ac}}}\skt{ ।}\pend

\pstart\versepar \ledsidenote{\passref{6.18.59}}\skt{हरेणुमात्रा वर्त्तिः स्याल्लेखनस्य प्रमाणतः । }\pend
\pstart\versepar \skt{}\edtext{\skt{प्रसादनस्य}}{\Afootnote{A K H\textsuperscript{pc}; \skt{प्रसादन}\ H\textsuperscript{ac}}}\skt{ }\edtext{\skt{चार्ध तु}}{\Afootnote{K H; \skt{चाध्यर्धा}\ A}}\skt{ }\edtext{\skt{द्विगुणा}}{\Afootnote{A K; \skt{द्विगुणे}\ H}}\skt{ रोपणस्य }\edtext{\skt{तु}}{\Afootnote{K H; \skt{च ॥}\ A}}\skt{ ।}\pend

\pstart\versepar \ledsidenote{\passref{6.18.60}}\skt{रसाञ्जनस्य मात्रा तु }\edtext{\skt{पिष्ट्वा}}{\Afootnote{K H; \skt{यथा}\ A}}\skt{ वर्त्तिमिता मता । }\pend
\pstart\versepar \skt{}\edtext{\skt{द्विस्त्रिश्चतुः}}{\lemma{\skt{चतुः}}\Afootnote{A K; \skt{श्चतुः}\ H}}\skt{ शलाकाश्च }\edtext{\edtext{\skt{चूर्ण्णस्याप्यनुपूर्वशः}}{\lemma{\skt{चूर्ण्णस्याप्य्}}\Afootnote{K H; \skt{चूर्णस्या}\ A}}}{\lemma{\skt{अनुपूर्वशः}}\Afootnote{K H; \skt{प्यनुपूर्वशः ॥}\ A}}\skt{ ।}\edtext{}{\lemma{\skt{मता । द्विस्त्रिश्}}\Afootnote{K; \skt{मता द्वित्रि}\ A H}}\pend

\pstart\versepar \ledsidenote{\passref{6.18.61}}\edtext{\edtext{\skt{तेषान्तुल्यगुणान्येव}}{\lemma{\skt{तुल्यगुणान्य्}}\Afootnote{K H; \skt{तुल्यगुणा}\ A}}}{\lemma{\skt{एव}}\Afootnote{K H; \skt{न्येव}\ A}}\skt{ }\edtext{\skt{विदध्याद्भाजनानपि}}{\Afootnote{K; \skt{विदध्याद्भाजनान्यपि ।}\ A; \skt{विदध्याद्भाजनानि वा ॥}\ H}}\skt{ । }\pend
\pstart\versepar \skt{सौवर्ण्णं }\edtext{\skt{रजतं}}{\Afootnote{K H; \skt{राजतं}\ A}}\skt{ }\edtext{\skt{शार्ङ्गं}}{\Afootnote{A K; \emph{om.}\ H}}\skt{ }\edtext{\edtext{\edtext{\skt{वैडूर्यङ्कांसमेव}}{\lemma{\skt{वैडूर्यङ्}}\Afootnote{K; \skt{ताम्रं वैदूर्य}\ A; \skt{सार्ग्गं वैडूर्यङ्}\ H}}}{\lemma{\skt{कांसम्}}\Afootnote{K H; \skt{कांस्य}\ A}}}{\lemma{\skt{एव}}\Afootnote{K H; \skt{जम् ॥}\ A}}\skt{ }\edtext{\skt{च}}{\Afootnote{K H; \emph{om.}\ A}}\skt{ ।}\pend

\pstart\versepar \ledsidenote{\passref{6.18.62}}\edtext{\skt{आयसञ्च}}{\Afootnote{K H; \skt{आयसानि}\ A}}\skt{ }\edtext{\skt{विधेयानि}}{\Afootnote{K H; \skt{च योज्यानि}\ A}}\skt{ शलाकाश्च यथाक्रमात् । }\pend
\pstart\versepar \skt{वक्त्रयोर्मुकुलाकारा }\edtext{\skt{कलायपरिमण्डला}}{\Afootnote{A K; \skt{कलायपरिमण्डलं ।}\ H}}\skt{ ।}\edtext{}{\lemma{\skt{यथाक्रमात् । वक्त्रयोर्}}\Afootnote{K H; \skt{यथाक्रमम् । वक्त्रयोर्}\ A}}\pend

\pstart\versepar \ledsidenote{\passref{6.18.63}}\skt{अष्टांगुला }\edtext{\skt{मध्ये}}{\Afootnote{K; \skt{तनुर्मध्ये}\ A; \skt{यता मध्ये}\ H}}\skt{ }\edtext{\skt{कृशा}}{\Afootnote{K H; \emph{om.}\ A}}\skt{ }\edtext{\skt{सुकृता}}{\Afootnote{A K H\textsuperscript{ac}; \skt{सु}\ H\textsuperscript{pc}}}\skt{ साधुनिग्रहा । }\pend
\pstart\versepar \skt{}\edtext{\edtext{\skt{औदुम्बर्यश्मजा}}{\lemma{\skt{औदुम्बर्य्}}\Afootnote{K H; \skt{औदुम्ब}\ A}}}{\lemma{\skt{अश्मजा}}\Afootnote{K; \skt{र्यश्मजा}\ A; \skt{अस्मजा}\ H}}\skt{ वापि }\edtext{\skt{ताररीतिहिता}}{\Afootnote{K; \skt{शारीरी वा हिता}\ A; \skt{तारारीती हिता}\ H}}\skt{ भवेत् ।}\pend

\pstart\versepar \ledsidenote{\passref{6.18.64}}\edtext{\skt{सुखासीनस्य}}{\Afootnote{K H; \skt{वामेनाक्षि विनिर्भुज्य हस्तेन सु}\ A}}\skt{ }\edtext{\skt{निर्भुज्य}}{\Afootnote{K; \emph{om.}\ A; \skt{निङ्गुज्य}\ H}}\skt{ }\edtext{\skt{नेत्रं}}{\Afootnote{K H; \emph{om.}\ A}}\skt{ }\edtext{\skt{चाततपाणिना}}{\Afootnote{K H; \skt{समाहितः ।}\ A}}\skt{ । }\pend
\pstart\versepar \skt{}\edtext{\skt{अदक्षिणं}}{\Afootnote{K H; \skt{शलाकया}\ A}}\skt{ दक्षिणेन }\edtext{\skt{क्षिपेत्कानीनमञ्जनं}}{\lemma{\skt{क्षिपेत्कानीनम्}}\Afootnote{A K; \skt{क्षिपेत्कानीम्}\ H}}\skt{ ।}\pend

\pstart\versepar \ledsidenote{\passref{6.18.65}}\edtext{\skt{अपाङ्गम्वा}}{\Afootnote{K H; \skt{आपाङ्ग्यं वा}\ A}}\skt{ }\edtext{\edtext{\skt{यथादोषमथवापि}}{\lemma{\skt{यथादोषम्}}\Afootnote{K; \skt{यथायोगं}\ A H}}}{\lemma{\skt{अथवापि}}\Afootnote{K H; \skt{कुर्याच्चापि}\ A}}\skt{ }\edtext{\skt{गतागतां}}{\Afootnote{K; \skt{गतागतम्}\ A H}}\skt{ । }\pend
\pstart\versepar \skt{}\edtext{\skt{वर्त्मावले पि}}{\Afootnote{K; \skt{वर्त्मोपलेपि}\ A; \skt{वर्त्म्मावलपि}\ H}}\skt{ }\edtext{\skt{वासां}}{\Afootnote{K; \skt{वा य}\ A; \skt{य}\ H}}\skt{ }\edtext{\edtext{\skt{वाप्यङ्गुल्यैव}}{\lemma{\skt{वाप्य्}}\Afootnote{K; \skt{त्त}\ A; \skt{च्चापि}\ H}}}{\lemma{\skt{अङ्गुल्यैव}}\Afootnote{K; \skt{दङ्गुल्यैव}\ A; \skt{अतुल्यैव}\ H}}\skt{ प्रयोजयेत् ।}\pend

\pstart\versepar \ledsidenote{\passref{6.18.66}}\edtext{\edtext{\skt{अन्त्यान्त्ययोर्वार्नाञ्ज्याद्धि}}{\lemma{\skt{अन्त्यान्त्ययोर्}}\Afootnote{K; \skt{अक्षि नास्यन्तयोर्}\ A; \skt{अड्यान्तयोर्}\ H}}}{\lemma{\skt{वार्नाञ्ज्याद्धि}}\Afootnote{K; \skt{अञ्ज्या}\ A; \skt{व्वानाञ्ज्याद्धि}\ H}}\skt{ }\edtext{\skt{बाधमानो}}{\Afootnote{K H; \skt{द्बाधमानो}\ A}}\skt{ }\edtext{\skt{क्षि वा}}{\Afootnote{K H; \skt{ऽपि वा}\ A}}\skt{ भिषक् । }\pend
\pstart\versepar \skt{न }\edtext{\skt{चानिर्गतदोषे}}{\Afootnote{K; \skt{चानिर्वान्तदोषे}\ A; \skt{भानिर्ग्गतदोषे}\ H}}\skt{ }\edtext{\skt{क्ष्णिधावनं}}{\Afootnote{K H; \skt{ऽक्ष्णि धावनं}\ A}}\skt{ संप्रयोजयेत् ।}\pend

\pstart \ledsidenote{\passref{6.18.67}}\edtext{\skt{दोषप्रतिनिवृत्तस्तु}}{\Afootnote{K H; \skt{दोषः प्रतिनिवृत्तः स}\ A}}\skt{ }\edtext{\edtext{\edtext{\skt{हन्याद्दृष्टेर्बलन्तथा}}{\lemma{\skt{हन्याद्दृष्टेर्}}\Afootnote{K H; \skt{न्हन्याद्दृष्टे}\ A}}}{\lemma{\skt{बलन्}}\Afootnote{K H; \skt{र्बलं}\ A}}}{\lemma{\skt{तथा}}\Afootnote{K H; \skt{तथा । गतदोषमपेताश्रु पश्येद्यत्सम्यगम्भसा ॥}\ A}}\skt{ ।}\pend

\pstart\versepar \ledsidenote{\passref{6.18.68}}\edtext{\edtext{\edtext{\skt{अधावनाद्बहिस्तिष्ठं}}{\lemma{\skt{अधावनाद्बहिस्}}\Afootnote{K; \skt{प्रक्षाल्याक्षि यथादो}\ A; \skt{अबाधमानाद्बहिस्}\ H}}}{\lemma{\skt{तिष्ठं}}\Afootnote{K; \skt{तिष्ठ}\ H}}\skt{ }\edtext{\skt{भूयः}}{\Afootnote{K; \skt{न्भूयः}\ H}}\skt{ }\edtext{\skt{सञ्जनयेद्भयं}}{\Afootnote{K; \skt{षं कार्यं प्रत्यञ्जनं ततः ।}\ A; \skt{सञ्जायते भयं ।}\ H}}}{\lemma{\skt{नाद्… सञ्}}\Afootnote{\emph{om.}\ A}}\skt{ । }\pend
\pstart\versepar \skt{श्रमोदावर्त्तरुदितमद्यक्रोधभयज्वरैः ।}\pend

\pstart\versepar \ledsidenote{\passref{6.18.69}}\edtext{\edtext{\skt{मलवेगविघातोष्णैश्चार्त्तानां}}{\lemma{\skt{मलवेगविघातोष्णैश्}}\Afootnote{K; \skt{वेगाघातशिरोदोषैश्}\ A; \skt{मलवेगाभिघातोष्णैश्}\ H}}}{\lemma{\skt{चार्त्तानां}}\Afootnote{A K H\textsuperscript{ac}; \skt{चार्त्ता}\ H\textsuperscript{pc}}}\skt{ }\edtext{\skt{नाञ्जनं}}{\Afootnote{K; \skt{नेष्यतेऽञ्जनम् ।}\ A; \skt{अ ञ्जनं}\ H\textsuperscript{ac}; \skt{ना नाञ्जनं}\ H\textsuperscript{pc}}}\skt{ }\edtext{\skt{हितं}}{\Afootnote{K H; \emph{om.}\ A}}\skt{ । }\pend
\pstart\versepar \skt{}\edtext{\skt{रागोपदेहौ}}{\Afootnote{K H; \skt{रागरु}\ A}}\skt{ }\edtext{\skt{तिमिरं}}{\Afootnote{K H; \skt{क्तिमिरास्राव}\ A}}\skt{ }\edtext{\skt{शूलं}}{\Afootnote{K H; \skt{शूल}\ A}}\skt{ }\edtext{\edtext{\skt{संरम्भमेव}}{\lemma{\skt{संरम्भम्}}\Afootnote{K H; \skt{संरम्भ}\ A}}\skt{ च}}{\lemma{\skt{एव च}}\Afootnote{K H; \skt{संभवात् ॥}\ A}}\skt{ ।}\pend

\pstart\versepar \ledsidenote{\passref{6.18.70}}\skt{निद्राक्षये क्रियाशक्तिं कालेष्वेषु करोति तत् । }\pend
\pstart\versepar \skt{शिरोरुग्धूपनं }\edtext{\skt{तप्ते}}{\Afootnote{K H; \skt{ते रा}\ A}}\skt{ }\edtext{\skt{रागाधीमन्थसम्भवः}}{\Afootnote{K H; \skt{गस्रावाधीमन्थसंभवम् ॥}\ A}}\skt{ ॥}\edtext{}{\lemma{\skt{कालेष्वेषु करोति तत् । शिरोरुग्धूपनं}}\Afootnote{K\textsuperscript{pc}; \skt{प्रवाते दृग्बलक्षयम् । रजोधूमह}\ A; \skt{शिरोरुग्धूपनं}\ K\textsuperscript{ac}; \skt{प्रतान्ते दिग्धलक्षणं । सिरोरुद्धूमसन्}\ H}}\pend

\pstart\versepar \ledsidenote{\passref{6.18.71}}\edtext{\skt{संरम्भशूलौनस्यान्ते}}{\Afootnote{A K; \skt{संरम्भभूलौनस्यान्ते}\ H}}\skt{ }\edtext{\skt{प्रतिश्यायं}}{\Afootnote{K H; \skt{शिरोरुजि}\ A}}\skt{ }\edtext{\skt{करोति च}}{\Afootnote{K H; \skt{शिरोरुजम् ।}\ A}}\skt{ । }\pend
\pstart\versepar \skt{}\edtext{\skt{शिरस्नाते}}{\Afootnote{K H; \skt{शिरःस्नाते}\ A}}\skt{ }\edtext{\skt{तिशीते च}}{\Afootnote{K H; \skt{ऽतिशीते च}\ A}}\skt{ रवावनुदिते }\edtext{\skt{पि}}{\Afootnote{K H; \skt{ऽपि}\ A}}\skt{ च ।}\pend

\pstart\versepar \ledsidenote{\passref{6.18.72}}\edtext{\skt{दोषस्थैर्यादपार्थ}}{\Afootnote{K; \skt{दोषस्थैर्यादपार्थं}\ A; \skt{दोषं स्थैर्यादपार्थं}\ H}}\skt{ स्याद्दोषोत्क्लेशं }\edtext{\skt{करोति वा}}{\Afootnote{K H; \skt{करोति च ।}\ A}}\skt{ । }\pend
\pstart\versepar \skt{अजीर्ण्णे }\edtext{\edtext{\skt{प्येवमेतत्स्यात्स्रोतोमार्गनिरोधनात्}}{\lemma{\skt{प्य्}}\Afootnote{K H; \skt{ऽप्य्}\ A}}}{\lemma{\skt{एतत्स्यात्स्रोतोमार्गनिरोधनात्}}\Afootnote{K H; \skt{एव स्यात्स्रोतोमार्गावरोधनात् ॥}\ A}}\skt{ ।}\pend

\pstart\versepar \ledsidenote{\passref{6.18.73}}\edtext{\skt{दोषवेगोदये}}{\Afootnote{A K; \skt{दोषवेगोदयो}\ H}}\skt{ }\edtext{\skt{दन्तं}}{\Afootnote{K H; \skt{दत्तं}\ A}}\skt{ }\edtext{\skt{कुर्यात्तांस्तानुपद्रवान्}}{\lemma{\skt{तानुपद्रवान्}}\Afootnote{K; \skt{स्तानुपद्रवान् ।}\ A; \skt{तानुद्रवान् ।}\ H}}\skt{ । }\pend
\pstart\versepar \skt{तस्मात्परिहरन्दोषानञ्जनं साधु योजयेत् ।}\pend

\pstart \ledsidenote{\passref{6.18.74}}\edtext{\skt{व्यापदश्च}}{\Afootnote{K H; \skt{लेखनस्य विशेषेण काल एष प्रकीर्तितः । व्यापदश्च}\ A}}\skt{ जयेदेताः }\edtext{\skt{सेकाश्च्योतनलेखनैः}}{\lemma{\skt{च्योतनलेखनैः}}\Afootnote{K H; \skt{च्योतनलेपनैः ॥}\ A}}\skt{ ।}\pend

\pstart\versepar \ledsidenote{\passref{6.18.75}}\edtext{\edtext{\skt{यथास्वन्धूमकवडैर्नस्यैश्च्यापि}}{\lemma{\skt{धूमकवडैर्}}\Afootnote{K H; \skt{धूमकवलैर्}\ A}}}{\lemma{\skt{च्यापि}}\Afootnote{K; \skt{चापि}\ A H}}\skt{ }\edtext{\skt{समुच्छ्रिताः}}{\Afootnote{K; \skt{समुत्थिताः ।}\ A; \skt{समुच्छ्रितां ॥}\ H}}\skt{ । }\pend
\pstart\versepar \skt{}\edtext{\skt{विषदं}}{\Afootnote{K; \skt{विशदं}\ A H}}\skt{ }\edtext{\skt{लघ्वनास्रावि}}{\Afootnote{A K; \skt{लघ्वनाम्वावि}\ H}}\skt{ क्रियापटु सुनिर्मलं ।}\pend

\pstart\versepar \ledsidenote{\passref{6.18.76}}\edtext{\edtext{\skt{प्रशान्तदोषन्नयनं}}{\lemma{\skt{प्रशान्तदोषन्}}\Afootnote{K; \skt{संशान्तोपद्र}\ A; \skt{संशान्तदोषं}\ H}}}{\lemma{\skt{नयनं}}\Afootnote{K H; \skt{वं नेत्रं}\ A}}\skt{ विरिक्तं सम्यगादिशेत् । }\pend
\pstart\versepar \skt{जिह्मं दारुणदुर्वर्णं }\edtext{\skt{स्रस्तरूक्षमतीव}}{\lemma{\skt{स्रस्तरूक्षम्}}\Afootnote{K H; \skt{स्रस्तं रूक्षम्}\ A}}\skt{ च ।}\pend

\pstart\versepar \ledsidenote{\passref{6.18.77}}\edtext{\skt{नेत्रे}}{\Afootnote{K; \skt{नेत्रं}\ A; \skt{नेत्र}\ H}}\skt{ }\edtext{\skt{विरेकातियोगोत्स्यन्दने}}{\Afootnote{K; \skt{विरेकातियोगे स्यन्दते}\ A; \skt{विकारे तियोगोत्स्यन्दने}\ H}}\skt{ }\edtext{\skt{चातिमात्रशः}}{\Afootnote{A K; \skt{चातिमात्रसः ॥}\ H}}\skt{ । }\pend
\pstart\versepar \skt{तत्र सन्तर्पणं कार्यं }\edtext{\skt{विधाणं}}{\Afootnote{K; \skt{विधानं}\ A H}}\skt{ }\edtext{\skt{वातनाशनं}}{\Afootnote{K H; \skt{चानिलापहम् ॥}\ A}}\skt{ ।}\pend

\pstart\versepar \ledsidenote{\passref{6.18.78}}\edtext{\skt{विरिक्तं}}{\Afootnote{K H; \skt{अक्षि मन्दविरिक्तं}\ A}}\skt{ }\edtext{\edtext{\skt{हीनमक्षि}}{\lemma{\skt{हीनम्}}\Afootnote{K H; \emph{om.}\ A}}}{\lemma{\skt{अक्षि}}\Afootnote{K H; \emph{om.}\ A}}\skt{ स्यादुदग्रतरदोषवत् । }\pend
\pstart\versepar \skt{कण्डूपदेहौ पैच्छिल्यं }\edtext{\skt{संरम्भश्चार्तिमेव}}{\lemma{\skt{चार्तिम्}}\Afootnote{K; \skt{चात्रम्}\ H}}\skt{ च । }\pend
\pstart\versepar \skt{}\edtext{\skt{धूमनस्याञ्जनैस्तत्र}}{\lemma{\skt{धूमनस्याञ्जनैस्}}\Afootnote{K H; \skt{त् । धूमनस्याञ्जनैस्}\ A}}\skt{ हितन्दोषावसेचनं ।}\edtext{}{\lemma{\skt{स्यादुदग्रतरदोषवत् । कण्डूपदेहौ}}\Afootnote{K H; \skt{स्यादुदग्रतरदोष}\ A}}\edtext{}{\lemma{\skt{वत्… एव}}\Afootnote{\emph{om.}\ A}}\pend

\pstart\versepar \ledsidenote{\passref{6.18.79}}\skt{स्नेहवर्ण्णबलोपेतं प्रसन्नन्दोषवर्जितं । }\pend
\pstart\versepar \skt{ज्ञेयं }\edtext{\skt{प्रसादनं}}{\Afootnote{K H; \skt{प्रसादने}\ A}}\skt{ सम्यगुपयुक्ते }\edtext{\skt{क्षि}}{\Afootnote{K; \skt{ऽक्षि}\ A H}}\skt{ }\edtext{\skt{निर्वृतं}}{\Afootnote{A K; \skt{निर्व्वृतिं ॥}\ H}}\skt{ ।}\pend

\pstart\versepar \ledsidenote{\passref{6.18.80}}\edtext{\skt{किञ्चिद्धीनविकारः}}{\Afootnote{K H; \skt{किञ्चिद्धीनविकारं}\ A}}\skt{ }\edtext{\skt{स्यात्तर्पणाद्विकृतादपि}}{\Afootnote{K H; \skt{स्यात्तर्पणाद्धि कृतादति ।}\ A}}\skt{ । }\pend
\pstart\versepar \skt{तत्र }\edtext{\skt{दोषहरं}}{\Afootnote{A K\textsuperscript{pc} H; \skt{दोषआहरं}\ K\textsuperscript{ac}}}\skt{ रूक्षं }\edtext{\skt{विधानं}}{\Afootnote{K H; \skt{भेषजं}\ A}}\skt{ शस्यते }\edtext{\skt{मृदु}}{\Afootnote{A K\textsuperscript{pc}; \skt{मृदुः}\ K\textsuperscript{ac} H}}\skt{ ।}\pend

\pstart \ledsidenote{\passref{6.18.81}}\skt{साधारणमपि ज्ञेयमेवं }\edtext{\skt{रोपणलक्षणं}}{\Afootnote{K H; \skt{रोपणलक्षणम् । प्रसादनवदाचष्टे तस्मिन्युक्तेऽतिभेषजम् ॥}\ A}}\skt{ ।}\pend

\pstart\versepar \ledsidenote{\passref{6.18.82}}\skt{स्नेहनं रोपणं वापि }\edtext{\skt{हीनयोगे}}{\Afootnote{K H; \skt{हीनयुक्त}\ A}}\skt{ }\edtext{\skt{त्वपार्थकं}}{\Afootnote{K H; \skt{मपार्थकम् ।}\ A}}\skt{ । }\pend
\pstart\versepar \skt{}\edtext{\skt{तस्मान्मात्राविदा}}{\Afootnote{K; \skt{कर्तव्यं मात्रया तस्मा}\ A; \skt{तस्मात्मात्राविदा}\ H}}\skt{ }\edtext{\edtext{\skt{कार्यमञ्जनं}}{\lemma{\skt{कार्यम्}}\Afootnote{K H; \emph{om.}\ A}}}{\lemma{\skt{अञ्जनं}}\Afootnote{K H; \skt{दञ्जनं}\ A}}\skt{ सिद्धिमिच्छता । }\pend
\pstart\versepar \skt{}\edtext{\skt{धूमैः}}{\Afootnote{K; \emph{om.}\ A; \skt{धूमै}\ H}}\skt{ }\edtext{\edtext{\skt{सनस्यैरथ}}{\lemma{\skt{सनस्यैर्}}\Afootnote{K H; \emph{om.}\ A}}\skt{ वा}}{\lemma{\skt{अथ वा}}\Afootnote{K H; \emph{om.}\ A}}\skt{ }\edtext{\edtext{\skt{जयेच्छ्लेष्माणमीरितं}}{\lemma{\skt{जयेच्छ्लेष्माणम्}}\Afootnote{K H; \emph{om.}\ A}}}{\lemma{\skt{ईरितं}}\Afootnote{K; \emph{om.}\ A; \skt{ईरितां ॥}\ H}}\skt{ ।}\pend

\pstart\versepar \ledsidenote{\passref{6.18.83}}\skt{पुटपाकक्रियाद्यासु }\edtext{\skt{क्रियास्वेषैव}}{\lemma{\skt{एषैव}}\Afootnote{K H; \skt{एषैव(क्रियास्वेकैव)}\ A}}\skt{ कल्पना । }\pend
\pstart\versepar \skt{सहस्रशो }\edtext{\skt{पि}}{\Afootnote{K; \skt{चाञ्ज}\ A H}}\skt{ }\edtext{\skt{कुर्वीत}}{\Afootnote{K; \skt{नेषु}\ A; \skt{ने च}\ H}}\skt{ }\edtext{\skt{बीजेनानेन}}{\Afootnote{K H; \skt{बीजेनोक्तेन}\ A}}\skt{ }\edtext{\skt{पूजितः}}{\Afootnote{K; \skt{पूजिताः ॥}\ A; \skt{पूजिताः ॥}\ H}}\skt{ ॥ }\pend
\pstart\versepar \skt{}\edtext{\edtext{\skt{अञ्जनन्नस्यकवलधूमपानमनुक्रमात्}}{\lemma{\skt{अञ्जनन्}}\Afootnote{K H; \emph{om.}\ A}}}{\lemma{\skt{नस्यकवलधूमपानम्}}\Afootnote{K H; \emph{om.}\ A}}\skt{ । }\pend
\pstart\versepar \skt{}\edtext{\skt{युञ्ज्याद्यथार्थमेतानि}}{\lemma{\skt{एतानि}}\Afootnote{K H; \emph{om.}\ A}}\skt{ }\edtext{\skt{पृथक्पृथगथापि}}{\Afootnote{K; \emph{om.}\ A; \skt{पृथगथापि}\ H}}\skt{ }\edtext{\skt{वा}}{\Afootnote{K H; \emph{om.}\ A}}\skt{ ।}\edtext{}{\lemma{\skt{अनुक्रमात् । युञ्ज्याद्यथार्थम्}}\Afootnote{K; \emph{om.}\ A; \skt{अनुक्रमात् । युज्याद्यथार्थम्}\ H}}\pend

\pstart\versepar \ledsidenote{\passref{6.18.84}}\skt{दृष्टेर्बलविवृध्यर्थं }\edtext{\skt{पाप्मरोगक्षयाय च}}{\Afootnote{K H; \skt{याप्यरोगक्षयाय च ।}\ A}}\skt{ । }\pend
\pstart\versepar \skt{}\edtext{\edtext{\skt{नृपार्हाण्यञ्जनाग्राणि}}{\lemma{\skt{नृपार्हाण्य्}}\Afootnote{K; \skt{राजार्हा}\ A; \skt{नृपार}\ H}}}{\lemma{\skt{अञ्जनाग्राणि}}\Afootnote{K; \skt{न्यञ्जनाग्र्याणि}\ A; \skt{हाञ्जनाग्राणि}\ H}}\skt{ निबोधेमान्यतः परं ।}\pend

\pstart\versepar \ledsidenote{\passref{6.18.85}}\edtext{\edtext{\skt{भागानष्टावञ्जन-स्य}}{\lemma{\skt{भागानष्टाव्}}\Afootnote{K H; \skt{अष्टौ भागा}\ A}}}{\lemma{\skt{अञ्जन-स्य}}\Afootnote{K H; \skt{नञ्जनस्य}\ A}}\skt{ }\edtext{\skt{नीलोत्पलसुगन्धिनः}}{\Afootnote{K H; \skt{नीलोत्पलसमत्विषः ।}\ A}}\skt{ । }\pend
\pstart\versepar \skt{}\edtext{\skt{ताम्रं}}{\Afootnote{K H; \skt{औडुम्बरं}\ A}}\skt{ }\edtext{\skt{सुवर्ण्णञ्च}}{\Afootnote{K H; \skt{शातकुम्भं}\ A}}\skt{ }\edtext{\skt{भवेद्रजतञ्च}}{\Afootnote{K H; \skt{राजतं च}\ A}}\skt{ }\edtext{\skt{समांशतः}}{\Afootnote{K; \skt{समासतः ॥}\ A; \skt{समांसतः ॥}\ H}}\skt{ ।}\pend

\pstart\versepar \ledsidenote{\passref{6.18.86}}\edtext{\skt{एकादशैतानङ्गांस्तु}}{\Afootnote{K; \skt{एकादशैतान्भागांस्तु}\ A; \skt{एकादशैतानंशास्तु}\ H}}\skt{ योजयेत्कुशलो भिषक् । }\pend
\pstart\versepar \skt{मूषाक्षिप्तन्तु }\edtext{\edtext{\skt{तद्धान्तमावृतं}}{\lemma{\skt{तद्धान्तम्}}\Afootnote{K; \skt{दाध्मात}\ A; \skt{तद्धा}\ H}}}{\lemma{\skt{आवृतं}}\Afootnote{K; \skt{मावृतं}\ A; \skt{म्या वृन्तं}\ H}}\skt{ जातवेदसि ।}\edtext{}{\lemma{\skt{भिषक् । मूषाक्षिप्तन्तु}}\Afootnote{H; \skt{भिषक् । मूषाक्षिप्तं त}\ A; \skt{भिषक् । मूषो क्षिप्तन्तु}\ K\textsuperscript{ac}; \skt{भिषक् । मूषाक्षिप्तन्तु}\ K\textsuperscript{pc}}}\pend

\pstart\versepar \ledsidenote{\passref{6.18.87}}\edtext{\skt{खदिरं}}{\Afootnote{K H; \skt{खदिरा}\ A}}\skt{ }\edtext{\skt{स्यन्दनाङ्गारैः}}{\Afootnote{K H; \skt{श्मन्तकाङ्गरै}\ A}}\skt{ }\edtext{\skt{गोशकृद्भिरथापि}}{\lemma{\skt{गोशकृद्भिर्}}\Afootnote{K H; \skt{र्गोशकृद्भिर्}\ A}}\skt{ वा । }\pend
\pstart\versepar \skt{गवां शकृद्रसे मूत्रे }\edtext{\skt{घृते}}{\Afootnote{K H; \emph{om.}\ A}}\skt{ }\edtext{\skt{दध्न्यथ}}{\lemma{\skt{अथ}}\Afootnote{K H; \skt{सर्पिषि}\ A}}\skt{ माक्षिके}\pend

\pstart\versepar \ledsidenote{\passref{6.18.88}}\edtext{\edtext{\edtext{\skt{तैलम्मद्यवसामज्जासर्वगन्धोदकेष्वपि}}{\lemma{\skt{तैलम्}}\Afootnote{K H; \skt{तैल}\ A}}}{\lemma{\skt{मद्यवसामज्जासर्वगन्धोदकेष्व्}}\Afootnote{K H; \skt{मद्यवसामज्जसर्वगन्धोदकेषु}\ A}}}{\lemma{\skt{अपि}}\Afootnote{K H; \skt{च ।}\ A}}\skt{ । }\pend
\pstart\versepar \skt{}\edtext{\skt{द्राक्षासितेक्षुत्रिफलारसेषु}}{\Afootnote{K H; \skt{द्राक्षारसेक्षुत्रिफलारसेषु}\ A}}\skt{ }\edtext{\skt{सुहिमेषु}}{\Afootnote{A K; \skt{विहिमेषु}\ H}}\skt{ च ।}\pend

\pstart\versepar \ledsidenote{\passref{6.18.89}}\edtext{\skt{सारिवादिकषायेषु}}{\Afootnote{K H; \skt{सारिवादिकषाये च कषाये}\ A}}\skt{ }\edtext{\skt{उत्पलादौ}}{\Afootnote{K; \skt{चोत्पलादि}\ A; \skt{उत्पलादेस्}\ H}}\skt{ }\edtext{\skt{तथैव च}}{\Afootnote{K H; \skt{के ।}\ A}}\skt{ । }\pend
\pstart\versepar \skt{}\edtext{\skt{निषेचयेत्पृथक्वैनं}}{\Afootnote{K; \skt{निषेचयेत्पृथक्चैनं}\ A; \skt{निषेचयेत्पुनस्त्वेनं}\ H}}\skt{ }\edtext{\skt{धान्तं}}{\Afootnote{K H; \skt{ध्मातं}\ A}}\skt{ }\edtext{\skt{धान्तं}}{\Afootnote{K H; \skt{ध्मातं}\ A}}\skt{ पुनः पुनः ।}\pend

\pstart\versepar \ledsidenote{\passref{6.18.90}}\skt{ततो }\edtext{\skt{न्तरिक्षे}}{\Afootnote{K H; \skt{ऽन्तरीक्षे}\ A}}\skt{ सप्ताहं }\edtext{\skt{प्लोताबद्धं}}{\Afootnote{K H; \skt{प्लोतबद्धं}\ A}}\skt{ स्थितं जले । }\pend
\pstart\versepar \skt{विशोष्य }\edtext{\skt{चूर्ण्णयेन्मुक्तां}}{\Afootnote{K; \skt{चूर्णयेन्मुक्तां}\ A; \skt{चूर्ण्णयेत्मुक्ता}\ H}}\skt{ }\edtext{\skt{स्फाटिकं}}{\Afootnote{K H; \skt{स्फटिकं}\ A}}\skt{ विद्रुमन्तथा ।}\pend

\pstart\versepar \ledsidenote{\passref{6.18.91}}\edtext{\skt{कालानुशारिवाञ्चापि}}{\lemma{\skt{कालानुशारिवाञ्}}\Afootnote{K H; \skt{कालानुसारिवां}\ A}}\skt{ }\edtext{\edtext{\skt{शुचिरावाप्य}}{\lemma{\skt{शुचिर्}}\Afootnote{K H; \skt{शुचि}\ A}}}{\lemma{\skt{आवाप्य}}\Afootnote{K H; \skt{रावाप्य}\ A}}\skt{ योगतः । }\pend
\pstart\versepar \skt{एतच्चूर्ण्णाञ्जनं श्रेष्ठं निहितं भाजने शुभे ।}\pend

\pstart\versepar \ledsidenote{\passref{6.18.92}}\edtext{\skt{वैडूर्यदन्तस्फटिकशङ्खशैलायसे}}{\Afootnote{K H; \skt{दन्तस्फटिकवैदूर्यशङ्खशैलासनो}\ A}}\skt{ }\edtext{\skt{नवे}}{\Afootnote{K H; \skt{द्भवे ।}\ A}}\skt{ । }\pend
\pstart\versepar \skt{}\edtext{\skt{शातकोम्भे}}{\Afootnote{K H; \skt{शातकुम्भे}\ A}}\skt{ }\edtext{\skt{थ}}{\Afootnote{K H; \skt{ऽथ}\ A}}\skt{ }\edtext{\skt{शार्ङ्गो वा}}{\Afootnote{K; \skt{शार्ङ्गे वा}\ A; \skt{शारङ्गो}\ H}}\skt{ राजते वा }\edtext{\skt{सुसंवृते}}{\Afootnote{K H; \skt{सुसंस्कृते ।}\ A}}\skt{ । }\pend
\pstart\versepar \skt{}\edtext{\skt{सहस्रपाकवत्पूज्यभूमिपाय}}{\Afootnote{K; \skt{सहस्रपाकवत्पूजां कृत्वा राज्ञः}\ A; \skt{सहस्रपावत्पूज्यभूमिपाय}\ H}}\skt{ }\edtext{\skt{प्रयोजयेत्}}{\Afootnote{A K; \skt{प्रदापयेत् ।}\ H}}\skt{ ।}\pend

\pstart\versepar \ledsidenote{\passref{6.18.93}}\skt{तेनाञ्जिताक्षो नृपतिर्भवेत्सर्वजनप्रियः । }\pend
\pstart\versepar \skt{अधृष्यः सर्वभूतानां दृष्टिरोगविवर्जितः ।}\pend

\pstart\versepar \ledsidenote{\passref{6.18.94}}\edtext{\edtext{\skt{कुष्ठश्चन्दनमेला}}{\lemma{\skt{कुष्ठश्}}\Afootnote{K; \skt{कुष्ठं}\ A; \skt{कुष्ठ}\ H}}\skt{ च}}{\lemma{\skt{चन्दनमेला च}}\Afootnote{K H; \skt{चन्दनमेलाश्च}\ A}}\skt{ पत्रं मधुकमञ्जनं । }\pend
\pstart\versepar \skt{}\edtext{\skt{मेषशृंग्याश्च}}{\Afootnote{K H; \skt{मेषशृङ्गस्य}\ A}}\skt{ पुष्पाणि वक्रं रत्नानि सप्त च ।}\pend

\pstart\versepar \ledsidenote{\passref{6.18.95}}\skt{उत्पलस्य बृहत्योश्च पद्मस्यापि च }\edtext{\skt{केसरं}}{\Afootnote{K H; \skt{केशरम् ।}\ A}}\skt{ । }\pend
\pstart\versepar \skt{}\edtext{\edtext{\skt{नागपुष्पमुशीराणि}}{\lemma{\skt{नागपुष्पम्}}\Afootnote{K H; \skt{नागपुष्प}\ A}}}{\lemma{\skt{उशीराणि}}\Afootnote{K H; \skt{मुशीराणि}\ A}}\skt{ }\edtext{\edtext{\skt{पिप्पल्यस्तुत्थमेव}}{\lemma{\skt{पिप्पल्यस्}}\Afootnote{K H; \skt{पिप्पली}\ A}}}{\lemma{\skt{एव}}\Afootnote{K H; \skt{उत्त}\ A}}\skt{ }\edtext{\skt{च}}{\Afootnote{K H; \skt{मम् ॥}\ A}}\skt{ ।}\pend

\pstart\versepar \ledsidenote{\passref{6.18.96}}\skt{कुक्कुटाण्डकपालानि }\edtext{\edtext{\skt{पथ्यान्दार्वीं}}{\lemma{\skt{पथ्यान्}}\Afootnote{K H; \skt{दार्वीं}\ A}}}{\lemma{\skt{दार्वीं}}\Afootnote{K H; \skt{पथ्यां}\ A}}\skt{ }\edtext{\skt{सरोचनं}}{\Afootnote{K H; \skt{सरोचनाम् ।}\ A}}\skt{ । }\pend
\pstart\versepar \skt{}\edtext{\skt{मरिचान्यक्षमज्जा च}}{\Afootnote{K; \skt{मरिचान्यक्षमज्जानं}\ A; \skt{मरिचान्यक्षरमज्जा च}\ H}}\skt{ }\edtext{\skt{तुल्या चः}}{\Afootnote{K; \skt{तुल्यां च}\ A; \skt{तुल्या च}\ H}}\skt{ }\edtext{\skt{गृहगोधिका}}{\Afootnote{K H; \skt{गृहगोपिकाम् ॥}\ A}}\skt{ ।}\pend

\pstart\versepar \ledsidenote{\passref{6.18.97}}\skt{कृत्वा सूक्ष्मं ततश्चूर्ण्णं न्यसेदभ्यर्च्य पूर्ववत् । }\pend
\pstart\versepar \skt{एतद्भद्रोदयं नाम }\edtext{\skt{सदा}}{\Afootnote{K H; \skt{सदै}\ A}}\skt{ }\edtext{\skt{वार्हति}}{\Afootnote{A K; \skt{र्द्दति}\ H}}\skt{ }\edtext{\skt{भूपतिः}}{\Afootnote{K H; \skt{भूमिपः ॥}\ A}}\skt{ ।}\pend

\pstart\versepar \ledsidenote{\passref{6.18.98}}\edtext{\skt{वक्रञ्च}}{\Afootnote{K H; \skt{वक्रं स}\ A}}\skt{ मरिचञ्चैव }\edtext{\skt{मांसी}}{\Afootnote{K H; \skt{मांसीं}\ A}}\skt{ शैलेयमेव च । }\pend
\pstart\versepar \skt{तुल्यांशानि }\edtext{\skt{समग्रैस्तैः}}{\Afootnote{K H; \skt{समानैस्तैः}\ A}}\skt{ }\edtext{\skt{समा}}{\Afootnote{K H; \skt{सम}\ A}}\skt{ }\edtext{\skt{चात्र}}{\Afootnote{K H; \skt{ग्रैश्च}\ A}}\skt{ मनःशिला ।}\pend

\pstart\versepar \ledsidenote{\passref{6.18.99}}\skt{पत्रस्य भागाश्चत्वारो द्विगुणं }\edtext{\skt{सर्वतो}}{\Afootnote{A K; \skt{सर्व्व}\ H}}\skt{ }\edtext{\skt{ञ्जनं}}{\Afootnote{A K; \skt{त्वंजनं ॥}\ H}}\skt{ । }\pend
\pstart\versepar \skt{तावच्च यष्टीमधुकं }\edtext{\skt{पूर्ववत्त्वेतदञ्जनं}}{\Afootnote{K H; \skt{पूर्ववच्चैतदञ्जनम् ॥}\ A}}\skt{ ॥}\pend

\pstart\versepar \ledsidenote{\passref{6.18.100}}\edtext{\skt{मनःशिला}}{\Afootnote{K; \skt{मनःशिलां}\ A H}}\skt{ देवकाष्ठं }\edtext{\skt{रजने}}{\Afootnote{K H; \skt{रजन्यौ}\ A}}\skt{ त्रिफलोषणं । }\pend
\pstart\versepar \skt{}\edtext{\skt{द्राक्षामधुकमञ्जिष्ठाः}}{\Afootnote{K; \skt{लाक्षालशुनमञ्जिष्ठा}\ A; \skt{द्राक्षालशुनमञ्जिष्ठा}\ H}}\skt{ }\edtext{\skt{सैन्धवेलाः}}{\Afootnote{K; \skt{सैन्धवैलाः}\ A H}}\skt{ }\edtext{\skt{सगैरिकाः}}{\Afootnote{K; \skt{समाक्षिकाः ॥}\ A; \skt{सशैरिकाः ।}\ H}}\skt{ ।}\pend

\pstart\versepar \ledsidenote{\passref{6.18.101}}\edtext{\skt{लोध्रं}}{\Afootnote{K H; \skt{रोध्रं}\ A}}\skt{ }\edtext{\skt{शावरकं}}{\Afootnote{K H; \skt{सावरकं}\ A}}\skt{ }\edtext{\edtext{\skt{चूर्ण्णमायसं}}{\lemma{\skt{चूर्ण्णम्}}\Afootnote{K H; \skt{चूर्ण}\ A}}}{\lemma{\skt{आयसं}}\Afootnote{K H; \skt{मायसं}\ A}}\skt{ }\edtext{\edtext{\skt{ताम्रजन्तथा}}{\lemma{\skt{ताम्रजन्}}\Afootnote{K H; \skt{ताम्रमे}\ A}}}{\lemma{\skt{तथा}}\Afootnote{K H; \skt{व च ।}\ A}}\skt{ । }\pend
\pstart\versepar \skt{}\edtext{\edtext{\skt{कालानुशारिवाञ्चापि}}{\lemma{\skt{कालानुशारिवाञ्}}\Afootnote{K H; \skt{कालानुसारिवां}\ A}}}{\lemma{\skt{चापि}}\Afootnote{K H; \skt{चैव}\ A}}\skt{ कुक्कुटाण्डदलानि च ।}\pend

\pstart\versepar \ledsidenote{\passref{6.18.102}}\edtext{\skt{पिष्ट्वा}}{\Afootnote{K H; \emph{om.}\ A}}\skt{ }\edtext{\skt{च्छागेन}}{\Afootnote{K H; \skt{तुल्यानि}\ A}}\skt{ पयसा }\edtext{\skt{गुडिकाङ्कारयेद्भिषक्}}{\lemma{\skt{गुडिकाङ्}}\Afootnote{K H; \skt{पिष्ट्वा गुटिकां}\ A}}\skt{ । }\pend
\pstart\versepar \skt{कण्डूतिमिरशुक्रार्मरक्तराज्युपशान्तये ।}\edtext{}{\lemma{\skt{कारयेद्भिषक् । कण्डूतिमिरशुक्रार्मरक्तराज्युपशान्तये}}\Afootnote{K H; \skt{कारयेद्बुधः । कण्डूतिमिरशुक्लार्मरक्तराज्युपशान्तये ॥}\ A}}\pend

\pstart\versepar \ledsidenote{\passref{6.18.103}}\edtext{\skt{कांसापमार्जनमषी}}{\Afootnote{K; \skt{कांस्यापमार्जनमसी}\ A; \skt{कांसापमार्ज्जनमसी}\ H}}\skt{ मधुकं सैन्धवन्तथा । }\pend
\pstart\versepar \skt{एरण्डमूलञ्च समं }\edtext{\skt{बृहत्योश्च}}{\Afootnote{K H; \skt{बृहत्यंश}\ A}}\skt{ }\edtext{\skt{द्विरंशतः}}{\Afootnote{K H; \skt{द्वयान्वितम् ॥}\ A}}\skt{ ।}\pend

\pstart\versepar \ledsidenote{\passref{6.18.104}}\skt{आजेन पयसा पिष्ट्वा ताम्रपात्रम्प्रलेपयेत् । }\pend
\pstart\versepar \skt{सप्तकृत्वस्तु ता वर्त्त्यः }\edtext{\skt{च्छायाशुष्का}}{\Afootnote{A K; \skt{श्छायाशुष्का}\ H}}\skt{ }\edtext{\skt{रुजापहा}}{\Afootnote{K H; \skt{रुजापहाः ॥}\ A}}\skt{ ।}\edtext{}{\lemma{\skt{प्रलेपयेत् । सप्तकृत्वस्तु ता}}\Afootnote{A K; \skt{प्रलेपयेत् ॥ सप्तकृत्यस्तु ता}\ H}}\pend

\pstart\versepar \ledsidenote{\passref{6.18.106}}\skt{रसक्रियाविधानेन यथोक्तविधिकोविदः । }\pend
\pstart\versepar \skt{}\edtext{\skt{पिण्डांजनानि}}{\Afootnote{A K; \skt{पिण्डाजनानि}\ H}}\skt{ }\edtext{\skt{मतिमां}}{\Afootnote{K; \emph{om.}\ A; \skt{मतिमा}\ H}}\skt{ }\edtext{\skt{विदध्यात्तु}}{\Afootnote{K; \skt{कुर्वीत}\ A; \skt{न्विदध्यात्तु}\ H}}\skt{ }\edtext{\edtext{\skt{यथागदमिति}}{\lemma{\skt{यथागदम्}}\Afootnote{K H; \skt{यथायोगमतन्}\ A}}}{\lemma{\skt{इति}}\Afootnote{K H; \skt{द्रितः ॥}\ A}}\skt{ ॥}\pend

\pstart \ledsidenote{\passref{6.18.trailer}}\edtext{\skt{॥ शालाक्ये १७ ॥ {\symfont ❈}}}{\Afootnote{K; \skt{इति सुश्रुतसंहितायामुत्तरतन्त्रान्तर्गते शालाक्यतन्त्रे क्रियाकल्पो नामाष्टदशोऽध्यायः ॥१८ ॥}\ A; \skt{॥ शालाक्ये क्रियाकल्पः सप्तदशमोऽध्यायः ॥}\ H}}\skt{ ॥}\pend

\endnumbering
\beginnumbering
\eledsection*{Adhyāya 19}
\pstart \ledsidenote{\passref{6.19.1}}\skt{अथातो }\edtext{\edtext{\skt{नयनाभिघातप्रतिषेधमध्यायम्व्याख्यास्यामः}}{\lemma{\skt{अध्यायम्}}\Afootnote{K H; \emph{om.}\ A}}}{\lemma{\skt{व्याख्यास्यामः}}\Afootnote{A H; \emph{om.}\ K}}\skt{ ॥}\pend

\pstart\versepar \ledsidenote{\passref{6.19.3ab}}\skt{अभ्याहते तु नयने बहुधा }\edtext{\skt{णराणां}}{\Afootnote{K; \skt{नराणां}\ A H}}\skt{ }\pend
\pstart\versepar \skt{संरम्भरागतुमुलासु रुजासु }\edtext{\skt{धीमां}}{\Afootnote{K; \skt{धीमान् ।}\ A H}}\skt{ ।}\pend

\pstart\versepar \ledsidenote{\passref{6.19.3cd}}\skt{दृष्टिप्रसादजननैस्तमुपक्रमेत }\pend
\pstart\versepar \skt{}\edtext{\skt{नस्यास्यलेपपरिषेवनतर्पणाद्यं}}{\Afootnote{K; \skt{नस्यास्यलेपपरिषेचनतर्पणाद्यमुक्तं पुनः क्षतजपित्तजशूलपथ्यम् ॥}\ A; \skt{नस्यास्यलेपपरिषेवणतर्प्पणाद्यं ॥}\ H}}\skt{ ॥}\edtext{}{\lemma{\skt{… न}}\Afootnote{\emph{om.}\ A}}\pend

\pstart\versepar \ledsidenote{\passref{6.19.4ab}}\edtext{\skt{उक्तं}}{\Afootnote{K; \skt{उक्ता}\ H}}\skt{ }\edtext{\skt{पुरा}}{\Afootnote{K; \skt{रा}\ H}}\skt{ क्षतजपित्तजशूलपथ्यं }\pend
\pstart\versepar \skt{दृष्टिप्रसादजननं }\edtext{\edtext{\skt{विधिमाशु}}{\lemma{\skt{विधिम्}}\Afootnote{K H; \skt{विधि}\ A}}}{\lemma{\skt{आशु}}\Afootnote{K H; \skt{माशु}\ A}}\skt{ }\edtext{\skt{कुर्यात्}}{\Afootnote{K; \skt{कुर्यात्स्निग्धैर्हिमैश्च मधुरैश्च तथा प्रयोगैः ।}\ A; \skt{हन्यात् ।}\ H}}\skt{ ।}\edtext{}{\lemma{\skt{… दृ}}\Afootnote{\emph{om.}\ A}}\pend

\pstart\versepar \ledsidenote{\passref{6.19.4cd}}\skt{स्निग्धैर्हिमैश्च मधुरैश्च तथा प्रयोगैः }\pend
\pstart\versepar \skt{}\edtext{\skt{स्वेदाग्निधूमभयशोकरुजाभितापैः}}{\Afootnote{K H; \skt{स्वेदाग्निधूमभयशोकरुजाभिघातैरभ्याहतामपि तथैव भिषक्चिकित्सेत् ॥}\ A}}\skt{ ।}\edtext{}{\lemma{\skt{… स्वे}}\Afootnote{\emph{om.}\ A}}\pend

\pstart\versepar \ledsidenote{\passref{6.19.5ab}}\edtext{\edtext{\skt{अभ्याहतामपि}}{\lemma{\skt{अभ्याहताम्}}\Afootnote{K H; \skt{सद्यो}\ A}}}{\lemma{\skt{अपि}}\Afootnote{K H; \skt{हते}\ A}}\skt{ }\edtext{\skt{तथा च}}{\Afootnote{K H; \skt{नयन}\ A}}\skt{ भिषक्चिकित्सेद् }\pend
\pstart\versepar \skt{वैशेषिकश्च }\edtext{\skt{परिषेचनलेपबस्ति}}{\Afootnote{\skt{वति दोषमवेक्ष्य कार्यः ।}\ A; \skt{परिषेचनलेपवस्ति ।}\ K; \skt{परिषेचनतर्प्पणाद्यैः ।}\ H}}\skt{ ।}\edtext{}{\lemma{\skt{भिषक्चिकित्सेद्वैशेषिकश्च}}\Afootnote{K; \skt{एष विधिस्तदूर्ध्वं स्यन्देरितो भ}\ A; \skt{भिषक्चिकित्सेत्मस्तिष्कमूर्द्ध्व}\ H}}\pend

\pstart \ledsidenote{\passref{6.19.5cd}}\edtext{\edtext{\skt{मस्तिष्कमूर्धपरिषेचनतर्पनाद्यैः}}{\lemma{\skt{मस्तिष्कम्}}\Afootnote{K; \skt{अभ्याहतं नयनमीषदथास्यबाष्पसंस्वेदितं भवति तन्निरुजं क्षणेन ॥}\ A}}}{\lemma{\skt{ऊर्धपरिषेचनतर्पनाद्यैः}}\Afootnote{K; \emph{om.}\ A}}\skt{ ।}\pend

\pstart\versepar \ledsidenote{\passref{6.19.8cd}}\edtext{\edtext{\skt{षट्सप्ततिर्नयनजा}}{\lemma{\skt{षट्सप्ततिर्}}\Afootnote{K; \skt{षट्सप्तति}\ A; \skt{षट्सप्तरिर्}\ H}}\skt{ य}}{\lemma{\skt{नयनजा य}}\Afootnote{K H; \skt{र्नयनजा य}\ A}}\skt{ इमे प्रदिष्टाः । }\pend
\pstart\versepar \skt{रोगा }\edtext{\edtext{\skt{भवन्त्यमहताम्महताञ्च}}{\lemma{\skt{भवन्त्य्}}\Afootnote{K H; \skt{भवन्}\ A}}}{\lemma{\skt{अमहताम्}}\Afootnote{K H; \skt{त्यमहतां}\ A}}\skt{ }\edtext{\skt{नेत्रे}}{\Afootnote{K H; \skt{तेभ्यः ॥}\ A}}\skt{ ।}\pend

\pstart\versepar \ledsidenote{\passref{6.19.9ab}}\edtext{\skt{स्तन्यप्रदोषजनितः}}{\Afootnote{K; \skt{स्तन्यप्रकोपकफमा}\ A; \skt{स्तन्यप्रयजनितः}\ H}}\skt{ कफरक्तदोषाद् । }\pend
\pstart\versepar \skt{बालाक्षिवर्त्मभय एव }\edtext{\skt{कुकूणको न्यस्}}{\Afootnote{K H; \skt{कुकूणकोऽन्यः ।}\ A}}\edtext{}{\lemma{\skt{कफरक्तदोषाद् । बालाक्षिवर्त्मभय}}\Afootnote{K\textsuperscript{pc}; \skt{रुतपित्तरक्तैर्बलाक्षिवर्त्मभव}\ A; \skt{कफरक्तदोषाद् । ब्लाक्षिवर्त्मभय}\ K\textsuperscript{ac}; \skt{कफरक्तदोषाद्बालाक्षिवर्त्म्मभव}\ H}}\pend

\pstart\versepar \ledsidenote{\passref{6.19.9cd}}\edtext{\edtext{\skt{तेनाक्षिकण्डूरथ}}{\lemma{\skt{तेनाक्षिकण्डूर्}}\Afootnote{K H; \skt{मृद्नाति ने}\ A}}}{\lemma{\skt{अथ}}\Afootnote{K H; \emph{om.}\ A}}\skt{ }\edtext{\skt{मर्दति}}{\Afootnote{K H; \skt{त्रमति}\ A}}\skt{ }\edtext{\skt{चाक्षिकूटो}}{\Afootnote{K H; \skt{कण्डुमथाक्षिकूटं}\ A}}\skt{ । }\pend
\pstart\versepar \skt{नासाललाटमपि }\edtext{\skt{चाथ}}{\Afootnote{K H; \skt{तेन}\ A}}\skt{ शिशुः स नित्यं}\pend

\pstart\versepar \ledsidenote{\passref{6.19.10ab}}\skt{सूर्यप्रभान्न सहते }\edtext{\edtext{\skt{मुहुराश्रवेच्च}}{\lemma{\skt{मुहुर्}}\Afootnote{K H; \emph{om.}\ A}}}{\lemma{\skt{आश्रवेच्च}}\Afootnote{K H; \skt{स्रवति प्रबद्धं}\ A}}\skt{ । }\pend
\pstart\versepar \skt{}\edtext{\skt{तस्याहते}}{\Afootnote{K H; \skt{तस्याहरे}\ A}}\skt{ }\edtext{\edtext{\skt{रुधिरमाशु}}{\lemma{\skt{रुधिरम्}}\Afootnote{K H; \skt{द्रुधिर}\ A}}}{\lemma{\skt{आशु}}\Afootnote{K H; \skt{माशु}\ A}}\skt{ }\edtext{\skt{विनिर्लिखेद्वा}}{\Afootnote{K H; \skt{विनिर्लिखेच्च ।}\ A}}\pend

\pstart\versepar \ledsidenote{\passref{6.19.10cd}}\skt{क्षौद्रायुतैश्च }\edtext{\skt{कटकैः}}{\Afootnote{K H; \skt{कटुभिः}\ A}}\skt{ }\edtext{\skt{प्रतिसारयीतः}}{\Afootnote{K; \skt{प्रतिसारयेत्तु}\ A; \skt{प्रतिसारयीत}\ H}}\skt{ । }\pend
\pstart\versepar \skt{मातुः }\edtext{\skt{शिशोरभिहतञ्च}}{\lemma{\skt{अभिहतञ्च}}\Afootnote{K H; \skt{अभिहितं च}\ A}}\skt{ विधिम्विदध्यात् ।}\pend

\pstart\versepar \ledsidenote{\passref{6.19.11}}\edtext{\skt{तं वामयेन्मधुकमाक्षिकसंयुतैश्च}}{\Afootnote{K H; \skt{तं वामयेत्तु मधुसैन्धवसंप्रयुक्तैः}\ A}}\skt{ }\pend
\pstart\versepar \skt{पीतम्पयः }\edtext{\skt{सह}}{\Afootnote{K H; \skt{खलु}\ A}}\skt{ फलैः }\edtext{\skt{खरमञ्जरोत्थैः}}{\Afootnote{K H; \skt{खरमञ्जरीणाम् ॥}\ A}}\skt{ ॥}\pend

\pstart\versepar \ledsidenote{\passref{6.19.12ab}}\edtext{\skt{स्यात्पिप्पलीलवणमाक्षिकसंयुतैर्वा}}{\Afootnote{K; \skt{स्यात्पिप्पलीलवणमाक्षिकसं}\ A H}}\skt{ }\pend
\pstart\versepar \skt{तोयैर्वमेत मधुमागधिकायुतैर्वा । }\pend
\pstart\versepar \skt{}\edtext{\skt{नैवं}}{\Afootnote{K H; \skt{नैनं}\ A}}\skt{ }\edtext{\skt{वमन्तमथ}}{\lemma{\skt{अथ}}\Afootnote{K H; \skt{अपि}\ A}}\skt{ }\edtext{\skt{वामयितुं}}{\Afootnote{A K\textsuperscript{pc} H; \skt{वायितुं}\ K\textsuperscript{ac}}}\skt{ यतेत ।}\edtext{}{\lemma{\skt{युतैर्… मधुमागधिकायु}}\Afootnote{\emph{om.}\ A}}\edtext{}{\lemma{\skt{युतैर्… मधुमागधिकायु}}\Afootnote{\emph{om.}\ H}}\pend

\pstart\versepar \ledsidenote{\passref{6.19.12cd}}\skt{दत्वा }\edtext{\skt{वचामदनदुग्धभुजि}}{\Afootnote{K H; \skt{वचामशनदुग्धभुजे}\ A}}\skt{ }\edtext{\skt{प्रयोज्य}}{\Afootnote{A K; \skt{प्रयोज्यम्}\ H}}\skt{ }\edtext{\skt{पूर्वन्}}{\Afootnote{K; \skt{मूर्ध्वं}\ A; \skt{ऊर्द्ध्वन्}\ H}}\skt{ }\pend
\pstart\versepar \skt{ततः फलयुतं वमनं विधिज्ञैः ।}\pend

\pstart\versepar \ledsidenote{\passref{6.19.13ab}}\edtext{\skt{जम्ब्वाम्रधात्रीवरुणेषु च}}{\Afootnote{K; \skt{जम्ब्वाम्रधात्र्यणुदलैः परि}\ A; \skt{जम्ब्वाम्रधातृतरुणेषु च}\ H}}\skt{ धावनार्थं }\pend
\pstart\versepar \skt{कार्यं }\edtext{\edtext{\skt{कषायमथ}}{\lemma{\skt{कषायम्}}\Afootnote{K H; \skt{कषाय}\ A}}}{\lemma{\skt{अथ}}\Afootnote{K H; \skt{मव}\ A}}\skt{ }\edtext{\skt{सेवनमेव}}{\lemma{\skt{सेवनम्}}\Afootnote{K H; \skt{सेचनम्}\ A}}\skt{ }\edtext{\skt{चापि}}{\Afootnote{A K; \skt{वापि ।}\ H}}\skt{ ।}\pend

\pstart\versepar \ledsidenote{\passref{6.19.13cd}}\edtext{\skt{आश्च्योतने पि च}}{\Afootnote{K H; \skt{आश्च्योतने च}\ A}}\skt{ }\edtext{\skt{घृतं }\edtext{\skt{ह्यमृतं}}{\lemma{\skt{अमृतं}}\Afootnote{K H; \skt{त्र घृतं}\ A}}}{\lemma{\skt{घृतं ह्य्}}\Afootnote{K H; \skt{हितम}\ A}}\skt{ }\edtext{\skt{विपक्वम्}}{\Afootnote{K H; \skt{गुडूची}\ A}}\skt{ }\pend
\pstart\versepar \skt{}\edtext{\skt{इष्टं}}{\Afootnote{K H; \skt{सिद्धं}\ A}}\skt{ }\edtext{\edtext{\skt{तथाहुरपि}}{\lemma{\skt{तथाहुर्}}\Afootnote{K H; \skt{तथाहु}\ A}}\skt{ वा}}{\lemma{\skt{अपि वा}}\Afootnote{K H; \skt{रपि च}\ A}}\skt{ }\edtext{\skt{त्रिफलाविपक्वं}}{\Afootnote{A K; \skt{तृफलाकृतञ्च ॥}\ H}}\skt{ ।}\pend

\pstart\versepar \ledsidenote{\passref{6.19.14ab}}\edtext{\skt{नेपालिशंखमरिचानि}}{\Afootnote{K H; \skt{नेपालजामरिचशङ्ख}\ A}}\skt{ }\edtext{\skt{रसाञ्जनञ्च}}{\Afootnote{K H; \skt{रसाञ्जनानि}\ A}}\skt{ }\pend
\pstart\versepar \skt{सिन्धुप्रसूतगुडमाक्षिकसंयुतानि ।}\pend

\pstart\versepar \ledsidenote{\passref{6.19.14cd}}\skt{स्यादञ्जनं }\edtext{\skt{मधुरसामधुताम्रकैर्वा}}{\Afootnote{K H; \skt{मधुरसामधुकाम्रकैर्वा}\ A}}\skt{ }\pend
\pstart\versepar \skt{कृष्णायसं }\edtext{\skt{मधुरसे}}{\Afootnote{K; \skt{घृतपयो मधु}\ A; \skt{मधुघृते}\ H}}\skt{ }\edtext{\skt{त्वथ}}{\Afootnote{K H; \emph{om.}\ A}}\skt{ वापि }\edtext{\skt{घृष्टं}}{\Afootnote{K H; \skt{दग्धम् ॥}\ A}}\skt{ ।}\pend

\pstart\versepar \ledsidenote{\passref{6.19.15ab}}\skt{व्योषम्पलाण्डु मधुकं लवणोत्तमञ्च }\pend
\pstart\versepar \skt{लाक्षाञ्च गैरिकयुतां }\edtext{\skt{गुडिकाञ्जनम्वा}}{\lemma{\skt{गुडिकाञ्जनम्}}\Afootnote{K H; \skt{गुटिकाञ्जनं}\ A}}\skt{ ।}\pend

\pstart\versepar \ledsidenote{\passref{6.19.15cd}}\skt{निम्बच्छदं }\edtext{\skt{मधुकदार्व्यथ}}{\lemma{\skt{अथ}}\Afootnote{K; \skt{स}\ A; \skt{सि}\ H}}\skt{ }\edtext{\skt{ताम्रलोध्रम्}}{\Afootnote{A K; \skt{ता प्रलोध्रम्}\ H}}\skt{ }\pend
\pstart\versepar \skt{इच्छन्ति चात्र भिषजोऽञ्जनमंशतुल्यं ।}\pend

\pstart\versepar \ledsidenote{\passref{6.19.16ab}}\skt{स्रोतोजशंखदधिसैन्धवमर्धपक्षम् }\pend
\pstart\versepar \skt{}\edtext{\skt{आभावितं तु}}{\Afootnote{K H; \skt{शु}\ A}}\skt{ }\edtext{\edtext{\skt{शिशुशुक्रतदञ्जनन्तु}}{\lemma{\skt{शिशुशुक्रतदञ्जनन्}}\Afootnote{K; \skt{क्रं शिशोर्नुदति भावितमञ्जने}\ A; \skt{शिशुशुक्रतञ्जनन्तु ।}\ H}}}{\lemma{\skt{तु}}\Afootnote{K; \skt{न ।}\ A; \emph{om.}\ H}}\skt{ ।}\pend

\pstart\versepar \ledsidenote{\passref{6.19.16cd}}\skt{स्यन्दे कफादभिहितं क्रममाचरेच्च }\pend
\pstart\versepar \skt{बालस्य }\edtext{\skt{योगकुशलः}}{\Afootnote{K H; \skt{रोगकुशलो}\ A}}\skt{ }\edtext{\skt{कुकुणं}}{\Afootnote{K; \skt{ऽक्षिगदं}\ A; \skt{कुकूणं}\ H}}\skt{ जिघांसुः ॥}\pend

\pstart \ledsidenote{\passref{6.19.17ab}}\edtext{\skt{समुद्रमिव}}{\lemma{\skt{समुद्रम्}}\Afootnote{K H; \skt{समुद्र}\ A}}\skt{ गम्भीरं नैव शक्यं }\edtext{\skt{चिकित्सितुं}}{\Afootnote{K H; \skt{चिकित्सितम् ।}\ A}}\skt{ ।}\pend

\pstart \ledsidenote{\passref{6.19.17cd}}\skt{वक्तुन्निरवशेषेण श्लोकानांमयुतैरपि ।}\pend

\pstart \ledsidenote{\passref{6.19.18ab}}\edtext{\skt{सहस्रैरपि वा }\edtext{\skt{प्युक्तमर्थमल्पमतिर्नरः}}{\lemma{\skt{उक्तम्}}\Afootnote{K H; \skt{क्तम्}\ A}}}{\lemma{\skt{अपि वा प्य्}}\Afootnote{K H; \skt{अपि वा प्रो}\ A}}\skt{ ।}\pend

\pstart \ledsidenote{\passref{6.19.18cd}}\skt{तर्कग्रन्थार्थरहितो }\edtext{\skt{नैव}}{\Afootnote{A K H\textsuperscript{pc}; \skt{नैवव}\ H\textsuperscript{ac}}}\skt{ }\edtext{\edtext{\skt{गृह्णात्यपण्डितः}}{\lemma{\skt{गृह्णात्य्}}\Afootnote{K H; \skt{गृह्णा}\ A}}}{\lemma{\skt{अपण्डितः}}\Afootnote{K H; \skt{त्यपण्डितः ॥}\ A}}\skt{ ।}\pend

\pstart \ledsidenote{\passref{6.19.19ab}}\edtext{\skt{यदिदं}}{\Afootnote{K H; \skt{तदिदं}\ A}}\skt{ बहुगूढार्थं चिकित्साबीजमीरितं ।}\pend

\pstart \ledsidenote{\passref{6.19.19cd}}\skt{कुशलेनाभिपन्नन्तद्बहुधाभिप्ररोहति ।}\pend

\pstart \ledsidenote{\passref{6.19.20ab}}\edtext{\skt{तस्मान्मतिमता}}{\Afootnote{K; \skt{तस्मान्मतिमता}\ A; \skt{तस्मात्मतिमता}\ H}}\skt{ नित्यं नानाशास्त्रार्थदर्शिना ।}\pend

\pstart \ledsidenote{\passref{6.19.20cd}}\edtext{\skt{सर्वमूह्यमगूढार्थं}}{\lemma{\skt{अगूढार्थं}}\Afootnote{K H; \skt{अगाधार्थं}\ A}}\skt{ }\edtext{\skt{शास्त्रमागमबुद्धिनेति}}{\lemma{\skt{आगमबुद्धिनेति}}\Afootnote{K H; \skt{आगमबुद्धिना ॥}\ A}}\skt{ ॥}\pend

\pstart \ledsidenote{\passref{6.19.trailer}}\edtext{\skt{शालाक्ये ॥ १८ ॥ {\symfont ❈}}}{\Afootnote{K; \skt{इति सुश्रुतसंहितायामुत्तरतन्त्रान्तर्गते शालाक्यतन्त्रे नयनाभिघातचिकित्सितं नामैकोनविंशोऽध्यायः ॥१९ ॥}\ A; \skt{शालाक्ये नयनाभिघातप्रतिषेधोऽष्टादशमो ध्यायः ॥}\ H}}\skt{ ॥}\pend

\endnumbering
\beginnumbering
\eledsection*{Adhyāya 20}
\pstart \ledsidenote{\passref{6.20.1}}\skt{अथातः }\edtext{\skt{कर्ण्णगतरोगविज्ञानीयं}}{\Afootnote{H; \skt{कर्णगतरोगविज्ञानीयमध्यायं}\ A; \skt{कर्ण्णगतरोगविज्ञानीयं व्याव्याख्यास्यामः ॥}\ K}}\skt{ }\edtext{\skt{व्याख्यास्यामः}}{\Afootnote{A H; \emph{om.}\ K}}\skt{ ॥}\pend

\pstart\versepar \ledsidenote{\passref{6.20.3}}\edtext{\skt{कर्ण्णशूलः}}{\Afootnote{K H; \skt{कर्णशूलं}\ A}}\skt{ प्रणादश्च बाधिर्यं क्ष्वेड एव च । }\pend
\pstart\versepar \skt{}\edtext{\skt{कर्ण्णाश्रावः}}{\Afootnote{K H; \skt{कर्णस्रावः}\ A}}\skt{ }\edtext{\skt{कर्ण्णकण्डू}}{\Afootnote{K; \skt{कर्णकण्डूः}\ A H}}\skt{ }\edtext{\skt{कर्ण्णवर्चस्तथैव}}{\lemma{\skt{कर्ण्णवर्चस्}}\Afootnote{A K; \skt{कर्ण्णवच्चस्}\ H}}\skt{ च ।}\pend

\pstart\versepar \ledsidenote{\passref{6.20.4}}\edtext{\skt{क्रिमिकर्ण्णः}}{\Afootnote{K H; \skt{कृमिकर्ण}\ A}}\skt{ }\edtext{\skt{प्रतीनाहे}}{\Afootnote{K; \skt{प्रतिनाहौ}\ A; \skt{प्रतीनाहौ}\ H}}\skt{ }\edtext{\edtext{\skt{विद्रधिर्विविधस्तथा}}{\lemma{\skt{विद्रधिर्}}\Afootnote{K H; \skt{विद्रधि}\ A}}}{\lemma{\skt{विविधस्}}\Afootnote{K H; \skt{र्द्विविधस्}\ A}}\skt{ । }\pend
\pstart\versepar \skt{कर्ण्णपाकः }\edtext{\edtext{\skt{पूतिकर्ण्णस्तथैवार्शश्चतुर्विधं}}{\lemma{\skt{पूतिकर्ण्णस्}}\Afootnote{K H; \skt{पूतिकर्ण}\ A}}}{\lemma{\skt{तथैवार्शश्}}\Afootnote{K H; \skt{स्तथैवार्शश्}\ A}}\skt{ ।}\pend

\pstart\versepar \ledsidenote{\passref{6.20.5}}\edtext{\skt{तथार्बुदं}}{\Afootnote{K; \skt{कर्णार्बुदं}\ A}}\skt{ सप्तविधं शोफश्चापि चतुर्विधः । }\pend
\pstart\versepar \skt{एते कर्ण्णगता रोगा }\edtext{\skt{विंशदष्टौ च}}{\Afootnote{K H; \skt{अष्टाविंशति}\ A}}\skt{ }\edtext{\skt{कीर्त्तिताः}}{\Afootnote{K H; \skt{रीरिताः ॥}\ A}}\skt{ ।}\edtext{}{\lemma{\skt{… ए}}\Afootnote{\emph{om.}\ H}}\pend

\pstart\versepar \ledsidenote{\passref{6.20.6}}\skt{समीरणः }\edtext{\skt{स्रोतगतो}}{\Afootnote{K; \skt{श्रोत्रगतो}\ A H}}\skt{ }\edtext{\skt{न्यथाचरन्समन्ततः}}{\Afootnote{K; \skt{ऽन्यथाचरः समन्ततः}\ A; \skt{ऽन्यथाचरन्समन्ततः}\ H}}\skt{ }\edtext{\edtext{\skt{शूलमतीव}}{\lemma{\skt{शूलम्}}\Afootnote{K H; \skt{शूल}\ A}}}{\lemma{\skt{अतीव}}\Afootnote{K H; \skt{मतीव}\ A}}\skt{ कर्ण्णयोः । }\pend
\pstart\versepar \skt{करोति दोषैश्च }\edtext{\edtext{\skt{यथास्वमावृतः}}{\lemma{\skt{यथास्वम्}}\Afootnote{K; \skt{यथास्व}\ A; \skt{यथाश्वम्}\ H}}\skt{ स}}{\lemma{\skt{आवृतः स}}\Afootnote{K H; \skt{मावृतः स}\ A}}\skt{ कर्ण्णशूलः कथितो }\edtext{\skt{दुरासदः}}{\Afootnote{K H; \skt{दुराचरः ॥}\ A}}\skt{ ।}\pend

\pstart\versepar \ledsidenote{\passref{6.20.7}}\edtext{\skt{यदा च}}{\Afootnote{K H; \skt{यदा तु}\ A}}\skt{ }\edtext{\skt{तिष्ठेत्तु}}{\Afootnote{K H; \skt{नाडीषु}\ A}}\skt{ }\edtext{\edtext{\skt{विमार्गमागतः}}{\lemma{\skt{विमार्गम्}}\Afootnote{K H; \skt{विमार्ग}\ A}}}{\lemma{\skt{आगतः}}\Afootnote{K H; \skt{मागतः}\ A}}\skt{ }\edtext{\skt{समेत्य}}{\Afootnote{K H\textsuperscript{pc}; \skt{स ए}\ A; \skt{समे}\ H\textsuperscript{ac}}}\skt{ }\edtext{\skt{नाडीं}}{\Afootnote{K H; \emph{om.}\ A}}\skt{ }\edtext{\skt{सहशब्दवाहिनीं}}{\Afootnote{K H; \skt{व शब्दाभिवहासु तिष्ठ}\ A}}\skt{ । }\pend
\pstart\versepar \skt{}\edtext{\skt{शृणोति}}{\Afootnote{K H; \skt{ति । श्र्णोति}\ A}}\skt{ }\edtext{\skt{शब्दान्विविधान्तथा}}{\lemma{\skt{तथा}}\Afootnote{K H; \skt{स्तदा}\ A}}\skt{ नरः }\edtext{\edtext{\skt{प्रणादमेनं}}{\lemma{\skt{प्रणादम्}}\Afootnote{K H; \skt{प्रणाद}\ A}}}{\lemma{\skt{एनं}}\Afootnote{K H; \skt{मेनं}\ A}}\skt{ कथयन्ति चामयं ।}\pend

\pstart\versepar \ledsidenote{\passref{6.20.8}}\skt{स एव }\edtext{\skt{शब्दाभिवहाः}}{\Afootnote{K H; \skt{शब्दानुवहा}\ A}}\skt{ }\edtext{\skt{सिरा}}{\Afootnote{K H; \skt{यदा}\ A}}\skt{ }\edtext{\skt{यदा}}{\Afootnote{K H; \skt{सिराः}\ A}}\skt{ कफानुयातो व्यनुसृत्य तिष्ठति । }\pend
\pstart\versepar \skt{तदा }\edtext{\skt{नरस्याप्रतिकारसेविनो}}{\Afootnote{A K; \skt{नरस्य प्रतिकारसेविनो}\ H}}\skt{ }\edtext{\skt{भवेद्धि}}{\Afootnote{K H; \skt{भवेत्तु}\ A}}\skt{ }\edtext{\edtext{\skt{बाधिर्यमसंशयं}}{\lemma{\skt{बाधिर्यम्}}\Afootnote{K H; \skt{बाधिर्य}\ A}}}{\lemma{\skt{असंशयं}}\Afootnote{K H; \skt{मसंशयं}\ A}}\skt{ }\edtext{\skt{खलुः}}{\Afootnote{K; \skt{खलु ॥}\ A H}}\skt{ ।}\pend

\pstart\versepar \ledsidenote{\passref{6.20.9}}\edtext{\skt{श्रमात्क्लमाद्वापि हि}}{\Afootnote{K; \skt{श्रमात्क्षया}\ A; \skt{श्रमाक्लमाद्वापि हि}\ H}}\skt{ }\edtext{\skt{रूक्षभोजिनः}}{\Afootnote{K H; \skt{द्रूक्षकषायभोजना}\ A}}\skt{ }\edtext{\skt{समीरणः}}{\Afootnote{K H; \skt{त्समीरणः}\ A}}\skt{ शब्दपथे }\edtext{\skt{व्यवस्थितः}}{\Afootnote{K H; \skt{प्रतिष्ठितः ।}\ A}}\skt{ । }\pend
\pstart\versepar \skt{}\edtext{\skt{विरिक्तमूर्धस्य}}{\Afootnote{K H; \skt{विरिक्तशीर्षस्य}\ A}}\skt{ }\edtext{\skt{नरस्य}}{\Afootnote{K H; \skt{च शीत}\ A}}\skt{ }\edtext{\skt{दारुणं}}{\Afootnote{K H; \skt{सेविनः}\ A}}\skt{ करोति हि }\edtext{\edtext{\skt{क्ष्वेलमतीव}}{\lemma{\skt{क्ष्वेलम्}}\Afootnote{K H; \skt{क्ष्वेड}\ A}}}{\lemma{\skt{अतीव}}\Afootnote{K H; \skt{मतीव}\ A}}\skt{ कर्ण्णयोः ।}\pend

\pstart\versepar \ledsidenote{\passref{6.20.10}}\skt{शिरोभिघातादथवा निमज्जतो जले }\edtext{\skt{प्रपाकादथवापि}}{\Afootnote{K H; \skt{प्रपाकादयवापि}\ A}}\skt{ विद्रधेः । }\pend
\pstart\versepar \skt{}\edtext{\skt{स्रवेद्धि}}{\Afootnote{K H; \skt{स्रवेत्तु}\ A}}\skt{ पूयं }\edtext{\skt{खलु}}{\Afootnote{K H; \skt{श्रव}\ A}}\skt{ }\edtext{\skt{यत्र}}{\Afootnote{K H; \skt{णोऽ}\ A}}\skt{ }\edtext{\skt{कर्ण्णतः स}}{\Afootnote{K H; \skt{निलावृतः स}\ A}}\skt{ कर्ण्णसंस्राव इति प्रकीर्त्तितः}\pend

\pstart\versepar \ledsidenote{\passref{6.20.11}}\skt{कफेन }\edtext{\skt{कण्डूपचितेन}}{\Afootnote{K H; \skt{कण्डूः प्रचितेन}\ A}}\skt{ कर्ण्णयोर्भृशम्भवेत्स्रोतसि }\edtext{\skt{कण्डुसंज्ञितः}}{\Afootnote{K H; \skt{कर्णसंज्ञिते ।}\ A}}\skt{ । }\pend
\pstart\versepar \skt{}\edtext{\skt{विशेषिते}}{\Afootnote{K; \skt{विशोषिते}\ A H}}\skt{ श्लेष्मणि पित्ततेजसा }\edtext{\skt{ध्रुवं}}{\Afootnote{K H; \skt{नृणां}\ A}}\skt{ भवेत्स्रोतसि कर्ण्णगूथकः ।}\pend

\pstart\versepar \ledsidenote{\passref{6.20.12}}\edtext{\skt{स कर्ण्णगूथो}}{\Afootnote{K H; \skt{स कर्णविट्को}\ A}}\skt{ द्रवतां यदा गतो }\edtext{\skt{विलापितो}}{\Afootnote{K H; \skt{विलायितो}\ A}}\skt{ }\edtext{\skt{घ्राणसुखं}}{\Afootnote{K H; \skt{घ्राणमुखं}\ A}}\skt{ प्रपद्यते । }\pend
\pstart\versepar \skt{तदा स कर्ण्णप्रतिनाहसंज्ञितो भवेद्विकारः शिरसो }\edtext{\skt{नुवेदकृत्}}{\Afootnote{K; \skt{ऽभितापनः ॥}\ A; \skt{ऽनुवेदकृत् ॥}\ H}}\skt{ ।}\pend

\pstart\versepar \ledsidenote{\passref{6.20.13}}\skt{यदा तु मूर्च्छन्त्यथवापि जन्तवः }\edtext{\edtext{\skt{सृजन्त्यपत्यान्यथवापि}}{\lemma{\skt{अपत्यान्य्}}\Afootnote{K H; \skt{अपत्या}\ A}}}{\lemma{\skt{अथवापि}}\Afootnote{K H; \skt{न्यथवापि}\ A}}\skt{ }\edtext{\skt{मक्षिका}}{\Afootnote{K H; \skt{मक्षिकाः ।}\ A}}\skt{ । }\pend
\pstart\versepar \skt{तदञ्जनत्वाच्छ्रवणो }\edtext{\skt{निरुध्यते}}{\Afootnote{K H; \skt{निरुच्यते}\ A}}\skt{ भिषग्भिराद्यैः }\edtext{\skt{क्रिमिकर्ण्ण}}{\Afootnote{K H; \skt{कृमिकर्ण}\ A}}\skt{ }\edtext{\edtext{\skt{इत्यथ}}{\lemma{\skt{इत्य्}}\Afootnote{K H; \skt{को ग}\ A}}}{\lemma{\skt{अथ}}\Afootnote{K H; \skt{दः ॥}\ A}}\skt{ ।}\pend

\pstart\versepar \ledsidenote{\passref{6.20.14}}\edtext{\skt{क्षताभिघातप्रथितस्तु}}{\Afootnote{K; \skt{क्षताभिघातप्रभवस्तु}\ A H}}\skt{ }\edtext{\skt{विद्रधि}}{\Afootnote{K; \skt{विद्रधिर्}\ A H}}\skt{ भवेत्तथा दोषकृतो परः पुनः । }\pend
\pstart\versepar \skt{}\edtext{\skt{सरक्तपीतारुणमस्रमास्रवेत्प्रतोदधूमायनदाहचोषवान्}}{\lemma{\skt{आस्रवेत्प्रतोदधूमायनदाहचोषवान्}}\Afootnote{K H; \skt{आस्रवेत्प्रतोदधूमायनदाहचोपवान् ॥}\ A}}\skt{ ।}\pend

\pstart\versepar \ledsidenote{\passref{6.20.15}}\skt{भवेत्प्रपाकः खलु पित्तकोपतो }\edtext{\skt{विकोथविक्लेदकरस्तु}}{\Afootnote{K H; \skt{विकोथविक्लेदकरश्च}\ A}}\skt{ कर्ण्णयोः । }\pend
\pstart\versepar \skt{}\edtext{\skt{स्रोतस्थिते}}{\Afootnote{K; \skt{स्थिते कफे}\ A; \skt{स्रोत्रस्थिते}\ H}}\skt{ }\edtext{\skt{श्लेष्मणि}}{\Afootnote{K H; \skt{स्रोतसि}\ A}}\skt{ }\edtext{\skt{पित्ततेजसा}}{\Afootnote{K H; \skt{पित्ततेजसः}\ A}}\skt{ }\edtext{\skt{विलीयमाने}}{\Afootnote{K H; \skt{विलाय्यमाने}\ A}}\skt{ }\edtext{\skt{भृशसंप्रतापिते}}{\Afootnote{K; \skt{भृशसंप्रतापवान् ॥}\ A; \skt{भृशसंप्रतापितः ॥}\ H}}\skt{ ।}\pend

\pstart\versepar \ledsidenote{\passref{6.20.16}}\skt{अवेदनो वाप्यथवा }\edtext{\skt{रुजान्वितो}}{\Afootnote{K H; \skt{सवेदनो}\ A}}\skt{ }\edtext{\skt{घणं}}{\Afootnote{K; \skt{घनं}\ A H}}\skt{ }\edtext{\skt{स्रवेत्पूतिकपूतिकर्ण्णतः}}{\Afootnote{K H; \skt{स्रवेत्पूति च पूतिकर्णकः ।}\ A}}\skt{ । }\pend
\pstart\versepar \skt{}\edtext{\edtext{\skt{प्रदिष्टलिङ्गान्यर्शांसि}}{\lemma{\skt{प्रदिष्टलिङ्गान्य्}}\Afootnote{K; \skt{प्रदिष्टलिङ्गान्य}\ A; \skt{प्रतिष्टलिंगान्य्}\ H}}}{\lemma{\skt{अर्शांसि}}\Afootnote{K H; \skt{रशांसि}\ A}}\skt{ }\edtext{\skt{यत्नतस्तथैव}}{\lemma{\skt{यत्नतस्}}\Afootnote{K H; \skt{तत्त्वतस्}\ A}}\skt{ }\edtext{\edtext{\edtext{\skt{शोफार्बुदलिङ्गमीरितमिति}}{\lemma{\skt{शोफार्बुदलिङ्गम्}}\Afootnote{K H; \skt{शोफार्बुदलिङ्ग}\ A}}}{\lemma{\skt{ईरितम्}}\Afootnote{K H; \skt{मीरितम् ।}\ A}}}{\lemma{\skt{इति}}\Afootnote{K; \skt{मया पुरस्तात्प्रसमीक्ष्य योजयेदिहैव तावत्प्रयतो भिषग्वरः ॥}\ A; \skt{इति ॥ {\symfont ❈} ॥}\ H}}\skt{ ॥}\pend

\pstart \ledsidenote{\passref{6.20.trailer}}\edtext{\skt{शालाक्ये १९ ॥ {\symfont ❈}}}{\Afootnote{K; \skt{इति सुश्रुतसंहितायामुत्तरतन्त्रान्तर्गते शालाक्यतन्त्रे कर्णगतरोगविज्ञानीयो नाम विंशतितमोऽध्यायः ॥२० ॥}\ A; \skt{शालाक्ये कर्ण्णगतरोगविज्ञानीयमेकोनविंशतितमो ध्यायः ॥ {\symfont ❈} ॥}\ H}}\skt{ ॥}\pend

\endnumbering
\beginnumbering
\eledsection*{Adhyāya 21}
\pstart \ledsidenote{\passref{6.21.1}}\skt{अथातः कर्णगतरोगप्रतिषेधं }\edtext{\skt{अध्यायं}}{\Afootnote{K H; \emph{om.}\ A}}\skt{ }\edtext{\skt{व्याख्यास्यामः}}{\Afootnote{A H; \emph{om.}\ K}}\skt{ ॥}\pend

\pstart\versepar \ledsidenote{\passref{6.21.4}}\skt{कर्णशूले प्रणादे च बाधिर्यक्ष्वेडयोरपि । }\pend
\pstart\versepar \skt{चतुर्ण्णामपि रोगाणां सामान्यं भेषजम्विदुः ।}\pend

\pstart\versepar \ledsidenote{\passref{6.21.5}}\skt{स्निग्धं वातहरैः }\edtext{\edtext{\skt{स्वेदैरुपस्निन्नं}}{\lemma{\skt{स्वेदैर्}}\Afootnote{K H; \skt{स्वेदै}\ A}}}{\lemma{\skt{उपस्निन्नं}}\Afootnote{K; \skt{र्नरं स्नेह}\ A; \skt{ऊपस्निन्नम्}\ H}}\skt{ विरेचितं । }\pend
\pstart\versepar \skt{नाडीस्वेदैरुपचरेत्पिण्डस्वेदैस्तथैव च ।}\pend

\pstart \ledsidenote{\passref{6.21.8}}\skt{मीनकुक्कुटलावानाम्मांसजैः }\edtext{\skt{पयसाथ}}{\Afootnote{K H; \skt{पयसापि}\ A}}\skt{ }\edtext{\skt{वा}}{\Afootnote{K H; \skt{वा । पिण्डैः स्वेदं च कुर्वीत कर्णशूलनिवारणम् ॥}\ A}}\skt{ ।}\pend

\pstart\versepar \ledsidenote{\passref{6.21.9}}\edtext{\skt{अश्वत्थपत्रवल्लीम्वा}}{\Afootnote{K H; \skt{अश्वत्थपत्रखल्लं वा}\ A}}\skt{ }\edtext{\skt{विधाय}}{\Afootnote{A K; \skt{विधार्य}\ H}}\skt{ }\edtext{\skt{बहुपुत्रिकां}}{\Afootnote{K H; \skt{बहुपत्रकम् ।}\ A}}\skt{ । }\pend
\pstart\versepar \skt{}\edtext{\skt{अभ्यज्याङ्गारसंपूर्ण्णं}}{\Afootnote{K; \skt{तदङ्गारैः सुसंपूर्णं}\ A; \skt{अभ्यज्यागारसंपूर्ण्णां}\ H}}\skt{ }\edtext{\skt{विदध्याच्छ्रवणे}}{\Afootnote{K; \skt{निदध्याछ्रवणो}\ A; \skt{निदध्याच्छ्रवणे}\ H}}\skt{ }\edtext{\skt{सुखं}}{\Afootnote{K H; \skt{परि ॥}\ A}}\skt{ ।}\pend

\pstart\versepar \ledsidenote{\passref{6.21.11}}\edtext{\skt{क्षौमगुग्गुल्वगरुभिः}}{\Afootnote{K; \skt{क्षौमगुग्गुल्वगुरुभिः}\ A H}}\skt{ }\edtext{\skt{सघृतैर्धूपयीत च}}{\lemma{\skt{धूपयीत च}}\Afootnote{K H; \skt{धूपयेच्च तम् ।}\ A}}\skt{ । }\pend
\pstart\versepar \skt{}\edtext{\skt{भुक्तोपरि}}{\Afootnote{K H; \skt{भक्तोपरि}\ A}}\skt{ हितं सर्पिर्बस्तिकर्म च पूजितं ।}\pend

\pstart\versepar \ledsidenote{\passref{6.21.12}}\edtext{\edtext{\skt{उष्णम्पयोनुपानम्वा}}{\lemma{\skt{उष्णम्}}\Afootnote{K H; \emph{om.}\ A}}}{\lemma{\skt{पयोनुपानम्वा}}\Afootnote{K H; \skt{निरन्नो नि}\ A}}\skt{ }\edtext{\edtext{\skt{त्रिरात्रम्पाययेद्घृतं}}{\lemma{\skt{त्रिरात्रम्}}\Afootnote{K H; \skt{शि तत्सर्पिः}\ A}}}{\lemma{\skt{पाययेद्घृतं}}\Afootnote{K H; \skt{पीत्वोपरि पिबेत्पयः ।}\ A}}\skt{ । }\pend
\pstart\versepar \skt{}\edtext{\skt{मूर्धबस्तिषु}}{\Afootnote{K; \skt{मूर्ध्वस्तिषु}\ A; \skt{मूर्द्ध्नि बस्तिषु}\ H}}\skt{ }\edtext{\skt{नस्ये च}}{\Afootnote{A K; \skt{नस्ये थ}\ H}}\skt{ मस्तिष्के परिषेचने ।}\pend

\pstart\versepar \ledsidenote{\passref{6.21.13}}\skt{शतपाकं बलातैलं }\edtext{\skt{प्रशंसंति च}}{\Afootnote{K H; \skt{प्रशस्तं चापि}\ A}}\skt{ भोजने । }\pend
\pstart\versepar \skt{}\edtext{\skt{कण्ठकारीमजाक्षीरे}}{\Afootnote{K; \skt{कण्टकारीमजाक्षीरे}\ A H}}\skt{ पक्त्वा क्षीरेण तेन }\edtext{\skt{तु}}{\Afootnote{K H; \skt{च ॥}\ A}}\skt{ ।}\pend

\pstart\versepar \ledsidenote{\passref{6.21.14}}\skt{विपचेत्कुक्कुटवसां }\edtext{\edtext{\skt{कर्ण्णयोस्तत्र}}{\lemma{\skt{कर्ण्णयोस्}}\Afootnote{K H; \skt{कर्णयो}\ A}}}{\lemma{\skt{तत्र}}\Afootnote{K H; \skt{स्तत्प्र}\ A}}\skt{ पूरणं । }\pend
\pstart\versepar \skt{तण्डुलीयकमूलानि }\edtext{\skt{फलमङ्कोठजं}}{\Afootnote{K H; \skt{फलमङ्कोलजं}\ A}}\skt{ तथा ।}\pend

\pstart\versepar \ledsidenote{\passref{6.21.15}}\edtext{\skt{अहिंसापेचुकामूलां}}{\Afootnote{K; \skt{अहिंसाकेन्दुकान्मूलं}\ A; \skt{अहिंसापेचुकामूलं}\ H}}\skt{ सरलन्देवदारु च । }\pend
\pstart\versepar \skt{लशुनं शृङ्गवेरञ्च तथा वंशावलेखनं ।}\pend

\pstart\versepar \ledsidenote{\passref{6.21.16}}\edtext{\skt{कल्कैरेषान्तथैवाम्ब्लैः}}{\lemma{\skt{तथैवाम्ब्लैः}}\Afootnote{K; \skt{तथाम्लैश्च}\ A; \skt{तथैवाम्लैः}\ H}}\skt{ }\edtext{\skt{पचेत्स्नेहचतुर्गुणं}}{\Afootnote{K H; \skt{पचेत्स्नेहं चतुर्विधम् ।}\ A}}\skt{ । }\pend
\pstart\versepar \skt{}\edtext{\edtext{\skt{वेदनानिग्रहार्थन्तद्धितं}}{\lemma{\skt{वेदनानिग्रहार्थन्}}\Afootnote{K H; \skt{वेदनायाः प्रशान्त्य}\ A}}}{\lemma{\skt{तद्धितं}}\Afootnote{K H; \skt{र्थं हितं}\ A}}\skt{ }\edtext{\skt{स्यात्कर्ण्णपूरणं}}{\Afootnote{K H; \skt{तत्कर्णपूरणम् ॥}\ A}}\skt{ ।}\pend

\pstart\versepar \ledsidenote{\passref{6.21.17}}\skt{लशुनार्द्रकशिग्रूणां }\edtext{\skt{मुरुंग्या}}{\Afootnote{K; \skt{मुरङ्ग्या}\ A H}}\skt{ }\edtext{\skt{मूलकस्य च}}{\Afootnote{A K; \skt{मूलमेव च ।}\ H}}\skt{ । }\pend
\pstart\versepar \skt{}\edtext{\skt{कदल्याश्च}}{\Afootnote{K H; \skt{कदल्याः स्व}\ A}}\skt{ रसः श्रेष्ठः }\edtext{\skt{सुखोष्णः}}{\Afootnote{K H; \skt{कदुष्णः}\ A}}\skt{ कर्ण्णपूरणे ।}\pend

\pstart\versepar \ledsidenote{\passref{6.21.18}}\edtext{\skt{क्षौद्रञ्च}}{\Afootnote{K H; \emph{om.}\ A}}\skt{ }\edtext{\skt{शृङ्गवेररञ्च}}{\Afootnote{K H; \skt{शृङ्गवेररसः}\ A}}\skt{ }\edtext{\skt{सैन्धवं}}{\Afootnote{K H; \skt{क्षौद्रं सैन्धवं}\ A}}\skt{ तैलमेव च । }\pend
\pstart\versepar \skt{}\edtext{\skt{कटुष्णः}}{\Afootnote{K; \skt{कदुष्णं}\ A; \skt{कटूष्णं}\ H}}\skt{ कर्ण्णयोर्देयमेतद्वा वेदनापहं ।}\pend

\pstart\versepar \ledsidenote{\passref{6.21.19}}\edtext{\skt{वंशावलेखसंयुक्तें}}{\Afootnote{K; \skt{वंशावलेखनायुक्ते}\ A; \skt{वर्शावलेखसंयुक्ते}\ H}}\skt{ मूत्रे }\edtext{\skt{वाजाविके}}{\Afootnote{K H\textsuperscript{pc}; \skt{चाजाविके}\ A; \skt{जाविके}\ H\textsuperscript{ac}}}\skt{ भिषक् । }\pend
\pstart\versepar \skt{सर्प्पिः पचेत्तेन कर्ण्णं पूरयेत्कर्ण्णशूलिनः ।}\pend

\pstart\versepar \ledsidenote{\passref{6.21.20}}\skt{महतः पञ्चमूलस्य }\edtext{\edtext{\skt{काण्डान्यष्टाङ्गुलानि}}{\lemma{\skt{काण्डान्य्}}\Afootnote{K H; \skt{काण्डम्}\ A}}\skt{ वा}}{\lemma{\skt{अष्टाङ्गुलानि वा}}\Afootnote{K H; \skt{अष्टादशाङ्गुलम् ।}\ A}}\skt{ । }\pend
\pstart\versepar \skt{}\edtext{\skt{क्षौमेणावेष्ट्य}}{\Afootnote{A K; \skt{क्षौमेणावेष्ट}\ H}}\skt{ संसिच्य तैलेनादीपयेत्ततः ।}\pend

\pstart\versepar \ledsidenote{\passref{6.21.21}}\skt{यत्तैलं च्यवते तेभ्यो }\edtext{\skt{दीप्तेभ्यो}}{\Afootnote{K H; \skt{धृतेभ्यो}\ A}}\skt{ भाजनोपरि । }\pend
\pstart\versepar \skt{}\edtext{\skt{एतत्तु}}{\Afootnote{K H; \skt{ज्ञेयं त}\ A}}\skt{ }\edtext{\skt{दीपिकातैलं}}{\Afootnote{K H; \skt{द्दीपिकातैलं}\ A}}\skt{ सद्यो गृह्णाति वेदनां ।}\pend

\pstart \ledsidenote{\passref{6.21.22}}\skt{कुर्यादेवं भद्रकाष्ठे कुष्ठे }\edtext{\skt{वासरलेपि वा}}{\Afootnote{K H; \skt{काष्ठे च सारले । मतिमान्दीपिकातैलं कर्णशूलनिबर्हणम् ॥}\ A}}\skt{ ।}\pend

\pstart\versepar \ledsidenote{\passref{6.21.28}}\edtext{\skt{मूत्रेष्वम्ब्लेषु}}{\lemma{\skt{अम्ब्लेषु}}\Afootnote{K; \skt{अम्लेषु}\ A H}}\skt{ वातघ्ने }\edtext{\skt{गणे च}}{\Afootnote{A K; \skt{गणेषु}\ H}}\skt{ }\edtext{\skt{कथिते}}{\Afootnote{K; \skt{क्वथिते}\ A H}}\skt{ भिषक् । }\pend
\pstart\versepar \skt{पचेच्चतुर्विधं स्नेहं }\edtext{\skt{पूरणार्थं तु}}{\Afootnote{K H; \skt{पूरणं तच्च}\ A}}\skt{ कर्ण्णयोः ।}\pend

\pstart \ledsidenote{\passref{6.21.29}}\edtext{\edtext{\skt{घृतस्नेहक्रियाश्चोक्ताः}}{\lemma{\skt{घृतस्नेहक्रियाश्}}\Afootnote{K H; \skt{एता एव क्रियाः}\ A}}}{\lemma{\skt{चोक्ताः}}\Afootnote{K; \skt{कुर्या}\ A; \skt{चोक्तं}\ H}}\skt{ }\edtext{\skt{शूले स्युः}}{\Afootnote{K H; \skt{त्पित्तघ्नैः}\ A}}\skt{ }\edtext{\skt{पित्तसंयुते}}{\Afootnote{K H; \skt{पित्तसंयुते । काकोल्यादौ दशक्षीरं तिक्तं चात्र हितं हविः ॥}\ A}}\skt{ ।}\pend

\pstart\versepar \ledsidenote{\passref{6.21.35}}\edtext{\skt{सर्षपेंगुदजाः}}{\Afootnote{K; \skt{सामा}\ A; \skt{सर्षपे गुदजाः}\ H}}\skt{ }\edtext{\skt{स्नेहः}}{\Afootnote{K H; \skt{न्यतो}\ A}}\skt{ }\edtext{\skt{सकफे}}{\Afootnote{K H; \skt{विशेषेण बा}\ A}}\skt{ }\edtext{\skt{पूरणं}}{\Afootnote{K H; \skt{धिर्ये पूरणं}\ A}}\skt{ }\edtext{\skt{हितं}}{\Afootnote{K H; \skt{शृणु ।}\ A}}\skt{ । }\pend
\pstart\versepar \skt{गवां मूत्रेण बिल्वानि पिष्ट्वा तैलं विपाचयेत् ।}\pend

\pstart\versepar \ledsidenote{\passref{6.21.36}}\edtext{\skt{सक्षारं}}{\Afootnote{K H; \skt{सजलं}\ A}}\skt{ }\edtext{\skt{तच्च}}{\Afootnote{K H; \skt{च स}\ A}}\skt{ }\edtext{\skt{वितरेद्बाधिर्ये}}{\Afootnote{K H; \skt{दुग्धं च बाधिर्ये}\ A}}\skt{ कर्ण्णपूरणं । }\pend
\pstart\versepar \skt{}\edtext{\skt{सितामधुकनिम्बैश्च}}{\Afootnote{K H; \skt{सितामधुकबिम्बीभिः}\ A}}\skt{ सिद्धम्वाजे }\edtext{\edtext{\skt{पयस्यथ}}{\lemma{\skt{पयस्य्}}\Afootnote{K H; \skt{पयस्य}\ A}}}{\lemma{\skt{अथ}}\Afootnote{K H; \skt{पि ॥}\ A}}\skt{ ।}\pend

\pstart\versepar \ledsidenote{\passref{6.21.37}}\edtext{\skt{उष्णं}}{\Afootnote{K H; \emph{om.}\ A}}\skt{ }\edtext{\skt{मथित्वा}}{\Afootnote{K H; \skt{बिम्}\ A}}\skt{ }\edtext{\skt{निःक्वाथे}}{\Afootnote{K H; \skt{बीक्वाथे}\ A}}\skt{ }\edtext{\skt{बेल्वशीतं}}{\Afootnote{K; \skt{विमथ्योष्णं शीतीभूतं}\ A; \skt{बैल्वशीतं}\ H}}\skt{ }\edtext{\skt{समुद्ध्रितं}}{\Afootnote{K; \skt{तदुद्धृतम् ।}\ A; \skt{समुद्धितं ।}\ H}}\skt{ । }\pend
\pstart\versepar \skt{पुनः पचेद्दशक्षीरं }\edtext{\skt{दत्वा}}{\Afootnote{K H; \skt{सिता}\ A}}\skt{ }\edtext{\skt{मधुकचन्दने}}{\Afootnote{K H; \skt{मधुकचन्दनैः ॥}\ A}}\skt{ ।}\pend

\pstart\versepar \ledsidenote{\passref{6.21.38}}\edtext{\skt{बिल्वाम्ब्लगाढं}}{\Afootnote{K; \skt{बिल्वाम्बुगाढं}\ A; \skt{बिल्वाम्लगाढं}\ H}}\skt{ तत्तैलं बाधिर्ये कर्ण्णपूरणं । }\pend
\pstart\versepar \skt{वक्ष्यते }\edtext{\skt{न्यत्\_ प्रतिश्याये}}{\Afootnote{K; \skt{यः प्रतिश्याये}\ A; \skt{ऽन्यः प्रतिश्याये}\ H}}\skt{ विधिः सो प्यत्र पूजितः ।}\pend

\pstart\versepar \ledsidenote{\passref{6.21.39}}\skt{वातव्याधिषु यश्चोक्तो }\edtext{\skt{विस्तरेण}}{\Afootnote{K H; \skt{विधिः स च}\ A}}\skt{ }\edtext{\skt{सुखावहः}}{\Afootnote{K H; \skt{हितो भवे}\ A}}\skt{ ॥ }\pend
\pstart\versepar \skt{}\edtext{\skt{कर्ण्णास्रावे}}{\Afootnote{K H; \skt{त् । कर्णस्रावे}\ A}}\skt{ पूतिकर्ण्णे तथैव }\edtext{\skt{क्रिमिकर्ण्णके}}{\Afootnote{K H; \skt{कृमिकर्णके ॥}\ A}}\skt{ ।}\pend

\pstart\versepar \ledsidenote{\passref{6.21.40}}\skt{समानं कर्म कुर्वीत }\edtext{\skt{योगां}}{\Afootnote{K; \skt{योगा}\ A H}}\skt{ }\edtext{\skt{वैशेषिकानपि}}{\Afootnote{K; \skt{न्वैशेषिकानपि}\ A H}}\skt{ । }\pend
\pstart\versepar \skt{शिरोविरेचनं चैव }\edtext{\skt{धूपणं}}{\Afootnote{K; \skt{धूपनं}\ A; \skt{लेपनं}\ H}}\skt{ पूरणं तथा ।}\pend

\pstart\versepar \ledsidenote{\passref{6.21.41}}\skt{प्रमार्जनन्धावनञ्च वीक्ष्य वीक्ष्यावचारयेत् । }\pend
\pstart\versepar \skt{राजवृक्षादितोयेन }\edtext{\skt{सुरसादिजलेण}}{\Afootnote{K; \skt{सुरसादिगणेन}\ A; \skt{सुरसादिजलेन}\ H}}\skt{ वा ।}\pend

\pstart\versepar \ledsidenote{\passref{6.21.42}}\skt{कर्ण्णप्रक्षालनं कार्यं }\edtext{\skt{चूर्ण्णैरेतैर्श्च}}{\lemma{\skt{एतैर्श्च}}\Afootnote{K H; \skt{एषां च}\ A}}\skt{ पूरयेत् । }\pend
\pstart\versepar \skt{चूर्ण्णपञ्ज्चकषायञ्च }\edtext{\edtext{\skt{कपित्थरसमेव}}{\lemma{\skt{कपित्थरसम्}}\Afootnote{K H; \skt{कपित्थरसयोजितम् ॥}\ A}}}{\lemma{\skt{एव}}\Afootnote{K H; \emph{om.}\ A}}\skt{ }\edtext{\skt{च}}{\Afootnote{K H; \emph{om.}\ A}}\skt{ ।}\edtext{}{\lemma{\skt{पूरयेत् । चूर्ण्णपञ्ज्चकषायञ्च}}\Afootnote{K H; \skt{पूरणम् । क्वाथं पञ्ज्चकषायं तु}\ A}}\pend

\pstart\versepar \ledsidenote{\passref{6.21.43}}\skt{कर्ण्णस्रावे प्रशंसन्ति पूरणं मधुना सह । }\pend
\pstart\versepar \skt{}\edtext{\skt{सर्जत्वक्चूर्ण्णसंयुक्तंः}}{\Afootnote{K; \skt{सर्जत्वक्चूर्णसंयुक्तः}\ A H}}\skt{ }\edtext{\skt{कर्प्पासीफलजो}}{\Afootnote{K H; \skt{कार्पासीफलजो}\ A}}\skt{ रसः ।}\pend

\pstart\versepar \ledsidenote{\passref{6.21.44}}\edtext{\skt{मधुना}}{\Afootnote{K H; \skt{योजितो मधुना}\ A}}\skt{ }\edtext{\skt{संश्रितः}}{\Afootnote{K H; \emph{om.}\ A}}\skt{ }\edtext{\skt{साधुः}}{\Afootnote{K H; \skt{वापि}\ A}}\skt{ }\edtext{\skt{कर्ण्णास्रावे}}{\Afootnote{K H; \skt{कर्णस्रावे}\ A}}\skt{ प्रशस्यते । }\pend
\pstart\versepar \skt{}\edtext{\skt{लाक्षासर्जरसौ}}{\Afootnote{K H; \skt{लाक्षा रसाञ्जनं}\ A}}\skt{ }\edtext{\skt{वापि}}{\Afootnote{K; \skt{सर्जश्}\ A; \skt{चापि}\ H}}\skt{ }\edtext{\skt{चूर्ण्णितौ}}{\Afootnote{K H; \skt{चूर्णितं}\ A}}\skt{ कर्ण्णपूरणं ।}\pend

\pstart\versepar \ledsidenote{\passref{6.21.45}}\edtext{\skt{सशैवलमहावृक्षजम्ब्वाम्रप्रसवायुतैः}}{\Afootnote{K H; \skt{सशैवलं महावृक्षजम्ब्वाम्रप्रसवायुतम् ।}\ A}}\skt{ । }\pend
\pstart\versepar \skt{}\edtext{\skt{कुलिंगक्षौद्रमाण्डूकैः}}{\Afootnote{K H; \skt{कुलीरक्षौद्रमण्डूकी}\ A}}\skt{ सिद्धं तैलञ्च पूजितं ।}\pend

\pstart\versepar \ledsidenote{\passref{6.21.46}}\edtext{\skt{तिन्दुकान्यभया}}{\Afootnote{\skt{तिन्दुकान्यभया}\ A K\textsuperscript{pc} H; \skt{तिन्दुकान्यतभया}\ K\textsuperscript{ac}}}\skt{ }\edtext{\skt{लोध्रं}}{\Afootnote{K H; \skt{रोध्रं}\ A}}\skt{ समङ्गामलकं }\edtext{\skt{तथा}}{\Afootnote{K H; \skt{मधु ।}\ A}}\skt{ । }\pend
\pstart\versepar \skt{पूरणञ्चात्र }\edtext{\skt{पथ्यां}}{\Afootnote{K; \skt{पथ्यं}\ A H}}\skt{ }\edtext{\edtext{\skt{स्यात्कपित्थरसमेव}}{\lemma{\skt{स्यात्कपित्थरसम्}}\Afootnote{K H; \skt{स्यात्कपित्थरसयोजितम् ॥}\ A}}}{\lemma{\skt{एव}}\Afootnote{K H; \emph{om.}\ A}}\skt{ }\edtext{\skt{च}}{\Afootnote{K H; \emph{om.}\ A}}\skt{ ।}\pend

\pstart\versepar \ledsidenote{\passref{6.21.47}}\edtext{\edtext{\skt{कपित्थाम्ररसञ्चापि}}{\lemma{\skt{कपित्थाम्ररसञ्}}\Afootnote{K H; \skt{रसमाम्रकपित्था}\ A}}}{\lemma{\skt{चापि}}\Afootnote{K H; \skt{नां}\ A}}\skt{ }\edtext{\skt{मधूकस्य}}{\Afootnote{K H; \skt{मधूकध}\ A}}\skt{ }\edtext{\skt{धवस्य च}}{\Afootnote{K H; \skt{वशालजम् ।}\ A}}\skt{ । }\pend
\pstart\versepar \skt{पूरणार्थं प्रशंसन्ति तैलं }\edtext{\skt{वार्त्तिविनाशनं}}{\Afootnote{K H; \skt{वा तैर्विपाचितम् ॥}\ A}}\skt{ ।}\pend

\pstart\versepar \ledsidenote{\passref{6.21.48}}\edtext{\skt{प्रियंगुमधुकाम्बष्ठाधातक्युत्पलपर्ण्णिनी}}{\Afootnote{K H; \skt{प्रियङ्गुमधुकाम्बष्ठाधातकीशिलपर्णिभिः ।}\ A}}\skt{ । }\pend
\pstart\versepar \skt{}\edtext{\skt{मञ्जिष्ठालोध्रलाक्षाभिः}}{\Afootnote{K H; \skt{मञ्जिष्ठालोभ्रलाक्षाभिः}\ A}}\skt{ कपित्थस्य रसेन }\edtext{\skt{च}}{\Afootnote{K H; \skt{वा ॥}\ A}}\skt{ ।}\pend

\pstart \ledsidenote{\passref{6.21.49}}\edtext{\skt{पचेत्तैलन्तदास्रावमभिगृह्णाति}}{\lemma{\skt{अभिगृह्णाति}}\Afootnote{K; \skt{अवगृह्णाति}\ A; \skt{अतिगृह्णाति}\ H}}\skt{ }\edtext{\skt{पूरणात्}}{\Afootnote{K H; \skt{पूरणात् । घृतं रसाञ्ज्नं नार्याः क्षीरेण मधुसंयुतम् ॥}\ A}}\skt{ ।}\pend

\pstart \ledsidenote{\passref{6.21.51}}\edtext{\skt{क्रिमिकर्ण्णविधानार्थं}}{\Afootnote{K; \skt{पूरणः पूतिकर्णस्य शमनो मधुसंयुतः । कृमिकर्णकनाशार्थं}\ A; \skt{क्रिमिकर्ण्णविनाशार्थं}\ H}}\skt{ }\edtext{\skt{क्रिमिघ्नं}}{\Afootnote{K H; \skt{कृमिघ्नं}\ A}}\skt{ }\edtext{\skt{योजयेद्विधिं}}{\Afootnote{A K\textsuperscript{pc} H; \skt{योजयेद्विध्ं ।}\ K\textsuperscript{ac}}}\skt{ ।}\pend

\pstart\versepar \ledsidenote{\passref{6.21.52}}\edtext{\skt{वार्ताकिधूमञ्च}}{\Afootnote{K H; \skt{वार्ताकुधूमश्च}\ A}}\skt{ हितः }\edtext{\skt{सर्षपस्नेह}}{\Afootnote{K H; \skt{सार्षपस्नेह}\ A}}\skt{ एव च । }\pend
\pstart\versepar \skt{}\edtext{\skt{पूरणं}}{\Afootnote{K H; \skt{कृमिघ्नं}\ A}}\skt{ हरितालेन गवाम्मूत्रयुतेन }\edtext{\skt{तु}}{\Afootnote{K H; \skt{च ॥}\ A}}\skt{ ।}\pend

\pstart\versepar \ledsidenote{\passref{6.21.53}}\edtext{\skt{धूपनं}}{\Afootnote{K H; \skt{गुग्गुलोः}\ A}}\skt{ कर्ण्णदौर्गन्ध्ये }\edtext{\skt{गुग्गुलुः}}{\Afootnote{K; \skt{धूपनं}\ A; \skt{गुग्गुलु}\ H}}\skt{ श्रेष्ठ }\edtext{\skt{उच्यते}}{\Afootnote{K H; \skt{मुच्यते ।}\ A}}\skt{ । }\pend
\pstart\versepar \skt{}\edtext{\skt{प्रर्च्छदनं}}{\Afootnote{K; \skt{छर्दनं}\ A; \skt{प्रच्छर्द्दनन्}\ H}}\skt{ }\edtext{\skt{धूमपानं}}{\Afootnote{K H; \skt{धूमपानं च}\ A}}\skt{ }\edtext{\skt{कवडस्य च}}{\Afootnote{K H; \skt{कवलस्य च}\ A}}\skt{ धारणं ।}\pend

\pstart\versepar \ledsidenote{\passref{6.21.54}}\edtext{\skt{कर्ण्णक्ष्वेडहितं}}{\Afootnote{K H; \skt{कर्णक्ष्वेडे हितं}\ A}}\skt{ तैलं }\edtext{\skt{सार्षपम्वापि}}{\lemma{\skt{वापि}}\Afootnote{K H; \skt{चैव}\ A}}\skt{ }\edtext{\skt{पूजितं}}{\Afootnote{K H; \skt{पूरणम् ।}\ A}}\skt{ । }\pend
\pstart\versepar \skt{विद्रधौ }\edtext{\skt{वापि}}{\Afootnote{K; \skt{चापि}\ A H}}\skt{ कुर्वीत विद्रध्युक्तं }\edtext{\skt{विधिं}}{\Afootnote{K H; \skt{चिकि}\ A}}\skt{ }\edtext{\skt{भिषक्}}{\Afootnote{K H; \skt{त्सितम् ॥}\ A}}\skt{ ।}\pend

\pstart\versepar \ledsidenote{\passref{6.21.55}}\edtext{\skt{क्लेदयित्वा तु}}{\Afootnote{K H; \skt{प्रक्लेद्य धीमां}\ A}}\skt{ }\edtext{\skt{तैलेन}}{\Afootnote{K H; \skt{स्तैलेन}\ A}}\skt{ स्वेदेन प्रविलाय्य च । }\pend
\pstart\versepar \skt{}\edtext{\skt{विशोधयेत्कर्ण्णगूथं}}{\Afootnote{K H; \skt{शोधयेत्कर्णविट्कं तु}\ A}}\skt{ भिषक्सम्यक्छलाकया ।}\pend

\pstart\versepar \ledsidenote{\passref{6.21.56}}\edtext{\skt{स्नेहः}}{\Afootnote{K; \skt{नाडी}\ A; \skt{स्नेह}\ H}}\skt{ }\edtext{\skt{स्वेदो}}{\Afootnote{K H; \skt{स्वेदो}\ A}}\skt{ }\edtext{\skt{थ}}{\Afootnote{K H; \skt{ऽथ}\ A}}\skt{ वमनं }\edtext{\skt{धूमो}}{\Afootnote{A K; \skt{धूपो}\ H}}\skt{ }\edtext{\skt{मूर्ध्नि}}{\Afootnote{K H; \skt{मूर्ध}\ A}}\skt{ विरेचनं । }\pend
\pstart\versepar \skt{विधिश्च }\edtext{\skt{कफहा}}{\Afootnote{K H; \skt{कफहृ}\ A}}\skt{ }\edtext{\skt{सर्वं}}{\Afootnote{K H; \skt{त्सर्वः}\ A}}\skt{ }\edtext{\skt{कर्ण्णकण्डूमपोहति}}{\lemma{\skt{अपोहति}}\Afootnote{A K; \skt{अतीष्यते ।}\ H}}\skt{ ॥}\pend

\pstart\versepar \ledsidenote{\passref{6.21.57}}\skt{अथ कर्ण्णप्रतीनाहे स्नेहस्वेदौ प्रयोजयेत् । }\pend
\pstart\versepar \skt{ततो विरिक्तशिरसः क्रियां प्राप्तां समाचरेत् ।}\pend

\pstart\versepar \ledsidenote{\passref{6.21.58}}\skt{कर्ण्णपाकस्य भैषज्यं कुर्यात्क्षतविसर्पवत् । }\pend
\pstart\versepar \skt{कर्ण्णच्छिद्रे }\edtext{\skt{वत्तमानं}}{\Afootnote{K; \skt{वर्तमानं}\ A H}}\skt{ }\edtext{\skt{किटूक्लेदकफादिकं}}{\Afootnote{K H\textsuperscript{pc}; \skt{कीटं क्लेदमलादि वा ॥}\ A; \skt{क टूक्लेदकफादिकं ।}\ H\textsuperscript{ac}}}\skt{ ।}\edtext{}{\lemma{\skt{कुर्यात्क्षतविसर्पवत् । कर्ण्णच्छिद्रे}}\Afootnote{K H; \skt{कुर्यात्पित्तविसर्पवत् । कर्णच्छिद्रे}\ A}}\pend

\pstart\versepar \ledsidenote{\passref{6.21.59}}\edtext{\skt{शृङ्गेणापहरेद्द्वीपानथवापि}}{\Afootnote{K; \skt{शृङ्गेणापहरेद्धीमानथवापि}\ A; \skt{शृङ्गेणापचरेद्धीमानथवापि}\ H}}\skt{ शलाकया । }\pend
\pstart\versepar \skt{शेषाणान्तु }\edtext{\skt{विकाराणा}}{\Afootnote{K; \skt{विकाराणां}\ A H}}\skt{ }\edtext{\edtext{\skt{प्राक्चिकित्सितमीतितं}}{\lemma{\skt{प्राक्चिकित्सितम्}}\Afootnote{K H; \skt{प्राक्चिकित्सित}\ A}}}{\lemma{\skt{ईतितं}}\Afootnote{K; \skt{मीरितम् ॥}\ A; \skt{ईरितमिति ॥}\ H}}\skt{ ॥}\pend

\pstart \ledsidenote{\passref{6.21.trailer}}\edtext{\skt{शालाक्ये}}{\Afootnote{K; \skt{इति सुश्रुतसंहितायामुत्तरतन्त्रान्तर्गते शालाक्यतन्त्रे कर्णगतरोगप्रतिषेधो नामैकविंशोऽध्यायः ॥२१ ॥}\ A; \skt{शालाक्ये कर्ण्णरोगप्रतिषेधो विंशतितमोऽध्यायः ॥ ॥}\ H}}\skt{ ॥ २० ॥}\pend

\pstart\versepar \ledsidenote{\passref{6.21.60}}\edtext{\skt{कफस्यन्दं}}{\Afootnote{K H; \emph{om.}\ A}}\skt{ }\edtext{\skt{रक्तस्यन्दं}}{\Afootnote{K H; \emph{om.}\ A}}\skt{ }\edtext{\skt{लेख्यं}}{\Afootnote{K H; \emph{om.}\ A}}\skt{ }\edtext{\skt{च्छ्वेद्यं च}}{\Afootnote{K; \emph{om.}\ A; \skt{च्छेद्यञ्च}\ H}}\skt{ }\edtext{\skt{पक्ष्मगं}}{\Afootnote{K H; \emph{om.}\ A}}\skt{ । }\pend
\pstart\versepar \skt{}\edtext{\skt{दृष्टिक्रियाविधाणञ्च}}{\Afootnote{K; \emph{om.}\ A; \skt{दृष्टिक्रियाविधानञ्च}\ H}}\skt{ }\edtext{\edtext{\skt{नेत्राघतस्तथैव}}{\lemma{\skt{नेत्राघतस्}}\Afootnote{K H; \emph{om.}\ A}}\skt{ च}}{\lemma{\skt{तथैव च}}\Afootnote{K H; \emph{om.}\ A}}\skt{ । }\pend
\pstart\versepar \skt{}\edtext{\skt{कर्ण्णरोगं}}{\Afootnote{K H; \emph{om.}\ A}}\skt{ }\edtext{\skt{प्रतीकारं}}{\Afootnote{K H; \emph{om.}\ A}}\skt{ }\edtext{\skt{प्रोक्तं वै}}{\Afootnote{K H; \emph{om.}\ A}}\skt{ }\edtext{\skt{द्वितीयो}}{\Afootnote{K H; \emph{om.}\ A}}\skt{ }\edtext{\skt{दश ॥ {\symfont ❈}}}{\Afootnote{K; \emph{om.}\ A; \skt{दश ॥}\ H}}\skt{ ॥ ० ॥}\pend

\endnumbering
\beginnumbering
\eledsection*{Adhyāya 22}
\pstart \ledsidenote{\passref{6.22.1}}\skt{अथातो }\edtext{\skt{घ्राणगतरोगविज्ञानीयं}}{\Afootnote{H; \skt{नासागतरोगविज्ञानीयमध्यायं}\ A; \skt{घ्राणगतरोगविज्ञानीयं व्याव्याख्यास्यामः ॥}\ K}}\skt{ }\edtext{\skt{व्याख्यास्यामः}}{\Afootnote{A H; \emph{om.}\ K}}\skt{ ॥}\pend

\pstart\versepar \ledsidenote{\passref{6.22.3}}\skt{अपीनसः }\edtext{\skt{पूतिनासं}}{\Afootnote{K H; \skt{पूतिनस्यं}\ A}}\skt{ नासापाकस्तथैव च । }\pend
\pstart\versepar \skt{तथा शोणितपित्तं च }\edtext{\skt{पूयशोणितमेव}}{\lemma{\skt{पूयशोणितम्}}\Afootnote{A K; \skt{पूयः शोणितम्}\ H}}\skt{ च ।}\pend

\pstart\versepar \ledsidenote{\passref{6.22.4}}\edtext{\skt{क्षवथुर्भ्रंशथुर्दीप्तिर्नासानाहः}}{\lemma{\skt{दीप्तिर्}}\Afootnote{K H; \skt{दीप्तो}\ A}}\skt{ परिस्रवः । }\pend
\pstart\versepar \skt{}\edtext{\skt{भवेच्छोषेण}}{\Afootnote{K H; \skt{नासाशोषेण}\ A}}\skt{ सहिता }\edtext{\skt{दशैकश्चेरिता}}{\lemma{\skt{दशैकश्}}\Afootnote{K H; \skt{दशैकाश्}\ A}}\skt{ गदाः ।}\pend

\pstart\versepar \ledsidenote{\passref{6.22.5}}\skt{चत्वार्यर्शांसि चत्वारः शोफाः सप्तार्बुदानि च । }\pend
\pstart\versepar \skt{प्रतिश्यायाश्च ये पञ्च वक्ष्यन्ते सचिकित्सिताः । }\pend
\pstart\versepar \skt{}\edtext{\skt{एकत्रिंशत्सहैतैस्तु}}{\Afootnote{K H; \skt{एकत्रिंशन्मितास्ते तु}\ A}}\skt{ नासारोगाः प्रकीर्त्तिताः ॥}\pend

\pstart\versepar \ledsidenote{\passref{6.22.6}}\edtext{\skt{आनाह्यते}}{\Afootnote{K; \skt{आनह्यते}\ A H}}\skt{ }\edtext{\skt{शुष्यति}}{\Afootnote{K H; \emph{om.}\ A}}\skt{ यस्य }\edtext{\skt{नासाः}}{\Afootnote{K; \emph{om.}\ A; \skt{नासा}\ H}}\skt{ }\edtext{\edtext{\edtext{\skt{प्रक्लेदमायात्यथ}}{\lemma{\skt{प्रक्लेदम्}}\Afootnote{K H; \emph{om.}\ A}}}{\lemma{\skt{आयात्य्}}\Afootnote{K H; \skt{वि}\ A}}}{\lemma{\skt{अथ}}\Afootnote{K H; \emph{om.}\ A}}\skt{ }\edtext{\skt{धूम्यते च}}{\Afootnote{K; \skt{धूप्यते च}\ A H}}\skt{ । }\pend
\pstart\versepar \skt{}\edtext{\skt{न}}{\Afootnote{K H; \skt{प्रक्लिद्यते शुष्यति चापि नासा । न}\ A}}\skt{ वेत्ति यो गन्धरसांश्च जन्तुर्जुष्टं व्यवस्येत्तमपीनसेन । }\pend

\pstart\versepar \ledsidenote{\passref{6.22.7}}\edtext{\skt{तञ्चानिलं}}{\Afootnote{K; \skt{तं चानिल}\ A H}}\skt{ श्लेष्मभवं विकारं }\edtext{\skt{ब्रूयात्प्रतिश्यायसमानलिङ्गं}}{\Afootnote{A K\textsuperscript{pc} H; \skt{ब्रूयात्प्रतिश्यायसमानि लिङ्गं ।}\ K\textsuperscript{ac}}}\skt{ । }\pend
\pstart\versepar \skt{दोषैर्विदग्धैर्गलतालुमूले }\edtext{\skt{संदूषितो}}{\Afootnote{K H; \skt{संवासितो}\ A}}\skt{ यस्य समीरणस्तु ।}\pend

\pstart\versepar \ledsidenote{\passref{6.22.8}}\skt{निरेति पूतिर्मुखनासिकाभ्यां तं पूतिनासं प्रवदन्ति रोगं । }\pend
\pstart\versepar \skt{घ्राणाश्रितं पित्तमरूंषि }\edtext{\skt{कुर्याद्यस्मिन्विपाके}}{\Afootnote{K H; \skt{कुर्याद्यस्मिन्विकारे}\ A}}\skt{ बलवांश्च पाकः ।}\pend

\pstart\versepar \ledsidenote{\passref{6.22.9}}\edtext{\skt{तन्नासिकापाक}}{\Afootnote{K H; \skt{तं नासिकापाकम्}\ A}}\skt{ इति }\edtext{\edtext{\skt{व्यवस्येद्विक्लेदकोथावपि}}{\lemma{\skt{व्यवस्येद्विक्लेदकोथाव्}}\Afootnote{K\textsuperscript{ac} H; \skt{व्यवस्येद्विक्लेदको}\ A K\textsuperscript{pc}}}}{\lemma{\skt{अपि}}\Afootnote{\skt{थावपि}\ A K\textsuperscript{pc}; \skt{अथ वापि}\ K\textsuperscript{ac} H}}\skt{ यत्र }\edtext{\skt{दृष्टौः}}{\Afootnote{K\textsuperscript{pc}; \skt{दृष्टौ ।}\ A; \emph{om.}\ K\textsuperscript{ac} H}}\skt{ । }\pend
\pstart\versepar \skt{चतुर्विधं द्विप्रभवं }\edtext{\skt{द्विमार्गं}}{\Afootnote{K H; \emph{om.}\ A}}\skt{ }\edtext{\skt{वक्ष्ये}}{\Afootnote{K H; \skt{वक्ष्या}\ A}}\skt{ }\edtext{\edtext{\skt{पुनश्चोपरि}}{\lemma{\skt{पुनश्}}\Afootnote{K H; \skt{मि भू}\ A}}}{\lemma{\skt{चोपरि}}\Afootnote{K H; \skt{यः खलु}\ A}}\skt{ रक्तपित्तं ।}\pend

\pstart\versepar \ledsidenote{\passref{6.22.10}}\skt{दोषैर्विदग्धैरथ वापि जन्तोर्ललाटदेशे }\edtext{\skt{भिहतस्य तैस्तैः}}{\Afootnote{K\textsuperscript{pc}; \skt{ऽभिहतस्य तैस्तु ।}\ A; \skt{भिहतस्य तैस्ताः ।}\ K\textsuperscript{ac}; \skt{ऽभिहतस्य तैस्तैः ।}\ H}}\skt{ । }\pend
\pstart\versepar \skt{नासा स्रवेत्पूयमसृग्विमिश्रं तं पूयरक्तं प्रवदन्ति रोगं ।}\pend

\pstart \ledsidenote{\passref{6.22.11}}\edtext{\skt{घ्राणाश्रितो}}{\Afootnote{K H; \skt{घ्राणाश्रिते}\ A}}\skt{ मर्मणि }\edtext{\skt{सम्प्रदुष्टो}}{\Afootnote{K H; \skt{संप्रदुष्टे}\ A}}\skt{ यस्यानिलो नासिकया निरेति ।}\pend

\pstart\versepar \ledsidenote{\passref{6.22.12}}\skt{कफानुयातो बहुशो }\edtext{\skt{तिशब्दस्तं}}{\Afootnote{K H; \skt{सशब्दस्तं}\ A}}\skt{ रोगमाहुः क्षवथुं }\edtext{\skt{विधिज्ञः}}{\Afootnote{K H; \skt{विधिज्ञाः ।}\ A}}\skt{ । }\pend
\pstart\versepar \skt{तीक्ष्णोपयोगादतिजिघ्रतो वा }\edtext{\skt{भावात्कटूनर्कनिरीक्षणाद्वा}}{\Afootnote{K H; \skt{भावान्कटूनर्कनिरीक्षणाद्वा ॥}\ A}}\skt{ ।}\pend

\pstart\versepar \ledsidenote{\passref{6.22.13}}\skt{सूत्रादिभिर्वा }\edtext{\skt{तरुणास्थिमर्मण्युद्घर्षिते}}{\lemma{\skt{उद्घर्षिते}}\Afootnote{K H; \skt{उद्घाटिते}\ A}}\skt{ }\edtext{\skt{न्यत् । क्षवथुर्निरेति}}{\lemma{\skt{न्यत् । क्षवथुर्}}\Afootnote{K; \skt{ऽन्यः क्षवथुर्}\ A; \skt{न्यः क्षवथुर्}\ H}}\skt{ । }\pend
\pstart\versepar \skt{}\edtext{\skt{प्रभ्रश्यते}}{\Afootnote{A K; \skt{प्रभ्रस्यते}\ H}}\skt{ }\edtext{\skt{नस्तुकयोर्हि}}{\Afootnote{K; \skt{नासिकयैव}\ A; \skt{नस्तुकाभ्यां हि}\ H}}\skt{ }\edtext{\skt{यस्य}}{\Afootnote{K H; \skt{यश्च}\ A}}\skt{ सान्द्रो विदग्धो लवणः }\edtext{\skt{कफस्तु}}{\lemma{\skt{कफस्}}\Afootnote{A K; \skt{कफास्}\ H}}\skt{ ।}\pend

\pstart\versepar \ledsidenote{\passref{6.22.14}}\skt{प्राक्सञ्चितो मूर्धनि }\edtext{\skt{सम्प्रतप्तस्तं}}{\Afootnote{K H; \skt{पित्ततप्तस्तं}\ A}}\skt{ }\edtext{\skt{भ्रंशथुं}}{\Afootnote{A K; \skt{भ्रंसथुं}\ H}}\skt{ }\edtext{\skt{रोगमुदाहरन्ति}}{\lemma{\skt{रोगम्}}\Afootnote{K H; \skt{व्याधिम्}\ A}}\skt{ । }\pend
\pstart\versepar \skt{}\edtext{\skt{घ्राणं}}{\Afootnote{K H; \skt{घ्राणे}\ A}}\skt{ भृशं }\edtext{\skt{दाहि तु}}{\Afootnote{K H; \skt{दा}\ A}}\skt{ }\edtext{\skt{यस्य}}{\Afootnote{K H; \skt{हस}\ A}}\skt{ }\edtext{\edtext{\skt{जन्त्वोर्वातश्च}}{\lemma{\skt{जन्त्वोर्}}\Afootnote{K; \skt{मन्वि}\ A; \skt{जन्तोर्}\ H}}}{\lemma{\skt{वातश्च}}\Afootnote{K H; \skt{ते तु}\ A}}\skt{ }\edtext{\skt{दीप्यन्निव}}{\Afootnote{K H; \skt{वि}\ A}}\skt{ }\edtext{\skt{निश्चरेच्च}}{\Afootnote{K H; \skt{निःसरेद्धूम इवेह वायुः ॥}\ A}}\skt{ ।}\pend

\pstart\versepar \ledsidenote{\passref{6.22.15}}\skt{नासा प्रदीप्ते }\edtext{\skt{नयनस्य}}{\Afootnote{K H; \skt{व च यस्य}\ A}}\skt{ }\edtext{\skt{चापि}}{\Afootnote{K H; \skt{जन्तोर्}\ A}}\skt{ व्याधिं तु }\edtext{\skt{दीप्तिं तमुदाहरन्ति}}{\lemma{\skt{दीप्तिं तम्}}\Afootnote{K H; \skt{तं दीप्तम्}\ A}}\skt{ । }\pend
\pstart\versepar \skt{कफावृतो }\edtext{\skt{वायुरुदारसंज्ञो}}{\lemma{\skt{उदारसंज्ञो}}\Afootnote{K H; \skt{उदानसंज्ञो}\ A}}\skt{ यदा स्वमार्गे }\edtext{\skt{विगुणस्तु}}{\Afootnote{K H; \skt{विगुणः स्थि}\ A}}\skt{ }\edtext{\skt{तिष्ठेत्}}{\Afootnote{K H; \skt{तः स्यात् ॥}\ A}}\skt{ ।}\pend

\pstart\versepar \ledsidenote{\passref{6.22.16}}\edtext{\skt{नासां}}{\Afootnote{K; \skt{घ्राणं}\ A; \skt{नासा}\ H}}\skt{ }\edtext{\edtext{\skt{वृणोत्यास्यगतः}}{\lemma{\skt{वृणोत्य्}}\Afootnote{K H; \skt{वृणोती}\ A}}}{\lemma{\skt{आस्यगतः}}\Afootnote{K H; \skt{व तदा}\ A}}\skt{ सरोगो नासाप्रतीनाह इति प्रदिष्टः । }\pend
\pstart\versepar \skt{}\edtext{\skt{अजस्रमच्छं}}{\lemma{\skt{अजस्रम्}}\Afootnote{A K; \skt{अजास्रम्}\ H}}\skt{ सलिलप्रकाशं }\edtext{\skt{यस्याविवर्ण्णं}}{\Afootnote{K H; \skt{यस्याविवर्णे}\ A}}\skt{ }\edtext{\skt{श्रवतीह}}{\Afootnote{K; \skt{स्रवतीह}\ A H}}\skt{ नासा ।}\pend

\pstart\versepar \ledsidenote{\passref{6.22.17}}\skt{रात्रौ }\edtext{\skt{विशेषेण च}}{\Afootnote{K H; \skt{विशेषेण हि}\ A}}\skt{ }\edtext{\skt{तत्तु}}{\Afootnote{K H; \skt{तं वि}\ A}}\skt{ }\edtext{\skt{रोगं}}{\Afootnote{K H; \skt{कारं}\ A}}\skt{ }\edtext{\skt{नासाप्रतिस्राव}}{\Afootnote{K H; \skt{नासापरिस्रावम्}\ A}}\skt{ इति व्यवस्येत् । }\pend
\pstart\versepar \skt{घ्राणाश्रिते }\edtext{\skt{स्रोतसि}}{\Afootnote{K H; \skt{श्लेष्मणि}\ A}}\skt{ }\edtext{\skt{मारुतेन}}{\Afootnote{K H; \skt{मारुतेन पित्तेन}\ A}}\skt{ गाढं }\edtext{\skt{प्रदीप्ते}}{\Afootnote{K H; \emph{om.}\ A}}\skt{ परिशोषिते च ।}\edtext{}{\lemma{\skt{व्यवस्येत् । घ्राणाश्रिते}}\Afootnote{A K; \skt{व्यवस्येत् ॥ घ्राणास्रिते}\ H}}\pend

\pstart\versepar \ledsidenote{\passref{6.22.18}}\edtext{\edtext{\skt{कृच्छ्राच्छ्वसेदूर्ध्वमधश्च}}{\lemma{\skt{कृच्छ्राच्छ्वसेदूर्ध्वम्}}\Afootnote{K H; \skt{समुच्छ्वसित्यूर्ध्व}\ A}}}{\lemma{\skt{अधश्च}}\Afootnote{K H; \skt{मधश्च}\ A}}\skt{ }\edtext{\edtext{\skt{जन्तुर्यस्मिं}}{\lemma{\skt{जन्तुर्}}\Afootnote{K H; \skt{कृच्छ्रा}\ A}}\skt{ स}}{\lemma{\skt{यस्मिं स}}\Afootnote{K; \skt{द्यस्तस्य}\ A; \skt{यस्मिन्स}\ H}}\skt{ नासापरिशोष उक्तः । }\pend
\pstart\versepar \skt{}\edtext{\skt{दोशैः}}{\Afootnote{K; \skt{दोषैस्}\ A H}}\skt{ }\edtext{\skt{समस्तैः}}{\Afootnote{K H; \skt{त्रिभिस्तैः}\ A}}\skt{ पृथगेकशश्च }\edtext{\skt{ब्रूयाद्यथार्शांसि}}{\Afootnote{K; \skt{ब्रूयात्तथार्शांसि}\ A H}}\skt{ तथैव }\edtext{\skt{शोफां}}{\Afootnote{K; \skt{शोफान्}\ A H}}\skt{ ।}\pend

\pstart\versepar \ledsidenote{\passref{6.22.19}}\edtext{\skt{स्वशास्त्रसिद्धान्तमवेक्ष्य}}{\lemma{\skt{स्वशास्त्रसिद्धान्तम्}}\Afootnote{K H; \skt{शालाक्यसिद्धान्तम्}\ A}}\skt{ चापि सर्वात्मकं }\edtext{\skt{सप्तविधार्बुदं}}{\Afootnote{K H; \skt{सप्तममर्बुदं}\ A}}\skt{ स्यात् । }\pend
\pstart\versepar \skt{रोगप्रतिश्याय }\edtext{\skt{इतीह}}{\Afootnote{K H; \skt{इहोप}\ A}}\skt{ }\edtext{\skt{यत्तु तं}}{\Afootnote{K H; \skt{दिष्टः स}\ A}}\skt{ वक्ष्यते }\edtext{\skt{पञ्चविधं}}{\Afootnote{K H; \skt{पञ्चविधः}\ A}}\skt{ पुरस्तात् ।}\edtext{}{\lemma{\skt{स्यात् । रोगप्रतिश्याय}}\Afootnote{K; \skt{तु । रोगः प्रतिश्याय}\ A; \skt{स्यात् ॥ रोगः प्रतिश्याय}\ H}}\pend

\pstart\versepar \ledsidenote{\passref{6.22.20}}\skt{नासास्रोतोगता }\edtext{\skt{रोगा}}{\Afootnote{K H; \skt{रोगास्}\ A}}\skt{ }\edtext{\skt{एकत्रिंशत्प्रकीर्त्तिताः}}{\Afootnote{K H; \skt{त्रिंशदेकश्च कीर्तिताः ।}\ A}}\skt{ । }\pend
\pstart\versepar \skt{}\edtext{\skt{स्रोतःपथे तु}}{\Afootnote{K; \skt{स्रोतःपथे य}\ A H}}\skt{ }\edtext{\skt{विपुलं}}{\Afootnote{K; \skt{द्विपुलं}\ A H}}\skt{ }\edtext{\skt{शोफव}}{\Afootnote{K H; \skt{कोशव}\ A}}\skt{ }\edtext{\skt{त्वर्बुदं}}{\Afootnote{K; \skt{च्चार्बुदं}\ A; \skt{त्त्वर्ब्बुदं}\ H}}\skt{ }\edtext{\skt{हितं}}{\Afootnote{K H; \skt{भवेत् ॥)}\ A}}\skt{ ।}\pend

\pstart \ledsidenote{\passref{6.22.21}}\skt{शोफास्तु }\edtext{\skt{शोफविज्ञाना}}{\Afootnote{A K; \skt{शोफविज्ञानान्}\ H}}\skt{ नासास्रोतोव्यवस्थिता }\edtext{\skt{इति}}{\Afootnote{K H; \skt{निदानेऽर्शांसि निर्दिष्टान्येवं तानि विभावयेत् ॥}\ A}}\skt{ ॥}\pend

\pstart \ledsidenote{\passref{6.22.trailer}}\edtext{\skt{शालाक्ये २१ ॥ {\symfont ❈}}}{\Afootnote{K; \skt{इति सुश्रुतसंहितायामुत्तरतन्त्रान्तर्गते शालाक्यतन्त्रे नासागतरोगविज्ञानीयो नाम द्वाविंशोऽध्यायः ॥२२ ॥}\ A; \skt{शालाक्ये घ्राणगतरोगविज्ञानीयमेकविंशतितमोऽध्यायः ॥ ॥}\ H}}\skt{ ॥}\pend

\endnumbering
\beginnumbering
\eledsection*{Adhyāya 23}
\pstart \ledsidenote{\passref{6.23.1}}\skt{अथातो }\edtext{\skt{घ्राणगतप्रतिषेधं}}{\Afootnote{K; \skt{नासागतरोगप्रतिषेधं}\ A; \skt{घ्राणगतरोगप्रतिषेधम्}\ H}}\skt{ }\edtext{\skt{व्याख्यास्यामः}}{\Afootnote{A; \skt{व्याव्याख्यास्यामः ॥}\ K; \skt{अध्यायं व्याख्यास्यामः ॥}\ H}}\skt{ ॥}\pend

\pstart\versepar \ledsidenote{\passref{6.23.3}}\edtext{\skt{अपीनसे}}{\Afootnote{K H; \skt{पूर्वोद्दिष्टे}\ A}}\skt{ }\edtext{\skt{पूतिनासे च}}{\Afootnote{K H; \skt{पूतिनस्ये च}\ A}}\skt{ जन्तोः }\edtext{\skt{स्नेहास्वेदौ}}{\Afootnote{K; \skt{स्नेहस्वेदौ}\ A H}}\skt{ च्छर्द्दनं स्रंसनञ्च । }\pend
\pstart\versepar \skt{}\edtext{\skt{हितं}}{\Afootnote{K H; \skt{युक्तं}\ A}}\skt{ }\edtext{\skt{भवेल्लघु}}{\Afootnote{K H; \skt{भक्तं}\ A}}\skt{ }\edtext{\skt{तीक्ष्णञ्च}}{\Afootnote{K H; \skt{तीक्ष्णम}\ A}}\skt{ }\edtext{\edtext{\skt{भक्तमुष्णन्तोयं}}{\lemma{\skt{भक्तम्}}\Afootnote{K; \skt{ल्पं लघु स्या}\ A; \skt{भक्षम्}\ H}}}{\lemma{\skt{उष्णन्}}\Afootnote{K\textsuperscript{pc} H; \skt{दुष्णं}\ A; \skt{उष्णा न्}\ K\textsuperscript{ac}}}\skt{ धूमपानञ्च }\edtext{\skt{कार्यं}}{\Afootnote{K H; \skt{काले ॥}\ A}}\skt{ ॥}\pend

\pstart\versepar \ledsidenote{\passref{6.23.4}}\edtext{\skt{कलिंगकं}}{\Afootnote{K H; \emph{om.}\ A}}\skt{ }\edtext{\skt{मरिचं}}{\Afootnote{K H; \emph{om.}\ A}}\skt{ हिंगु }\edtext{\skt{लाक्षा}}{\Afootnote{K H; \skt{व्योषं वत्सकाख्यं शिवाटी लाक्षा}\ A}}\skt{ }\edtext{\skt{शेफालिका}}{\Afootnote{K H; \skt{बीजं}\ A}}\skt{ }\edtext{\skt{सुरसा}}{\Afootnote{K H; \skt{सौरभं}\ A}}\skt{ कट्फलञ्च । }\pend
\pstart\versepar \skt{}\edtext{\skt{वचा}}{\Afootnote{K H; \skt{उग्रा}\ A}}\skt{ }\edtext{\skt{कुष्ठं}}{\Afootnote{A K; \skt{कुष्ठ}\ H}}\skt{ }\edtext{\skt{तिक्ष्णगन्धा}}{\Afootnote{K; \skt{तीक्ष्णगन्धा}\ A H}}\skt{ }\edtext{\skt{विडंगहितं}}{\Afootnote{K; \skt{विडङ्गं श्रेष्ठं}\ A; \skt{विडंगां हितं}\ H}}\skt{ नित्यं चावपीडे }\edtext{\skt{कणा}}{\Afootnote{K H; \skt{करञ्}\ A}}\skt{ }\edtext{\skt{च}}{\Afootnote{K H; \skt{जम् ॥}\ A}}\skt{ ।}\pend

\pstart\versepar \ledsidenote{\passref{6.23.5}}\edtext{\skt{एभिर्द्रव्यैः}}{\lemma{\skt{एभिर्}}\Afootnote{K H; \skt{एतैर्}\ A}}\skt{ सार्षपं }\edtext{\skt{चापि}}{\Afootnote{K H; \skt{मूत्रयुक्तं}\ A}}\skt{ तैलं }\edtext{\skt{पचेद्धीमान्नस्य}}{\Afootnote{K H; \skt{धीमान्नस्य}\ A}}\skt{ हेतोः }\edtext{\skt{समूत्रं}}{\Afootnote{K H; \skt{पचेत ।}\ A}}\skt{ । }\pend
\pstart\versepar \skt{नासापाके }\edtext{\skt{पित्तहरं}}{\Afootnote{K H; \skt{पित्तहृत्सं}\ A}}\skt{ }\edtext{\skt{विधाणं}}{\Afootnote{K; \skt{विधानं}\ A H}}\skt{ कार्यं }\edtext{\skt{सर्वम्बाह्यमभ्यन्तरञ्च}}{\lemma{\skt{अभ्यन्तरञ्}}\Afootnote{K H; \skt{आभ्यन्तरं}\ A}}\skt{ ।}\pend

\pstart\versepar \ledsidenote{\passref{6.23.6}}\edtext{\skt{हरेद्रक्तं}}{\Afootnote{K H; \skt{हृत्वा रक्तं}\ A}}\skt{ }\edtext{\skt{क्षीरवृक्षत्वचञ्च}}{\Afootnote{K; \skt{क्षीरवृक्षत्वचश्च}\ A; \skt{क्षत्व च}\ H}}\skt{ }\edtext{\skt{योज्याः}}{\Afootnote{K H; \skt{साज्याः}\ A}}\skt{ }\edtext{\skt{सेके}}{\Afootnote{K H; \skt{सेका}\ A}}\skt{ }\edtext{\skt{सघृताश्च}}{\Afootnote{K H; \skt{योजनीयाश्च}\ A}}\skt{ }\edtext{\skt{प्रदेहाः}}{\Afootnote{K H; \skt{लेपाः ।}\ A}}\skt{ । }\pend
\pstart\versepar \skt{}\edtext{\skt{वक्ष्याम्यूर्ध्वं}}{\lemma{\skt{वक्ष्याम्य्}}\Afootnote{A K\textsuperscript{pc} H; \skt{वक्ष्य्म्य्}\ K\textsuperscript{ac}}}\skt{ }\edtext{\skt{रक्तपित्तोपशान्तिं}}{\Afootnote{K H; \skt{रक्तपीतोपशान्तिं}\ A}}\skt{ }\edtext{\skt{नाडीवत्स्यात्पूयरक्तचिकित्स्यं}}{\Afootnote{K; \skt{नाडीवत्स्यात्पूयरक्ते चिकित्सा ॥}\ A; \skt{नाडीवत्स्यात्पूयरक्तं चिकित्सां ॥}\ H}}\skt{ ।}\pend

\pstart\versepar \ledsidenote{\passref{6.23.7}}\skt{वान्ते }\edtext{\skt{सम्यक्चावपीडाश्च}}{\Afootnote{K H; \skt{सम्यक्चावपीडं व}\ A}}\skt{ }\edtext{\edtext{\skt{योज्यास्तीक्ष्णा}}{\lemma{\skt{योज्यास्}}\Afootnote{K H; \skt{दन्ति}\ A}}}{\lemma{\skt{तीक्ष्णा}}\Afootnote{K H; \skt{तीक्ष्णं}\ A}}\skt{ }\edtext{\skt{धूमाः}}{\Afootnote{K H; \skt{धूमं}\ A}}\skt{ }\edtext{\skt{सोधनं}}{\Afootnote{K; \skt{शोधनं}\ A H}}\skt{ चात्र }\edtext{\skt{नस्यं}}{\Afootnote{K H; \skt{नस्यम् । क्षेप्यं नस्यं}\ A}}\skt{ । }\pend
\pstart\versepar \skt{}\edtext{\skt{वैरेचनं}}{\Afootnote{K H; \skt{मूर्धवैरेच}\ A}}\skt{ }\edtext{\skt{प्रथमे}}{\Afootnote{K H; \emph{om.}\ A}}\skt{ }\edtext{\skt{नस्ततस्तु}}{\Afootnote{K H; \skt{नीयै}\ A}}\skt{ }\edtext{\skt{नाड्या}}{\Afootnote{K H; \skt{र्नाड्या}\ A}}\skt{ चूर्ण्णं क्षवथौ }\edtext{\skt{भ्रंशथौ}}{\Afootnote{A K; \skt{भ्रन्सथौ}\ H}}\skt{ च ।}\pend

\pstart\versepar \ledsidenote{\passref{6.23.8}}\edtext{\edtext{\edtext{\skt{कुर्यात्स्वेदाञ्छिरसश्चानिलघ्नां}}{\lemma{\skt{कुर्यात्स्वेदाञ्}}\Afootnote{K; \skt{कुर्यात्स्वेदा}\ A; \skt{कुर्यात्स्वेदाश्}\ H}}}{\lemma{\skt{छिरसश्}}\Afootnote{K H; \skt{न्मूर्ध्नि वा}\ A}}}{\lemma{\skt{चानिलघ्नां}}\Afootnote{K; \skt{तामयघ्ना}\ A; \skt{चानिलघ्नं}\ H}}\skt{ }\edtext{\skt{दद्या}}{\Afootnote{K; \skt{न्स्निग्धान्}\ A; \skt{स्निग्धौ}\ H}}\skt{ }\edtext{\skt{धूमानन्यदन्यद्धितञ्च}}{\Afootnote{K; \skt{धूमान्यद्यन्यद्धितं च ।}\ A; \skt{धूमानन्यद्धितञ्च ।}\ H}}\skt{ । }\pend
\pstart\versepar \skt{दीप्ते रोगे }\edtext{\skt{पैत्तिकञ्च}}{\Afootnote{K H; \skt{पैत्तिकं सं}\ A}}\skt{ }\edtext{\skt{विधिं}}{\Afootnote{K H; \skt{विधानं}\ A}}\skt{ }\edtext{\skt{कुर्यान्मधुरं}}{\Afootnote{K H; \skt{कुर्यात्सर्वं स्वादु य}\ A}}\skt{ }\edtext{\skt{शीतलञ्च}}{\lemma{\skt{शीतलञ्}}\Afootnote{K H; \skt{च्छीतलं}\ A}}\skt{ ।}\pend

\pstart\versepar \ledsidenote{\passref{6.23.9}}\skt{नासानाहे स्नेहपानं }\edtext{\skt{प्रधाणं}}{\Afootnote{K; \skt{प्रधानं}\ A H}}\skt{ }\edtext{\skt{स्निग्धां}}{\Afootnote{K H; \skt{स्निग्धा}\ A}}\skt{ }\edtext{\skt{धूमां}}{\Afootnote{K H; \skt{धूमा}\ A}}\skt{ }\edtext{\skt{मूर्ध्नि}}{\Afootnote{K H; \skt{मूर्ध}\ A}}\skt{ बस्तिश्च नित्यं । }\pend
\pstart\versepar \skt{बलातैलं }\edtext{\skt{सर्वथा}}{\Afootnote{K H; \skt{सर्वथै}\ A}}\skt{ }\edtext{\edtext{\edtext{\skt{चोपयोज्यम्वातव्याधावुक्तमन्यच्च}}{\lemma{\skt{चोपयोज्यम्}}\Afootnote{K H; \skt{वोपयोज्यं}\ A}}}{\lemma{\skt{उक्तम्}}\Afootnote{K H; \emph{om.}\ A}}\skt{ य}}{\lemma{\skt{अन्यच्च य}}\Afootnote{K; \skt{अन्यदुक्तं च य}\ A; \skt{अन्यच्च यः}\ H}}\skt{ }\edtext{\skt{स्यात्}}{\Afootnote{K H; \skt{द्यत् ॥}\ A}}\skt{ ।}\pend

\pstart\versepar \ledsidenote{\passref{6.23.10}}\skt{नासास्रावे }\edtext{\edtext{\skt{घ्राणतश्चूर्ण्णमुक्तन्नाड्या}}{\lemma{\skt{चूर्ण्णम्}}\Afootnote{K H; \skt{चूर्ण}\ A}}}{\lemma{\skt{उक्तन्}}\Afootnote{K H; \skt{मुक्तं}\ A}}\skt{ }\edtext{\skt{देयं ये}}{\Afootnote{K H; \skt{देयं यो}\ A}}\skt{ }\edtext{\skt{वपीडाश्च}}{\Afootnote{K; \skt{ऽवपीडश्च}\ A; \skt{ऽवपीडाश्च}\ H}}\skt{ }\edtext{\skt{पथ्या}}{\Afootnote{K H; \skt{तीक्ष्णः ।}\ A}}\skt{ । }\pend
\pstart\versepar \skt{}\edtext{\edtext{\edtext{\skt{तीक्ष्णान्धूमान्देवदार्व्यार्द्रकाणां}}{\lemma{\skt{तीक्ष्णान्}}\Afootnote{K H; \skt{तीक्ष्णं}\ A}}}{\lemma{\skt{धूमान्}}\Afootnote{K H; \skt{धूमं}\ A}}}{\lemma{\skt{देवदार्व्यार्द्रकाणां}}\Afootnote{K; \skt{देवदार्वग्निकाभ्यां}\ A; \skt{देवदार्व्व्या र्द्रकानां}\ H\textsuperscript{ac}; \skt{देवदा र्वार्द्रकानां}\ H\textsuperscript{pc}}}\skt{ }\edtext{\skt{मांसञ्चाजं}}{\lemma{\skt{चाजं}}\Afootnote{K H; \skt{वाजं}\ A}}\skt{ }\edtext{\skt{हितमात्रादिशन्ति}}{\Afootnote{K; \skt{युक्तमत्रादिशन्ति ॥}\ A; \skt{हितमत्रादिशन्ति ॥}\ H}}\skt{ ।}\pend

\pstart\versepar \ledsidenote{\passref{6.23.11}}\skt{नासाशोषे क्षीरसर्पिः प्रधानं सिद्धं तैलं }\edtext{\skt{चानुकल्पेन}}{\Afootnote{K H; \skt{चाणुकल्पेन}\ A}}\skt{ नस्यं । }\pend
\pstart\versepar \skt{सर्पिःपानं भोजनं जाङ्गलैश्च }\edtext{\skt{स्नेहे}}{\Afootnote{K; \skt{स्नेहः}\ A H}}\skt{ }\edtext{\skt{स्वेदौ}}{\Afootnote{K H; \skt{स्वेदः}\ A}}\skt{ }\edtext{\edtext{\skt{स्नेहिकश्चात्र}}{\lemma{\skt{स्नेहिकश्}}\Afootnote{K; \skt{स्नैहिकश्}\ A H}}}{\lemma{\skt{चात्र}}\Afootnote{K; \skt{चापि}\ A H}}\skt{ धूमः ।}\pend

\pstart \ledsidenote{\passref{6.23.12}}\edtext{\skt{शेषां}}{\Afootnote{K; \skt{शेषा}\ A H}}\skt{ }\edtext{\skt{रोगां}}{\Afootnote{K; \skt{न्रोगा}\ A H}}\skt{ }\edtext{\skt{घ्राणजां}}{\Afootnote{K; \skt{न्घ्राणजा}\ A H}}\skt{ }\edtext{\skt{सन्नियच्छेद्यथा}}{\Afootnote{K; \skt{न्सन्नियच्छेदुक्तं}\ A; \skt{न्सन्नियच्छेत्यथा}\ H}}\skt{ तेषां }\edtext{\skt{कीर्त्तितं}}{\Afootnote{K H; \skt{यद्यथा}\ A}}\skt{ }\edtext{\edtext{\skt{संविधाणमिति}}{\lemma{\skt{संविधाणम्}}\Afootnote{K; \skt{संविधानम्}\ A H}}}{\lemma{\skt{इति}}\Afootnote{K H; \emph{om.}\ A}}\skt{ ॥ ॥}\pend

\pstart \ledsidenote{\passref{6.23.trailer}}\edtext{\skt{शालाक्ये}}{\Afootnote{K; \skt{इति सुश्रुतसंहितायामुत्तरतन्त्रान्तर्गते शालाक्यतन्त्रे नासागतरोगप्रतिषेधो नाम त्रयोविंशोऽध्यायः ॥२३ ॥}\ A; \skt{शालाक्ये घ्राणगतरोगप्रतिषेधो द्वाविंशतिमः ॥ ॥}\ H}}\skt{ २२ ॥ ० ।}\pend

\endnumbering
\beginnumbering
\eledsection*{Adhyāya 24}
\endnumbering
\beginnumbering
\eledsection*{Adhyāya 25}
\pstart \ledsidenote{\passref{6.25.1}}\skt{अथातः }\edtext{\skt{शिरोगतरोगविज्ञानीयं}}{\Afootnote{K; \skt{शिरोरोगविज्ञानीयम्}\ A; \skt{शिरोगतरोगविज्ञानी}\ H}}\skt{ }\edtext{\skt{अध्यायं}}{\Afootnote{A K; \skt{यं}\ H}}\skt{ }\edtext{\skt{व्याख्यास्यामः}}{\Afootnote{A H; \skt{व्या व्याख्यास्यामः ॥}\ K}}\skt{ ॥}\pend

\pstart\versepar \ledsidenote{\passref{6.25.3}}\skt{शिरो }\edtext{\skt{रुजानि}}{\Afootnote{K H; \skt{रुजति}\ A}}\skt{ }\edtext{\skt{मर्त्यानाम्वातपित्तकफैस्त्रिभिः}}{\lemma{\skt{वातपित्तकफैस्त्रिभिः}}\Afootnote{K; \skt{वातपित्तकफैस्त्रिभिः ।}\ A H}}\skt{ । }\pend
\pstart\versepar \skt{सन्निपातेन रक्तेन क्षयेण क्रिमिभिस्तथा ।}\pend

\pstart\versepar \ledsidenote{\passref{6.25.4}}\edtext{\skt{सूर्यावर्तावेभेदाभ्यां}}{\Afootnote{K; \skt{सूर्यावर्तानन्तवातार्धावभेदक}\ A; \skt{सूर्यावर्तावभेदाभ्यां}\ H}}\skt{ }\edtext{\skt{शंखकेन}}{\Afootnote{K H; \skt{शङ्खकैः । ए}\ A}}\skt{ }\edtext{\skt{तथैव च}}{\Afootnote{K H; \skt{का}\ A}}\skt{ । }\pend
\pstart\versepar \skt{}\edtext{\skt{दशप्रकारावस्थास्य}}{\Afootnote{K H; \skt{दशप्रकारस्य}\ A}}\skt{ लक्षणं सम्प्रवक्ष्यते ।}\pend

\pstart\versepar \ledsidenote{\passref{6.25.5}}\skt{यस्यानिमित्तं शिरसो }\edtext{\skt{रुजास्तु}}{\Afootnote{K H; \skt{रुजश्च}\ A}}\skt{ भवन्ति तीव्रा }\edtext{\skt{निवितातिमात्रं}}{\Afootnote{K H; \skt{निशि चातिमात्रम् ।}\ A}}\skt{ । }\pend
\pstart\versepar \skt{}\edtext{\skt{वन्धोपतापैश्च}}{\Afootnote{K H; \skt{बन्धोपतापैश्च}\ A}}\skt{ }\edtext{\skt{भवेद्भिशेषः}}{\Afootnote{K; \skt{भवेद्विशेषः}\ A; \skt{भवेद्भिशेषः}\ H}}\skt{ }\edtext{\skt{शिरोभिघातः स}}{\Afootnote{K H; \skt{शिरोभितापः स}\ A}}\skt{ समीरणेन ।}\pend

\pstart\versepar \ledsidenote{\passref{6.25.6}}\edtext{\skt{यस्योष्णमङ्गारनिभं}}{\Afootnote{K H; \skt{यस्योष्णमङ्गारचितं}\ A}}\skt{ यथैव }\edtext{\skt{दह्येच्छिरो}}{\Afootnote{K; \skt{दह्येत}\ A; \skt{दह्येच्छि{\symfont ¦}छिरो}\ H}}\skt{ }\edtext{\skt{धूप्यति}}{\Afootnote{K H; \skt{धूप्येत शि}\ A}}\skt{ }\edtext{\skt{मात्रं}}{\Afootnote{K H; \skt{रोक्सिनासम् ।}\ A}}\skt{ । }\pend
\pstart\versepar \skt{शीतेन रात्रौ च }\edtext{\skt{भवेद्विशेषः}}{\Afootnote{K; \skt{भवेद्विशेषः}\ A H}}\skt{ शिरोभितापः स तु पित्तकोपात् ।}\pend

\pstart\versepar \ledsidenote{\passref{6.25.7}}\edtext{\skt{शिरो}}{\Afootnote{K H; \skt{शिरो}\ A}}\skt{ }\edtext{\skt{भवेत्यस्य}}{\Afootnote{K H; \skt{गलं यस्य}\ A}}\skt{ कफोपदिग्धं गुरु }\edtext{\skt{प्रतिस्तवमथो}}{\Afootnote{K; \skt{प्रतिष्टब्धमथो}\ A; \skt{प्रतिस्तव्वमथो}\ H}}\skt{ हिमञ्च । }\pend
\pstart\versepar \skt{}\edtext{\skt{शूनाक्षिकूटां}}{\Afootnote{K; \skt{शूनाक्षिकूटं}\ A H}}\skt{ वदनञ्च यस्य }\edtext{\skt{शिराभिघातः}}{\Afootnote{K H; \skt{शिरोभितापः}\ A}}\skt{ }\edtext{\skt{सकप्रकोपात्}}{\Afootnote{\skt{स कफप्रकोपात् ॥}\ A; \skt{सककोपात् ।}\ K; \skt{कफप्रकोपात् ॥}\ H}}\skt{ ।}\pend

\pstart\versepar \ledsidenote{\passref{6.25.8}}\edtext{\skt{शिरोभिघाते}}{\Afootnote{K H; \skt{शिरोभितापे}\ A}}\skt{ }\edtext{\skt{त्रितयप्रवृत्त}}{\Afootnote{K H; \skt{त्रितयप्रवृत्ते स}\ A}}\skt{ र्वाणि लिङ्गानि }\edtext{\skt{मुहुर्भवन्ति}}{\Afootnote{K; \skt{समुद्भवन्ति ।}\ A; \skt{मुहुर्भवन्ति ॥}\ H}}\skt{ । }\pend
\pstart\versepar \skt{}\edtext{\skt{रक्तात्मकं}}{\Afootnote{K H; \skt{रक्तात्मकः}\ A}}\skt{ }\edtext{\skt{पित्तसमानलिङ्गं}}{\Afootnote{K H; \skt{पित्तसमानलिङ्गः}\ A}}\skt{ स्पर्शासहत्वं शिरसो भवेच्च ।}\pend

\pstart\versepar \ledsidenote{\passref{6.25.9}}\edtext{\skt{असृग्वसाश्लेष्मसमीरणानां}}{\Afootnote{K; \skt{वसाबलासक्षतसंभवानां}\ A; \skt{असृग्वसाश्लेष्मसमीरणां}\ H}}\skt{ }\edtext{\skt{शिरोगतानांमिहसंक्षयेण}}{\Afootnote{K; \skt{शिरोगतानामिह संक्षयेण}\ A H}}\skt{ । }\pend
\pstart\versepar \skt{}\edtext{\skt{क्षतप्रवृत्तिः}}{\Afootnote{K H; \skt{क्षयप्रवृत्तः}\ A}}\skt{ }\edtext{\skt{शिरसोभितापः}}{\Afootnote{K H\textsuperscript{pc}; \skt{शिरसोऽभितापः}\ A; \skt{शिरभितापः}\ H\textsuperscript{ac}}}\skt{ कष्टो }\edtext{\skt{भवेदुग्रजातिमात्रं}}{\Afootnote{K H; \skt{भवेदुग्ररुजोऽतिमात्रम् ॥}\ A}}\skt{ ।}\pend

\pstart\versepar \ledsidenote{\passref{6.25.10}}\edtext{\edtext{\skt{संस्वेदनश्च्छर्दनधूमनस्यै}}{\lemma{\skt{संस्वेदनश्}}\Afootnote{K H; \skt{संस्वेदन}\ A}}}{\lemma{\skt{च्छर्दनधूमनस्यै}}\Afootnote{K; \skt{च्छर्दनधूमनस्यैर्}\ A; \skt{च्छर्द्दनधूमनस्यैर्}\ H}}\skt{ }\edtext{\edtext{\skt{रसृश्च}}{\lemma{\skt{रसृश्}}\Afootnote{K; \skt{असृ}\ A H}}}{\lemma{\skt{च}}\Afootnote{K; \skt{ग्विमोक्षैश्च}\ A H}}\skt{ }\edtext{\skt{विवृद्धिमे}}{\Afootnote{K; \skt{विवृद्धिमेति}\ A H}}\skt{ }\pend
\pstart\versepar \skt{निस्तुद्यते यस्य }\edtext{\skt{शिरोतिमात्रं}}{\Afootnote{K H; \skt{शिरोऽतिमात्रं}\ A}}\skt{ }\edtext{\skt{सभक्षमाणं}}{\Afootnote{K H; \skt{संभक्ष्यमाणं}\ A}}\skt{ }\edtext{\skt{स्फुरतीव}}{\Afootnote{K; \skt{स्फुटतीव}\ A; \skt{स्फुततीव}\ H}}\skt{ }\edtext{\skt{चापि}}{\Afootnote{K H; \skt{चान्तः ॥}\ A}}\skt{ ।}\pend

\pstart\versepar \ledsidenote{\passref{6.25.11}}\edtext{\skt{ग्राणाच्च}}{\Afootnote{K H; \skt{घ्राणाच्च}\ A}}\skt{ गच्छेत्सलिलं सरक्तं }\edtext{\skt{शिरोभिघातः}}{\Afootnote{K H; \skt{शिरोभितापः}\ A}}\skt{ }\edtext{\skt{क्रिमिभिः स}}{\Afootnote{K H; \skt{कृमिभिः स}\ A}}\skt{ घोरः । }\pend
\pstart\versepar \skt{}\edtext{\skt{सूर्योदये या}}{\Afootnote{K H; \skt{सूर्योदयं या}\ A}}\skt{ }\edtext{\edtext{\skt{मन्दमक्षिभ्रुवं}}{\lemma{\skt{मन्दम्}}\Afootnote{K; \skt{प्रति मन्दमन्द}\ A; \skt{प्रति मन्दमन्दम्}\ H}}}{\lemma{\skt{अक्षिभ्रुवं}}\Afootnote{K H; \skt{मक्षिभ्रुवं}\ A}}\skt{ }\edtext{\skt{वेदनयातिगाढं}}{\Afootnote{K; \skt{रुक्समुपैति गाढम् ॥}\ A; \skt{वेदनस्यातिगाढं ॥}\ H}}\skt{ ।}\pend

\pstart\versepar \ledsidenote{\passref{6.25.12}}\skt{विवर्धते चांशुमता सहैव सूर्यापवृत्तौ विनिवर्तते च । }\pend
\pstart\versepar \skt{शीतेन शान्तिं लभते कदाचिदुष्णेन जन्तुः }\edtext{\skt{सुखमाप्नुयाद्वा}}{\lemma{\skt{आप्नुयाद्वा}}\Afootnote{K H; \skt{आप्नुयाच्च ॥}\ A}}\skt{ ।}\pend

\pstart \ledsidenote{\passref{6.25.13}}\edtext{\skt{आवर्तसंज्ञाः स तु}}{\Afootnote{K H; \skt{तं भास्करावर्तमुदाहरन्}\ A}}\skt{ }\edtext{\skt{सूर्यपूर्वा }\edtext{\skt{व्याधिर्स्मृतः}}{\Afootnote{K; \skt{व्याधिर्स्मृतः}\ H}}\skt{ पित्तसमीरणाभ्यां}}{\lemma{\skt{सूर्यपूर्वा… पित्त}}\Afootnote{\emph{om.}\ A}}\skt{ ।}\pend

\pstart \ledsidenote{\passref{6.25.15cd}}\edtext{\skt{यस्योक्त}}{\Afootnote{K H; \skt{यस्योत्त}\ A}}\skt{ }\edtext{\skt{सार्थागंमतीव}}{\Afootnote{K; \skt{माङ्गार्धमतीव}\ A; \skt{सार्थागमतीव}\ H}}\skt{ }\edtext{\skt{जन्ताः}}{\Afootnote{K H; \skt{जन्तोः}\ A}}\skt{ }\edtext{\skt{सम्भेदतोदभ्रुममोहशूलैः}}{\Afootnote{K H; \skt{संभेदतोदभ्रमशूलजुष्टम् ॥}\ A}}\skt{ ।}\pend

\pstart\versepar \ledsidenote{\passref{6.25.16}}\edtext{\skt{पक्षाद्दशाहादथवा}}{\Afootnote{K; \skt{पक्षाद्दशाहादथवाप्य्}\ A; \skt{पक्षाद्दशाहादपि चाप्य्}\ H}}\skt{ }\edtext{\skt{कदाचित्तस्यावभेदं}}{\Afootnote{K; \skt{अकस्मात्तस्यार्धमेदं}\ A; \skt{अकस्मात्तस्यार्द्धभेदन्}\ H}}\skt{ त्रितयाह्यवस्येत् । }\pend
\pstart\versepar \skt{शंखाश्रितो वायुरुदीर्ण्णवेगः कृतानुयात्रः }\edtext{\skt{कफरक्तपित्तैः}}{\Afootnote{K H; \skt{कफपित्तरक्तैः ॥}\ A}}\skt{ ।}\edtext{}{\lemma{\skt{त्रितयाह्यवस्येत् । शंखाश्रितो}}\Afootnote{K; \skt{त्रितयाद्व्यवस्येत् । शङ्खाश्रितो}\ A; \skt{त्रितया ह्यवस्येत् ॥ शंखाशितो}\ H}}\pend

\pstart\versepar \ledsidenote{\passref{6.25.17}}\edtext{\skt{रुजा}}{\Afootnote{K H; \skt{रुजः}\ A}}\skt{ }\edtext{\skt{सुतीव्रां}}{\Afootnote{K H; \skt{सुतीव्राः}\ A}}\skt{ }\edtext{\skt{प्रतनातिशीर्ष}}{\Afootnote{K; \skt{प्रतनोति मूर्ध्नि}\ A; \skt{प्रतसोतोतिशीर्ष}\ H}}\skt{ }\edtext{\skt{विशेषतञ्चापि}}{\lemma{\skt{विशेषतञ्}}\Afootnote{K H; \skt{विशेषतश्}\ A}}\skt{ हि }\edtext{\skt{संखयास्तु}}{\Afootnote{K; \skt{शङ्खयोस्तु ।}\ A; \skt{तं खयोस्तु ॥}\ H}}\skt{ । }\pend
\pstart\versepar \skt{}\edtext{\skt{सुकष्टमेतं}}{\Afootnote{K; \skt{सुकष्टमेनं}\ A; \skt{सुकष्टमेत}\ H}}\skt{ }\edtext{\skt{खलुशंखसंज्ञां}}{\Afootnote{K; \skt{खलु शङ्खकाख्यं}\ A; \skt{त्खलुशंखसंज्ञां}\ H}}\skt{ महर्षयो वेदविदः पुराणाः ।}\pend

\pstart\versepar \ledsidenote{\passref{6.25.18}}\skt{व्याधिं वदन्त्युद्गतमृत्युकल्पं }\pend
\pstart\versepar \skt{भिषक्सहस्रैरपि }\edtext{\skt{दुर्निवारमिति}}{\lemma{\skt{इति}}\Afootnote{K H; \emph{om.}\ A}}\skt{ ॥}\pend

\pstart \ledsidenote{\passref{6.25.trailer}}\edtext{\skt{शालाक्ये ॥ २३॥ {\symfont ❈}}}{\Afootnote{K; \skt{इति सुश्रुतसंहितायामुत्तरतन्त्रान्तर्गते शालाक्यतन्त्रे शिरोरोगविज्ञानीयो नाम पञ्चविंशोऽध्यायः ॥२५ ॥}\ A; \skt{शालाक्ये शिरोगतरोगविज्ञानीयं त्रयोविंशतिमोऽध्यायः ॥ ॥}\ H}}\skt{ ॥}\pend

\endnumbering
\beginnumbering
\eledsection*{Adhyāya 26}
\pstart \ledsidenote{\passref{6.26.1}}\skt{अथातः }\edtext{\edtext{\skt{शिरोगतरोगप्रतिषेधमध्यायं}}{\lemma{\skt{शिरोगतरोगप्रतिषेधम्}}\Afootnote{K H; \skt{शिरोरोगप्रतिषेधं}\ A}}}{\lemma{\skt{अध्यायं}}\Afootnote{K H; \emph{om.}\ A}}\skt{ }\edtext{\skt{व्याख्यास्यामः}}{\Afootnote{A H; \skt{व्या व्याख्यास्यामः ॥}\ K}}\skt{ ॥}\pend

\pstart\versepar \ledsidenote{\passref{6.26.3}}\edtext{\skt{वातव्याधिविधिं}}{\Afootnote{K; \skt{वातव्याधिविधिः}\ A; \skt{वातव्याधिविधि}\ H}}\skt{ }\edtext{\skt{कुर्यात्च्छिरोरोगे}}{\Afootnote{K; \skt{कार्यः शिरोरोगे}\ A; \skt{कुर्यात्च्छिरोरोगे}\ H}}\skt{ }\edtext{\skt{निलात्मके}}{\Afootnote{K H; \skt{ऽनिलात्मके ।}\ A}}\skt{ । }\pend
\pstart\versepar \skt{}\edtext{\skt{पयोनुपानं}}{\Afootnote{K; \skt{पयो नुपानं}\ A; \skt{पयोऽनुपानं}\ H}}\skt{ सेवेत घृतन्तैलमथापि वा ।}\pend

\pstart\versepar \ledsidenote{\passref{6.26.4}}\edtext{\skt{मुद्गान्माषां}}{\Afootnote{K; \skt{मुद्गा}\ A; \skt{मुद्गान्माषा}\ H}}\skt{ }\edtext{\skt{कुलंत्थाम्श्च}}{\Afootnote{K; \skt{न्कुलत्थान्माषांश्च}\ A; \skt{न्कुलत्थाञ्श्च}\ H}}\skt{ }\edtext{\skt{खादेद्वानिशि}}{\Afootnote{K; \skt{खादेच्च निशि}\ A; \skt{खादेद्वा निशि}\ H}}\skt{ }\edtext{\skt{केवलां}}{\Afootnote{K H; \skt{केवला}\ A}}\skt{ । }\pend
\pstart\versepar \skt{}\edtext{\skt{कट्वाष्णां}}{\Afootnote{K; \skt{न् । कटूष्णांश्}\ A; \skt{कटूष्णा}\ H}}\skt{ }\edtext{\skt{समसर्पिष्कानुष्णं}}{\Afootnote{K; \skt{च ससर्पिष्कानुष्णं}\ A; \skt{न्समसर्प्पिष्कानुष्णञ्}\ H}}\skt{ }\edtext{\skt{चानुपिवेत्पयः}}{\Afootnote{K H; \skt{चानु पयः पिबेत् ॥}\ A}}\skt{ ।}\pend

\pstart\versepar \ledsidenote{\passref{6.26.5}}\edtext{\skt{पिवेद्वा}}{\Afootnote{K H; \skt{पिबेद्वा}\ A}}\skt{ पयसा तैलं }\edtext{\edtext{\skt{तत्कालश्चापि}}{\lemma{\skt{तत्कालश्}}\Afootnote{K; \skt{तत्कल्कं}\ A; \skt{तत्कालम्}\ H}}}{\lemma{\skt{चापि}}\Afootnote{K; \skt{वापि}\ A H}}\skt{ मानवः । }\pend
\pstart\versepar \skt{वातघ्नसिद्धैः क्षीरैश्च }\edtext{\skt{सुखोष्णैर्ल्लेपयेच्छिरः}}{\Afootnote{K; \skt{सुखोष्णैः सेकमाचरेत् ॥}\ A; \skt{सुखोष्णैर्ल्लेपयेच्छिरः ॥}\ H}}\skt{ ।}\pend

\pstart \ledsidenote{\passref{6.26.6}}\edtext{\skt{स्विन्नैर्वा}}{\Afootnote{K; \skt{तत्सिद्धैः पायसैर्वापि सुखोष्णैर्लेपयेच्छिरः । स्विन्नैर्वा}\ A; \skt{स्निन्नैर्व्वा}\ H}}\skt{ }\edtext{\skt{मत्स्यापिशितैः}}{\Afootnote{K H; \skt{मत्स्यपिशितैः}\ A}}\skt{ कृशरैर्वा ससैन्धवैः ।}\pend

\pstart\versepar \ledsidenote{\passref{6.26.7}}\skt{चन्दनोत्पलकुष्ठैर्वा }\edtext{\skt{स्निग्धैर्म्मार्गधिवायुतैः}}{\Afootnote{K H; \skt{सुश्लक्ष्णैर्मगधायुतैः ।}\ A}}\skt{ ॥ }\pend
\pstart\versepar \skt{}\edtext{\skt{स्विन्नस्य}}{\Afootnote{K; \skt{स्निग्धस्य}\ A; \skt{स्निन्नस्य}\ H}}\skt{ }\edtext{\skt{नस्य}}{\Afootnote{K H; \skt{तैलं नस्यं}\ A}}\skt{ }\edtext{\skt{तैलस्यात्कुलीररससाधितं}}{\Afootnote{K H; \skt{स्यात्कुलीररससाधितम् ॥}\ A}}\skt{ ।}\pend

\pstart\versepar \ledsidenote{\passref{6.26.8}}\skt{वरुणादौ गणेक्षुण्णे }\edtext{\skt{क्षीरमेवोदकं}}{\Afootnote{K; \skt{क्षीरमर्धोदकं}\ A H}}\skt{ पचेत् । }\pend
\pstart\versepar \skt{}\edtext{\skt{क्षीरावसिष्टन्तच्छीतंमथित्वासारमुद्वरेत्}}{\lemma{\skt{तच्छीतंमथित्वासारमुद्वरेत्}}\Afootnote{K; \skt{च तन्मथ्यं शीतं सारमुपाहरेत् ॥}\ A; \skt{तच्छीतस्मथित्वासारमुद्वरेत् ।}\ H}}\skt{ ।}\edtext{}{\lemma{\skt{पचेत् । क्षीरावसिष्टन्}}\Afootnote{K; \skt{पचेत् । क्षीरशेषं}\ A; \skt{पचेत् ॥ क्षीरावलिष्टन्}\ H}}\pend

\pstart\versepar \ledsidenote{\passref{6.26.9}}\edtext{\skt{तद्वा}}{\Afootnote{K H; \skt{ततो}\ A}}\skt{ मधुरकैः सिद्धं }\edtext{\skt{नस्यतः}}{\Afootnote{K H; \skt{नस्ये त}\ A}}\skt{ }\edtext{\skt{पूजितं}}{\Afootnote{K H; \skt{त्पूजितं}\ A}}\skt{ हविः । }\pend
\pstart\versepar \skt{}\edtext{\skt{अस्मिं}}{\Afootnote{K; \skt{तस्मि}\ A; \skt{अस्मि}\ H}}\skt{ }\edtext{\skt{विपक्वं}}{\Afootnote{K; \skt{न्विपक्वे}\ A; \skt{न्विपक्वं}\ H}}\skt{ }\edtext{\skt{क्षीरे च}}{\Afootnote{K H; \skt{क्षीरे तु}\ A}}\skt{ पेयं सर्पिः सशर्करं ।}\pend

\pstart\versepar \ledsidenote{\passref{6.26.10}}\skt{धूमं चास्य यथाकालं }\edtext{\skt{स्नेहिकं}}{\Afootnote{K; \skt{स्नैहिकं}\ A H}}\skt{ योजयेद्भिषक् । }\pend
\pstart\versepar \skt{पानाभ्यञ्जननस्येषु }\edtext{\skt{बस्तिकर्मणि}}{\Afootnote{\skt{बसिकर्मणि}\ A; \skt{वस्तिकर्मणि}\ K H}}\skt{ सेचने ।}\edtext{}{\lemma{\skt{योजयेद्भिषक् । पानाभ्यञ्जननस्येषु}}\Afootnote{A K; \skt{योजयेद्भिषल् ॥ पानाभ्यञ्जननस्येषु}\ H}}\pend

\pstart\versepar \ledsidenote{\passref{6.26.11}}\edtext{\edtext{\skt{त्रिवृतश्चात्र}}{\lemma{\skt{त्रिवृतश्}}\Afootnote{K; \skt{विदध्यात्त्रैवृतं}\ A; \skt{त्रिवृत्तं}\ H}}}{\lemma{\skt{चात्र}}\Afootnote{K H; \emph{om.}\ A}}\skt{ }\edtext{\skt{देयं}}{\Afootnote{K; \skt{धी}\ A; \skt{पेयं}\ H}}\skt{ }\edtext{\skt{स्याद्वलातैलमथापि}}{\lemma{\skt{स्याद्वलातैलम्}}\Afootnote{K H; \skt{मान्बलातैलम्}\ A}}\skt{ वा । }\pend
\pstart\versepar \skt{भोजयेच्च रसैः स्निग्धैः पयोभिर्वा सुसंस्कृतैः ।}\pend

\pstart\versepar \ledsidenote{\passref{6.26.12}}\edtext{\skt{पित्तरक्तसमुत्थानं}}{\Afootnote{K H; \skt{पित्तरक्तसमुत्थानौ}\ A}}\skt{ }\edtext{\skt{शिरोरोगं}}{\Afootnote{K H; \skt{शिरोरोगौ}\ A}}\skt{ निवाचयेत् । }\pend
\pstart\versepar \skt{शिरः }\edtext{\skt{प्रलयैः}}{\Afootnote{K H; \skt{रोलेपैः}\ A}}\skt{ }\edtext{\skt{सघृतैः}}{\Afootnote{K H; \skt{ससर्पिष्कैः}\ A}}\skt{ परिषेकैश्च शीतलैः ।}\edtext{}{\lemma{\skt{निवाचयेत् । शिरः}}\Afootnote{K; \skt{निवारयेत् । शि}\ A; \skt{निधाचयेत् ॥ शिरः}\ H}}\pend

\pstart\versepar \ledsidenote{\passref{6.26.13}}\edtext{\skt{क्षीरेक्षुरसधान्याम्लमस्तुक्षौद्रसितायुतैः}}{\Afootnote{K H; \skt{क्षीरेक्षुरसधान्याम्लमस्तुक्षौद्रसिताजलैः ।}\ A}}\skt{ । }\pend
\pstart\versepar \skt{}\edtext{\skt{नलवश्जामेवाल}}{\Afootnote{K; \skt{नलवञ्जुलकह्लार}\ A; \skt{नलवंजुलवाल}\ H}}\skt{ }\edtext{\skt{चन्दनो}}{\Afootnote{K H; \skt{चन्दनो}\ A}}\skt{ }\edtext{\skt{शीरपद्मकैः}}{\Afootnote{K H; \skt{त्पलपद्मकैः ॥}\ A}}\pend

\pstart\versepar \ledsidenote{\passref{6.26.14}}\edtext{\skt{शंखगैरिककल्हार}}{\Afootnote{K; \skt{वंशशैवलयष्ट्याह्व}\ A; \skt{शंखगैर्ककल्हार}\ H}}\skt{ मुस्ताम्भोरुहसंयुतैः । }\pend
\pstart\versepar \skt{शिरः प्रलेपैः }\edtext{\skt{सघृतैः}}{\Afootnote{K H; \skt{सघृतै}\ A}}\skt{ }\edtext{\skt{सर्व्वैर्वा}}{\Afootnote{K; \skt{र्वैस}\ A; \skt{सर्व्व वा}\ H}}\skt{ }\edtext{\skt{पयसायुतैः}}{\Afootnote{K; \skt{र्पैश्च तथाविधैः ॥}\ A; \skt{ओअयसायुतैः ॥}\ H}}\pend

\pstart\versepar \ledsidenote{\passref{6.26.15}}\edtext{\skt{मुखालेपैश्च}}{\Afootnote{K H; \emph{om.}\ A}}\skt{ }\edtext{\edtext{\skt{मधुरैर्नस्यकर्म्मादिभेषजं}}{\lemma{\skt{मधुरैर्}}\Afootnote{K; \skt{मधुरैश्च}\ A; \skt{मधुरै}\ H}}}{\lemma{\skt{नस्यकर्म्मादिभेषजं}}\Afootnote{K; \skt{मुखालैपैर्नस्यकर्मभिरेव च ।}\ A; \skt{र्न्नस्त कर्म्मादिभेषजं ॥}\ H}}\skt{ । }\pend
\pstart\versepar \skt{आस्थापनैर्विरेकैश्च }\edtext{\skt{पत्थैश्च}}{\Afootnote{K H; \skt{पथ्यैश्च}\ A}}\skt{ }\edtext{\skt{स्नेहबस्तिभिः}}{\Afootnote{A; \skt{स्नेहवस्तिभिः ।}\ K H}}\skt{ ।}\pend

\pstart\versepar \ledsidenote{\passref{6.26.16}}\skt{क्षीरसर्पिर्हितन्नस्यं वसा वा }\edtext{\skt{जाङ्गलाः}}{\Afootnote{K H; \skt{जाङ्गला}\ A}}\skt{ शुभाः । }\pend
\pstart\versepar \skt{उत्पलादिविपक्वेन }\edtext{\skt{क्षीरेणास्थापनं}}{\Afootnote{A K; \skt{क्षीरेणास्थापन्}\ H}}\skt{ }\edtext{\skt{तथा}}{\Afootnote{K H; \skt{हितम् ॥}\ A}}\skt{ ।}\pend

\pstart\versepar \ledsidenote{\passref{6.26.17}}\skt{भोजनं जाङ्गलरसैः सर्पिषा चानुवासनं । }\pend
\pstart\versepar \skt{}\edtext{\skt{मधुरं}}{\Afootnote{K H; \skt{मधुरैः}\ A}}\skt{ }\edtext{\edtext{\skt{क्षीरसर्पिश्च}}{\lemma{\skt{क्षीरसर्पिश्}}\Afootnote{K H; \skt{क्षीरसर्पिस्तु}\ A}}}{\lemma{\skt{च}}\Afootnote{K H; \skt{स्नेहने च}\ A}}\skt{ }\edtext{\skt{सक्षौद्रं}}{\Afootnote{K H; \skt{स}\ A}}\skt{ }\edtext{\skt{शर्करायुतं}}{\Afootnote{K H; \skt{शर्करम् ॥}\ A}}\skt{ ।}\pend

\pstart\versepar \ledsidenote{\passref{6.26.18}}\edtext{\skt{पित्तरक्तसमुद्दिष्टं}}{\Afootnote{K H; \skt{पित्तरक्तघ्नमुद्दिष्टं}\ A}}\skt{ }\edtext{\skt{यच्चान्यदपि}}{\Afootnote{K H; \skt{यच्चान्यदपि}\ A}}\skt{ }\edtext{\skt{तद्विधं}}{\Afootnote{K H; \skt{तद्धितम् ।}\ A}}\skt{ ॥ }\pend
\pstart\versepar \skt{कफोत्थितं }\edtext{\skt{शिरोरोगञ्जयेत्कफविनाशनः}}{\Afootnote{K; \skt{शिरोरोगं जयेत्कफनिवारणैः ॥}\ A; \skt{शिरोरोगञ्जयेत्कफविनाशनैः ॥}\ H}}\skt{ ॥}\pend

\pstart \ledsidenote{\passref{6.26.19ab}}\edtext{\edtext{\edtext{\edtext{\skt{शिरोविरेकैर्वमनैस्तीक्ष्णैर्गण्डुषधारणैः}}{\lemma{\skt{शिरोविरेकैर्}}\Afootnote{A K; \skt{शिरोविरेकै}\ H}}}{\lemma{\skt{वमनैस्}}\Afootnote{A K; \skt{र्व्वमनै}\ H}}}{\lemma{\skt{तीक्ष्णैर्}}\Afootnote{A K; \skt{स्तीक्ष्णै}\ H}}}{\lemma{\skt{गण्डुषधारणैः}}\Afootnote{K; \skt{गण्डूषधारणैः ।}\ A; \skt{र्गण्डुषधारणैः ॥}\ H}}\skt{ ।}\pend

\pstart\versepar \ledsidenote{\passref{6.26.20}}\edtext{\skt{मधुकसारेणा}}{\Afootnote{K; \skt{शिरो मधूकसारेण}\ A; \skt{मधुकसारेण}\ H}}\skt{ }\edtext{\skt{शिरः}}{\Afootnote{K H; \emph{om.}\ A}}\skt{ स्निग्धं }\edtext{\skt{चास्य}}{\Afootnote{K H; \skt{चापि}\ A}}\skt{ विरेचयेत् । }\pend
\pstart\versepar \skt{इङ्गुदस्य }\edtext{\skt{धवत्वग्नि}}{\Afootnote{K H; \skt{त्वचा वापि}\ A}}\skt{ मेषशृङ्गस्य }\edtext{\skt{वायुनः}}{\Afootnote{K; \skt{वा भिषक् ॥}\ A; \skt{वायुनं ॥}\ H}}\skt{ ।}\pend

\pstart\versepar \ledsidenote{\passref{6.26.21}}\edtext{\skt{एभेनेव}}{\Afootnote{K; \skt{आभ्यामेव}\ A; \skt{एभेव}\ H}}\skt{ }\edtext{\skt{कृतावर्त्तिधूमपाने}}{\Afootnote{K H; \skt{कृतां वर्तिं धूमपाने}\ A}}\skt{ प्रयोजयेत् । }\pend
\pstart\versepar \skt{कफघ्नां }\edtext{\skt{वमनाञ्जात्र}}{\Afootnote{K; \skt{यं कट्फ}\ A; \skt{वमनञ्जात्र}\ H}}\skt{ }\edtext{\skt{जिघ्रच्चूर्ण्णञ्च}}{\Afootnote{K; \skt{लचूर्णं च}\ A; \skt{जिघ्रे चूर्ण्णञ्ज}\ H}}\skt{ }\edtext{\skt{कट्फलं}}{\Afootnote{K H; \skt{कवलाश्च कफापहाः ॥}\ A}}\skt{ ।}\edtext{}{\lemma{\skt{प्रयोजयेत् । कफघ्नां}}\Afootnote{K; \skt{प्रयोजयेत् । घ्रे}\ A; \skt{प्रयोजयेत् । कफघ्रेम्}\ H}}\pend

\pstart\versepar \ledsidenote{\passref{6.26.22}}\edtext{\skt{सरलाकाष्ठशार्ङ्गष्ठादेवकाष्ठैः}}{\Afootnote{K H; \skt{सरलाकुष्ठशार्ङ्गेष्टादेवकाष्ठैः}\ A}}\skt{ सरोहिषैः । }\pend
\pstart\versepar \skt{क्षारपिष्टैः सलवणैः }\edtext{\skt{सुखोष्णैर्ल्लेपयेच्छिरः}}{\Afootnote{A K; \skt{सुखोष्णैर्ल्लेपयेच्छिरः ॥}\ H}}\skt{ ।}\pend

\pstart\versepar \ledsidenote{\passref{6.26.23}}\edtext{\skt{यवषष्ठिकयोश्चान्नं}}{\lemma{\skt{यवषष्ठिकयोश्}}\Afootnote{K H; \skt{यवषष्टिकयोश्}\ A}}\skt{ व्योषक्षारसमायुतं । }\pend
\pstart\versepar \skt{}\edtext{\skt{पटोलनिम्वकौलत्थीमात्रावद्भोजयेद्रसैः}}{\Afootnote{K; \skt{पटोलमुद्गकौलत्थैर्मात्रावद्भोजयेद्रसैः ॥}\ A; \skt{पटोलमुद्गकौलत्थैर्म्मात्रावद्भोजयेद्रसैः ॥}\ H}}\skt{ ॥}\pend

\pstart\versepar \ledsidenote{\passref{6.26.24}}\edtext{\skt{शिरोरोगत्रिदोषोत्थे}}{\Afootnote{K H; \skt{शिरोरोगे त्रिदोषोत्थे}\ A}}\skt{ त्रिदोषघ्नो विधिर्हितः । }\pend
\pstart\versepar \skt{सर्पिः }\edtext{\skt{पानविशेषेण}}{\Afootnote{K H; \skt{पानं विशेषेण}\ A}}\skt{ पुराणम्वा दिशन्ति हि ॥}\pend

\pstart\versepar \ledsidenote{\passref{6.26.25}}\skt{क्षयजे }\edtext{\skt{क्षयमासाध्य}}{\Afootnote{K H; \skt{क्षयमासाद्य}\ A}}\skt{ कर्तव्यो }\edtext{\skt{वृंहणो}}{\Afootnote{K H; \skt{बृंहणो}\ A}}\skt{ विधिः । }\pend
\pstart\versepar \skt{पाने नस्ये च सर्पिः }\edtext{\skt{स्यद्वातघ्नमधुरैः}}{\Afootnote{K H; \skt{स्याद्वातघ्नमधुरैः}\ A}}\skt{ }\edtext{\skt{श्रितं}}{\Afootnote{K H; \skt{शृतम् ॥}\ A}}\skt{ ।}\pend

\pstart\versepar \ledsidenote{\passref{6.26.26}}\skt{क्षयकासापहं चात्र सर्पिः पथ्यतमं }\edtext{\skt{सदा}}{\Afootnote{K H; \skt{विदुः ।}\ A}}\skt{ ॥ }\pend
\pstart\versepar \skt{}\edtext{\edtext{\skt{क्रिमिभिर्भक्षमानस्य}}{\lemma{\skt{क्रिमिभिर्}}\Afootnote{K H; \skt{कृमिभिर्}\ A}}}{\lemma{\skt{भक्षमानस्य}}\Afootnote{K; \skt{भक्ष्यमाणस्य}\ A; \skt{भक्षमाणस्य}\ H}}\skt{ }\edtext{\skt{वक्षते}}{\Afootnote{K; \skt{वक्ष्यते}\ A; \skt{वक्ष}\ H}}\skt{ शिरसः }\edtext{\skt{क्रिया}}{\Afootnote{A K; \skt{क्रियां ॥}\ H}}\skt{ ।}\pend

\pstart\versepar \ledsidenote{\passref{6.26.27}}\edtext{\skt{शोणितं}}{\Afootnote{K H; \skt{नस्ये हि शोणितं}\ A}}\skt{ }\edtext{\skt{नस्यतो}}{\Afootnote{K H; \emph{om.}\ A}}\skt{ दद्यात्तेन मूर्च्छन्ति जन्तवः । }\pend
\pstart\versepar \skt{}\edtext{\skt{मत्ताः}}{\Afootnote{A K; \skt{मत्ता}\ H}}\skt{ शोणितगन्धेन }\edtext{\skt{समयान्ति}}{\Afootnote{K H; \skt{समायान्ति}\ A}}\skt{ }\edtext{\skt{ततस्तुतः}}{\Afootnote{K H; \skt{यतस्ततः ॥}\ A}}\skt{ ।}\pend

\pstart\versepar \ledsidenote{\passref{6.26.28}}\skt{तेषां }\edtext{\skt{निर्हरणां}}{\Afootnote{K; \skt{निर्हरणं}\ A; \skt{निह्र्हरणं}\ H}}\skt{ कार्यं ततः }\edtext{\skt{शीर्षविरेचनैः}}{\Afootnote{K H; \skt{मूर्धविरेचनैः ।}\ A}}\skt{ । }\pend
\pstart\versepar \skt{}\edtext{\skt{ह्रस्वशिग्रुकवीजैश्च}}{\Afootnote{K H; \skt{ह्रस्वशिग्रुकबीजैर्वा}\ A}}\skt{ }\edtext{\skt{कांसनीलसमायुतैः}}{\Afootnote{K H; \skt{कांस्यनीलीसमायुतैः ॥}\ A}}\skt{ ।}\pend

\pstart\versepar \ledsidenote{\passref{6.26.29}}\skt{कृमिघ्नैरवपीडैश्च मूत्रपिष्टैरुपाचरेत् । }\pend
\pstart\versepar \skt{पूतिमत्स्यकृतां }\edtext{\skt{धूमां}}{\Afootnote{K H; \skt{धूमा}\ A}}\skt{ }\edtext{\skt{क्रिमिघ्नाञ्श्च}}{\Afootnote{K H; \skt{न्कृमिघ्नांश्च}\ A}}\skt{ प्रयोजयेत् ।}\edtext{}{\lemma{\skt{उपाचरेत् । पूतिमत्स्यकृतां}}\Afootnote{K; \skt{उपाचरेत् । पूतिमत्स्ययुतान्}\ A; \skt{उपाचरेत् । पूतिमत्स्यं कृतां}\ H}}\pend

\pstart\versepar \ledsidenote{\passref{6.26.30}}\skt{भोजनानि }\edtext{\skt{क्रिमिघ्नानि}}{\Afootnote{K H; \skt{कृमिघ्नानि}\ A}}\skt{ पानानि }\edtext{\skt{विविधाणि च}}{\Afootnote{K; \skt{विविधानि च ।}\ A H}}\skt{ ॥ }\pend
\pstart\versepar \skt{सूर्यावर्तेविधातव्यं नस्यकर्मादिभेषजं ।}\pend

\pstart \ledsidenote{\passref{6.26.31cd}}\skt{तथा }\edtext{\skt{भवेद्देह}}{\Afootnote{K H; \skt{र्धभेदके}\ A}}\skt{ व्याधौ }\edtext{\skt{प्राप्तमन्यच्च}}{\Afootnote{K; \skt{प्राप्तमन्यच्च}\ A; \skt{प्र्प्तमन्यच्च}\ H}}\skt{ यद्भवेत् ।}\pend

\pstart \ledsidenote{\passref{6.26.32ab}}\edtext{\skt{शिरीषमूलकफलैरवपीडनयोर्हितः}}{\lemma{\skt{अवपीडनयोर्}}\Afootnote{K H; \skt{अवपीडोऽनयोर्}\ A}}\skt{ ॥}\pend

\pstart \ledsidenote{\passref{6.26.33}}\edtext{\skt{अवभेदहितो}}{\Afootnote{K; \skt{अवपीडो हितश्}\ A; \skt{अर्द्वभेदहितश्}\ H}}\skt{ }\edtext{\skt{वास्याद्व चा}}{\Afootnote{K; \skt{चात्र वचा}\ A; \skt{चस्याद्व चा}\ H}}\skt{ }\edtext{\skt{पिप्पलिसंयुतं}}{\Afootnote{K H; \skt{मागधिकायुतः । मधुकेनाव्पीडो वा मधुना सह संयुतः ॥}\ A}}\skt{ ।}\pend

\pstart\versepar \ledsidenote{\passref{6.26.35}}\edtext{\skt{शारिवोत्पलकुष्ठानि}}{\Afootnote{K H; \skt{सारिवोत्पलकुष्ठानि}\ A}}\skt{ मधुकं चाम्लपेषितं । }\pend
\pstart\versepar \skt{}\edtext{\skt{सर्पिः}}{\Afootnote{K H; \skt{सर्पि}\ A}}\skt{ }\edtext{\skt{क्षौद्रयुतो}}{\Afootnote{K H; \skt{स्तैलयुतो}\ A}}\skt{ लेपो द्वयोरपि सुखावहः ।}\pend

\pstart \ledsidenote{\passref{6.26.37cd}}\skt{आहारश्च विधातव्यो }\edtext{\skt{वातपित्तप्रणाशनः}}{\Afootnote{K H; \skt{वातपित्तविनाशनः ॥}\ A}}\skt{ ।}\pend

\pstart\versepar \ledsidenote{\passref{6.26.38}}\edtext{\skt{मधुशीर्दोषसंयाव}}{\Afootnote{K H; \skt{मधुमस्तकसंयाव}\ A}}\skt{ }\edtext{\skt{हविः}}{\Afootnote{K H; \skt{घृत}\ A}}\skt{ }\edtext{\skt{पूरैश्च यः}}{\Afootnote{K H; \skt{पूरैश्च भो}\ A}}\skt{ }\edtext{\skt{क्रमः}}{\Afootnote{K H; \skt{जनम् ।}\ A}}\skt{ ॥ }\pend
\pstart\versepar \skt{क्षीरसर्पिः प्रशंसन्ति }\edtext{\skt{नस्यः}}{\Afootnote{K H; \skt{नस्ये}\ A}}\skt{ }\edtext{\skt{पानञ्च}}{\Afootnote{K H; \skt{पाने च}\ A}}\skt{ शंखके ।}\pend

\pstart\versepar \ledsidenote{\passref{6.26.39}}\edtext{\skt{उदकानां}}{\Afootnote{K H; \skt{जाङ्गलानां}\ A}}\skt{ }\edtext{\skt{रसैः}}{\Afootnote{A K; \skt{रसै।}\ H}}\skt{ }\edtext{\skt{स्निग्धैराहारश्चात्र}}{\lemma{\skt{स्निग्धैराहारश्}}\Afootnote{A K; \skt{स्निग्धैराहारश्}\ H}}\skt{ शस्यते । }\pend
\pstart\versepar \skt{}\edtext{\skt{शतावरीं}}{\Afootnote{A K; \skt{शतावरी}\ H}}\skt{ }\edtext{\skt{कृष्णतिलां}}{\Afootnote{K H; \skt{तिला}\ A}}\skt{ }\edtext{\skt{मधुकं}}{\Afootnote{K H; \skt{न्कृष्णान्मधुकं}\ A}}\skt{ }\edtext{\skt{उत्पलं}}{\Afootnote{K H; \skt{नीलमुत्पलम् ॥}\ A}}\skt{ ।}\pend

\pstart\versepar \ledsidenote{\passref{6.26.40}}\edtext{\edtext{\skt{दूर्वापुनर्न्नवाम्वापि}}{\lemma{\skt{दूर्वापुनर्न्नवाम्}}\Afootnote{K H; \skt{दूर्वां पुनर्नवां}\ A}}}{\lemma{\skt{वापि}}\Afootnote{K H; \skt{चैव}\ A}}\skt{ }\edtext{\skt{लेपं}}{\Afootnote{K H; \skt{लेपे}\ A}}\skt{ साध्ववचारयेत् । }\pend
\pstart\versepar \skt{महासुगन्धमथवा पालिन्दीञ्चाम्लपेषितां ।}\edtext{}{\lemma{\skt{साध्ववचारयेत् । महासुगन्धम्}}\Afootnote{K H; \skt{साध्ववचारयेत् । महासुगन्धाम्}\ A}}\pend

\pstart\versepar \ledsidenote{\passref{6.26.41}}\edtext{\skt{शीतां श्वा पि}}{\Afootnote{K; \skt{शीतांश्चात्र}\ A; \skt{शीतं शा पि}\ H}}\skt{ }\edtext{\skt{परीषेकां}}{\Afootnote{K H; \skt{परीषेका}\ A}}\skt{ }\edtext{\skt{प्रदेहांश्चात्रावतारयेत्}}{\lemma{\skt{प्रदेहांश्}}\Afootnote{K H; \skt{न्प्रदेहांश्}\ A}}\skt{ । }\pend
\pstart\versepar \skt{सूर्यावर्तापहश्चास्मिन्नवपीडं }\edtext{\skt{प्रयोजयेत्}}{\Afootnote{K H; \skt{प्रयोजयेत् । अवपीडश्च देयोऽत्र सूर्यावर्तनिवारणः ॥}\ A}}\skt{ ।}\edtext{}{\lemma{\skt{चात्रावतारयेत्… नव}}\Afootnote{\emph{om.}\ A}}\edtext{}{\lemma{\skt{चात्रावतारयेत् । सूर्यावर्तापहश्}}\Afootnote{K; \skt{चावतार{\symfont ¦} येत् । सूर्यावर्तापहश्}\ H}}\pend

\pstart\versepar \ledsidenote{\passref{6.26.42}}\edtext{\skt{क्षयक्रिमिकृतौ}}{\Afootnote{K H; \skt{कृमिक्षयकृतौ}\ A}}\skt{ हित्वा शिरोरोगेषु वुद्धिमान् । }\pend
\pstart\versepar \skt{मधुतैलसमायुक्तैः }\edtext{\skt{शिरांस्यभिविरेचयेत्}}{\Afootnote{K H; \skt{शिरांस्यतिविरेचयेत् ॥}\ A}}\skt{ ।}\edtext{}{\lemma{\skt{वुद्धिमान् । मधुतैलसमायुक्तैः}}\Afootnote{K H; \skt{बुद्धिमान् । मधुतैलसमायुक्तैः}\ A}}\pend

\pstart\versepar \ledsidenote{\passref{6.26.43}}\skt{पश्चात्सर्षपतैलेन ततो }\edtext{\skt{नस्यः}}{\Afootnote{K H; \skt{नस्यं}\ A}}\skt{ प्रशस्यते । }\pend
\pstart\versepar \skt{न चेच्छान्तिम्व्रजन्त्येवं }\edtext{\skt{स्निग्धस्विन्नन्ततो}}{\lemma{\skt{स्निग्धस्विन्नन्}}\Afootnote{K H; \skt{स्निग्धस्विन्नांस्}\ A}}\skt{ भिषक् ।}\edtext{}{\lemma{\skt{प्रशस्यते । न}}\Afootnote{K H; \skt{प्रयोजयेत् । न}\ A}}\pend

\pstart\versepar \ledsidenote{\passref{6.26.44}}\edtext{\skt{एतानुपक्रमेत्सम्यक्शिराणां}}{\Afootnote{K H; \skt{पश्चादुपाचरेत्सम्यक्सिराणाम्}\ A}}\skt{ }\edtext{\skt{मृदुदारणैः}}{\Afootnote{K H; \skt{अथ मोक्षणैः ।}\ A}}\skt{ ॥ }\pend
\pstart\versepar \skt{षट्सप्ततिर्नेत्ररोगा }\edtext{\skt{विंशदष्टौरकर्ण्णजाः}}{\Afootnote{K; \skt{दशाष्टादश कर्णजाः ॥}\ A; \skt{विंशदष्टौकर्ण्णजाः ।}\ H}}\skt{ ।}\pend

\pstart\versepar \ledsidenote{\passref{6.26.45}}\skt{एकत्रिंशद्घ्राणगताः }\edtext{\skt{शिरोरोगा}}{\Afootnote{\skt{शिरस्येका}\ A; \skt{शिरोर्गा}\ K H}}\skt{ दशैव तु ॥ }\pend
\pstart\versepar \skt{}\edtext{\skt{इति ह}}{\Afootnote{K H; \skt{इति}\ A}}\skt{ }\edtext{\skt{विस्तरणः}}{\Afootnote{K H; \skt{विस्तरतो}\ A}}\skt{ }\edtext{\skt{प्रोक्ताः}}{\Afootnote{K H; \skt{दृष्टाः(दिष्टाः)}\ A}}\skt{ }\edtext{\skt{सलक्षणा}}{\Afootnote{K H; \skt{सलक्षण}\ A}}\skt{ चिकित्सिताः ।}\pend

\pstart\versepar \ledsidenote{\passref{6.26.46}}\skt{संहितायामभिहिताः पञ्चषष्टिमुखामयाः । }\pend
\pstart\versepar \skt{}\edtext{\skt{एतावन्तो}}{\Afootnote{A K; \emph{om.}\ H}}\skt{ }\edtext{\edtext{\skt{यथास्थूलमुत्तमाङ्गगता}}{\lemma{\skt{उत्तम्}}\Afootnote{K H; \skt{उत्त}\ A}}}{\lemma{\skt{आङ्गगता}}\Afootnote{K H; \skt{माङ्गगता}\ A}}\skt{ गदाः ॥ }\pend
\pstart\versepar \skt{अस्मिं }\edtext{\skt{छास्त्रस्तु}}{\Afootnote{K H; \skt{छास्त्रे}\ A}}\skt{ }\edtext{\skt{निपतिताः}}{\Afootnote{K H; \skt{निगदिताः}\ A}}\skt{ }\edtext{\skt{संख्यारूपचिकित्सितैरिति}}{\lemma{\skt{इति}}\Afootnote{K H; \emph{om.}\ A}}\skt{ ॥}\edtext{}{\lemma{\skt{… ए}}\Afootnote{\emph{om.}\ A}}\pend

\pstart \ledsidenote{\passref{6.26.trailer}}\edtext{\skt{शालाक्ये ॥ २४ ॥ {\symfont ❈}}}{\Afootnote{K; \skt{इति सुश्रुतसंहितायामुत्तरतन्त्रान्तर्गते शालाक्यतन्त्रे शिरोरोगप्रतिषेधो नाम षड्विंशोऽध्यायः ॥ २६ ॥}\ A; \skt{शालाक्ये शिरोगतरोगप्रतिषेधश्चतुर्व्विंशतितमोऽध्यायः ॥ ॥}\ H}}\skt{ ॥}\pend

\endnumbering
